% Options for packages loaded elsewhere
\PassOptionsToPackage{unicode}{hyperref}
\PassOptionsToPackage{hyphens}{url}
\documentclass[
  british,
  a4paper,
]{book}
\usepackage{xcolor}
\usepackage{amsmath,amssymb}
\setcounter{secnumdepth}{5}
\usepackage{iftex}
\ifPDFTeX
  \usepackage[T1]{fontenc}
  \usepackage[utf8]{inputenc}
  \usepackage{textcomp} % provide euro and other symbols
\else % if luatex or xetex
  \usepackage{unicode-math} % this also loads fontspec
  \defaultfontfeatures{Scale=MatchLowercase}
  \defaultfontfeatures[\rmfamily]{Ligatures=TeX,Scale=1}
\fi
\usepackage{lmodern}
\ifPDFTeX\else
  % xetex/luatex font selection
\fi
% Use upquote if available, for straight quotes in verbatim environments
\IfFileExists{upquote.sty}{\usepackage{upquote}}{}
\IfFileExists{microtype.sty}{% use microtype if available
  \usepackage[]{microtype}
  \UseMicrotypeSet[protrusion]{basicmath} % disable protrusion for tt fonts
}{}
\makeatletter
\@ifundefined{KOMAClassName}{% if non-KOMA class
  \IfFileExists{parskip.sty}{%
    \usepackage{parskip}
  }{% else
    \setlength{\parindent}{0pt}
    \setlength{\parskip}{6pt plus 2pt minus 1pt}}
}{% if KOMA class
  \KOMAoptions{parskip=half}}
\makeatother
\usepackage{color}
\usepackage{fancyvrb}
\newcommand{\VerbBar}{|}
\newcommand{\VERB}{\Verb[commandchars=\\\{\}]}
\DefineVerbatimEnvironment{Highlighting}{Verbatim}{commandchars=\\\{\}}
% Add ',fontsize=\small' for more characters per line
\usepackage{framed}
\definecolor{shadecolor}{RGB}{248,248,248}
\newenvironment{Shaded}{\begin{snugshade}}{\end{snugshade}}
\newcommand{\AlertTok}[1]{\textcolor[rgb]{0.94,0.16,0.16}{#1}}
\newcommand{\AnnotationTok}[1]{\textcolor[rgb]{0.56,0.35,0.01}{\textbf{\textit{#1}}}}
\newcommand{\AttributeTok}[1]{\textcolor[rgb]{0.13,0.29,0.53}{#1}}
\newcommand{\BaseNTok}[1]{\textcolor[rgb]{0.00,0.00,0.81}{#1}}
\newcommand{\BuiltInTok}[1]{#1}
\newcommand{\CharTok}[1]{\textcolor[rgb]{0.31,0.60,0.02}{#1}}
\newcommand{\CommentTok}[1]{\textcolor[rgb]{0.56,0.35,0.01}{\textit{#1}}}
\newcommand{\CommentVarTok}[1]{\textcolor[rgb]{0.56,0.35,0.01}{\textbf{\textit{#1}}}}
\newcommand{\ConstantTok}[1]{\textcolor[rgb]{0.56,0.35,0.01}{#1}}
\newcommand{\ControlFlowTok}[1]{\textcolor[rgb]{0.13,0.29,0.53}{\textbf{#1}}}
\newcommand{\DataTypeTok}[1]{\textcolor[rgb]{0.13,0.29,0.53}{#1}}
\newcommand{\DecValTok}[1]{\textcolor[rgb]{0.00,0.00,0.81}{#1}}
\newcommand{\DocumentationTok}[1]{\textcolor[rgb]{0.56,0.35,0.01}{\textbf{\textit{#1}}}}
\newcommand{\ErrorTok}[1]{\textcolor[rgb]{0.64,0.00,0.00}{\textbf{#1}}}
\newcommand{\ExtensionTok}[1]{#1}
\newcommand{\FloatTok}[1]{\textcolor[rgb]{0.00,0.00,0.81}{#1}}
\newcommand{\FunctionTok}[1]{\textcolor[rgb]{0.13,0.29,0.53}{\textbf{#1}}}
\newcommand{\ImportTok}[1]{#1}
\newcommand{\InformationTok}[1]{\textcolor[rgb]{0.56,0.35,0.01}{\textbf{\textit{#1}}}}
\newcommand{\KeywordTok}[1]{\textcolor[rgb]{0.13,0.29,0.53}{\textbf{#1}}}
\newcommand{\NormalTok}[1]{#1}
\newcommand{\OperatorTok}[1]{\textcolor[rgb]{0.81,0.36,0.00}{\textbf{#1}}}
\newcommand{\OtherTok}[1]{\textcolor[rgb]{0.56,0.35,0.01}{#1}}
\newcommand{\PreprocessorTok}[1]{\textcolor[rgb]{0.56,0.35,0.01}{\textit{#1}}}
\newcommand{\RegionMarkerTok}[1]{#1}
\newcommand{\SpecialCharTok}[1]{\textcolor[rgb]{0.81,0.36,0.00}{\textbf{#1}}}
\newcommand{\SpecialStringTok}[1]{\textcolor[rgb]{0.31,0.60,0.02}{#1}}
\newcommand{\StringTok}[1]{\textcolor[rgb]{0.31,0.60,0.02}{#1}}
\newcommand{\VariableTok}[1]{\textcolor[rgb]{0.00,0.00,0.00}{#1}}
\newcommand{\VerbatimStringTok}[1]{\textcolor[rgb]{0.31,0.60,0.02}{#1}}
\newcommand{\WarningTok}[1]{\textcolor[rgb]{0.56,0.35,0.01}{\textbf{\textit{#1}}}}
\usepackage{longtable,booktabs,array}
\newcounter{none} % for unnumbered tables
\usepackage{calc} % for calculating minipage widths
% Correct order of tables after \paragraph or \subparagraph
\usepackage{etoolbox}
\makeatletter
\patchcmd\longtable{\par}{\if@noskipsec\mbox{}\fi\par}{}{}
\makeatother
% Allow footnotes in longtable head/foot
\IfFileExists{footnotehyper.sty}{\usepackage{footnotehyper}}{\usepackage{footnote}}
\makesavenoteenv{longtable}
\usepackage{graphicx}
\makeatletter
\newsavebox\pandoc@box
\newcommand*\pandocbounded[1]{% scales image to fit in text height/width
  \sbox\pandoc@box{#1}%
  \Gscale@div\@tempa{\textheight}{\dimexpr\ht\pandoc@box+\dp\pandoc@box\relax}%
  \Gscale@div\@tempb{\linewidth}{\wd\pandoc@box}%
  \ifdim\@tempb\p@<\@tempa\p@\let\@tempa\@tempb\fi% select the smaller of both
  \ifdim\@tempa\p@<\p@\scalebox{\@tempa}{\usebox\pandoc@box}%
  \else\usebox{\pandoc@box}%
  \fi%
}
% Set default figure placement to htbp
\def\fps@figure{htbp}
\makeatother
\ifLuaTeX
\usepackage[bidi=basic,shorthands=off]{babel}
\else
\usepackage[bidi=default,shorthands=off]{babel}
\fi
\ifLuaTeX
  \usepackage{selnolig} % disable illegal ligatures
\fi
\setlength{\emergencystretch}{3em} % prevent overfull lines
\providecommand{\tightlist}{%
  \setlength{\itemsep}{0pt}\setlength{\parskip}{0pt}}
\usepackage[]{natbib}
\bibliographystyle{apalike}
\usepackage{microtype}
\usepackage{booktabs}
\usepackage{longtable}
\usepackage{float}
\usepackage{xcolor}
\definecolor{linkblue}{HTML}{1565C0}
% hyperref is loaded by pandoc's default LaTeX template AFTER this
% header-includes block, so \hypersetup here would error with
% "Undefined control sequence". Pass options forward so they apply
% when hyperref actually loads.
\PassOptionsToPackage{colorlinks=true,linkcolor=linkblue,citecolor=linkblue,urlcolor=linkblue}{hyperref}
% This volume is theorem-heavy (Maekawa, Kawasaki, isometry arguments).
\usepackage{amsmath}
\usepackage{amssymb}

\setcounter{secnumdepth}{4}
\setcounter{tocdepth}{3}

% LaTeX fallbacks for the Harrer-style HTML callout boxes
% ({block, type='boxempty'} etc. — see style/style.css). PDF degrades to
% plain quoted blocks; HTML keeps the styled boxes.
\newenvironment{boxempty}{\begin{quote}}{\end{quote}}
\newenvironment{boxinfo}{\begin{quote}}{\end{quote}}
\newenvironment{boximportant}{\begin{quote}}{\end{quote}}
\newenvironment{boxreport}{\begin{quote}}{\end{quote}}
\newenvironment{boxquestion}{\begin{quote}}{\end{quote}}

\usepackage{fancyhdr}
\setlength{\headheight}{25pt}
\pagestyle{fancy}
\fancyhead{}
\fancyhead[C]{\small\nouppercase{\leftmark}}
\fancyfoot{}
\fancyfoot[C]{\scriptsize R. Heller · Date: 2025-01-01 · Last changed: 2026-07-31\hfill\thepage}
\fancypagestyle{plain}{\fancyhf{}\renewcommand{\headrulewidth}{0pt}\fancyfoot[C]{\scriptsize R. Heller · Date: 2025-01-01 · Last changed: 2026-07-31\hfill\thepage}}


% Wrap long code and URLs within the printable text area.
\usepackage{fvextra}
\fvset{breaklines=true,breakanywhere=true}
\usepackage{xurl}
\usepackage{bookmark}
\IfFileExists{xurl.sty}{\usepackage{xurl}}{} % add URL line breaks if available
\urlstyle{same}
\hypersetup{
  pdftitle={Folding in R},
  pdfauthor={R. Heller},
  pdflang={en-GB},
  pdfkeywords={dimension reduction, manifold learning, origami, benchmarking, R},
  hidelinks,
  pdfcreator={LaTeX via pandoc}}

\title{Folding in R}
\usepackage{etoolbox}
\makeatletter
\providecommand{\subtitle}[1]{% add subtitle to \maketitle
  \apptocmd{\@title}{\par {\large #1 \par}}{}{}
}
\makeatother
\subtitle{Origami, Manifolds, and Dimension Reduction}
\author{R. Heller}
\date{2026-11-09}

\begin{document}
\maketitle

{
\setcounter{tocdepth}{2}
\tableofcontents
}
\listoffigures
\chapter*{Preface}\label{preface}
\addcontentsline{toc}{chapter}{Preface}

\emph{R. Heller · Date: 2026-07-29 · Last changed: 2026-11-09}

A rigid origami crease pattern generates a manifold whose true low-dimensional
structure is known exactly, in closed form. Rigid folding bends the sheet along
creases only; the facets stay flat. The folded surface is therefore isometric to
the flat sheet, and geodesic distances measured along the surface are unchanged
by folding. The flat crease pattern is not an approximation of the correct
two-dimensional embedding. It \emph{is} the correct embedding --- the answer key.

That single property is what this book is built on.

Standard benchmarks for dimension reduction --- the Swiss roll, the S-curve, the
severed sphere --- give one fixed geometry and a visual impression of success.
A crease pattern gives an exact reconstruction error and a parameterised family
of test problems. The fold angle \(\theta\) runs continuously from a flat sheet to
a fully compressed one, so every method's performance becomes a curve rather
than a single number, and no configuration can be cherry-picked.

Chari and Pachter have argued that two-dimensional embeddings of
high-dimensional biological data are systematically misleading. The usual
rebuttal is anecdotal. On a manifold where the truth is known exactly, the
distortion can be \emph{measured} instead of argued about --- and, more usefully, the
evaluation metrics themselves can be checked against truth. That is the reason
this book exists.

\section*{What this book is not}\label{what-this-book-is-not}
\addcontentsline{toc}{section}{What this book is not}

It is not a treatment of rigidity theory, computational origami design, or deep
generative modelling. It is not a tutorial on \texttt{bookdown}, \texttt{make}, or project
tooling. Where a section could not be tied back to the ground-truth claim above,
it was cut rather than kept for completeness.

\section*{Companion package}\label{companion-package}
\addcontentsline{toc}{section}{Companion package}

The companion R package \href{https://github.com/r-heller/foldbenchR}{foldbenchR} is publicly available. This book renders with the helper scripts versioned alongside its chapters, so the package is not required to build the book or inspect its recorded benchmarks.

\chapter*{Impressum}\label{impressum}
\addcontentsline{toc}{chapter}{Impressum}

\emph{R. Heller · Date: 2026-07-29 · Last changed: 2026-07-29}

\emph{Folding in R: Origami, Manifolds, and Dimension Reduction}

R. Heller\\
ORCID: \href{https://orcid.org/0000-0001-8006-9742}{0000-0001-8006-9742}

Source repository: \url{https://github.com/r-heller/folding-in-r}\\
Online edition: \url{https://r-heller.github.io/folding-in-r/}

Prose is licensed under \href{https://creativecommons.org/licenses/by/4.0/}{CC BY 4.0}.
Code is licensed under the MIT License; see \texttt{LICENSE-CODE.md} in the repository.

This is a living document. The version you are reading was rendered on
2026-11-09.

\chapter*{Acknowledgments}\label{acknowledgments}
\addcontentsline{toc}{chapter}{Acknowledgments}

\emph{R. Heller · Date: 2026-07-29 · Last changed: 2026-07-29}

This book stands on two bodies of work that developed independently and rarely
cite each other: the mathematics of rigid origami, and the practice of
non-linear dimension reduction. Bringing them together is the only original
move here; both halves were built by others.

The R package ecosystem this book depends on --- \texttt{bookdown}, \texttt{Rtsne}, \texttt{umap},
\texttt{igraph}, \texttt{vegan}, \texttt{torch} and their dependencies --- is maintained largely by
volunteers. The bibliography in \texttt{packages.bib} is generated automatically so
that those authors are cited by name rather than thanked in the abstract.

\chapter*{How to use this book}\label{how-to-use-this-book}
\addcontentsline{toc}{chapter}{How to use this book}

\emph{R. Heller · Date: 2026-07-29 · Last changed: 2026-08-31}

\section*{Standing methodological rules}\label{standing-methodological-rules}
\addcontentsline{toc}{section}{Standing methodological rules}

Three rules apply throughout and are not restated in each chapter.

\begin{enumerate}
\def\labelenumi{\arabic{enumi}.}
\item
  \textbf{Every stochastic result is reported across at least 20 seeds, and across
  more where the spread requires it.} t-SNE and UMAP are seed-dependent; a
  single run is an anecdote. Twenty is a floor, not a design: the count each
  method needs is measured, by inverting the 0.02 reportable difference against
  that method's within-cell spread, and a method needing more gets more.
  \texttt{scripts/measure-seed-budget.R} does the measuring and
  \texttt{data/processed/seed-budget.rds} is its record. Seeds are fixed in
  \texttt{\_common.R} as \texttt{BENCH\_SEEDS} and disclosed with every figure.

  Where an artefact is produced below the floor --- a pilot, or an arm whose cost
  forbids it --- the count is named in the text that reports it, and
  \texttt{scripts/check-artefact-producers.R} will not let it pass unnamed.
\item
  \textbf{Every \(\theta\)-dependent result is reported as a curve, never as a single
  number.} Fold angle is a continuous difficulty parameter. Reporting one
  value invites cherry-picking, and hides exactly the regime changes the
  benchmark exists to expose.
\item
  \textbf{Every figure carries \texttt{fig.alt}.} No result is encoded by colour alone.
\end{enumerate}

\section*{What you need}\label{what-you-need}
\addcontentsline{toc}{section}{What you need}

R 4.1 or newer and the packages pinned in \texttt{renv.lock}. Run \texttt{renv::restore()}
once after cloning. There is nothing else to install: the pattern
constructors, the sampler and the evaluation metrics are plain scripts in \texttt{R/},
sourced by \texttt{\_common.R} when the book renders. The code that produces every
result in this book is in the same repository as the book, at the same commit.

The Chapter \ref{results} benchmark grid takes hours to compute. It is
precomputed and committed under \texttt{data/processed/}; \texttt{scripts/run-benchmark-grid.R}
is kept as the provenance record rather than run during the build.

\section*{How the chapters are shaped}\label{how-the-chapters-are-shaped}
\addcontentsline{toc}{section}{How the chapters are shaped}

Every body chapter uses the same nine sections, in the same order, with the same
anchors. The anchors are namespaced by a short chapter mnemonic --- \texttt{geom-},
\texttt{intro-}, \texttt{eval-} --- because pandoc slugs headings by their text, and twelve
chapters each carrying a heading called ``Results'' would otherwise collide and
silently renumber each other's cross-references.

\begin{verbatim}
## What this chapter answers  {#<mn>-question}
## Setup                      {#<mn>-setup}
## Background                 {#<mn>-background}
## <free-wording core>        {#<mn>-core}
## Results                    {#<mn>-results}
## Diagnostics                {#<mn>-diagnostics}
## Where it fails             {#<mn>-limits}
## Reproduce this             {#<mn>-reproduce}
## Further reading            {#<mn>-reading}
\end{verbatim}

Only the core section is freely worded --- it is the one that carries the
chapter's actual argument, and calling it ``Core'' in every chapter would be
worse than useless. Everything else is fixed, and \texttt{scripts/lint-chapters.R}
fails the build on a chapter that is missing a slot, has them out of order, or
anchors one to the wrong namespace.

Read it as a contract with you rather than a template for us. \emph{Where it fails}
is in every chapter because every method in this book fails somewhere, and a
chapter that could not fill that section would not be finished. \emph{Reproduce this}
names the artefact and the script that produced it, so any number in the chapter
can be traced to the run that generated it.

\section*{Callouts}\label{callouts}
\addcontentsline{toc}{section}{Callouts}

Five callout classes are defined in \texttt{style/style.css}, each with a LaTeX
fallback so they survive into the PDF.

\begin{itemize}
\tightlist
\item
  \texttt{boxinfo} --- context a reader can skip without losing the thread.
\item
  \texttt{boximportant} --- a result that is easy to misread, or a default that will bite.
\item
  \texttt{boxquestion} --- something to try, with the answer available in the sources.
\item
  \texttt{boxreport} --- an aside on provenance, attribution or the state of the
  literature.
\item
  \texttt{boxempty} --- an unadorned block, no icon and no colour. It carries the
  suggested citation in \hyperref[citing]{Citing this book}.
\end{itemize}

Write them as knitr block chunks:

\begin{verbatim}
\begin{boximportant}
umap::umap() defaults to preserve.seed = TRUE.
\end{boximportant}
\end{verbatim}

Never as a raw \texttt{\textless{}div\ class="boxinfo"\textgreater{}} and never as a pandoc fenced div
(\texttt{:::\ boxinfo}). Those two forms are not what the stylesheet targets: they
render as unstyled markup in HTML and disappear entirely in PDF, where the
boxes are LaTeX environments declared in \texttt{style/preamble.tex}. That is how a
sibling volume shipped eighty unstyled boxes.

\section*{How to read it}\label{how-to-read-it}
\addcontentsline{toc}{section}{How to read it}

Part I builds the object. Part II tests four families of method against it.
Part III is the benchmark proper and contains the book's main contribution ---
Chapter \ref{evaluation}, where the evaluation metrics themselves are checked
against exact truth. Part IV asks whether any of it changes a real decision.

Readers who only want the benchmark can start at Chapter \ref{benchmarks} and
refer back to Chapter \ref{generative} for the generative model.

\part{Folding}\label{part-folding}

\chapter{Introduction}\label{intro}

\emph{R. Heller · Date: 2026-07-29 · Last changed: 2026-09-07}

\emph{Why should a crease pattern tell you anything about a gene expression matrix?}

\section{What this chapter answers}\label{intro-question}

It should not, and this book never argues that it does. A gene expression matrix
is not a folded sheet of paper, and no chapter here claims a resemblance between
them.

What a crease pattern has is a property no real dataset has. Its correct
two-dimensional structure is known exactly, in closed form, before a single
point is sampled. Rigid folding bends a flat sheet along its creases while the
facets stay flat and unstretched, so the folded surface is isometric to the
sheet it came from. The flat pattern is not an approximation to the folded
surface's intrinsic geometry. It \textbf{is} that geometry.

One thing follows, and the book is built on it: an evaluator can be handed the
right answer and asked whether it recognises it. A metric that ranks the correct
embedding below a plainly wrong one where the correct embedding is \emph{known} has
not earned trust where it is not. That is a necessary condition rather than a
resemblance, and it is the only transfer to real data this book asserts.

\begin{figure}

{\centering \includegraphics[width=0.9\linewidth,alt={Two panels side by side. Panel A is a flat Miura-ori crease pattern of sixteen parallelogram facets, each in a different colour, with two facets greyed out. Panel B shows the same sheet folded into a corrugated three-dimensional surface, with the facets in matching colours. The greyed facets in A are the ones hidden behind folds in B.}]{01-introduction_files/figure-latex/intro-fold-figure-1} 

}

\caption{A flat crease pattern and the same sheet folded. Turning the folded sheet changes which facets you can see; it does not change the pattern.}\label{fig:intro-fold-figure}
\end{figure}

The pattern in panel A is the same however panel B is turned. What changes is
how much of it any one view exposes --- and a dimension-reduction method sees
something closer to B than to A. The gap between the two is what the rest of
this book measures.

\section{The argument this book joins}\label{intro-background}

Chari and Pachter argued that the two-dimensional embeddings on which
single-cell biology has come to depend distort their data badly enough that
conclusions drawn from the pictures are unsafe \citep{chari2023specious}. The reply
from Lause, Berens and Kobak was not that the distortion is small. It was that
the case against the embeddings rested on metrics that ``focused on preservation
of distances, where 2D PCA was unsurprisingly the best'' \citep{lause2024elephant} ---
that the disagreement had moved from the embeddings to the instruments used to
judge them.

That framing is not this book's, and saying so belongs here rather than in a
footnote. It is the published position of one of the two parties to the dispute.
What this book adds is not the position but a measurement of it: how far each
evaluator sits from an evaluator that has the answer key.

Machado, Behrisch and Telea supply the other half of the reason to look
\citep{machado2025necessary}. They construct projections that score well on
trustworthiness and continuity while displaying structure the data does not
contain. That is the false-\emph{positive} direction, and it needs no ground truth:
to show that a bad embedding scores high, it is enough to know the embedding is
bad. The complement --- that the \emph{right} embedding scores low --- cannot be shown
without an answer key. It is rarely shown for that reason, and it is what
Chapter \ref{evaluation} does.

\begin{Shaded}
\begin{Highlighting}[]
\NormalTok{audit }\OtherTok{\textless{}{-}} \FunctionTok{read\_run}\NormalTok{(}\StringTok{"data/processed/evaluator{-}audit.rds"}\NormalTok{, }\StringTok{"d7982cf7eb3afe7d"}\NormalTok{)}
\NormalTok{prov  }\OtherTok{\textless{}{-}} \FunctionTok{attr}\NormalTok{(audit, }\StringTok{"provenance"}\NormalTok{)}

\CommentTok{\# knitr rounds an inline numeric to options(digits), which turns 4e{-}16 into 0.}
\CommentTok{\# These return character, so a small number survives the trip into prose.}
\CommentTok{\# trimws() because formatC() pads "fg" output to a common width, and a padded}
\CommentTok{\# number lands in the prose with its leading spaces intact.}
\NormalTok{fmt }\OtherTok{\textless{}{-}} \ControlFlowTok{function}\NormalTok{(x, }\AttributeTok{d =} \DecValTok{3}\DataTypeTok{L}\NormalTok{) }\FunctionTok{trimws}\NormalTok{(}\FunctionTok{formatC}\NormalTok{(}\FunctionTok{signif}\NormalTok{(x, d), }\AttributeTok{format =} \StringTok{"fg"}\NormalTok{, }\AttributeTok{digits =}\NormalTok{ d))}
\NormalTok{sci }\OtherTok{\textless{}{-}} \ControlFlowTok{function}\NormalTok{(x, }\AttributeTok{d =} \DecValTok{2}\DataTypeTok{L}\NormalTok{) }\FunctionTok{formatC}\NormalTok{(x, }\AttributeTok{format =} \StringTok{"e"}\NormalTok{, }\AttributeTok{digits =}\NormalTok{ d }\SpecialCharTok{{-}} \DecValTok{1}\DataTypeTok{L}\NormalTok{)}

\NormalTok{amb }\OtherTok{\textless{}{-}}\NormalTok{ audit}\SpecialCharTok{$}\NormalTok{inversions[audit}\SpecialCharTok{$}\NormalTok{inversions}\SpecialCharTok{$}\NormalTok{reference }\SpecialCharTok{==} \StringTok{"ambient"} \SpecialCharTok{\&}
\NormalTok{                          audit}\SpecialCharTok{$}\NormalTok{inversions}\SpecialCharTok{$}\NormalTok{n }\SpecialCharTok{==} \FunctionTok{max}\NormalTok{(prov}\SpecialCharTok{$}\NormalTok{n), ]}
\NormalTok{inverts    }\OtherTok{\textless{}{-}} \FunctionTok{tapply}\NormalTok{(amb}\SpecialCharTok{$}\NormalTok{inverted, amb}\SpecialCharTok{$}\NormalTok{metric, mean)}
\NormalTok{trust\_span }\OtherTok{\textless{}{-}} \FunctionTok{range}\NormalTok{(}\FunctionTok{c}\NormalTok{(amb}\SpecialCharTok{$}\NormalTok{truth[amb}\SpecialCharTok{$}\NormalTok{metric }\SpecialCharTok{==} \StringTok{"trust"}\NormalTok{],}
\NormalTok{                      amb}\SpecialCharTok{$}\NormalTok{pca2[amb}\SpecialCharTok{$}\NormalTok{metric }\SpecialCharTok{==} \StringTok{"trust"}\NormalTok{]))}

\NormalTok{perfect }\OtherTok{\textless{}{-}}\NormalTok{ audit}\SpecialCharTok{$}\NormalTok{scores[audit}\SpecialCharTok{$}\NormalTok{scores}\SpecialCharTok{$}\NormalTok{reference }\SpecialCharTok{==} \StringTok{"ambient"} \SpecialCharTok{\&}
\NormalTok{                          audit}\SpecialCharTok{$}\NormalTok{scores}\SpecialCharTok{$}\NormalTok{candidate }\SpecialCharTok{==} \StringTok{"truth"} \SpecialCharTok{\&}
\NormalTok{                          audit}\SpecialCharTok{$}\NormalTok{scores}\SpecialCharTok{$}\NormalTok{k }\SpecialCharTok{==}\NormalTok{ K\_DEFAULT }\SpecialCharTok{\&}
\NormalTok{                          audit}\SpecialCharTok{$}\NormalTok{scores}\SpecialCharTok{$}\NormalTok{n }\SpecialCharTok{==} \FunctionTok{max}\NormalTok{(prov}\SpecialCharTok{$}\NormalTok{n), ]}
\NormalTok{ceiling }\OtherTok{\textless{}{-}} \FunctionTok{aggregate}\NormalTok{(}\FunctionTok{cbind}\NormalTok{(knn, qnx) }\SpecialCharTok{\textasciitilde{}}\NormalTok{ theta, perfect, mean)}
\end{Highlighting}
\end{Shaded}

The design is one sentence. Hand each evaluator the exact chart as though it
were a submission, alongside a two-component PCA of the same folded cloud, and
record which of the two it prefers. Across the audit's ambient-referenced cells
at the larger sample size, continuity ranks the wrong answer above the right one
in 41.8 per cent of them. Trustworthiness almost
never inverts, and only because it barely moves: over the whole sweep it rates
every candidate in the audit, the exact chart and the wrong one alike, at
0.984 or above. At the last fold angle, kNN preservation rates
the \emph{exact} unfolding 0.58 and \(Q_{NX}\) rates it
0.714. A reader shown the first of those would call the
embedding a partial failure. It is the answer.

\textbf{The law that chapter set out to state does not exist, and this introduction
will not promise it.} The plan of record specified a threshold in the fold
angle: rank metrics are invariant to a monotone contraction of distances until
neighbour identity changes, so they should invert only past some critical
\(\theta\). They do invert, and the mechanism is the one predicted. The threshold
is not. It slides with the neighbourhood size across most of the available range
(Figure \ref{fig:eval-fig-threshold}), so there is no threshold in \(\theta\)
alone. What survives is a threshold surface in \(\theta\) and the neighbourhood
\emph{fraction} \(k/n\), and it survives for continuity only. That is a smaller finding
than the one budgeted for it, and Chapter \ref{evaluation} reports it as the
smaller thing it is.

None of this arbitrates the dispute above. Espadoto and colleagues survey what
dimension-reduction evaluations actually report \citep{espadoto2021survey}; the claim
in this book is about those reports, and it requires no position on what a
two-dimensional embedding is for.

\section{The answer key}\label{intro-core}

The construction is one pattern family and it is not new. Miura's report
describes the tessellation this book folds as ``the isometric transfer of an
infinite plane subject to biaxial shortening'' \citep{miura1985method}, which is the
whole claim in the words of the person who first wrote it down. Chapter
\ref{folding-geometry} derives the folding rather than citing it, because the
derivation is what makes the answer key exact instead of approximate.

\begin{Shaded}
\begin{Highlighting}[]
\NormalTok{p }\OtherTok{\textless{}{-}} \FunctionTok{miura\_ori}\NormalTok{(}\AttributeTok{nx =}\NormalTok{ PATTERN\_GRID}\SpecialCharTok{$}\NormalTok{miura[[}\StringTok{"nx"}\NormalTok{]], }\AttributeTok{ny =}\NormalTok{ PATTERN\_GRID}\SpecialCharTok{$}\NormalTok{miura[[}\StringTok{"ny"}\NormalTok{]])}
\NormalTok{m }\OtherTok{\textless{}{-}} \FunctionTok{sample\_manifold}\NormalTok{(p, }\AttributeTok{theta =}\NormalTok{ FIG\_INTRO\_THETA, }\AttributeTok{n =}\NormalTok{ N\_DEFAULT,}
                     \AttributeTok{seed =}\NormalTok{ BENCH\_SEEDS[}\DecValTok{1}\NormalTok{])}

\NormalTok{iso  }\OtherTok{\textless{}{-}} \FunctionTok{facet\_isometry\_error}\NormalTok{(p, FIG\_INTRO\_THETA)}
\NormalTok{pca2 }\OtherTok{\textless{}{-}}\NormalTok{ stats}\SpecialCharTok{::}\FunctionTok{prcomp}\NormalTok{(m}\SpecialCharTok{$}\NormalTok{X, }\AttributeTok{rank. =}\NormalTok{ EMBED\_DIM)}\SpecialCharTok{$}\NormalTok{x[, }\FunctionTok{seq\_len}\NormalTok{(EMBED\_DIM)]}
\NormalTok{err  }\OtherTok{\textless{}{-}} \FunctionTok{c}\NormalTok{(}\AttributeTok{chart =} \FunctionTok{reconstruction\_error}\NormalTok{(m}\SpecialCharTok{$}\NormalTok{truth, m}\SpecialCharTok{$}\NormalTok{truth),}
          \AttributeTok{pca   =} \FunctionTok{reconstruction\_error}\NormalTok{(pca2, m}\SpecialCharTok{$}\NormalTok{truth))}
\NormalTok{qn   }\OtherTok{\textless{}{-}} \FunctionTok{c}\NormalTok{(}\AttributeTok{chart =} \FunctionTok{qnx}\NormalTok{(m}\SpecialCharTok{$}\NormalTok{truth, }\FunctionTok{reference\_dist}\NormalTok{(m, }\StringTok{"chart"}\NormalTok{), }\AttributeTok{K =}\NormalTok{ K\_QNX),}
          \AttributeTok{pca   =} \FunctionTok{qnx}\NormalTok{(pca2,    }\FunctionTok{reference\_dist}\NormalTok{(m, }\StringTok{"chart"}\NormalTok{), }\AttributeTok{K =}\NormalTok{ K\_QNX))}
\end{Highlighting}
\end{Shaded}

Two numbers fix what ``answer key'' means here. Folding preserves every pairwise
distance within every facet to 4.4e-16 --- floating-point noise on a closed
form, not a tolerance anyone chose. And handing the chart back as though it were
a submission scores 5.7e-16 against itself under the book's
headline error, where a two-component PCA of the same folded cloud scores
0.0775.

\textbf{\(\theta\) is a parameter, not an angle.} It runs on the unit interval: zero is
the flat sheet, one is flat-folded, and the dihedral angle at a crease is
derived from it rather than being another name for it. Because it is continuous
and the geometry moves monotonically along it, every result in this book is a
curve rather than a number --- standing rule 2, and the reason a benchmark with
one fixed geometry cannot produce the observations Part III reports.

\textbf{Two headline numbers per cell, and the pair is deliberate.} The first is a
normalised Procrustes RMSE: the isometry-recovery score, admitting translation,
rotation, reflection and one isotropic scale, and charging for everything else.
The second is \(Q_{NX}\) computed with the exact chart as its reference geometry,
which is invariant under any diffeomorphism of the plane and therefore neutral
between preserving distances and preserving neighbourhoods. On the sample above
the two disagree about PCA: an error of 0.0775 against a
\(Q_{NX}\) of 0.899, where the exact chart scores
1 on the second and 5.7e-16 on the first.
The gap between them, as a function of \(\theta\), is a measurement of how much
observed error is genuine structural loss and how much is harmless
reparameterisation. Reporting only the first would build the book's scoring
function out of one side of the dispute it claims to measure.

\textbf{Folding preserves one kind of distance and destroys the other.} Distances
along the surface are unchanged, exactly; distances through the surrounding
space contract by a factor that depends only on the dihedral angle. Every method
in Part II consumes the second kind and is scored against the first. That
mismatch is not a flaw in the benchmark, it is the benchmark's subject --- and it
comes with a cost worth knowing before Part III: a crease has reach exactly
zero, so the sampling guarantees usually invoked to justify these methods do not
apply to this book's manifolds at all --- not weakly, not approximately. Chapter
\ref{folding-geometry} says what replaces them.

\section{What is claimed, and what is not}\label{intro-limits}

\textbf{Most of the methods under test never claimed what the headline metric
measures.} Isomap consumes geodesic distances, and its assumption is satisfied
exactly on a crease pattern \citep{tenenbaum2000isomap}; it is the one method in Part
II trying to do the thing a Procrustes score scores. t-SNE optimises a
divergence between neighbour probability distributions \citep{vandermaaten2008tsne},
UMAP a cross-entropy on a fuzzy simplicial set \citep{mcinnes2018umap}, and neither
claims to reproduce distances at all --- a t-SNE embedding's very extent is a
function of its perplexity and its iteration count. Local linear embedding, the
Laplacian eigenmap and an autoencoder are in the same position. Scoring them
against isometric truth and reporting that they lose is not a finding about
them; it is a category error dressed as a measurement. The second headline
number exists so that a method preserving neighbourhoods without preserving
distances can be seen doing it, and each Part II chapter restates the point for
the method it covers.

\begin{Shaded}
\begin{Highlighting}[]
\NormalTok{e1  }\OtherTok{\textless{}{-}} \FunctionTok{read\_run}\NormalTok{(}\StringTok{"data/processed/e1{-}controlled.rds"}\NormalTok{, }\StringTok{"01d82cfe8f45bb44"}\NormalTok{)}
\NormalTok{ctl }\OtherTok{\textless{}{-}} \FunctionTok{local}\NormalTok{(\{                       }\CommentTok{\# both difficulty axes inside the joint range}
\NormalTok{  s }\OtherTok{\textless{}{-}}\NormalTok{ e1[}\FunctionTok{is.finite}\NormalTok{(e1}\SpecialCharTok{$}\NormalTok{rmse), ]; s}\SpecialCharTok{$}\NormalTok{lg }\OtherTok{\textless{}{-}} \FunctionTok{log}\NormalTok{(s}\SpecialCharTok{$}\NormalTok{gs)}
\NormalTok{  s[s}\SpecialCharTok{$}\NormalTok{lg }\SpecialCharTok{\textgreater{}=} \FunctionTok{max}\NormalTok{(}\FunctionTok{tapply}\NormalTok{(s}\SpecialCharTok{$}\NormalTok{lg, s}\SpecialCharTok{$}\NormalTok{kind, min)) }\SpecialCharTok{\&}\NormalTok{ s}\SpecialCharTok{$}\NormalTok{lg }\SpecialCharTok{\textless{}=} \FunctionTok{min}\NormalTok{(}\FunctionTok{tapply}\NormalTok{(s}\SpecialCharTok{$}\NormalTok{lg, s}\SpecialCharTok{$}\NormalTok{kind, max)) }\SpecialCharTok{\&}
\NormalTok{    s}\SpecialCharTok{$}\NormalTok{np }\SpecialCharTok{\textgreater{}=} \FunctionTok{max}\NormalTok{(}\FunctionTok{tapply}\NormalTok{(s}\SpecialCharTok{$}\NormalTok{np, s}\SpecialCharTok{$}\NormalTok{kind, min)) }\SpecialCharTok{\&}\NormalTok{ s}\SpecialCharTok{$}\NormalTok{np }\SpecialCharTok{\textless{}=} \FunctionTok{min}\NormalTok{(}\FunctionTok{tapply}\NormalTok{(s}\SpecialCharTok{$}\NormalTok{np, s}\SpecialCharTok{$}\NormalTok{kind, max)), ]}
\NormalTok{\})}
\NormalTok{w }\OtherTok{\textless{}{-}} \FunctionTok{reshape}\NormalTok{(ctl[, }\FunctionTok{c}\NormalTok{(}\StringTok{"kind"}\NormalTok{, }\StringTok{"label"}\NormalTok{, }\StringTok{"seed"}\NormalTok{, }\StringTok{"method"}\NormalTok{, }\StringTok{"rmse"}\NormalTok{)],}
             \AttributeTok{idvar =} \FunctionTok{c}\NormalTok{(}\StringTok{"kind"}\NormalTok{, }\StringTok{"label"}\NormalTok{, }\StringTok{"seed"}\NormalTok{), }\AttributeTok{timevar =} \StringTok{"method"}\NormalTok{,}
             \AttributeTok{direction =} \StringTok{"wide"}\NormalTok{)}
\NormalTok{w}\SpecialCharTok{$}\NormalTok{spread }\OtherTok{\textless{}{-}} \FunctionTok{abs}\NormalTok{(w}\SpecialCharTok{$}\NormalTok{rmse.pca }\SpecialCharTok{{-}}\NormalTok{ w}\SpecialCharTok{$}\NormalTok{rmse.isomap)}
\NormalTok{sp    }\OtherTok{\textless{}{-}} \FunctionTok{tapply}\NormalTok{(w}\SpecialCharTok{$}\NormalTok{spread, w}\SpecialCharTok{$}\NormalTok{kind, mean)}
\NormalTok{ratio }\OtherTok{\textless{}{-}}\NormalTok{ sp[[}\StringTok{"classic"}\NormalTok{]] }\SpecialCharTok{/}\NormalTok{ sp[[}\StringTok{"crease"}\NormalTok{]]}

\FunctionTok{set.seed}\NormalTok{(}\DecValTok{11}\NormalTok{)                         }\CommentTok{\# cluster bootstrap, resampling whole settings}
\NormalTok{ci }\OtherTok{\textless{}{-}} \FunctionTok{quantile}\NormalTok{(}\FunctionTok{replicate}\NormalTok{(}\DecValTok{2000}\NormalTok{, \{}
\NormalTok{  r }\OtherTok{\textless{}{-}} \FunctionTok{vapply}\NormalTok{(}\FunctionTok{split}\NormalTok{(w, w}\SpecialCharTok{$}\NormalTok{kind), }\ControlFlowTok{function}\NormalTok{(d) \{}
\NormalTok{    l }\OtherTok{\textless{}{-}} \FunctionTok{sample}\NormalTok{(}\FunctionTok{unique}\NormalTok{(d}\SpecialCharTok{$}\NormalTok{label), }\FunctionTok{length}\NormalTok{(}\FunctionTok{unique}\NormalTok{(d}\SpecialCharTok{$}\NormalTok{label)), }\AttributeTok{replace =} \ConstantTok{TRUE}\NormalTok{)}
    \FunctionTok{mean}\NormalTok{(}\FunctionTok{unlist}\NormalTok{(}\FunctionTok{lapply}\NormalTok{(l, }\ControlFlowTok{function}\NormalTok{(x) d}\SpecialCharTok{$}\NormalTok{spread[d}\SpecialCharTok{$}\NormalTok{label }\SpecialCharTok{==}\NormalTok{ x])))}
\NormalTok{  \}, }\FunctionTok{numeric}\NormalTok{(}\DecValTok{1}\NormalTok{))}
\NormalTok{  r[[}\StringTok{"classic"}\NormalTok{]] }\SpecialCharTok{/}\NormalTok{ r[[}\StringTok{"crease"}\NormalTok{]]}
\NormalTok{\}), }\FunctionTok{c}\NormalTok{(}\FloatTok{0.025}\NormalTok{, }\FloatTok{0.975}\NormalTok{))}

\NormalTok{fit   }\OtherTok{\textless{}{-}} \FunctionTok{lm}\NormalTok{(spread }\SpecialCharTok{\textasciitilde{}} \FunctionTok{I}\NormalTok{(kind }\SpecialCharTok{==} \StringTok{"classic"}\NormalTok{), }\AttributeTok{data =}\NormalTok{ w)   }\CommentTok{\# cluster{-}robust by setting}
\NormalTok{X     }\OtherTok{\textless{}{-}} \FunctionTok{model.matrix}\NormalTok{(fit); u }\OtherTok{\textless{}{-}} \FunctionTok{residuals}\NormalTok{(fit); bread }\OtherTok{\textless{}{-}} \FunctionTok{solve}\NormalTok{(}\FunctionTok{crossprod}\NormalTok{(X))}
\NormalTok{meat  }\OtherTok{\textless{}{-}} \FunctionTok{Reduce}\NormalTok{(}\StringTok{\textasciigrave{}}\AttributeTok{+}\StringTok{\textasciigrave{}}\NormalTok{, }\FunctionTok{lapply}\NormalTok{(}\FunctionTok{unique}\NormalTok{(w}\SpecialCharTok{$}\NormalTok{label), }\ControlFlowTok{function}\NormalTok{(l) \{}
\NormalTok{  i }\OtherTok{\textless{}{-}}\NormalTok{ w}\SpecialCharTok{$}\NormalTok{label }\SpecialCharTok{==}\NormalTok{ l}
  \FunctionTok{tcrossprod}\NormalTok{(}\FunctionTok{crossprod}\NormalTok{(X[i, , }\AttributeTok{drop =} \ConstantTok{FALSE}\NormalTok{], u[i]))}
\NormalTok{\}))}
\NormalTok{tval  }\OtherTok{\textless{}{-}} \FunctionTok{coef}\NormalTok{(fit)[}\DecValTok{2}\NormalTok{] }\SpecialCharTok{/} \FunctionTok{sqrt}\NormalTok{((bread }\SpecialCharTok{\%*\%}\NormalTok{ meat }\SpecialCharTok{\%*\%}\NormalTok{ bread)[}\DecValTok{2}\NormalTok{, }\DecValTok{2}\NormalTok{])}
\NormalTok{pval  }\OtherTok{\textless{}{-}} \DecValTok{2} \SpecialCharTok{*} \FunctionTok{pt}\NormalTok{(}\SpecialCharTok{{-}}\FunctionTok{abs}\NormalTok{(tval), }\AttributeTok{df =} \FunctionTok{length}\NormalTok{(}\FunctionTok{unique}\NormalTok{(w}\SpecialCharTok{$}\NormalTok{label)) }\SpecialCharTok{{-}} \DecValTok{2}\NormalTok{)}
\end{Highlighting}
\end{Shaded}

\textbf{The Swiss roll discriminates better than this benchmark does.} The claim this
project set out to make was that a crease pattern is a \emph{harder} test than the
classical surfaces. Experiment E1 was run before Chapter \ref{comparison} was
drafted, precisely so that the prose could not outlive a refutation, and it
refuted it. Matched on both axes of difficulty --- branch separation and ambient
non-planarity --- over 9 crease
settings and 5 Swiss rolls,
the mean gap between PCA and Isomap is 0.574 on the classical
side against 0.113 on the crease side: a ratio of
5.06, with a cluster bootstrap interval over settings of
1.7 to 9.8 and a cluster-robust \(p\) of
0.02. Inside the overlap the rankings are identical --- Isomap first
on both families, the linear methods behind it on both --- so nothing flips.
Crease patterns are not a harder benchmark. They are an easier one that tells
methods apart about five times less well, and a reader choosing between the two
for a straight comparison of methods should choose the roll. E1 pre-registered
three outcomes and none of them occurred; the replacement claim is in this book
because it fits the evidence, not because its antecedent was met.

\textbf{The difficulty axis stops short of the failure it was expected to reach.}
Folding corrupts the neighbourhood graph a method consumes, and Chapter
\ref{generative} measures the corruption growing along the whole sweep. What it
does not find, anywhere on the committed grid, is a single graph edge joining
facets two or more cells apart --- the sheet-to-sheet short circuit that
classically breaks Isomap \citep{balasubramanian2002isomap} lives past the end of the
axis, because the sweep stops before the sheet begins collapsing onto itself.
This benchmark measures the graded corruption that precedes that failure. It
does not reach the failure.

\textbf{Exactness alone is not the differentiator.} An arc-length-charted Swiss roll
also has an exact isometric key, and it is published prior art: Schoeneman and
colleagues noticed that the customary answer key is not the isometric one and
replaced the Archimedean spiral with a unit-speed one \citep{schoeneman2017error}.
Claiming exact ground truth as the novelty would be wrong. What a crease pattern
has that theirs does not is reach exactly zero rather than positive, so that it
sits outside the smoothness theory rather than at its edge; intrinsic dimension
above two while the chart stays in closed form, which is what makes an
irreducible-loss bound computable at all; and a difficulty parameter whose
critical value can be derived rather than found by sweeping and looking.

\textbf{One family, and two documented negative results.} Two further pattern
families were specified and both were withdrawn on evidence, which Chapter
\ref{benchmarks} reports as results. A benchmark built on one family cannot
support claims about pattern variety, and none are made.

\section{How the book is organised}\label{intro-roadmap}

\textbf{Part I builds the object.} Chapter \ref{folding-geometry} derives the
folding and proves what it preserves and what it destroys, including a
proposition that an ambient-referenced stress must rank the exact unfolding
below a two-component PCA --- a theorem rather than a simulation. Chapter
\ref{generative} turns the surface into data: sampling, noise models, and the
four measures of difficulty that the \(\theta\) axis moves.

\textbf{Part II tests the methods one family at a time}, each against the same
answer key. Chapter \ref{linear} takes the linear projections, Chapter
\ref{geodesic} the geodesic methods and the short-circuit failure they are
prone to, Chapter \ref{neighbor} the neighbour embeddings, and Chapter
\ref{autoencoders} the learned folds --- where the honest result is that the
autoencoder row is declared unavailable rather than quietly dropped.

\textbf{Part III is the benchmark.} Chapter \ref{benchmarks} builds the manifolds
and the irreducible-loss bound that says how much of an error was unavoidable.
Chapter \ref{evaluation} runs the known-answer test on the evaluators
themselves. Chapter \ref{results} reads the full grid, and asks whether the
benchmark separates methods before asking which method wins. Chapter
\ref{comparison} puts the construction against the Swiss roll, the S-curve and
the severed sphere, and reports that it loses on discrimination.

\textbf{Part IV spends the result.} Chapter \ref{application} uses the grid to
choose a method and its settings for a real dataset, under a rule fixed in code
before the grid existed, and reports what happens.

Two appendices carry the reference material: \ref{notation} fixes every symbol,
and \ref{datasets} documents the committed artefacts. Every artefact is read
with its digest, so a regenerated run is a build failure rather than a silently
changed sentence, and the three standing rules the chapters do not restate are
set out in \emph{How to use this book}.

Not every reader wants the whole path. Someone who came for the geometry can
read Chapters \ref{folding-geometry} and \ref{generative} and stop; someone
who came for the numbers can start at Chapter \ref{results} and work backwards
through the two chapters it depends on; and someone deciding whether to use this
benchmark at all should read Chapter \ref{comparison} first, because that is
where the case against it is made.

\section{Further reading}\label{intro-reading}

Where to start depends on what you came for. For the origami, Hull's
\emph{Origametry} \citep{hull2020origametry} is the single best entry point and covers
both the classical flat-foldability theorems this book uses as checks. For the
methods and for the rank metrics Part III audits, Lee and Verleysen
\citep{lee2007nonlinear} is the textbook treatment of everything Part II fits.

For the dispute, read Chari and Pachter \citep{chari2023specious} and the reply from
Lause, Berens and Kobak \citep{lause2024elephant} in that order; they are short, and
between them they define the question this book answers with measurement rather
than with argument. Espadoto and colleagues \citep{espadoto2021survey} survey what
evaluations report in practice, and Machado, Behrisch and Telea
\citep{machado2025necessary} is the direct complement to Chapter \ref{evaluation}.
For the nearest prior benchmark with an exact answer key, Schoeneman and
colleagues \citep{schoeneman2017error} --- read alongside Chapter \ref{comparison},
which measures this construction against theirs and reports that theirs
separates methods better.

Miura's own report \citep{miura1985method} is worth finding for the sentence quoted
above, and each chapter closes with a \emph{Further reading} section of its own that
goes deeper on its own subject than this one can.

\chapter{The geometry of folding}\label{folding-geometry}

\emph{R. Heller · Date: 2026-07-29 · Last changed: 2026-08-31}

\emph{What does folding preserve, and what does it destroy?}

\section{What this chapter answers}\label{geom-question}

A sheet of paper folded along creases is a surface in space. Flattened out again
it is a region of the plane. The two are the same paper, and that is the whole
of this book's method: the flat region is not a good approximation to the folded
surface's true two-dimensional structure, it \emph{is} that structure, written down
in closed form before any data are generated.

This chapter establishes what that sameness does and does not mean. Folding
preserves every distance measured \textbf{along} the surface and destroys a great
many measured \textbf{through} the space around it, and almost every difficulty a
dimension-reduction method has with a folded sheet is a consequence of that one
sentence. A method sees the second kind of distance and is scored against the
first.

Nothing here is asserted. Every claim in this chapter is a check in
\texttt{tests/testthat/}, and every number is computed while the page is built.

\section{Setup}\label{geom-setup}

\begin{Shaded}
\begin{Highlighting}[]
\NormalTok{p }\OtherTok{\textless{}{-}} \FunctionTok{miura\_ori}\NormalTok{(}\AttributeTok{nx =} \DecValTok{6}\DataTypeTok{L}\NormalTok{, }\AttributeTok{ny =} \DecValTok{6}\DataTypeTok{L}\NormalTok{, }\AttributeTok{alpha =}\NormalTok{ pi }\SpecialCharTok{/} \DecValTok{3}\NormalTok{)}
\NormalTok{f }\OtherTok{\textless{}{-}} \FunctionTok{fold}\NormalTok{(p, }\AttributeTok{theta =} \FloatTok{0.6}\NormalTok{)}
\end{Highlighting}
\end{Shaded}

\texttt{miura\_ori()} returns a flat crease pattern; \texttt{fold()} places it in space at a
folding parameter \(\theta\). Both are in \texttt{R/}, sourced when this book renders,
and the \hyperref[notation]{notation appendix} fixes the symbols.

\section{Background}\label{geom-background}

\subsection{What a crease pattern is, and what it is not}\label{what-a-crease-pattern-is-and-what-it-is-not}

A crease pattern is a planar straight-line graph on a region of paper: vertices,
straight creases between them, and the facets those creases enclose. It carries
one more piece of information per crease, an \textbf{assignment} --- mountain or valley
--- which says which way the paper turns there.

The literature on origami mathematics is largely about which patterns fold and
how to design one that folds to a given shape \citep{lang2011origami, demaine2007folding}. This book needs almost none of that. It needs one family
that provably folds, and it needs the folding written as a map it can evaluate
at any \(\theta\).

\begin{boxinfo}
\textbf{Not the Huzita--Hatori axioms.} A reader who has met origami
mathematics before will expect the construction axioms --- the seven
operations that say which points and lines a folder can produce with a
sequence of folds. They are not what makes a crease pattern a manifold,
and they play no part in this book. They describe what can be
\emph{constructed}; the question here is what a given pattern, however
obtained, \emph{is} once folded. Hull's \emph{Origametry}
{[}@hull2020origametry{]} covers both and is the place to start if the
constructive side is what you came for.
\end{boxinfo}

\subsection{Rigid folding}\label{rigid-folding}

The foldings this book uses are \textbf{rigid}: the facets stay flat and unstretched,
and all the motion happens at the creases. That is a strong restriction and it
is the source of everything that follows. A folding that let facets bend would
still produce a surface, but its flat pattern would no longer be its unfolding,
and the answer key would be gone.

Two classical conditions constrain which flat patterns can fold flat at all. At
every interior vertex, \textbf{Maekawa's theorem} says the numbers of mountain and
valley creases differ by exactly two, and \textbf{Kawasaki's theorem} says the
alternating sum of the sector angles around the vertex is zero
\citep{hull2020origametry}. Both are necessary and neither is sufficient. Chapter
\ref{benchmarks} returns to how little they guarantee; this chapter uses them
as checks rather than as foundations, which is the only honest way to use a
necessary condition.

The attribution here is tangled and it is better to say so than to pick a side.
Kawasaki's condition appears independently in work by Justin and by Murata, and
no citation resolves cleanly to a single first statement. \emph{Origametry} is the
honest reference until one does.

\subsection{Two kinds of isometry}\label{two-kinds-of-isometry}

The word \emph{isometry} does two jobs and the difference between them is this book's
crux.

An \textbf{intrinsic} isometry preserves distances measured along the surface --- the
length of the shortest path an ant could walk without leaving the paper. Rigid
folding is an intrinsic isometry, exactly and by construction: the facets do not
stretch, so a path drawn on the flat sheet has the same length after folding.

An \textbf{ambient} isometry preserves distances measured through the surrounding
space --- straight lines through the air. Folding is emphatically \emph{not} an ambient
isometry. Bring two points on either side of a crease together by folding, and
the straight-line distance between them collapses while the path along the paper
does not change at all.

Every method in Part II consumes ambient distance or something built from it.
The answer key is intrinsic. That mismatch is not a flaw in the benchmark; it is
the benchmark's subject.

\subsection{Curvature, and where it is not}\label{curvature-and-where-it-is-not}

The natural wrong intuition about a folded sheet is that its curvature has been
pushed into the vertices. It is worth writing against rather than merely
avoiding, because it is wrong in a way that would undermine the whole
construction if it were right.

Discrete Gaussian curvature at a vertex of a polyhedral surface is its angle
defect: \(2\pi\) minus the sum of the sector angles meeting there. For a folded
flat sheet those angles are the flat pattern's angles, unchanged by folding, and
they sum to exactly \(2\pi\) because the vertex was an interior point of a flat
piece of paper. \textbf{The Gaussian curvature is identically zero, at every vertex
and on every facet.} Were it otherwise the sheet would not be developable and
the flat pattern would not be its unfolding.

What does concentrate on the creases is \emph{extrinsic}: the second fundamental
form, which is a measure supported on the one-skeleton rather than a function.
The surface is intrinsically flat and extrinsically as bent as one likes.

For manifold learning the consequence that matters is a different quantity
entirely. The \textbf{reach} of a submanifold \citep{federer1959curvature} is the largest
distance within which every point of the ambient space has a unique nearest
point on the surface. Sampling theory for manifold learning is generally stated
in terms of it \citep{niyogi2008homology}, and a crease has reach exactly zero: a
point just outside the fold is equidistant from both facets. So the standard
sampling guarantees do not apply here at all --- not weakly, not approximately.
This book's manifolds sit outside the hypotheses of the theory usually invoked
to justify the methods being tested on them, which is worth knowing before Part
III reports that some of those methods struggle.

\section{The isometry, in five steps}\label{geom-core}

\subsection{One: the folding is written down, not solved for}\label{one-the-folding-is-written-down-not-solved-for}

The Miura-ori is a tessellation of congruent parallelograms \citep{miura1985method}.
Its flat vertex \((i, j)\) sits at

\[x = i\,a + (j \bmod 2)\, b \cos\alpha, \qquad y = j\, b \sin\alpha,\]

where \(a\) and \(b\) are the cell side lengths and \(\alpha\) its acute angle. The
\((j \bmod 2)\) term is the entire pattern: it makes the transverse crease lines
zigzag with period two instead of running straight, which is what turns a plain
parallelogram lattice into something that folds.

The folded placement is obtained by \emph{positing} a form and solving for the
constants rather than by rotating facets about creases and hoping the result
closes. Look for

\[x = i L_x + (j \bmod 2)\, d_x, \qquad y = j L_y, \qquad z = (i \bmod 2)\, H .\]

Two facet side vectors follow immediately, \(e_1 = (L_x, 0, \pm H)\) for the
horizontal edge of flat length \(a\), and \(e_2 = (\pm d_x, L_y, 0)\) for the zigzag
edge of flat length \(b\). Because both are shared by opposite sides, every facet
is a parallelogram and therefore planar with no further conditions. Isometry
then needs only the two edge lengths and the angle between them:

\[L_x^2 + H^2 = a^2, \qquad d_x^2 + L_y^2 = b^2, \qquad L_x d_x = ab\cos\alpha .\]

Three equations in four unknowns: a one-parameter family, which is the rigid
folding. Writing \(L_x = a\cos\varphi\) and \(H = a\sin\varphi\) solves the first
identically and leaves \(d_x = b\cos\alpha/\cos\varphi\) and
\(L_y = b\sqrt{1 - \cos^2\alpha/\cos^2\varphi}\).

\(L_y\) is real only while \(\varphi \le \alpha\), and vanishes exactly at
\(\varphi = \alpha\), which is the flat-folded state. So \(\varphi\) runs on
\([0, \alpha]\) and the natural normalised parameter is \(\theta = \varphi/\alpha\).

\textbf{\(\theta\) is a parameter, not an angle.} It runs on \([0, 1]\): zero is the flat
sheet and one is flat-folded. The dihedral angle \(\rho\) at a crease is \emph{derived}
from it and is never a synonym for it --- a distinction this project got wrong
once and now states in three places.

\subsection{Two: the isometry is exact, not converged}\label{two-the-isometry-is-exact-not-converged}

Because the placement was derived by imposing edge lengths and facet angles, the
isometry holds by construction rather than to a solver's tolerance. That is a
claim about algebra, and algebra is exactly the thing that might be wrong, so it
is measured.

\begin{Shaded}
\begin{Highlighting}[]
\NormalTok{worst\_iso }\OtherTok{\textless{}{-}} \FunctionTok{max}\NormalTok{(}\FunctionTok{vapply}\NormalTok{(}\FunctionTok{seq}\NormalTok{(}\DecValTok{0}\NormalTok{, }\DecValTok{1}\NormalTok{, }\AttributeTok{by =} \FloatTok{0.05}\NormalTok{),}
                        \ControlFlowTok{function}\NormalTok{(t) }\FunctionTok{facet\_isometry\_error}\NormalTok{(p, t), }\FunctionTok{numeric}\NormalTok{(}\DecValTok{1}\NormalTok{)))}
\end{Highlighting}
\end{Shaded}

\texttt{facet\_isometry\_error()} compares \textbf{all pairwise distances} between a facet's
vertices, folded against flat, not just its edges: an implementation that
preserved every edge of a parallelogram while shearing it would pass an edge
test and would not be a rigid folding. Over the whole \(\theta\) sweep the worst
departure on this pattern is 0, which is floating-point
noise on a closed form rather than a tolerance anyone chose.

\begin{figure}

{\centering \includegraphics[width=0.9\linewidth,alt={Two panels. The left panel is a flat Miura-ori crease pattern: a grid of congruent parallelograms whose vertical crease lines zigzag, with creases drawn in two colours and two line styles for mountain and valley. The right panel shows the same sheet folded into a corrugated surface, seen from above and to one side, with the near facets drawn over the far ones.}]{02-folding-geometry_files/figure-latex/geom-pattern-1} 

}

\caption{The flat crease pattern and the same sheet folded. Mountain and valley creases are distinguished by colour and by line style together, so the distinction survives greyscale printing.}\label{fig:geom-pattern}
\end{figure}

Panel A of Figure \ref{fig:geom-pattern} is the answer key. Panel B is what a
dimension-reduction method is shown. Everything in Part III is a measurement of
the gap between them.

\begin{figure}

{\centering \includegraphics[width=0.9\linewidth,alt={Three panels showing the same Miura-ori sheet folded to increasing degrees, left to right. In the leftmost the sheet is nearly flat and wide; in the middle it is visibly corrugated; in the rightmost it has contracted to roughly half its width and the corrugations are deep. All three are drawn at the same scale.}]{02-folding-geometry_files/figure-latex/geom-sweep-1} 

}

\caption{One pattern at three points of its one-parameter family. The sheet contracts in both directions as it folds; the flat chart behind it does not change at all.}\label{fig:geom-sweep}
\end{figure}

Figure \ref{fig:geom-sweep} is that family at three values of \(\theta\), drawn
at one scale. Two things are visible and both matter later. The sheet contracts
in \textbf{both} directions at once, which is why fold angle is a difficulty
parameter rather than a nuisance: it moves the ambient geometry continuously
while leaving the chart exactly where it was. And the contraction is not
uniform --- the far end of the range does most of it, which is why the \(\theta\)
grid in Part III is even in \(\theta\) and uneven in everything a method cares
about.

The derivation above is for one family. Nothing in this book requires a second,
and two candidates for one were withdrawn on evidence rather than dropped
quietly; Chapter \ref{benchmarks} reports both as results. What a reader should
take from that is the standard a folding has to meet here: a closed form, an
exact isometry, and a fold that reaches every crease family. Two of the three are
easy.

\subsection{Three: what folding does to distance along the surface}\label{three-what-folding-does-to-distance-along-the-surface}

Nothing. That is the point, and it needs no further argument once the facets are
congruent: a path on the sheet crosses the same facets in the same places at any
\(\theta\), so its length is the same. Geodesic distance between two sample points
is therefore exactly the distance between their positions in the flat chart,
subject to one caveat that is stated in \texttt{geom-limits} rather than hidden here.

\subsection{Four: what folding does to distance through space}\label{four-what-folding-does-to-distance-through-space}

Here folding does a great deal, and the amount is exact.

Take two points on opposite sides of a single crease, each at perpendicular
distance \(d\) from it and at the same position along it. Their chart distance is
\(2d\) at any fold angle. Their ambient distance, once the crease has closed to a
dihedral angle \(\rho\), is \(2d\sin(\rho/2)\). So

\[\frac{d_A}{d_U} = \sin(\rho/2),\]

which depends on \(\rho\) \textbf{alone}. Not on \(d\), not on where along the crease the
pair sits. That independence is the substance of the claim; a ratio that drifted
with distance would mean the relation was an approximation.

\begin{Shaded}
\begin{Highlighting}[]
\NormalTok{ratio\_err }\OtherTok{\textless{}{-}} \FunctionTok{max}\NormalTok{(}\FunctionTok{vapply}\NormalTok{(}\FunctionTok{seq}\NormalTok{(pi, }\FloatTok{0.05}\NormalTok{, }\AttributeTok{length.out =} \DecValTok{40}\NormalTok{), }\ControlFlowTok{function}\NormalTok{(rho) \{}
  \FunctionTok{max}\NormalTok{(}\FunctionTok{vapply}\NormalTok{(}\FunctionTok{c}\NormalTok{(}\FloatTok{0.1}\NormalTok{, }\DecValTok{1}\NormalTok{, }\FloatTok{7.3}\NormalTok{), }\ControlFlowTok{function}\NormalTok{(d) \{}
\NormalTok{    P }\OtherTok{\textless{}{-}} \FunctionTok{c}\NormalTok{( d }\SpecialCharTok{*} \FunctionTok{sin}\NormalTok{(rho }\SpecialCharTok{/} \DecValTok{2}\NormalTok{), }\DecValTok{0}\NormalTok{, d }\SpecialCharTok{*} \FunctionTok{cos}\NormalTok{(rho }\SpecialCharTok{/} \DecValTok{2}\NormalTok{))}
\NormalTok{    Q }\OtherTok{\textless{}{-}} \FunctionTok{c}\NormalTok{(}\SpecialCharTok{{-}}\NormalTok{d }\SpecialCharTok{*} \FunctionTok{sin}\NormalTok{(rho }\SpecialCharTok{/} \DecValTok{2}\NormalTok{), }\DecValTok{0}\NormalTok{, d }\SpecialCharTok{*} \FunctionTok{cos}\NormalTok{(rho }\SpecialCharTok{/} \DecValTok{2}\NormalTok{))}
    \FunctionTok{abs}\NormalTok{(}\FunctionTok{sqrt}\NormalTok{(}\FunctionTok{sum}\NormalTok{((P }\SpecialCharTok{{-}}\NormalTok{ Q)}\SpecialCharTok{\^{}}\DecValTok{2}\NormalTok{)) }\SpecialCharTok{/}\NormalTok{ (}\DecValTok{2} \SpecialCharTok{*}\NormalTok{ d) }\SpecialCharTok{{-}} \FunctionTok{sin}\NormalTok{(rho }\SpecialCharTok{/} \DecValTok{2}\NormalTok{))}
\NormalTok{  \}, }\FunctionTok{numeric}\NormalTok{(}\DecValTok{1}\NormalTok{)))}
\NormalTok{\}, }\FunctionTok{numeric}\NormalTok{(}\DecValTok{1}\NormalTok{)))}
\end{Highlighting}
\end{Shaded}

Over forty dihedral angles and three separations the worst departure from
\(\sin(\rho/2)\) is 0 --- one unit in the last place, which
is what ``exact'' means in double precision.

A flat sheet has \(\rho = \pi\) and contracts nothing. Every fold contracts, and
since \(\sin(\rho/2) < 1\) for every \(\rho < \pi\), \textbf{ambient distance is strictly
less than chart distance for every \(\theta > 0\).}

\subsection{Five: the floor under ``strictly''}\label{five-the-floor-under-strictly}

That last sentence is false as written, and the chapter says so rather than
leaving a reader to trip over it in Part III.

Because \(1 - \sin(\rho/2) \approx (\pi - \rho)^2/8\) near the flat state, the
contraction is quadratic in the fold angle while machine epsilon is not. There
is therefore a fold angle below which the contraction exists mathematically and
is not representable.

\begin{Shaded}
\begin{Highlighting}[]
\NormalTok{gap        }\OtherTok{\textless{}{-}} \ControlFlowTok{function}\NormalTok{(g) }\DecValTok{1} \SpecialCharTok{{-}} \FunctionTok{sin}\NormalTok{((pi }\SpecialCharTok{{-}}\NormalTok{ g) }\SpecialCharTok{/} \DecValTok{2}\NormalTok{)          }\CommentTok{\# g = pi {-} rho}
\NormalTok{floor\_pred }\OtherTok{\textless{}{-}} \FunctionTok{sqrt}\NormalTok{(}\DecValTok{8} \SpecialCharTok{*}\NormalTok{ .Machine}\SpecialCharTok{$}\NormalTok{double.eps }\SpecialCharTok{/} \DecValTok{2}\NormalTok{)}
\NormalTok{theta\_min  }\OtherTok{\textless{}{-}} \FunctionTok{min}\NormalTok{(THETA\_GRID[THETA\_GRID }\SpecialCharTok{\textgreater{}} \DecValTok{0}\NormalTok{])}
\end{Highlighting}
\end{Shaded}

The predicted floor is 0 radians, and it behaves:
\texttt{gap()} returns 0 at a fold angle of \(10^{-3}\) and exactly
0 at \(10^{-8}\). The smallest non-zero \(\theta\) this book's grid uses
is 0.05, which is above the floor by a wide margin --- and that margin is
the only reason Part III may treat the contraction as strict.

\section{Results}\label{geom-results}

Three consequences follow, and they are the reason this chapter comes before the
methods rather than after them.

\textbf{The answer key is not what the methods reproduce.} A method that consumes
ambient distance is being asked to reproduce a configuration that folding has
already contracted. Scoring it against the intrinsic chart is a legitimate
question --- it is the question of whether the method recovers the \emph{surface} --- but
it is not the question the method's loss function asks.

\textbf{The exact answer scores badly on the wrong metric, provably.} This deserves
stating as a proposition rather than demonstrating with a simulation.

\begin{quote}
\textbf{Proposition.} Let stress be measured against ambient distance, and let the
candidate embeddings be the exact chart \(\mathcal{U}\) and the two-component
PCA of the folded cloud. Then for every \(\theta > 0\),
\(\mathrm{stress}(\mathcal{U}) > \mathrm{stress}(\mathrm{PCA})\).
\end{quote}

The proof is one line given step four. Ambient-referenced stress measures
reproduction of \(d_A\); the contraction gives \(d_A < d_U\) strictly; and
two-component PCA is by construction the least-squares-optimal \emph{linear}
approximant to that same ambient configuration \citep{pearson1901lines, eckart1936approximation}. The exact unfolding is not competing to reproduce
ambient distance and loses to something that is.

An evaluator that ranks the true answer below a plainly wrong one is not
malfunctioning. It is answering the question it was asked. Chapter
\ref{evaluation} measures how far that goes on real evaluators, and finds it
goes further than this proposition requires.

\textbf{Difficulty has a knob.} Because \(\theta\) is continuous and the contraction is
monotone in it, every result in this book is a curve rather than a number. That
is standing rule 2, and it is not a stylistic preference. Chapter
\ref{benchmarks} reports a case where two method families swap places partway
along the \(\theta\) axis, so that the average over \(\theta\) names a different
winner from the curve at seven of its nine points. A benchmark with one fixed
geometry cannot produce that observation, because it has nowhere to put it.

This is the specific sense in which a crease pattern is a better instrument than
a Swiss roll, and it is worth separating from the claim it is often confused
with. It is not that folded sheets resemble real data more closely. It is that
the truth is known \textbf{and} the difficulty is continuously adjustable, so a
method's performance is a function rather than a point, and the shape of that
function is a property of the method. The critique that two-dimensional
embeddings distort more than their users admit \citep{chari2023specious, kobak2019art} is an argument this book answers with measurement, and
measurement of that kind needs an axis to measure along.

\begin{figure}

{\centering \includegraphics[width=0.9\linewidth,alt={A line chart. The horizontal axis is the dihedral angle from zero to pi; the vertical axis is the ratio of ambient to geodesic distance, from zero to one. A smooth curve rises from the origin to one at pi, steeply at first and flattening as it approaches the right-hand end. Measured points lie exactly on the curve.}]{02-folding-geometry_files/figure-latex/geom-curve-1} 

}

\caption{The ambient-to-geodesic distance ratio across a single crease, against the dihedral angle. The curve is $\sin(\rho/2)$; the points are measured from folded coordinates. A flat sheet sits at the right-hand end and contracts nothing.}\label{fig:geom-curve}
\end{figure}

\section{Diagnostics}\label{geom-diagnostics}

The most useful material in this chapter is about two things this project got
wrong. Both were caught, both are cheap to state, and both generalise past
origami.

\subsection{An isometry test cannot tell a fold from a pleat}\label{an-isometry-test-cannot-tell-a-fold-from-a-pleat}

A second pattern family, the Yoshimura, shipped for three days with a folding
that was not one. It folded its horizontal creases and left every diagonal at
\(\rho = \pi\) --- an accordion pleat wearing a diamond tessellation's crease
pattern.

It passed the facet-isometry test perfectly, because \textbf{a pleat is a rigid
folding}. It passed Kawasaki, because the flat pattern was never the problem.
Neither test could have caught it: both are properties a correct implementation
has, and a pleat has them too.

What caught it was deriving the mountain/valley assignment from the \emph{folded}
dihedral angles and testing Maekawa on the result. A pleat leaves whole crease
families flat, and a vertex with unfolded creases cannot balance:

\begin{Shaded}
\begin{Highlighting}[]
\NormalTok{a  }\OtherTok{\textless{}{-}} \FunctionTok{crease\_assignment}\NormalTok{(p)}
\NormalTok{iv }\OtherTok{\textless{}{-}} \FunctionTok{interior\_vertices}\NormalTok{(p)}
\NormalTok{maekawa }\OtherTok{\textless{}{-}} \FunctionTok{vapply}\NormalTok{(iv, }\ControlFlowTok{function}\NormalTok{(v) \{}
\NormalTok{  sel }\OtherTok{\textless{}{-}}\NormalTok{ (p}\SpecialCharTok{$}\NormalTok{creases}\SpecialCharTok{$}\NormalTok{i }\SpecialCharTok{==}\NormalTok{ v }\SpecialCharTok{|}\NormalTok{ p}\SpecialCharTok{$}\NormalTok{creases}\SpecialCharTok{$}\NormalTok{j }\SpecialCharTok{==}\NormalTok{ v) }\SpecialCharTok{\&}\NormalTok{ a }\SpecialCharTok{!=} \StringTok{"B"}
  \FunctionTok{abs}\NormalTok{(}\FunctionTok{sum}\NormalTok{(a[sel] }\SpecialCharTok{==} \StringTok{"M"}\NormalTok{) }\SpecialCharTok{{-}} \FunctionTok{sum}\NormalTok{(a[sel] }\SpecialCharTok{==} \StringTok{"V"}\NormalTok{))}
\NormalTok{\}, }\FunctionTok{numeric}\NormalTok{(}\DecValTok{1}\NormalTok{))}
\end{Highlighting}
\end{Shaded}

On this Miura all 25 interior vertices satisfy
\(\lvert M - V \rvert = 2\). On the pleat, twelve of sixteen
failed.

\subsection{A test that a correct implementation passes is not a test a wrong one fails}\label{a-test-that-a-correct-implementation-passes-is-not-a-test-a-wrong-one-fails}

That is the moral, and the same chapter has to record that this book learned it
twice.

The derived assignment above was itself wrong for four days, in a way Maekawa
could not see. Maekawa counts labels at a vertex; it is therefore invariant
under swapping every mountain for every valley, so it passes on a labelling, on
that labelling's exact inverse, and on a third that agrees with neither on half
the creases. All three were live in this repository at once, and the one the
figures drew was correct on half the interior creases.

\begin{Shaded}
\begin{Highlighting}[]
\NormalTok{inverted }\OtherTok{\textless{}{-}} \FunctionTok{ifelse}\NormalTok{(a }\SpecialCharTok{==} \StringTok{"M"}\NormalTok{, }\StringTok{"V"}\NormalTok{, }\FunctionTok{ifelse}\NormalTok{(a }\SpecialCharTok{==} \StringTok{"V"}\NormalTok{, }\StringTok{"M"}\NormalTok{, }\StringTok{"B"}\NormalTok{))}
\NormalTok{maekawa\_inv }\OtherTok{\textless{}{-}} \FunctionTok{vapply}\NormalTok{(iv, }\ControlFlowTok{function}\NormalTok{(v) \{}
\NormalTok{  sel }\OtherTok{\textless{}{-}}\NormalTok{ (p}\SpecialCharTok{$}\NormalTok{creases}\SpecialCharTok{$}\NormalTok{i }\SpecialCharTok{==}\NormalTok{ v }\SpecialCharTok{|}\NormalTok{ p}\SpecialCharTok{$}\NormalTok{creases}\SpecialCharTok{$}\NormalTok{j }\SpecialCharTok{==}\NormalTok{ v) }\SpecialCharTok{\&}\NormalTok{ inverted }\SpecialCharTok{!=} \StringTok{"B"}
  \FunctionTok{abs}\NormalTok{(}\FunctionTok{sum}\NormalTok{(inverted[sel] }\SpecialCharTok{==} \StringTok{"M"}\NormalTok{) }\SpecialCharTok{{-}} \FunctionTok{sum}\NormalTok{(inverted[sel] }\SpecialCharTok{==} \StringTok{"V"}\NormalTok{))}
\NormalTok{\}, }\FunctionTok{numeric}\NormalTok{(}\DecValTok{1}\NormalTok{))}
\end{Highlighting}
\end{Shaded}

The inverted labelling --- every mountain a valley --- satisfies Maekawa at
25 of 25 interior vertices, which is all of
them. The check that finally settled it computes the assignment a second time by
an independent route, from facet normals oriented against the ambient frame
rather than against the mesh's own winding, and compares.

\subsection{A first-order flex at the flat state proves nothing}\label{a-first-order-flex-at-the-flat-state-proves-nothing}

The third instance, recorded here because it is the same error in a different
costume. At \(\theta = 0\) every \(z\) column of a pattern's rigidity Jacobian
vanishes, so the count of non-trivial infinitesimal flexes is \(V - 3\) for \textbf{any}
planar pattern whatsoever --- including patterns with no rigid folding at all. A
first-order flex at the flat state is a property of being flat, not of being
foldable. Second-order analysis and a finite continuation are what decide, and
Chapter \ref{benchmarks} records what they decided.

\section{Where it fails}\label{geom-limits}

\subsection{The chart is the geodesic only inside a convex region}\label{the-chart-is-the-geodesic-only-inside-a-convex-region}

Step three claimed that geodesic distance equals chart distance. That is true
when the straight segment between two chart points stays on the paper, and a
Miura's unfolded outline is \textbf{not convex} --- its edge is a zigzag with teeth. For
a pair straddling a tooth the straight chart segment leaves the sheet, the true
geodesic must detour, and the chart distance is a \emph{lower bound} rather than an
equality.

\begin{Shaded}
\begin{Highlighting}[]
\NormalTok{m\_in  }\OtherTok{\textless{}{-}} \FunctionTok{sample\_manifold}\NormalTok{(p, }\AttributeTok{theta =} \FloatTok{0.6}\NormalTok{, }\AttributeTok{n =} \DecValTok{400}\DataTypeTok{L}\NormalTok{, }\AttributeTok{seed =} \DecValTok{1001}\DataTypeTok{L}\NormalTok{)}
\NormalTok{m\_out }\OtherTok{\textless{}{-}} \FunctionTok{sample\_manifold}\NormalTok{(p, }\AttributeTok{theta =} \FloatTok{0.6}\NormalTok{, }\AttributeTok{n =} \DecValTok{400}\DataTypeTok{L}\NormalTok{, }\AttributeTok{seed =} \DecValTok{1001}\DataTypeTok{L}\NormalTok{, }\AttributeTok{boundary =} \ConstantTok{TRUE}\NormalTok{)}
\NormalTok{exit\_in  }\OtherTok{\textless{}{-}} \FunctionTok{chart\_exit\_fraction}\NormalTok{(m\_in,  p)}
\NormalTok{exit\_out }\OtherTok{\textless{}{-}} \FunctionTok{chart\_exit\_fraction}\NormalTok{(m\_out, p)}
\end{Highlighting}
\end{Shaded}

\texttt{chart\_exit\_fraction()} measures how often it happens by sampling pairs and
walking the segment between them. Drawing strictly inside a convex sub-rectangle
gives an exit fraction of 0; including the boundary strip
gives 0.013. Every artefact in this book that the answer key is
read against is drawn the first way, and the producer that generates the
irreducible-loss suite asserts the zero rather than trusting this paragraph.

\subsection{Self-intersection}\label{self-intersection}

Nothing in the algebra of step one prevents a folded placement from passing
through itself, and a self-intersecting surface would invalidate every ambient
distance measured on it.

\begin{Shaded}
\begin{Highlighting}[]
\NormalTok{gap\_at }\OtherTok{\textless{}{-}} \FunctionTok{vapply}\NormalTok{(}\FunctionTok{c}\NormalTok{(}\FloatTok{0.3}\NormalTok{, }\FloatTok{0.6}\NormalTok{, }\FloatTok{0.9}\NormalTok{), }\ControlFlowTok{function}\NormalTok{(t) }\FunctionTok{facet\_gap}\NormalTok{(p, t)}\SpecialCharTok{$}\NormalTok{gap, }\FunctionTok{numeric}\NormalTok{(}\DecValTok{1}\NormalTok{))}
\end{Highlighting}
\end{Shaded}

\texttt{facet\_gap()} computes the exact minimum distance between non-adjacent facets ---
polygon to polygon, analytically, not by sampling points on them. At \(\theta\) of
0.3, 0.6 and 0.9 it is 0.851,
0.786 and 0.526: falling as the sheet
closes, as it must, and positive throughout.

\begin{figure}

{\centering \includegraphics[width=0.9\linewidth,alt={A line chart. The horizontal axis is the folding parameter from zero to one; the vertical axis is the minimum distance between non-adjacent facets. The curve starts high at the left and falls steadily, approaching but not reaching zero at the right-hand end. A vertical dashed line near the right marks the grid's largest fold angle, where the clearance is still clearly positive.}]{02-folding-geometry_files/figure-latex/geom-clearance-1} 

}

\caption{Minimum clearance between non-adjacent facets, against the folding parameter. The surface approaches itself as it folds and never touches; the dashed line marks the largest fold angle the benchmark grid uses.}\label{fig:geom-clearance}
\end{figure}

Figure \ref{fig:geom-clearance} is the whole range. The clearance falls
smoothly and stays positive, and the largest \(\theta\) the benchmark grid uses is
marked: at 0.95 the sheet still has 0.396
of clearance. That is the check behind an assumption every ambient distance in
this book depends on, and it is a check rather than an argument because the
algebra of step one does not forbid self-intersection --- it simply does not
produce one for this family in this range.

\subsection{Sampling, and what a finite sample of a zero-reach surface is}\label{sampling-and-what-a-finite-sample-of-a-zero-reach-surface-is}

The reach remark in \texttt{geom-background} has a practical edge that belongs here.
Every guarantee about recovering a manifold from a finite sample is stated in
terms of the density of the sample relative to the reach \citep{niyogi2008homology},
and a crease has reach zero, so no such guarantee is available at any sample
size. What replaces it in this book is not a weaker theorem but a measurement:
the sampled surface is checked for how close distant parts of it come, in units
of the local spacing, and that ratio is reported alongside every result that
depends on it. Isomap's failure mode on tightly folded surfaces is precisely a
short circuit between sheets that the sampling cannot distinguish
\citep{tenenbaum2000isomap, balasubramanian2002isomap}, and the quantity that
predicts it is a property of the sample rather than of the surface.

\subsection{One family}\label{one-family}

This book ships one pattern family that folds. Two others were specified and
both were withdrawn on evidence, which Chapter \ref{benchmarks} reports as
results rather than as gaps. A benchmark built on a single family cannot support
claims about pattern variety, and this book does not make any.

\subsection{What the chapter does not settle}\label{what-the-chapter-does-not-settle}

Whether a folded sheet is a \emph{good} model of the manifolds that arise in real
data is not a geometric question and is not answered here. Chapter
\ref{comparison} measures what this benchmark discriminates that the classical
ones do not, and Chapter \ref{application} takes the result to a real dataset.
The honest position in between is that a crease pattern is a manifold whose
truth is known exactly, which is a property no real dataset has and no
previously used benchmark provides.

\section{Reproduce this}\label{geom-reproduce}

Every claim above is an assertion in the test suite, and the suite runs on every
push:

\begin{itemize}
\tightlist
\item
  \texttt{tests/testthat/test-folding.R} --- facet isometry across the \(\theta\) sweep and
  the \(\alpha\) range, no self-intersection, the dihedral contract.
\item
  \texttt{tests/testthat/test-patterns.R} --- congruence, Kawasaki, Maekawa, vertex
  degree.
\item
  \texttt{tests/testthat/test-crease-assignment.R} --- the mountain/valley derivation
  against an independent criterion, and the demonstration that Maekawa cannot
  discriminate.
\item
  \texttt{tests/testthat/test-contraction.R} --- \(\sin(\rho/2)\) and the floating-point
  floor.
\end{itemize}

Run them with \texttt{Rscript\ tests/testthat.R} from the repository root.

\section{Further reading}\label{geom-reading}

Hull's \emph{Origametry} \citep{hull2020origametry} is the single best entry point to the
mathematics of flat-foldability and is the reference for both classical theorems
above. Demaine and O'Rourke \citep{demaine2007folding} is the standard treatment of
the computational side, and Lang \citep{lang2011origami} of design. Miura's original
report \citep{miura1985method} introduces the pattern this book uses; Schenk and Guest
\citep{schenk2013geometry} give the mechanics of folded shell structures and are the
natural next step for a reader interested in the surfaces rather than the
sheets.

On the manifold-learning side, Federer's reach \citep{federer1959curvature} and its
role in sampling guarantees \citep{niyogi2008homology} are what the zero-reach remark
above refers to, and Lee and Verleysen \citep{lee2007nonlinear} is the textbook
treatment of the methods Part II takes up.

\chapter{Folding as a generative model}\label{generative}

\emph{R. Heller · Date: 2026-07-29 · Last changed: 2026-09-07}

\emph{How do you turn one pattern into a family of test problems?}

\section{What this chapter answers}\label{gen-question}

Chapter \ref{folding-geometry} established one object: a crease pattern whose
flat chart is its exact two-dimensional structure. One object is not a
benchmark. A benchmark is a \emph{family} of problems ordered by difficulty, and this
chapter is where the single sheet becomes that family.

The generative model has four scalars and three switches. Three of the scalars
--- the cell sides \(a\) and \(b\) and the acute angle \(\alpha\) --- fix the chart. The
fourth, \(\theta\), moves the sheet through space and does not touch the chart at
all. The switches are the sample size, the noise model and the seed. Every row
of every artefact this book ships is one setting of those seven things, and
every curve in Parts II and III is a sweep over the fourth scalar with the other
six held.

This chapter is the reference for that model. It derives the folded placement in
closed form and verifies the derivation here rather than citing one; it says
exactly how points are drawn and what a careless draw would do to the density;
it names the three noise models and no others; and it measures what \(\theta\)
does to the \(k\)-nearest-neighbour graph a method actually consumes. Chapter
\ref{benchmarks} builds constructions on top of this model, and every chapter
from \ref{linear} onward reads an artefact produced by it.

\section{Setup}\label{gen-setup}

\begin{Shaded}
\begin{Highlighting}[]
\NormalTok{p  }\OtherTok{\textless{}{-}} \FunctionTok{miura\_ori}\NormalTok{(}\AttributeTok{nx =} \DecValTok{6}\DataTypeTok{L}\NormalTok{, }\AttributeTok{ny =} \DecValTok{6}\DataTypeTok{L}\NormalTok{, }\AttributeTok{a =} \DecValTok{1}\NormalTok{, }\AttributeTok{b =} \DecValTok{1}\NormalTok{, }\AttributeTok{alpha =}\NormalTok{ pi }\SpecialCharTok{/} \DecValTok{3}\NormalTok{)}
\NormalTok{th }\OtherTok{\textless{}{-}} \FunctionTok{max}\NormalTok{(THETA\_GRID)}
\NormalTok{s  }\OtherTok{\textless{}{-}} \FunctionTok{sample\_manifold}\NormalTok{(p, }\AttributeTok{theta =}\NormalTok{ th, }\AttributeTok{n =}\NormalTok{ N\_DEFAULT, }\AttributeTok{seed =} \DecValTok{1001}\DataTypeTok{L}\NormalTok{)}
\end{Highlighting}
\end{Shaded}

\texttt{p} is the pattern the benchmark grid uses --- \texttt{PATTERN\_GRID} names its size and
\texttt{R/constants.R} explains why there is one entry --- \texttt{th} the largest fold the grid
reaches, and \texttt{s} a sample of 800 points drawn from the sheet at that
fold. Nothing in this chapter reads a committed artefact; every number below is
computed while the page is built. The \hyperref[notation]{notation appendix} fixes the
symbols.

\section{Background}\label{gen-background}

Dimension reduction is usually benchmarked against a \emph{collection}: a fixed list
of datasets, some real and some synthetic, each scored once and averaged over
\citep{espadoto2021survey}. The synthetic members of such a collection --- the Swiss
roll, the S-curve, the severed sphere --- are already generative in the weak sense
that they come with a parameterisation, so a ground-truth chart exists
\citep{tenenbaum2000isomap, lee2007nonlinear}. What they mostly lack is a knob: a
continuous parameter that makes the \emph{same} problem harder while leaving the
answer exactly where it was.

A generative model that can carry a benchmark has to supply four things.

\textbf{An exact chart.} Not a low-dimensional summary that scores well, but the
coordinates the data were generated from. A rigid folding gives this for free,
which is the whole reason the book uses one.

\textbf{A difficulty axis that moves the ambient geometry and nothing else.} Folding
moves every point in \(\mathbb{R}^3\) and moves no point of the chart. That is the
property \(\theta\) has and a dataset collection does not: two cells of the grid
can differ in difficulty while their answer keys are literally the same matrix.

\textbf{A noise model that is not confounded with the axis.} This one is easy to get
wrong. A noise scale keyed to the folded object would shrink as the sheet
closes, so a fixed setting would become a rising noise level along the \(\theta\)
sweep and the two axes could not be separated afterwards. The scale here is a
fraction of the \emph{chart} diameter, which does not move.

\textbf{A seed, and a sampler that respects it.} Every stochastic result in this book
is reported across 20 seeds (standing rule 1), which requires a sampler
whose draw is a function of its seed and whose call does not disturb the
caller's random stream.

One thing the model cannot supply is a sampling guarantee. The theory that says
how many points are enough to recover a submanifold is stated in terms of its
reach \citep{niyogi2008homology}, and Chapter \ref{folding-geometry} established
that a crease has reach exactly zero. Nothing here is covered by that theory, so
sampling adequacy is \emph{measured} on every sample instead --- which is what
\texttt{gen-results} does.

\section{The Miura-ori in four scalars}\label{gen-core}

Four scalars go in and four come out, and the two sets are not the same four.

Going in are \(a\), \(b\), \(\alpha\) and \(\theta\). The first three fix the flat
pattern: vertex \((i, j)\) of the tessellation sits at

\[x = i\,a + (j \bmod 2)\, b \cos\alpha, \qquad y = j\, b \sin\alpha,\]

which is Miura's parallelogram tessellation \citep{miura1985method} and is all of
\texttt{miura\_ori()}. The counts \(n_x\) and \(n_y\) set how many cells there are, not what
shape they have. \(\theta\) appears nowhere in that expression, and that absence is
the model's load-bearing property.

Coming out are the four metric constants of the folded placement. Chapter
\ref{folding-geometry} derived them; the derivation is repeated here in one
line because this chapter is where a reader is sent to find it. Posit a folded
placement of the same combinatorial form,

\[x = i L_x + (j \bmod 2)\, d_x, \qquad y = j L_y, \qquad z = (i \bmod 2)\, H,\]

and impose isometry on one facet. Its two side vectors are
\(e_1 = (L_x, 0, \pm H)\) of flat length \(a\) and \(e_2 = (\pm d_x, L_y, 0)\) of flat
length \(b\), so preserving both lengths and the angle between them gives

\[L_x^2 + H^2 = a^2, \qquad d_x^2 + L_y^2 = b^2, \qquad L_x d_x = ab\cos\alpha .\]

Three equations, four unknowns, one free parameter. Writing \(L_x = a\cos\varphi\)
and \(H = a\sin\varphi\) solves the first identically and leaves
\(d_x = b\cos\alpha/\cos\varphi\) and
\(L_y = b\sqrt{1 - \cos^2\alpha/\cos^2\varphi}\). That \(L_y\) is real only while
\(\varphi \le \alpha\) and vanishes at \(\varphi = \alpha\), which is the
flat-folded state --- so \(\varphi\) runs on \([0, \alpha]\) and \(\theta = \varphi/\alpha\)
is the normalised parameter \texttt{fold()} takes.

\subsection{The derivation, checked}\label{the-derivation-checked}

Those three conditions are what the placement was built to satisfy, so testing
them on the formula would be circular. What is not circular is reading the four
constants back \emph{out} of the coordinates \texttt{fold()} returns and testing the
conditions on those.

\begin{Shaded}
\begin{Highlighting}[]
\NormalTok{scalars\_of }\OtherTok{\textless{}{-}} \ControlFlowTok{function}\NormalTok{(pattern, theta) \{}
\NormalTok{  r   }\OtherTok{\textless{}{-}}\NormalTok{ pattern}\SpecialCharTok{$}\NormalTok{params}
\NormalTok{  vid }\OtherTok{\textless{}{-}} \ControlFlowTok{function}\NormalTok{(i, j) j }\SpecialCharTok{*}\NormalTok{ (r}\SpecialCharTok{$}\NormalTok{nx }\SpecialCharTok{+} \DecValTok{1}\DataTypeTok{L}\NormalTok{) }\SpecialCharTok{+}\NormalTok{ i }\SpecialCharTok{+} \DecValTok{1}\DataTypeTok{L}        \CommentTok{\# (i, j) {-}\textgreater{} vertex row}
\NormalTok{  V   }\OtherTok{\textless{}{-}} \FunctionTok{unname}\NormalTok{(}\FunctionTok{fold}\NormalTok{(pattern, theta)}\SpecialCharTok{$}\NormalTok{vertices3)}
  \FunctionTok{c}\NormalTok{(}\AttributeTok{Lx =}\NormalTok{ V[}\FunctionTok{vid}\NormalTok{(}\DecValTok{1}\DataTypeTok{L}\NormalTok{, }\DecValTok{0}\DataTypeTok{L}\NormalTok{), }\DecValTok{1}\DataTypeTok{L}\NormalTok{] }\SpecialCharTok{{-}}\NormalTok{ V[}\FunctionTok{vid}\NormalTok{(}\DecValTok{0}\DataTypeTok{L}\NormalTok{, }\DecValTok{0}\DataTypeTok{L}\NormalTok{), }\DecValTok{1}\DataTypeTok{L}\NormalTok{],}
    \AttributeTok{H  =}\NormalTok{ V[}\FunctionTok{vid}\NormalTok{(}\DecValTok{1}\DataTypeTok{L}\NormalTok{, }\DecValTok{0}\DataTypeTok{L}\NormalTok{), }\DecValTok{3}\DataTypeTok{L}\NormalTok{] }\SpecialCharTok{{-}}\NormalTok{ V[}\FunctionTok{vid}\NormalTok{(}\DecValTok{0}\DataTypeTok{L}\NormalTok{, }\DecValTok{0}\DataTypeTok{L}\NormalTok{), }\DecValTok{3}\DataTypeTok{L}\NormalTok{],}
    \AttributeTok{dx =}\NormalTok{ V[}\FunctionTok{vid}\NormalTok{(}\DecValTok{0}\DataTypeTok{L}\NormalTok{, }\DecValTok{1}\DataTypeTok{L}\NormalTok{), }\DecValTok{1}\DataTypeTok{L}\NormalTok{] }\SpecialCharTok{{-}}\NormalTok{ V[}\FunctionTok{vid}\NormalTok{(}\DecValTok{0}\DataTypeTok{L}\NormalTok{, }\DecValTok{0}\DataTypeTok{L}\NormalTok{), }\DecValTok{1}\DataTypeTok{L}\NormalTok{],}
    \AttributeTok{Ly =}\NormalTok{ V[}\FunctionTok{vid}\NormalTok{(}\DecValTok{0}\DataTypeTok{L}\NormalTok{, }\DecValTok{1}\DataTypeTok{L}\NormalTok{), }\DecValTok{2}\DataTypeTok{L}\NormalTok{] }\SpecialCharTok{{-}}\NormalTok{ V[}\FunctionTok{vid}\NormalTok{(}\DecValTok{0}\DataTypeTok{L}\NormalTok{, }\DecValTok{0}\DataTypeTok{L}\NormalTok{), }\DecValTok{2}\DataTypeTok{L}\NormalTok{])}
\NormalTok{\}}
\NormalTok{conditions }\OtherTok{\textless{}{-}} \ControlFlowTok{function}\NormalTok{(pattern, theta) \{}
\NormalTok{  r }\OtherTok{\textless{}{-}}\NormalTok{ pattern}\SpecialCharTok{$}\NormalTok{params}
\NormalTok{  z }\OtherTok{\textless{}{-}} \FunctionTok{scalars\_of}\NormalTok{(pattern, theta)}
  \FunctionTok{c}\NormalTok{(z[[}\StringTok{"Lx"}\NormalTok{]]}\SpecialCharTok{\^{}}\DecValTok{2} \SpecialCharTok{+}\NormalTok{ z[[}\StringTok{"H"}\NormalTok{]]}\SpecialCharTok{\^{}}\DecValTok{2}  \SpecialCharTok{{-}}\NormalTok{ r}\SpecialCharTok{$}\NormalTok{a}\SpecialCharTok{\^{}}\DecValTok{2}\NormalTok{,}
\NormalTok{    z[[}\StringTok{"dx"}\NormalTok{]]}\SpecialCharTok{\^{}}\DecValTok{2} \SpecialCharTok{+}\NormalTok{ z[[}\StringTok{"Ly"}\NormalTok{]]}\SpecialCharTok{\^{}}\DecValTok{2} \SpecialCharTok{{-}}\NormalTok{ r}\SpecialCharTok{$}\NormalTok{b}\SpecialCharTok{\^{}}\DecValTok{2}\NormalTok{,}
\NormalTok{    z[[}\StringTok{"Lx"}\NormalTok{]] }\SpecialCharTok{*}\NormalTok{ z[[}\StringTok{"dx"}\NormalTok{]]     }\SpecialCharTok{{-}}\NormalTok{ r}\SpecialCharTok{$}\NormalTok{a }\SpecialCharTok{*}\NormalTok{ r}\SpecialCharTok{$}\NormalTok{b }\SpecialCharTok{*} \FunctionTok{cos}\NormalTok{(r}\SpecialCharTok{$}\NormalTok{alpha))}
\NormalTok{\}}
\NormalTok{alphas }\OtherTok{\textless{}{-}} \FunctionTok{seq}\NormalTok{(MIURA\_ALPHA\_RANGE[}\DecValTok{1}\NormalTok{], MIURA\_ALPHA\_RANGE[}\DecValTok{2}\NormalTok{], }\AttributeTok{length.out =} \DecValTok{14}\DataTypeTok{L}\NormalTok{)}
\NormalTok{worst  }\OtherTok{\textless{}{-}} \FunctionTok{max}\NormalTok{(}\FunctionTok{vapply}\NormalTok{(alphas, }\ControlFlowTok{function}\NormalTok{(A) \{}
\NormalTok{  pa }\OtherTok{\textless{}{-}} \FunctionTok{miura\_ori}\NormalTok{(}\DecValTok{4}\DataTypeTok{L}\NormalTok{, }\DecValTok{4}\DataTypeTok{L}\NormalTok{, }\AttributeTok{a =} \DecValTok{1}\NormalTok{, }\AttributeTok{b =} \FloatTok{1.3}\NormalTok{, }\AttributeTok{alpha =}\NormalTok{ A)}
  \FunctionTok{max}\NormalTok{(}\FunctionTok{abs}\NormalTok{(}\FunctionTok{vapply}\NormalTok{(THETA\_GRID, }\ControlFlowTok{function}\NormalTok{(t) }\FunctionTok{conditions}\NormalTok{(pa, t), }\FunctionTok{numeric}\NormalTok{(}\DecValTok{3}\NormalTok{))))}
\NormalTok{\}, }\FunctionTok{numeric}\NormalTok{(}\DecValTok{1}\NormalTok{)))}
\NormalTok{ends }\OtherTok{\textless{}{-}} \FunctionTok{rbind}\NormalTok{(}\AttributeTok{flat =} \FunctionTok{scalars\_of}\NormalTok{(p, }\DecValTok{0}\NormalTok{), }\AttributeTok{folded =} \FunctionTok{scalars\_of}\NormalTok{(p, }\DecValTok{1}\NormalTok{))}
\end{Highlighting}
\end{Shaded}

Over fourteen opening angles spanning \texttt{MIURA\_ALPHA\_RANGE} and the whole
\(\theta\) grid, the worst violation of any of the three conditions is
6.7e-16 --- floating-point noise on a
closed form, not a solver tolerance. The endpoints behave too: at \(\theta = 0\)
the four constants are
1, 0, 0.5, 0.866, which is the flat pattern
recovered exactly rather than in a limit, and at \(\theta = 1\) they are
0.5, 0.866, 1, 0 --- \(L_y\) has gone to zero
and the sheet has no extent in \(y\) at all.

The four constants also say what folding does to the sheet as an object. \(L_x\)
and \(L_y\) both fall while \(H\) and \(d_x\) rise, so the sheet's overall extent
shrinks in \emph{both} in-plane directions at once while it gains depth in the third
--- the auxetic behaviour the Miura is usually presented for, arriving here as a
consequence of the three conditions rather than as a design goal. Nothing else
in the model depends on \(n_x\) and \(n_y\) beyond the count of cells: a larger
tessellation is a larger sheet made of the same cell, which is why crease count
is a difficulty knob in its own right and is swept separately rather than folded
into \(\theta\).

\begin{boxreport}
\textbf{This derivation is not quoted from anywhere.} Schenk and Guest
{[}@schenk2013geometry{]} give unit-cell quantities for the Miura under
a different angle convention, and their parameterisation and this one
are \emph{related work rather than the same statement}. No map between
the two is asserted here, because none has been verified in this
repository, and a reconciliation asserted without one would be exactly
the citation-by-osmosis this book is trying to avoid. The algebra above
stands on the check in \texttt{gen-scalars} and on
\texttt{tests/testthat/test-folding.R}, and on nothing else.
\end{boxreport}

\begin{figure}

{\centering \includegraphics[width=0.9\linewidth,alt={Six panels in two rows of three, each showing a folded Miura-ori sheet drawn in projection. The top row uses a shallow opening angle and the bottom row a steep one, so the two rows have visibly different cell shapes. Left to right within each row the sheet goes from flat, through half folded, to nearly collapsed, contracting in both directions as it goes. Mountain and valley creases are drawn in two colours and two line styles.}]{03-generative-model_files/figure-latex/gen-fig-family-1} 

}

\caption{The generative family. Rows are two opening angles $\alpha$; columns are three values of the folding parameter. The chart differs between rows and is identical along each row: $\theta$ moves the sheet and never the answer key.}\label{fig:gen-fig-family}
\end{figure}

Figure \ref{fig:gen-fig-family} is the family. Reading down a column changes the
chart; reading across a row does not change it by one coordinate. A benchmark
cell is a position in that array.

\begin{figure}

{\centering \includegraphics[width=0.9\linewidth,alt={Four small line charts side by side, one for each constant. Reading left to right: the horizontal edge projection falls from one to about 0.55, the corrugation height rises from zero to about 0.84, the zigzag offset rises from one half to about 0.92, and the row spacing falls from about 0.87 to about 0.4. All four bend most towards the right-hand end. Measured points lie on every line.}]{03-generative-model_files/figure-latex/gen-fig-scalars-1} 

}

\caption{The four metric constants of the folded placement against the folding parameter, for the pattern of `gen-setup`. Lines are the closed form; points are the same constants read back out of the coordinates `fold()` returns.}\label{fig:gen-fig-scalars}
\end{figure}

Figure \ref{fig:gen-fig-scalars} is the one-parameter family written out. Two
features of it matter downstream. The constants are smooth and monotone in
\(\theta\), so difficulty is ordered; and they are markedly \emph{not} linear in it, so
the last fifth of the sweep does most of the geometric work. Both facts turn up
again in \texttt{gen-results}.

\section{Sampling and noise}\label{gen-sampling}

The chart is the sample space. \texttt{sample\_manifold()} draws points in the flat
pattern, records them as \texttt{truth}, and carries each one into \(\mathbb{R}^3\)
through the rigid motion of the facet it landed on to give \texttt{X}. The two matrices
are the same points seen twice, and nothing is interpolated between folded
vertices --- the per-facet map is built from a frame and then \emph{checked} against the
corners, so a folding that stopped being rigid would fail rather than quietly
return positions no point of the surface occupies.

Drawing in the chart means drawing uniformly over area, and that is where a
naive implementation goes wrong. \texttt{.fan\_triangles()} splits every facet into a
fan from its first vertex and returns the area of each triangle; the fan is
checked rather than assumed, since it tiles a convex polygon and not a
re-entrant one. \texttt{.draw\_chart\_points()} then picks a triangle with probability
proportional to its area and maps a pair of uniforms into it --- with a square
root on the first.

\begin{Shaded}
\begin{Highlighting}[]
\NormalTok{tri }\OtherTok{\textless{}{-}} \FunctionTok{.fan\_triangles}\NormalTok{(p}\SpecialCharTok{$}\NormalTok{vertices, p}\SpecialCharTok{$}\NormalTok{facets)          }\CommentTok{\# internal: see R/sampling.R}
\NormalTok{draw\_chart }\OtherTok{\textless{}{-}} \ControlFlowTok{function}\NormalTok{(m, root) \{                    }\CommentTok{\# root = FALSE is the naive draw}
\NormalTok{  k  }\OtherTok{\textless{}{-}} \FunctionTok{sample.int}\NormalTok{(}\FunctionTok{nrow}\NormalTok{(tri}\SpecialCharTok{$}\NormalTok{idx), m, }\AttributeTok{replace =} \ConstantTok{TRUE}\NormalTok{, }\AttributeTok{prob =}\NormalTok{ tri}\SpecialCharTok{$}\NormalTok{area)}
\NormalTok{  va }\OtherTok{\textless{}{-}}\NormalTok{ p}\SpecialCharTok{$}\NormalTok{vertices[tri}\SpecialCharTok{$}\NormalTok{idx[k, }\DecValTok{1}\DataTypeTok{L}\NormalTok{], ]}
\NormalTok{  vb }\OtherTok{\textless{}{-}}\NormalTok{ p}\SpecialCharTok{$}\NormalTok{vertices[tri}\SpecialCharTok{$}\NormalTok{idx[k, }\DecValTok{2}\DataTypeTok{L}\NormalTok{], ]}
\NormalTok{  vc }\OtherTok{\textless{}{-}}\NormalTok{ p}\SpecialCharTok{$}\NormalTok{vertices[tri}\SpecialCharTok{$}\NormalTok{idx[k, }\DecValTok{3}\DataTypeTok{L}\NormalTok{], ]}
\NormalTok{  r1 }\OtherTok{\textless{}{-}}\NormalTok{ stats}\SpecialCharTok{::}\FunctionTok{runif}\NormalTok{(m); }\ControlFlowTok{if}\NormalTok{ (root) r1 }\OtherTok{\textless{}{-}} \FunctionTok{sqrt}\NormalTok{(r1)}
\NormalTok{  r2 }\OtherTok{\textless{}{-}}\NormalTok{ stats}\SpecialCharTok{::}\FunctionTok{runif}\NormalTok{(m)}
  \FunctionTok{list}\NormalTok{(}\AttributeTok{u =}\NormalTok{ (}\DecValTok{1} \SpecialCharTok{{-}}\NormalTok{ r1) }\SpecialCharTok{*}\NormalTok{ va }\SpecialCharTok{+}\NormalTok{ (r1 }\SpecialCharTok{*}\NormalTok{ (}\DecValTok{1} \SpecialCharTok{{-}}\NormalTok{ r2)) }\SpecialCharTok{*}\NormalTok{ vb }\SpecialCharTok{+}\NormalTok{ (r1 }\SpecialCharTok{*}\NormalTok{ r2) }\SpecialCharTok{*}\NormalTok{ vc, }\AttributeTok{r1 =}\NormalTok{ r1)}
\NormalTok{\}}
\NormalTok{cell\_cv }\OtherTok{\textless{}{-}} \ControlFlowTok{function}\NormalTok{(u, }\AttributeTok{bins =} \DecValTok{36}\DataTypeTok{L}\NormalTok{) \{}
\NormalTok{  n }\OtherTok{\textless{}{-}} \FunctionTok{as.vector}\NormalTok{(}\FunctionTok{table}\NormalTok{(}\FunctionTok{cut}\NormalTok{(u[, }\DecValTok{1}\DataTypeTok{L}\NormalTok{], bins), }\FunctionTok{cut}\NormalTok{(u[, }\DecValTok{2}\DataTypeTok{L}\NormalTok{], bins)))}
\NormalTok{  stats}\SpecialCharTok{::}\FunctionTok{sd}\NormalTok{(n[n }\SpecialCharTok{\textgreater{}} \DecValTok{0}\DataTypeTok{L}\NormalTok{]) }\SpecialCharTok{/} \FunctionTok{mean}\NormalTok{(n[n }\SpecialCharTok{\textgreater{}} \DecValTok{0}\DataTypeTok{L}\NormalTok{])}
\NormalTok{\}}
\FunctionTok{set.seed}\NormalTok{(}\DecValTok{7}\DataTypeTok{L}\NormalTok{)}
\NormalTok{good }\OtherTok{\textless{}{-}} \FunctionTok{draw\_chart}\NormalTok{(}\DecValTok{20000}\DataTypeTok{L}\NormalTok{, }\AttributeTok{root =} \ConstantTok{TRUE}\NormalTok{)}
\NormalTok{bad  }\OtherTok{\textless{}{-}} \FunctionTok{draw\_chart}\NormalTok{(}\DecValTok{20000}\DataTypeTok{L}\NormalTok{, }\AttributeTok{root =} \ConstantTok{FALSE}\NormalTok{)}
\NormalTok{apex }\OtherTok{\textless{}{-}} \FunctionTok{c}\NormalTok{(}\AttributeTok{area =} \FunctionTok{mean}\NormalTok{(good}\SpecialCharTok{$}\NormalTok{r1 }\SpecialCharTok{\textless{}=} \FloatTok{0.5}\NormalTok{), }\AttributeTok{naive =} \FunctionTok{mean}\NormalTok{(bad}\SpecialCharTok{$}\NormalTok{r1 }\SpecialCharTok{\textless{}=} \FloatTok{0.5}\NormalTok{))}
\NormalTok{cv   }\OtherTok{\textless{}{-}} \FunctionTok{c}\NormalTok{(}\AttributeTok{area =} \FunctionTok{cell\_cv}\NormalTok{(good}\SpecialCharTok{$}\NormalTok{u),      }\AttributeTok{naive =} \FunctionTok{cell\_cv}\NormalTok{(bad}\SpecialCharTok{$}\NormalTok{u))}
\NormalTok{area\_spread }\OtherTok{\textless{}{-}} \FunctionTok{diff}\NormalTok{(}\FunctionTok{range}\NormalTok{(tri}\SpecialCharTok{$}\NormalTok{area)) }\SpecialCharTok{/} \FunctionTok{mean}\NormalTok{(tri}\SpecialCharTok{$}\NormalTok{area)}
\end{Highlighting}
\end{Shaded}

Without the square root, the radial coordinate is uniform where the area element
is not: the fraction of a triangle within half-scale of its apex is one quarter
by area, and the naive draw puts 0.502 of its points
there against 0.251 for the corrected one. The apex of
every fan is a \emph{vertex of the crease pattern}, so the naive draw piles density
onto exactly the features the benchmark is about. Counted on a square grid over
the chart, the cell-count coefficient of variation is
0.855 against 0.283 --- a density
gradient roughly 3 times as strong, and
every neighbourhood metric in Part II reads a density gradient as structure.

The area weighting itself is a different correction and on this family it does
nothing: every Miura facet is congruent, so the spread of triangle areas is
1.2e-15, and area-proportional
and facet-uniform draws coincide. That is a fact about the pattern, not about the sampler, and it is
listed in \texttt{gen-limits} as something this book does not test.

\begin{figure}

{\centering \includegraphics[width=0.9\linewidth,alt={Two square scatter panels of chart points over the same region of a Miura pattern. Panel A on the left is an even wash of points with no visible structure. Panel B on the right shows the same region with distinctly darker clusters arranged on a regular lattice, one at the corner vertex of every cell, and lighter gaps between them.}]{03-generative-model_files/figure-latex/gen-fig-density-1} 

}

\caption{The same triangles, drawn from twice. Panel A is the sampler's draw, uniform over area. Panel B is the naive barycentric draw without the square root, which concentrates points at the fan apex of every triangle — that is, at the crease-pattern vertices.}\label{fig:gen-fig-density}
\end{figure}

\texttt{boundary\ =\ FALSE}, which is the default and what every committed artefact uses,
keeps only points at least half a median crease length from the sheet edge, by
rejection from the same uniform draw. Rejection rather than re-triangulation is
deliberate: survivors of a uniform draw are uniform over whatever region
survives, however awkward its shape.

Then noise, and there are exactly three models. \texttt{"none"} leaves \texttt{X} alone.
\texttt{"ambient"} adds isotropic Gaussian displacement in \(\mathbb{R}^3\), applied
\emph{after} folding so that it moves points off the surface. \texttt{"outlier"} throws a
fixed fraction of the points a fixed distance off it. There is no on-manifold
noise model: displacement inside the chart is a reparameterisation of the same
surface rather than measurement error, and \texttt{sample\_manifold()} rejects any type
it does not know.

\begin{Shaded}
\begin{Highlighting}[]
\NormalTok{NOISE }\OtherTok{\textless{}{-}} \FunctionTok{list}\NormalTok{(}\AttributeTok{none    =} \FunctionTok{list}\NormalTok{(}\AttributeTok{type =} \StringTok{"none"}\NormalTok{,    }\AttributeTok{sd =} \DecValTok{0}\NormalTok{),}
              \AttributeTok{ambient =} \FunctionTok{list}\NormalTok{(}\AttributeTok{type =} \StringTok{"ambient"}\NormalTok{, }\AttributeTok{sd =} \FloatTok{0.02}\NormalTok{),}
              \AttributeTok{outlier =} \FunctionTok{list}\NormalTok{(}\AttributeTok{type =} \StringTok{"outlier"}\NormalTok{, }\AttributeTok{sd =} \FloatTok{0.05}\NormalTok{))}
\NormalTok{theta\_lo }\OtherTok{\textless{}{-}}\NormalTok{ THETA\_GRID[}\DecValTok{3}\DataTypeTok{L}\NormalTok{]}
\NormalTok{theta\_hi }\OtherTok{\textless{}{-}}\NormalTok{ THETA\_GRID[}\DecValTok{19}\DataTypeTok{L}\NormalTok{]}
\NormalTok{noise\_at }\OtherTok{\textless{}{-}} \ControlFlowTok{function}\NormalTok{(theta) \{}
\NormalTok{  base }\OtherTok{\textless{}{-}} \FunctionTok{sample\_manifold}\NormalTok{(p, theta, N\_DEFAULT, }\AttributeTok{seed =} \DecValTok{1001}\DataTypeTok{L}\NormalTok{)}
  \FunctionTok{do.call}\NormalTok{(rbind, }\FunctionTok{lapply}\NormalTok{(}\FunctionTok{names}\NormalTok{(NOISE), }\ControlFlowTok{function}\NormalTok{(nm) \{}
\NormalTok{    z }\OtherTok{\textless{}{-}} \FunctionTok{sample\_manifold}\NormalTok{(p, theta, N\_DEFAULT, }\AttributeTok{noise =}\NormalTok{ NOISE[[nm]], }\AttributeTok{seed =} \DecValTok{1001}\DataTypeTok{L}\NormalTok{)}
\NormalTok{    d }\OtherTok{\textless{}{-}} \FunctionTok{sqrt}\NormalTok{(}\FunctionTok{rowSums}\NormalTok{((z}\SpecialCharTok{$}\NormalTok{X }\SpecialCharTok{{-}}\NormalTok{ base}\SpecialCharTok{$}\NormalTok{X)}\SpecialCharTok{\^{}}\DecValTok{2}\NormalTok{))}
    \FunctionTok{data.frame}\NormalTok{(}\AttributeTok{theta =}\NormalTok{ theta, }\AttributeTok{noise =}\NormalTok{ nm, }\AttributeTok{moved =} \FunctionTok{sum}\NormalTok{(d }\SpecialCharTok{\textgreater{}} \DecValTok{0}\NormalTok{),}
               \AttributeTok{rms =} \FunctionTok{sqrt}\NormalTok{(}\FunctionTok{mean}\NormalTok{(d}\SpecialCharTok{\^{}}\DecValTok{2}\NormalTok{)), }\AttributeTok{worst =} \FunctionTok{max}\NormalTok{(d),}
               \AttributeTok{truth\_fixed =} \FunctionTok{identical}\NormalTok{(z}\SpecialCharTok{$}\NormalTok{truth, base}\SpecialCharTok{$}\NormalTok{truth),}
               \AttributeTok{extent =} \FunctionTok{max}\NormalTok{(stats}\SpecialCharTok{::}\FunctionTok{dist}\NormalTok{(base}\SpecialCharTok{$}\NormalTok{X)))}
\NormalTok{  \}))}
\NormalTok{\}}
\NormalTok{nz }\OtherTok{\textless{}{-}} \FunctionTok{rbind}\NormalTok{(}\FunctionTok{noise\_at}\NormalTok{(theta\_lo), }\FunctionTok{noise\_at}\NormalTok{(theta\_hi))}
\FunctionTok{set.seed}\NormalTok{(}\DecValTok{11}\DataTypeTok{L}\NormalTok{); before }\OtherTok{\textless{}{-}}\NormalTok{ stats}\SpecialCharTok{::}\FunctionTok{runif}\NormalTok{(}\DecValTok{1}\NormalTok{)}
\FunctionTok{set.seed}\NormalTok{(}\DecValTok{11}\DataTypeTok{L}\NormalTok{); }\FunctionTok{invisible}\NormalTok{(}\FunctionTok{sample\_manifold}\NormalTok{(p, th, }\DecValTok{200}\DataTypeTok{L}\NormalTok{, }\AttributeTok{seed =} \DecValTok{5}\DataTypeTok{L}\NormalTok{))}
\NormalTok{after }\OtherTok{\textless{}{-}}\NormalTok{ stats}\SpecialCharTok{::}\FunctionTok{runif}\NormalTok{(}\DecValTok{1}\NormalTok{)}
\NormalTok{stream\_kept }\OtherTok{\textless{}{-}} \FunctionTok{identical}\NormalTok{(before, after)}
\end{Highlighting}
\end{Shaded}

The \texttt{ambient} row moves all 800 points; the
\texttt{outlier} row moves 40 of them, by exactly
3.97 each, which is half the chart
diameter. \texttt{truth} is byte-identical to the clean draw in every row --- the answer
key does not move when the data get worse.

\section{Results: what theta does}\label{gen-results}

\(\theta\) changes no chart coordinate. What it changes is the graph a method
builds before it does anything else.

\begin{Shaded}
\begin{Highlighting}[]
\NormalTok{k }\OtherTok{\textless{}{-}}\NormalTok{ K\_DEFAULT}
\NormalTok{graph\_at }\OtherTok{\textless{}{-}} \ControlFlowTok{function}\NormalTok{(theta, }\AttributeTok{n =}\NormalTok{ N\_DEFAULT, }\AttributeTok{seed =} \DecValTok{1001}\DataTypeTok{L}\NormalTok{, }\AttributeTok{pattern =}\NormalTok{ p) \{}
\NormalTok{  z  }\OtherTok{\textless{}{-}} \FunctionTok{sample\_manifold}\NormalTok{(pattern, theta, n, }\AttributeTok{seed =}\NormalTok{ seed)}
\NormalTok{  dA }\OtherTok{\textless{}{-}} \FunctionTok{as.matrix}\NormalTok{(stats}\SpecialCharTok{::}\FunctionTok{dist}\NormalTok{(z}\SpecialCharTok{$}\NormalTok{X));     }\FunctionTok{diag}\NormalTok{(dA) }\OtherTok{\textless{}{-}} \ConstantTok{Inf}
\NormalTok{  dU }\OtherTok{\textless{}{-}} \FunctionTok{as.matrix}\NormalTok{(stats}\SpecialCharTok{::}\FunctionTok{dist}\NormalTok{(z}\SpecialCharTok{$}\NormalTok{truth)); }\FunctionTok{diag}\NormalTok{(dU) }\OtherTok{\textless{}{-}} \ConstantTok{Inf}
\NormalTok{  nn }\OtherTok{\textless{}{-}} \FunctionTok{t}\NormalTok{(}\FunctionTok{apply}\NormalTok{(dA, }\DecValTok{1}\DataTypeTok{L}\NormalTok{, }\ControlFlowTok{function}\NormalTok{(r) }\FunctionTok{order}\NormalTok{(r)[}\FunctionTok{seq\_len}\NormalTok{(k)]))}
\NormalTok{  rk }\OtherTok{\textless{}{-}} \FunctionTok{t}\NormalTok{(}\FunctionTok{apply}\NormalTok{(dU, }\DecValTok{1}\DataTypeTok{L}\NormalTok{, rank, }\AttributeTok{ties.method =} \StringTok{"first"}\NormalTok{))}
\NormalTok{  i  }\OtherTok{\textless{}{-}} \FunctionTok{rep}\NormalTok{(}\FunctionTok{seq\_len}\NormalTok{(n), k)}
\NormalTok{  j  }\OtherTok{\textless{}{-}} \FunctionTok{as.vector}\NormalTok{(nn)}
  \FunctionTok{list}\NormalTok{(}\AttributeTok{sample =}\NormalTok{ z, }\AttributeTok{i =}\NormalTok{ i, }\AttributeTok{j =}\NormalTok{ j, }\AttributeTok{chart\_rank =}\NormalTok{ rk[}\FunctionTok{cbind}\NormalTok{(i, j)])}
\NormalTok{\}}
\NormalTok{sweep\_theta }\OtherTok{\textless{}{-}} \FunctionTok{do.call}\NormalTok{(rbind, }\FunctionTok{lapply}\NormalTok{(THETA\_GRID, }\ControlFlowTok{function}\NormalTok{(t) \{}
\NormalTok{  g }\OtherTok{\textless{}{-}} \FunctionTok{graph\_at}\NormalTok{(t)}
  \FunctionTok{data.frame}\NormalTok{(}\AttributeTok{theta =}\NormalTok{ t,}
             \AttributeTok{false\_nb   =} \FunctionTok{mean}\NormalTok{(g}\SpecialCharTok{$}\NormalTok{chart\_rank }\SpecialCharTok{\textgreater{}}\NormalTok{ k),}
             \AttributeTok{worst\_rank =} \FunctionTok{max}\NormalTok{(g}\SpecialCharTok{$}\NormalTok{chart\_rank),}
             \AttributeTok{sci        =} \FunctionTok{short\_circuit\_index}\NormalTok{(g}\SpecialCharTok{$}\NormalTok{sample),}
             \AttributeTok{gs         =} \FunctionTok{branch\_gap}\NormalTok{(g}\SpecialCharTok{$}\NormalTok{sample)}\SpecialCharTok{$}\NormalTok{ratio,}
             \AttributeTok{clearance  =} \FunctionTok{facet\_gap}\NormalTok{(p, t)}\SpecialCharTok{$}\NormalTok{gap)}
\NormalTok{\}))}
\NormalTok{top  }\OtherTok{\textless{}{-}} \FunctionTok{graph\_at}\NormalTok{(th)}
\NormalTok{U    }\OtherTok{\textless{}{-}}\NormalTok{ top}\SpecialCharTok{$}\NormalTok{sample}\SpecialCharTok{$}\NormalTok{truth}
\NormalTok{fls  }\OtherTok{\textless{}{-}}\NormalTok{ top}\SpecialCharTok{$}\NormalTok{chart\_rank }\SpecialCharTok{\textgreater{}}\NormalTok{ k}
\NormalTok{fcol }\OtherTok{\textless{}{-}}\NormalTok{ (top}\SpecialCharTok{$}\NormalTok{sample}\SpecialCharTok{$}\NormalTok{facet }\SpecialCharTok{{-}} \DecValTok{1}\DataTypeTok{L}\NormalTok{) }\SpecialCharTok{\%\%}\NormalTok{ p}\SpecialCharTok{$}\NormalTok{params}\SpecialCharTok{$}\NormalTok{nx     }\CommentTok{\# facet column and row}
\NormalTok{frow }\OtherTok{\textless{}{-}}\NormalTok{ (top}\SpecialCharTok{$}\NormalTok{sample}\SpecialCharTok{$}\NormalTok{facet }\SpecialCharTok{{-}} \DecValTok{1}\DataTypeTok{L}\NormalTok{) }\SpecialCharTok{\%/\%}\NormalTok{ p}\SpecialCharTok{$}\NormalTok{params}\SpecialCharTok{$}\NormalTok{nx}
\NormalTok{wider }\OtherTok{\textless{}{-}} \FunctionTok{abs}\NormalTok{(U[top}\SpecialCharTok{$}\NormalTok{i, }\DecValTok{1}\DataTypeTok{L}\NormalTok{] }\SpecialCharTok{{-}}\NormalTok{ U[top}\SpecialCharTok{$}\NormalTok{j, }\DecValTok{1}\DataTypeTok{L}\NormalTok{]) }\SpecialCharTok{\textgreater{}} \FunctionTok{abs}\NormalTok{(U[top}\SpecialCharTok{$}\NormalTok{i, }\DecValTok{2}\DataTypeTok{L}\NormalTok{] }\SpecialCharTok{{-}}\NormalTok{ U[top}\SpecialCharTok{$}\NormalTok{j, }\DecValTok{2}\DataTypeTok{L}\NormalTok{])}
\NormalTok{mistakes }\OtherTok{\textless{}{-}} \FunctionTok{c}\NormalTok{(}
  \AttributeTok{n\_false    =} \FunctionTok{sum}\NormalTok{(fls),}
  \AttributeTok{n\_same     =} \FunctionTok{sum}\NormalTok{((top}\SpecialCharTok{$}\NormalTok{sample}\SpecialCharTok{$}\NormalTok{facet[top}\SpecialCharTok{$}\NormalTok{i] }\SpecialCharTok{==}\NormalTok{ top}\SpecialCharTok{$}\NormalTok{sample}\SpecialCharTok{$}\NormalTok{facet[top}\SpecialCharTok{$}\NormalTok{j])[fls]),}
  \AttributeTok{n\_one\_cell =} \FunctionTok{sum}\NormalTok{((}\FunctionTok{abs}\NormalTok{(fcol[top}\SpecialCharTok{$}\NormalTok{i] }\SpecialCharTok{{-}}\NormalTok{ fcol[top}\SpecialCharTok{$}\NormalTok{j]) }\SpecialCharTok{\textless{}=} \DecValTok{1}\DataTypeTok{L} \SpecialCharTok{\&}
                    \FunctionTok{abs}\NormalTok{(frow[top}\SpecialCharTok{$}\NormalTok{i] }\SpecialCharTok{{-}}\NormalTok{ frow[top}\SpecialCharTok{$}\NormalTok{j]) }\SpecialCharTok{\textless{}=} \DecValTok{1}\DataTypeTok{L}\NormalTok{)[fls]),}
  \AttributeTok{across\_f   =} \FunctionTok{round}\NormalTok{(}\DecValTok{100} \SpecialCharTok{*} \FunctionTok{mean}\NormalTok{(wider[fls])),}
  \AttributeTok{across\_h   =} \FunctionTok{round}\NormalTok{(}\DecValTok{100} \SpecialCharTok{*} \FunctionTok{mean}\NormalTok{(wider[}\SpecialCharTok{!}\NormalTok{fls])))}
\end{Highlighting}
\end{Shaded}

Four quantities, all computed on the same samples, all monotone, and none of
them linear.

\textbf{The graph acquires false neighbours.} An ambient \(k\)-nearest-neighbour edge
is \emph{false} when its far endpoint is not among the point's \(k\) nearest neighbours
in the chart. At \(\theta = 0\) the ambient and chart graphs are identical and the
false fraction is 0. At \(\theta = 0.95\) it is
0.302: about three edges in ten of the graph a
method consumes join points that are not neighbours on the paper. The worst
single edge joins a point to its 168th chart
neighbour.

\textbf{Edges contract.} \texttt{short\_circuit\_index()} is the median ambient
nearest-neighbour distance over the median true intrinsic distance between the
same pairs, so it is one on a faithful graph. It falls from
1 to 0.787.

\textbf{Distant sheets approach, but not far enough.} Balasubramanian and Schwartz's
ratio \(g/s\) --- the closest approach between surface-distant sampled points, in
units of the median nearest-neighbour spacing --- is the quantity that predicts
Isomap's topological instability \citep{balasubramanian2002isomap}. It falls from
20.4 to 5.84, which
is a large fall and still leaves the sheets nearly six sample spacings apart.

The false-neighbour count above is measured at \texttt{K\_DEFAULT} against the \emph{chart},
and it is worth saying so because the book uses two other pairings and they are
easy to conflate. \(Q_{NX}\) is scored at \texttt{K\_QNX} against the chart, because that
is the answer key and the co-ranking measure is noisy at small \(K\);
trustworthiness, continuity and kNN preservation are scored at \texttt{K\_DEFAULT}
against \emph{ambient} distance, because they ask whether an embedding preserved the
neighbourhoods a method could actually see. Three combinations of neighbourhood
size and reference geometry, each answering a different question, all named in
\texttt{R/constants.R} so that they cannot drift.

\begin{figure}

{\centering \includegraphics[width=0.9\linewidth,alt={Four line charts in a two by two arrangement, each against the folding parameter from zero to 0.95. Top left, the false-neighbour fraction rises from zero, slowly at first and then sharply, to about 0.3. Top right, the short-circuit index falls from one to about 0.79. Bottom left, the branch-gap ratio falls from about 20 to about 6. Bottom right, the minimum facet clearance falls from about 0.87 to about 0.4. Every curve bends most in its rightmost quarter.}]{03-generative-model_files/figure-latex/gen-fig-difficulty-1} 

}

\caption{Four difficulty measures over the benchmark grid's $\theta$ sweep, computed on one seed at $n = 800$. All four are monotone and none is linear: the last quarter of the axis does most of the work.}\label{fig:gen-fig-difficulty}
\end{figure}

Figure \ref{fig:gen-fig-difficulty} is why standing rule 2 exists. A single
number reported at one \(\theta\) can be read off any point of those curves, and
the curves are shallow over the first half of the axis and steep over the last
quarter. Averaging over \(\theta\) would compress the only axis the benchmark has;
Chapter \ref{results} reports a case where two method families change places
partway along it, which no summary at a fixed geometry can show.

Where the false edges sit is as informative as how many there are.

\begin{figure}

{\centering \includegraphics[width=0.9\linewidth,alt={Two panels showing the flat Miura chart with short line segments drawn over it. In panel A the segments form a dense, even mesh filling the facets, with visible gaps running along the zigzag crease lines. In panel B the segments are longer and lie almost entirely in narrow zigzag bands that follow those same crease lines, leaving the interiors of the facets empty.}]{03-generative-model_files/figure-latex/gen-fig-graph-1} 

}

\caption{The ambient $k$-nearest-neighbour graph at the grid's largest fold, drawn on the chart it should have respected. Panel A holds the edges that are also chart neighbours; panel B holds the false ones. The false edges are not scattered over the sheet — they trace the crease lines, and every one of them straddles a fold.}\label{fig:gen-fig-graph}
\end{figure}

Figure \ref{fig:gen-fig-graph} draws both halves of that graph on the chart, and
the two halves occupy different parts of it. Honest edges fill the facets and
thin out along the creases; false edges do the reverse and lie in narrow bands on
the crease lines. Of the 2416 false edges,
0 have both endpoints in the same facet --- inside a facet,
ambient distance \emph{is} chart distance --- so every false neighbour straddles a
fold. The mistakes are also strictly local: all
2416 of them join facets no more than one cell apart in
either direction. There is a mild directional tilt on top of that,
62\% of false edges being longer across the corrugation
than along it against 49\% of honest ones, but the
localisation is the dominant fact and the direction is not.

That localisation is why difficulty here is not simply ``more noise''. Additive
noise degrades every neighbourhood alike; folding leaves the interior of every
facet exact and corrupts a codimension-one set of neighbourhoods completely.
That is a structured failure, and whether a method is robust to it is a
different question from whether it tolerates noise --- which is why the benchmark
grid crosses the two axes instead of collapsing them.

\section{Diagnostics}\label{gen-diagnostics}

\textbf{\(\theta\) is not a dihedral angle.} It is the fraction of the way to
flat-folded, on \([0, 1]\), and \texttt{fold()} refuses anything outside that. The
dihedral angles are derived and returned in \texttt{rho}; at \(\theta = 0.95\) the
creases of \texttt{p} span 0.504 to
3.14 radians, the upper end being the boundary
edges, which do not fold at all. A flat sheet has \(\rho = \pi\), not zero.
Conflating the two is an error this project made once and corrected in three
places.

\textbf{\texttt{noise\$sd} is a fraction of the chart diameter.} Not of the folded object.
The \texttt{ambient} row displaces points by an RMS of
0.272 at \(\theta = 0.1\) and
0.272 at \(\theta = 0.9\) --- the
same number --- while the folded cloud's own extent falls from
6.19 to
3.93. Had the scale been keyed to
the folded object, the noise axis and the fold axis would move together and no
result could be attributed to either.

\textbf{In the \texttt{outlier} row, \texttt{sd} is a contamination fraction.} The grid passes
\texttt{sd\ =\ 0.05} there and means five per cent of the points, not a standard
deviation; the displacement magnitude is fixed, not random. This is a genuine
trap in the interface and it is documented in \texttt{R/sampling.R} rather than fixed,
because the artefacts were computed under it.

\textbf{A seed reproduces a draw without costing the caller their stream.}
\texttt{sample\_manifold()} saves \texttt{.Random.seed}, replaces it, and restores it on exit,
so drawing a sample inside a loop does not silently re-align the loop's own
randomness. Two draws from one stream, with a sample taken between them, agree:
TRUE. Called without a seed it touches \texttt{set.seed()}
not at all and records \texttt{NA}, so a returned object never claims a seed that would
not reproduce it.

\section{Where it fails}\label{gen-limits}

\textbf{The classical short circuit lives past the end of the grid.} \texttt{gen-results}
found no edge joining facets more than one cell apart anywhere on \texttt{THETA\_GRID}.
That is not because folding cannot produce one.

\begin{Shaded}
\begin{Highlighting}[]
\NormalTok{cells\_apart }\OtherTok{\textless{}{-}} \ControlFlowTok{function}\NormalTok{(g, pattern) \{}
\NormalTok{  cl }\OtherTok{\textless{}{-}}\NormalTok{ (g}\SpecialCharTok{$}\NormalTok{sample}\SpecialCharTok{$}\NormalTok{facet }\SpecialCharTok{{-}} \DecValTok{1}\DataTypeTok{L}\NormalTok{) }\SpecialCharTok{\%\%}\NormalTok{  pattern}\SpecialCharTok{$}\NormalTok{params}\SpecialCharTok{$}\NormalTok{nx}
\NormalTok{  rw }\OtherTok{\textless{}{-}}\NormalTok{ (g}\SpecialCharTok{$}\NormalTok{sample}\SpecialCharTok{$}\NormalTok{facet }\SpecialCharTok{{-}} \DecValTok{1}\DataTypeTok{L}\NormalTok{) }\SpecialCharTok{\%/\%}\NormalTok{ pattern}\SpecialCharTok{$}\NormalTok{params}\SpecialCharTok{$}\NormalTok{nx}
  \FunctionTok{pmax}\NormalTok{(}\FunctionTok{abs}\NormalTok{(cl[g}\SpecialCharTok{$}\NormalTok{i] }\SpecialCharTok{{-}}\NormalTok{ cl[g}\SpecialCharTok{$}\NormalTok{j]), }\FunctionTok{abs}\NormalTok{(rw[g}\SpecialCharTok{$}\NormalTok{i] }\SpecialCharTok{{-}}\NormalTok{ rw[g}\SpecialCharTok{$}\NormalTok{j]))}
\NormalTok{\}}
\NormalTok{beyond }\OtherTok{\textless{}{-}} \FunctionTok{do.call}\NormalTok{(rbind, }\FunctionTok{lapply}\NormalTok{(}\FunctionTok{c}\NormalTok{(}\FunctionTok{max}\NormalTok{(THETA\_GRID), }\FloatTok{0.99}\NormalTok{, }\FloatTok{0.999}\NormalTok{), }\ControlFlowTok{function}\NormalTok{(t) \{}
\NormalTok{  g }\OtherTok{\textless{}{-}} \FunctionTok{graph\_at}\NormalTok{(t)}
  \FunctionTok{data.frame}\NormalTok{(}\AttributeTok{theta =}\NormalTok{ t,}
             \AttributeTok{gs    =} \FunctionTok{branch\_gap}\NormalTok{(g}\SpecialCharTok{$}\NormalTok{sample)}\SpecialCharTok{$}\NormalTok{ratio,}
             \AttributeTok{sci   =} \FunctionTok{short\_circuit\_index}\NormalTok{(g}\SpecialCharTok{$}\NormalTok{sample),}
             \AttributeTok{jumps =} \FunctionTok{mean}\NormalTok{(}\FunctionTok{cells\_apart}\NormalTok{(g, p) }\SpecialCharTok{\textgreater{}=} \DecValTok{2}\DataTypeTok{L}\NormalTok{))}
\NormalTok{\}))}
\end{Highlighting}
\end{Shaded}

Past the grid the ratio \(g/s\) collapses --- 5.84 at
\(\theta = 0.95\), 1.37 at
\(\theta = 0.999\) --- and edges joining facets two or more cells
apart appear: 0 of the graph at the grid's edge,
0.174 at the last value. The grid stops at
0.95 because past it the sheet is collapsing onto itself and the cell measures a degenerate object rather than
a harder instance of the same one. The consequence is worth stating plainly:
\textbf{this benchmark's difficulty axis does not reach the regime that classically
breaks Isomap.} What it measures instead is the graded corruption of the
neighbourhood graph that precedes it.

\textbf{Area weighting is correct and untested by this family.} Every Miura facet is
congruent, so area-proportional and facet-uniform sampling agree to
1.2e-15 and no result in this
book distinguishes them. The weighting is there because a family with unequal facets would need it, and this
book ships no such family --- see Chapter \ref{folding-geometry} on the two that
were withdrawn.

\textbf{Trimming the boundary changes the region, not just the count.} With
\texttt{boundary\ =\ FALSE} the facets nearest the sheet edge keep a small share of their
points: across the 36 facets of \texttt{p} the occupancy of \texttt{s} runs
from 3 to
40. The survivors are uniform
over the region that survives, which is the guarantee that was wanted, but that
region is not the whole sheet.

\textbf{Three scalars are a cell, not a family of patterns.} Varying \(a\), \(b\) and
\(\alpha\) changes the cell shape and never the combinatorics: the same
49 vertices and 36 facets, connected the
same way. This model produces a one-parameter deformation of one tessellation,
so nothing here supports a claim about pattern variety.

\textbf{The chart is the geodesic only inside a convex region.} The caveat is
Chapter \ref{folding-geometry}'s and it applies to every sample drawn here;
\texttt{boundary\ =\ FALSE} is what keeps it negligible, and \texttt{chart\_exit\_fraction()} is
what measures it rather than assuming it.

\section{Reproduce this}\label{gen-reproduce}

Every claim above is an assertion in the suite, which runs on every push:

\begin{itemize}
\tightlist
\item
  \texttt{tests/testthat/test-sampling.R} --- area-uniformity against an independent
  shoelace computation, points lying in the facet they claim and on the folded
  surface, seed reproducibility, and the RNG-stream restore.
\item
  \texttt{tests/testthat/test-folding.R} --- facet isometry over the \(\theta\) sweep and
  the \(\alpha\) range, the derived dihedral contract, and the absence of
  self-intersection.
\item
  \texttt{tests/testthat/test-patterns.R} --- congruence, Kawasaki and Maekawa.
\end{itemize}

Run them with \texttt{Rscript\ tests/testthat.R} from the repository root. Every figure
and every number in this chapter is recomputed at render time from \texttt{R/} alone,
so a change to the sampler or the folding map shows up here as a changed page
rather than as a stale one.

The model's settings are not scattered: \texttt{THETA\_GRID}, \texttt{N\_DEFAULT},
\texttt{PATTERN\_GRID}, \texttt{MIURA\_ALPHA\_RANGE}, \texttt{K\_DEFAULT} and \texttt{K\_QNX} are all in
\texttt{R/constants.R}, and \texttt{scripts/run-benchmark-grid.R} is the producer that turns
them into \texttt{data/processed/benchmark-grid.rds}. Chapter \ref{linear} opens by
reading that artefact with its digest, and Appendix \ref{datasets} is its
codebook.

\section{Further reading}\label{gen-reading}

Miura's original report \citep{miura1985method} introduces the tessellation; Schenk
and Guest \citep{schenk2013geometry} treat the folded sheet as a mechanical
metamaterial and are the natural next step for a reader who wants the unit-cell
geometry from the engineering side. Hull's \emph{Origametry} \citep{hull2020origametry} is
the reference for flat-foldability and for why a pattern folds at all.

On the benchmarking side, Espadoto and colleagues \citep{espadoto2021survey} survey
how dimension reduction is usually evaluated and make the case for scale that
this book answers with a difficulty axis instead; Machado and colleagues
\citep{machado2025necessary} argue that projection quality metrics are necessary and
not sufficient, which is the question Chapter \ref{evaluation} takes up with an
exact answer key in hand. Lee and Verleysen \citep{lee2007nonlinear} remains the
textbook treatment of the methods Part II tests against this model.

\part{Methods under test}\label{part-methods-under-test}

\chapter{Linear projections and where they break}\label{linear}

\emph{R. Heller · Date: 2026-07-29 · Last changed: 2026-09-07}

\emph{When is a projection just a shadow?}

\section{What this chapter answers}\label{lin-question}

A linear method asks one question: which flat slice of the ambient space keeps
most of the cloud? On a folded sheet the answer is a shadow. The sheet is
intrinsically two-dimensional, so an exact two-dimensional answer exists ---
Chapter \ref{folding-geometry} writes it down in closed form --- but folding puts
it out of a projection's reach, because unfolding is not a linear map. This
chapter measures how far out of reach, and establishes three things the rest of
the book runs on.

\textbf{The reading convention.} Every chapter from here on opens by reading a
committed artefact together with the digest it was written against --- that is
\texttt{read\_run()}, and this is the chapter that explains what it refuses.

\textbf{The baseline.} PCA's reconstruction-error curve over \(\theta\) is the reference
every later method is judged against. It is not a straw man: PCA is exact on the
flat sheet, and on this grid it is the best method of all until the fold is deep
enough for Isomap to overtake it --- which happens later than the ordering in
Chapter \ref{results} might suggest, and later still on \(Q_{NX}\) than on
reconstruction error. Chapter \ref{geodesic} locates both crossings.

\textbf{The coincidence.} PCA and classical MDS are not two methods. The registry
keeps both because showing that they coincide is worth a figure, and Chapter
\ref{results} must not report them as two results.

None of this argues against linear projection. Linear methods are the honest
baseline here precisely because they cannot represent a fold and say so plainly,
rather than returning a picture with structure the data do not have.

\section{Setup}\label{lin-setup}

\begin{Shaded}
\begin{Highlighting}[]
\NormalTok{grid }\OtherTok{\textless{}{-}} \FunctionTok{read\_run}\NormalTok{(}\StringTok{"data/processed/benchmark{-}grid.rds"}\NormalTok{, }\StringTok{"2483a468d82d4d64"}\NormalTok{)}
\NormalTok{lin  }\OtherTok{\textless{}{-}} \FunctionTok{subset}\NormalTok{(grid, method }\SpecialCharTok{\%in\%} \FunctionTok{c}\NormalTok{(}\StringTok{"pca"}\NormalTok{, }\StringTok{"mds"}\NormalTok{))}
\NormalTok{p    }\OtherTok{\textless{}{-}} \FunctionTok{miura\_ori}\NormalTok{(}\AttributeTok{nx =} \DecValTok{6}\DataTypeTok{L}\NormalTok{, }\AttributeTok{ny =} \DecValTok{6}\DataTypeTok{L}\NormalTok{, }\AttributeTok{alpha =}\NormalTok{ pi }\SpecialCharTok{/} \DecValTok{3}\NormalTok{)}
\end{Highlighting}
\end{Shaded}

\texttt{read\_run()} reads a committed \texttt{.rds} and refuses three things. An artefact with
\textbf{no provenance} cannot be traced to the code that produced it, so no number may
be reported from it. One written by a \textbf{\texttt{-\/-quick} run} is a smoke test at
\texttt{N\_QUICK} points, not the book's evidence. And one whose \textbf{digest} is not the
digest named here is not the run these sentences were read off --- the chunks would
re-run, the figures would redraw, and the prose around them would quietly keep
the old values.

The digest is taken with the provenance block stripped. Provenance carries a
commit, a timestamp and an elapsed time, so hashing it would fire the check on
every regeneration, and a check that always fires is one that gets deleted. What
is hashed is the answer.

\begin{Shaded}
\begin{Highlighting}[]
\CommentTok{\# The third refusal. The message hands over the digest it found, so recording a}
\CommentTok{\# new one is easier than not recording it.}
\FunctionTok{read\_run}\NormalTok{(}\StringTok{"data/processed/benchmark{-}grid.rds"}\NormalTok{, }\StringTok{"0000000000000000"}\NormalTok{)}
\CommentTok{\#\textgreater{} Error: data/processed/benchmark{-}grid.rds is not the run this chapter was}
\CommentTok{\#\textgreater{}   written against.}
\CommentTok{\#\textgreater{}   expected digest: 0000000000000000}
\CommentTok{\#\textgreater{}   found:           2483a468d82d4d64}
\end{Highlighting}
\end{Shaded}

\section{Background}\label{lin-background}

\subsection{The projection}\label{the-projection}

PCA is the least-squares fit of a \(d\)-dimensional affine subspace to a cloud of
points \citep{pearson1901lines}, and equivalently the rotation onto the directions of
greatest variance \citep{hotelling1933components}. Both descriptions are one singular
value decomposition, and Eckart--Young \citep{eckart1936approximation} says the
truncated decomposition is \emph{optimal} among approximations of that rank rather
than merely good. That matters twice below: it is why PCA is the strongest thing
a linear method can be, and it is what \texttt{metric\_floor()} is built from.

\texttt{embed\_pca()} computes it with \texttt{prcomp()}, an SVD of the centred data matrix. The
plan this book was written from preferred that route to \texttt{princomp()}'s
eigendecomposition of the covariance on accuracy grounds. The argument does not
reproduce here and is not made: \texttt{lin-core} compares the two and they agree to the
last representable place. \texttt{prcomp()} is a convention, not a result.

\subsection{The distance version}\label{the-distance-version}

Classical multidimensional scaling \citep{torgerson1952mds} starts from a matrix of
distances, double-centres their squares, and takes the leading eigenvectors. It
is not Kruskal's multidimensional scaling \citep{kruskal1964fit}, which minimises a
stress functional numerically; the classical version is a closed-form
eigenproblem, and \texttt{embed\_mds()} is \texttt{cmdscale()} on the ambient distances.

For Euclidean distances the double-centred matrix of squares \emph{is} the Gram matrix
of the centred coordinates, whose eigenvectors are the principal directions. So
classical MDS on Euclidean distances is PCA --- a textbook identity
\citep{lee2007nonlinear}. What the next section adds is a measurement of it: ``provably
identical'' and ``identical in this implementation, at this tolerance, on this
data'' are different claims, and only the second can be checked.

\subsection{What is being scored}\label{what-is-being-scored}

\texttt{reconstruction\_error()} is the normalised full Procrustes RMSE. It quotients out
translation, rotation, reflection and one isotropic scale \citep{gower1975procrustes}
and nothing else, so zero means the embedding is the chart up to a similarity and
one means it explains none of the chart's variance at all.

\texttt{metric\_floor()} says when zero is even available: it is the tail of the chart's
own eigenvalue spectrum, the Eckart--Young residual of forcing a chart of
intrinsic dimension \(p\) into \(d < p\). Here chart and target are both
two-dimensional, so the \texttt{floor} column of this grid is zero in every one of its
rows. Nothing below is loss the problem imposes; all of it belongs to the method.

\section{A shadow, and two names for it}\label{lin-core}

\subsection{Exact on the flat sheet}\label{exact-on-the-flat-sheet}

\begin{Shaded}
\begin{Highlighting}[]
\NormalTok{thetas }\OtherTok{\textless{}{-}}\NormalTok{ THETA\_COARSE[}\FunctionTok{c}\NormalTok{(}\DecValTok{1}\DataTypeTok{L}\NormalTok{, }\DecValTok{3}\DataTypeTok{L}\NormalTok{, }\DecValTok{5}\DataTypeTok{L}\NormalTok{)]}
\NormalTok{samp   }\OtherTok{\textless{}{-}} \FunctionTok{lapply}\NormalTok{(thetas, }\ControlFlowTok{function}\NormalTok{(t)}
  \FunctionTok{sample\_manifold}\NormalTok{(p, }\AttributeTok{theta =}\NormalTok{ t, }\AttributeTok{n =}\NormalTok{ N\_DEFAULT, }\AttributeTok{seed =}\NormalTok{ BENCH\_SEEDS[}\DecValTok{1}\NormalTok{]))}
\NormalTok{Y }\OtherTok{\textless{}{-}} \FunctionTok{lapply}\NormalTok{(samp, }\ControlFlowTok{function}\NormalTok{(s) }\FunctionTok{embed}\NormalTok{(}\StringTok{"pca"}\NormalTok{, s))}
\end{Highlighting}
\end{Shaded}

\texttt{embed()} takes the sample object, not a matrix. A method may consume the ambient
cloud, a neighbourhood graph built from it, or geodesics, and the registry
records which; handing it a bare matrix would throw away the distinction the
whole benchmark is about.

\begin{Shaded}
\begin{Highlighting}[]
\NormalTok{err0       }\OtherTok{\textless{}{-}} \FunctionTok{reconstruction\_error}\NormalTok{(Y[[}\DecValTok{1}\NormalTok{]], samp[[}\DecValTok{1}\NormalTok{]]}\SpecialCharTok{$}\NormalTok{truth)}
\NormalTok{floor0     }\OtherTok{\textless{}{-}} \FunctionTok{metric\_floor}\NormalTok{(samp[[}\DecValTok{1}\NormalTok{]])}
\NormalTok{flat\_cells }\OtherTok{\textless{}{-}} \FunctionTok{subset}\NormalTok{(grid, method }\SpecialCharTok{==} \StringTok{"pca"} \SpecialCharTok{\&}\NormalTok{ theta }\SpecialCharTok{==} \DecValTok{0} \SpecialCharTok{\&}\NormalTok{ noise }\SpecialCharTok{==} \StringTok{"none"}\NormalTok{)}

\NormalTok{sv     }\OtherTok{\textless{}{-}} \FunctionTok{vapply}\NormalTok{(samp, }\ControlFlowTok{function}\NormalTok{(s) }\FunctionTok{prcomp}\NormalTok{(s}\SpecialCharTok{$}\NormalTok{X)}\SpecialCharTok{$}\NormalTok{sdev, }\FunctionTok{numeric}\NormalTok{(}\DecValTok{3}\NormalTok{))}
\NormalTok{thick  }\OtherTok{\textless{}{-}}\NormalTok{ sv[}\DecValTok{3}\NormalTok{, ] }\SpecialCharTok{/}\NormalTok{ sv[}\DecValTok{1}\NormalTok{, ]          }\CommentTok{\# the corrugation, relative to the sheet}
\NormalTok{aspect }\OtherTok{\textless{}{-}}\NormalTok{ sv[}\DecValTok{2}\NormalTok{, ] }\SpecialCharTok{/}\NormalTok{ sv[}\DecValTok{1}\NormalTok{, ]          }\CommentTok{\# of the two directions PCA keeps}
\NormalTok{u      }\OtherTok{\textless{}{-}} \FunctionTok{svd}\NormalTok{(}\FunctionTok{scale}\NormalTok{(samp[[}\DecValTok{1}\NormalTok{]]}\SpecialCharTok{$}\NormalTok{truth, }\AttributeTok{scale =} \ConstantTok{FALSE}\NormalTok{))}\SpecialCharTok{$}\NormalTok{d}
\NormalTok{chart\_aspect }\OtherTok{\textless{}{-}}\NormalTok{ u[}\DecValTok{2}\NormalTok{] }\SpecialCharTok{/}\NormalTok{ u[}\DecValTok{1}\NormalTok{]          }\CommentTok{\# of the chart, which folding never touches}
\end{Highlighting}
\end{Shaded}

At \(\theta = 0\) the sheet lies in a plane, PCA finds that plane, and the recovery
is exact. On this sample the error is
3.8e-16 against a floor of 0,
and over the whole clean flat row of the committed grid --- 20
seeds --- the worst is
1.4e-15, with
\(Q_{NX}\) exactly 1. That is not a good score. It is
floating-point rounding on a problem whose answer is a rigid motion.

It is worth pausing on how unusual that is. At \(\theta = 0\) PCA and classical MDS
are the only methods on this grid that recover the chart exactly; every method in
Chapters \ref{geodesic} and \ref{neighbor} posts a visible error on a sheet
that has not been folded at all.

\begin{figure}

{\centering \includegraphics[width=0.9\linewidth,alt={Three panels side by side, each a scatter of embedded points coloured along a purple-to-yellow ramp by one exact chart coordinate, with five thin dark level-set paths drawn over them. In the leftmost, at fold angle zero, the colour runs in straight vertical bands and the five paths are straight, evenly spaced and parallel. The middle panel at half the sweep looks much the same. In the rightmost, at the largest fold angle, the whole patch is smaller in both directions and the paths have bowed outwards into curves; they remain ordered and never cross.}]{04-linear-projections_files/figure-latex/lin-shadow-1} 

}

\caption{The two-component PCA of the same sheet at three fold angles, coloured by the exact chart coordinate along the corrugation. The grey paths are level sets of that coordinate: an embedding that recovered the chart draws them nested and non-crossing, and one that folded the chart back onto itself draws two of them on top of each other.}\label{fig:lin-shadow}
\end{figure}

Figure \ref{fig:lin-shadow} is the chapter's title in three panels. The left
panel is the chart, recovered. The right panel is a shadow: the sheet has
contracted along its corrugation and the projection has contracted with it,
because a projection is the one thing that cannot put back a direction the data
no longer extend into. Read the grey level sets rather than the colour. They stay
ordered and never cross --- PCA does not tear the chart or fold it back on itself ---
but they bow, and a bowed chart is a distortion no similarity transform can undo.

\subsection{What the projection discards}\label{what-the-projection-discards}

Folding lifts the sheet out of its plane, and the projection's difficulty is that
the lift is real. The third principal direction of the ambient cloud \emph{is} the
corrugation, and its spread relative to the leading direction rises from
0 on the flat sheet to 0.115 and
0.295 at the two folded angles of Figure
\ref{fig:lin-shadow}. Two components cannot keep it, and by Eckart--Young no
other pair of directions would keep more of the cloud, so the residual is not a
choice anybody made.

Discarding it would be harmless if what remained were the chart. It is not, and
that is the mechanism the reconstruction error picks up. The sheet does not
contract equally in its two chart directions: at these three fold angles the pair
PCA keeps has aspect ratio 0.861, 0.931 and
0.691, moving first up and then down, while the chart's own
aspect ratio is 0.861 at every one of them because folding
never touches the chart. What the projection returns is therefore an
\emph{anisotropically} rescaled chart --- and full Procrustes admits translation,
rotation, reflection and exactly one \textbf{isotropic} scale. Anisotropy is precisely
what the alignment cannot absorb, which is why the score climbs even though the
level sets in Figure \ref{fig:lin-shadow} never cross.

\subsection{One computation, four names}\label{one-computation-four-names}

\begin{Shaded}
\begin{Highlighting}[]
\NormalTok{w }\OtherTok{\textless{}{-}} \FunctionTok{reshape}\NormalTok{(lin[, }\FunctionTok{c}\NormalTok{(}\StringTok{"theta"}\NormalTok{, }\StringTok{"noise"}\NormalTok{, }\StringTok{"seed"}\NormalTok{, }\StringTok{"method"}\NormalTok{, }\StringTok{"rmse"}\NormalTok{)],}
             \AttributeTok{idvar =} \FunctionTok{c}\NormalTok{(}\StringTok{"theta"}\NormalTok{, }\StringTok{"noise"}\NormalTok{, }\StringTok{"seed"}\NormalTok{), }\AttributeTok{timevar =} \StringTok{"method"}\NormalTok{,}
             \AttributeTok{direction =} \StringTok{"wide"}\NormalTok{)}
\NormalTok{d\_eps  }\OtherTok{\textless{}{-}} \FunctionTok{abs}\NormalTok{(w}\SpecialCharTok{$}\NormalTok{rmse.pca }\SpecialCharTok{{-}}\NormalTok{ w}\SpecialCharTok{$}\NormalTok{rmse.mds) }\SpecialCharTok{/}\NormalTok{ .Machine}\SpecialCharTok{$}\NormalTok{double.eps}
\NormalTok{rank\_d }\OtherTok{\textless{}{-}} \FunctionTok{max}\NormalTok{(}\FunctionTok{abs}\NormalTok{(lin}\SpecialCharTok{$}\NormalTok{qnx[lin}\SpecialCharTok{$}\NormalTok{method }\SpecialCharTok{==} \StringTok{"pca"}\NormalTok{] }\SpecialCharTok{{-}}\NormalTok{ lin}\SpecialCharTok{$}\NormalTok{qnx[lin}\SpecialCharTok{$}\NormalTok{method }\SpecialCharTok{==} \StringTok{"mds"}\NormalTok{]))}

\NormalTok{mid    }\OtherTok{\textless{}{-}}\NormalTok{ samp[[}\DecValTok{2}\NormalTok{]]}
\NormalTok{d\_cmds }\OtherTok{\textless{}{-}} \FunctionTok{reconstruction\_error}\NormalTok{(}\FunctionTok{embed}\NormalTok{(}\StringTok{"mds"}\NormalTok{, mid), Y[[}\DecValTok{2}\NormalTok{]])}
\NormalTok{d\_prin }\OtherTok{\textless{}{-}} \FunctionTok{reconstruction\_error}\NormalTok{(}\FunctionTok{princomp}\NormalTok{(mid}\SpecialCharTok{$}\NormalTok{X)}\SpecialCharTok{$}\NormalTok{scores[, }\DecValTok{1}\SpecialCharTok{:}\DecValTok{2}\NormalTok{], Y[[}\DecValTok{2}\NormalTok{]])}
\NormalTok{d\_eig  }\OtherTok{\textless{}{-}} \FunctionTok{reconstruction\_error}\NormalTok{(}
  \FunctionTok{scale}\NormalTok{(mid}\SpecialCharTok{$}\NormalTok{X, }\AttributeTok{scale =} \ConstantTok{FALSE}\NormalTok{) }\SpecialCharTok{\%*\%} \FunctionTok{eigen}\NormalTok{(}\FunctionTok{cov}\NormalTok{(mid}\SpecialCharTok{$}\NormalTok{X))}\SpecialCharTok{$}\NormalTok{vectors[, }\DecValTok{1}\SpecialCharTok{:}\DecValTok{2}\NormalTok{], Y[[}\DecValTok{2}\NormalTok{]])}
\end{Highlighting}
\end{Shaded}

Across all 1200 cells of the grid, PCA and classical MDS differ in
reconstruction error by at most 12.4 units of machine
epsilon, and 185 of those cells are bit-identical. On \(Q_{NX}\)
the difference is 0 exactly --- the two methods return configurations
whose rank orders agree perfectly.

The same holds for the routes into PCA itself. At \(\theta\) of 0.5,
\texttt{cmdscale()}, \texttt{princomp()} and an eigendecomposition of the covariance matrix all
land on \texttt{prcomp()}'s answer at Procrustes distances of
9.4e-16,
2.6e-16 and
1.5e-16. Four names, one computation.
The accuracy argument for \texttt{prcomp()} has nowhere to show up on data of this
conditioning.

\begin{figure}

{\centering \includegraphics[width=0.9\linewidth,alt={Two panels. Panel A is a scatter of an embedded point cloud in which every filled grey dot sits exactly inside an open orange circle, so the two configurations are indistinguishable. Panel B is a histogram whose horizontal axis runs from zero to about thirteen units of machine epsilon; the bars are tallest at the left, near zero, and fall away steadily to the right.}]{04-linear-projections_files/figure-latex/lin-same-1} 

}

\caption{PCA and classical MDS on the same ambient cloud. Panel A overplots one cell: filled points are the PCA embedding, open circles the Procrustes-aligned MDS embedding, and every circle is centred on a point. Panel B is the difference in reconstruction error over every cell of the committed grid, in units of machine epsilon.}\label{fig:lin-same}
\end{figure}

Figure \ref{fig:lin-same} is the reason both entries survive in the registry and
the reason neither may be counted twice. Chapter \ref{results} reports one
linear result, not two, and \texttt{R/methods.R} says so at the definition of
\texttt{embed\_mds()} rather than leaving it to be remembered.

\section{Results}\label{lin-results}

\begin{Shaded}
\begin{Highlighting}[]
\NormalTok{pca   }\OtherTok{\textless{}{-}} \FunctionTok{subset}\NormalTok{(grid, method }\SpecialCharTok{==} \StringTok{"pca"}\NormalTok{)}
\NormalTok{curve }\OtherTok{\textless{}{-}} \FunctionTok{do.call}\NormalTok{(rbind, }\FunctionTok{lapply}\NormalTok{(}\FunctionTok{split}\NormalTok{(pca, }\FunctionTok{list}\NormalTok{(pca}\SpecialCharTok{$}\NormalTok{noise, pca}\SpecialCharTok{$}\NormalTok{theta)),}
  \ControlFlowTok{function}\NormalTok{(s) }\FunctionTok{data.frame}\NormalTok{(}\AttributeTok{noise =}\NormalTok{ s}\SpecialCharTok{$}\NormalTok{noise[}\DecValTok{1}\NormalTok{], }\AttributeTok{theta =}\NormalTok{ s}\SpecialCharTok{$}\NormalTok{theta[}\DecValTok{1}\NormalTok{],}
                         \AttributeTok{mean =} \FunctionTok{mean}\NormalTok{(s}\SpecialCharTok{$}\NormalTok{rmse), }\AttributeTok{lo =} \FunctionTok{min}\NormalTok{(s}\SpecialCharTok{$}\NormalTok{rmse),}
                         \AttributeTok{hi =} \FunctionTok{max}\NormalTok{(s}\SpecialCharTok{$}\NormalTok{rmse))))}
\NormalTok{curve}\SpecialCharTok{$}\NormalTok{noise }\OtherTok{\textless{}{-}} \FunctionTok{factor}\NormalTok{(curve}\SpecialCharTok{$}\NormalTok{noise, }\AttributeTok{levels =} \FunctionTok{c}\NormalTok{(}\StringTok{"none"}\NormalTok{, }\StringTok{"ambient"}\NormalTok{, }\StringTok{"outlier"}\NormalTok{))}
\NormalTok{curve }\OtherTok{\textless{}{-}}\NormalTok{ curve[}\FunctionTok{order}\NormalTok{(curve}\SpecialCharTok{$}\NormalTok{noise, curve}\SpecialCharTok{$}\NormalTok{theta), ]}

\NormalTok{clean  }\OtherTok{\textless{}{-}} \FunctionTok{subset}\NormalTok{(curve, noise }\SpecialCharTok{==} \StringTok{"none"} \SpecialCharTok{\&}\NormalTok{ theta }\SpecialCharTok{\textgreater{}} \DecValTok{0}\NormalTok{)}
\NormalTok{expo   }\OtherTok{\textless{}{-}} \FunctionTok{coef}\NormalTok{(}\FunctionTok{lm}\NormalTok{(}\FunctionTok{log}\NormalTok{(mean) }\SpecialCharTok{\textasciitilde{}} \FunctionTok{log}\NormalTok{(theta), }\FunctionTok{subset}\NormalTok{(clean, theta }\SpecialCharTok{\textless{}=} \FloatTok{0.5}\NormalTok{)))[[}\DecValTok{2}\NormalTok{]]}
\NormalTok{step   }\OtherTok{\textless{}{-}}\NormalTok{ clean}\SpecialCharTok{$}\NormalTok{mean[}\SpecialCharTok{{-}}\DecValTok{1}\NormalTok{] }\SpecialCharTok{/}\NormalTok{ clean}\SpecialCharTok{$}\NormalTok{mean[}\SpecialCharTok{{-}}\FunctionTok{nrow}\NormalTok{(clean)]}
\NormalTok{spread }\OtherTok{\textless{}{-}} \FunctionTok{max}\NormalTok{((curve}\SpecialCharTok{$}\NormalTok{hi }\SpecialCharTok{{-}}\NormalTok{ curve}\SpecialCharTok{$}\NormalTok{lo)[curve}\SpecialCharTok{$}\NormalTok{noise }\SpecialCharTok{==} \StringTok{"none"}\NormalTok{])}
\NormalTok{at     }\OtherTok{\textless{}{-}} \ControlFlowTok{function}\NormalTok{(nz, t) curve}\SpecialCharTok{$}\NormalTok{mean[curve}\SpecialCharTok{$}\NormalTok{noise }\SpecialCharTok{==}\NormalTok{ nz }\SpecialCharTok{\&}\NormalTok{ curve}\SpecialCharTok{$}\NormalTok{theta }\SpecialCharTok{==}\NormalTok{ t]}
\end{Highlighting}
\end{Shaded}

\begin{figure}

{\centering \includegraphics[width=0.9\linewidth,alt={A line chart with the folding parameter from zero to nearly one on the horizontal axis and reconstruction error from zero to about 0.65 on the vertical. Three curves, distinguished by colour and line style. The lowest, for no noise, starts at zero and rises smoothly with increasing slope to about 0.21. The middle one, ambient noise, starts near 0.12 and is almost flat over the first quarter before rising to about 0.30. The highest, outliers, starts near 0.36 and behaves the same way, reaching about 0.62. All three are monotone with no crossing, and every one steepens over the last part of the sweep; the shaded bands around them are very narrow.}]{04-linear-projections_files/figure-latex/lin-curve-fig-1} 

}

\caption{The reference curve. Two-component PCA's reconstruction error against the chart, over the fold sweep, for the three noise models. The line is the mean over the twenty benchmark seeds and the band is their full range; the band is narrower than the line except at the last fold angle. The floor is zero everywhere on this grid, so the whole of each curve is recoverable structure that the projection did not recover.}\label{fig:lin-curve-fig}
\end{figure}

Figure \ref{fig:lin-curve-fig} has no feature in it, and that is the result
rather than a disappointment. There is no threshold, no knee and no regime
change: PCA degrades smoothly and monotonically from exact to poor, and a
benchmark that reported a single fold angle would have found nothing here to
report.

The degradation obeys a law. Over the first half of the sweep the clean curve is
a power of \(\theta\) with exponent 1.99 --- quadratic, to three
figures. Chapter \ref{folding-geometry} derives why: ambient distance across a
crease contracts by \(\sin(\rho/2)\), and \(1 - \sin(\rho/2)\) is quadratic in the
departure from flatness. The error PCA accumulates is that contraction,
integrated over the sheet.

Past the middle the curve steepens. The step from the second-largest to the
largest fold angle multiplies the error by 1.49
against a median step of 1.24, which is the same endgame
acceleration Chapter \ref{folding-geometry} shows in the sheet's contraction.
Clean, the error reaches 0.048 at the middle of the
sweep and 0.212 at its end.

Seeds contribute almost nothing. The widest range across the twenty benchmark
seeds at any clean fold angle is 0.011, and it occurs at the last
one; everywhere else the band is thinner than the line drawn over it. The curve
is a property of the geometry and not of the draw, which is what makes it usable
as a reference.

Noise moves the curve up without changing its shape. At \(\theta = 0\) the error is
already 0.12 under ambient Gaussian noise and
0.361 under outliers, on a sheet that is flat and whose
floor is zero. Those intercepts are the noise, not the fold; the fold's
contribution is what rides on top of them.

\begin{Shaded}
\begin{Highlighting}[]
\NormalTok{prod }\OtherTok{\textless{}{-}} \FunctionTok{read\_run}\NormalTok{(}\StringTok{"data/processed/product{-}grid.rds"}\NormalTok{, }\StringTok{"bb8ade76d88b95a3"}\NormalTok{)}
\NormalTok{pr2  }\OtherTok{\textless{}{-}} \FunctionTok{subset}\NormalTok{(prod, family }\SpecialCharTok{==} \StringTok{"crease"} \SpecialCharTok{\&}\NormalTok{ d }\SpecialCharTok{==}\NormalTok{ EMBED\_DIM }\SpecialCharTok{\&}\NormalTok{ status }\SpecialCharTok{==} \StringTok{"ok"}\NormalTok{)}
\NormalTok{exc  }\OtherTok{\textless{}{-}} \FunctionTok{aggregate}\NormalTok{(}\FunctionTok{cbind}\NormalTok{(rmse, floor, excess) }\SpecialCharTok{\textasciitilde{}}\NormalTok{ method, pr2, median)}
\NormalTok{exc  }\OtherTok{\textless{}{-}}\NormalTok{ exc[}\FunctionTok{order}\NormalTok{(exc}\SpecialCharTok{$}\NormalTok{excess), ]}
\end{Highlighting}
\end{Shaded}

One number from Chapter \ref{benchmarks}'s product construction belongs here,
because a chapter with this title should admit where linear projection
emphatically does not break. Pairing two folded sheets gives a chart of intrinsic
dimension four whose floor in two dimensions is strictly positive --- median
0.514 here --- so a method can be scored by how far
\emph{above} the floor it lands. PCA's median excess is
0.00508. It is the smallest excess
of any method in this book, below Isomap's
0.0155, on a problem where the
optimum is known exactly. When the obstruction is dimension rather than
curvature, the optimal linear map is very nearly the optimal map.

\section{Diagnostics}\label{lin-diagnostics}

Two diagnostics a practitioner would reach for do not see the fold, and both
would report success on the right-hand end of Figure \ref{fig:lin-curve-fig}.

\begin{Shaded}
\begin{Highlighting}[]
\NormalTok{scree }\OtherTok{\textless{}{-}} \FunctionTok{do.call}\NormalTok{(rbind, }\FunctionTok{lapply}\NormalTok{(THETA\_GRID, }\ControlFlowTok{function}\NormalTok{(t) \{}
\NormalTok{  v }\OtherTok{\textless{}{-}} \FunctionTok{vapply}\NormalTok{(BENCH\_SEEDS[}\DecValTok{1}\SpecialCharTok{:}\DecValTok{5}\NormalTok{], }\ControlFlowTok{function}\NormalTok{(sd) \{}
\NormalTok{    s   }\OtherTok{\textless{}{-}} \FunctionTok{sample\_manifold}\NormalTok{(p, }\AttributeTok{theta =}\NormalTok{ t, }\AttributeTok{n =}\NormalTok{ N\_DEFAULT, }\AttributeTok{seed =}\NormalTok{ sd)}
\NormalTok{    lam }\OtherTok{\textless{}{-}} \FunctionTok{prcomp}\NormalTok{(s}\SpecialCharTok{$}\NormalTok{X)}\SpecialCharTok{$}\NormalTok{sdev}\SpecialCharTok{\^{}}\DecValTok{2}
    \FunctionTok{c}\NormalTok{(}\FunctionTok{sqrt}\NormalTok{(}\FunctionTok{sum}\NormalTok{(lam[}\SpecialCharTok{{-}}\FunctionTok{seq\_len}\NormalTok{(EMBED\_DIM)]) }\SpecialCharTok{/} \FunctionTok{sum}\NormalTok{(lam)),}
      \FunctionTok{reconstruction\_error}\NormalTok{(}\FunctionTok{embed}\NormalTok{(}\StringTok{"pca"}\NormalTok{, s), s}\SpecialCharTok{$}\NormalTok{truth))}
\NormalTok{  \}, }\FunctionTok{numeric}\NormalTok{(}\DecValTok{2}\NormalTok{))}
  \FunctionTok{data.frame}\NormalTok{(}\AttributeTok{theta =}\NormalTok{ t, }\AttributeTok{ambient =} \FunctionTok{mean}\NormalTok{(v[}\DecValTok{1}\NormalTok{, ]), }\AttributeTok{chart =} \FunctionTok{mean}\NormalTok{(v[}\DecValTok{2}\NormalTok{, ]))}
\NormalTok{\}))}
\end{Highlighting}
\end{Shaded}

\textbf{The scree plot does not collapse.} The plan of record predicted that variance
explained would collapse as the fold deepens. It does not. At the largest fold
angle the first two components still hold
94.3 per cent of the ambient
variance --- a fall of under six percentage points across the entire sweep, while
the chart error grows by more than two orders of magnitude. A reader choosing the
number of components from a scree plot would be right about the dimension and
uninformed about the geometry.

\begin{figure}

{\centering \includegraphics[width=0.9\linewidth,alt={A line chart over the folding parameter with two curves, distinguished by colour and line style. Both start at zero and rise. The upper curve, the ambient scree residual, rises roughly linearly to about 0.24. The lower curve, the chart reconstruction error, hugs the axis at first and then accelerates, reaching about 0.21 at the right-hand end, so the gap between them is wide on the left and narrow on the right.}]{04-linear-projections_files/figure-latex/lin-scree-fig-1} 

}

\caption{The residual a scree plot shows against the error the answer key shows, over the same sweep. Both are dimensionless and both are monotone, but they are not proxies for each other: the scree residual is measured against the ambient cloud and the reconstruction error against the chart, and their ratio is not constant.}\label{fig:lin-scree-fig}
\end{figure}

Figure \ref{fig:lin-scree-fig} shows why the two cannot be substituted. The
scree residual overstates the chart error by a factor of
16.7 at the smallest non-zero fold
angle and 1.14 at
the largest: monotone in each other on this one axis, and nowhere near
proportional.

\textbf{Trustworthiness does not see it either.} The grid scores trustworthiness and
continuity at \texttt{K\_DEFAULT} against \emph{ambient} distance and \(Q_{NX}\) at \texttt{K\_QNX}
against the \emph{chart}. Put the two conventions side by side:

\begin{Shaded}
\begin{Highlighting}[]
\NormalTok{tc }\OtherTok{\textless{}{-}} \FunctionTok{aggregate}\NormalTok{(}\FunctionTok{cbind}\NormalTok{(trust, cont, qnx, knn) }\SpecialCharTok{\textasciitilde{}}\NormalTok{ theta,}
                \FunctionTok{subset}\NormalTok{(grid, method }\SpecialCharTok{==} \StringTok{"pca"} \SpecialCharTok{\&}\NormalTok{ noise }\SpecialCharTok{==} \StringTok{"none"}\NormalTok{), mean)}
\NormalTok{worst }\OtherTok{\textless{}{-}}\NormalTok{ tc[}\FunctionTok{nrow}\NormalTok{(tc), ]}
\end{Highlighting}
\end{Shaded}

At the largest fold angle PCA scores 0.992 on
trustworthiness and 0.997 on continuity, while \(Q_{NX}\)
against the answer key is 0.676 and \(k\)-nearest-neighbour
preservation is 0.505. An embedding that has lost half its
neighbours by the chart's reckoning is rated near-perfect by the ambient one.
Neither number is wrong; they answer different questions, which is Chapter
\ref{evaluation}'s subject.

\section{Where it fails}\label{lin-limits}

\textbf{A projection cannot unfold, and this is not a tuning problem.} The optimal
subspace is already what PCA returns \citep{eckart1936approximation}, so no choice of
plane avoids the trade-off in \texttt{lin-core}: keep the corrugation and spend a
component on it, or discard it and accept an anisotropically squashed chart.
Every later method in Part II earns its keep by consuming something other than
ambient distance.

\textbf{No robustness.} The outlier intercept in Figure \ref{fig:lin-curve-fig} is a
covariance matrix being dragged by displaced points. PCA has no defence against
it and none is offered here; the book reports the cost rather than repairing it,
so that later methods meet a plain implementation.

\textbf{The floor excuses nothing on this grid.} \texttt{metric\_floor()} is zero in every
cell, so the whole curve is loss the projection is responsible for. That flips on
the product construction, where the floor is large and PCA's excess above it is
the smallest in the book. The two results are not in tension: a linear method is
close to optimal exactly when the obstruction is dimension, and far from it when
the obstruction is curvature.

\textbf{One family, one size, one \(\alpha\).} Everything above is measured on a single
Miura tessellation at one opening angle, so the ordering of methods at a given
fold angle is a property of this grid; Chapter \ref{results} is where it meets
anything else. A benchmark on one family cannot support a claim about pattern
variety, and this chapter makes none.

\section{Reproduce this}\label{lin-reproduce}

The grid is committed, not generated at render time.
\texttt{scripts/run-benchmark-grid.R} produces it and declares the artefact path in its
header; \texttt{scripts/check-artefact-producers.R} checks that every artefact the book
specifies has a producer. Re-running that script and reproducing its numbers
leaves the digest in \texttt{lin-setup} alone; getting different numbers changes it, and
the render then stops with both digests in the message.

The rest is checked by the suite, which runs on every push:

\begin{itemize}
\tightlist
\item
  \texttt{test-methods.R} --- \emph{``the ambient-metric methods recover a flat sheet exactly''}
  and \emph{``PCA and classical MDS agree, because they are the same thing''}.
\item
  \texttt{test-metrics.R} --- \emph{``a 2-D chart has no irreducible loss at d = 2''}, \emph{``the
  floor is attained, exactly, by the top-d principal projection''}, and \emph{``the
  floor does not move with theta or with noise''}.
\item
  \texttt{test-artefacts.R} --- \emph{``a chapter cannot render against a run it was not written
  against''}.
\end{itemize}

Run them with \texttt{Rscript\ tests/testthat.R} from the repository root.

\section{Further reading}\label{lin-reading}

Pearson \citep{pearson1901lines} and Hotelling \citep{hotelling1933components} are the two
original readings of PCA and are still the clearest statement of why they are the
same thing; Eckart and Young \citep{eckart1936approximation} supply the optimality
that makes it a baseline rather than a heuristic. Torgerson
\citep{torgerson1952mds} is classical scaling's source and Kruskal \citep{kruskal1964fit}
the stress-minimising alternative it is routinely confused with. Gower
\citep{gower1975procrustes} is the reference for the alignment the score is built on;
Chapter \ref{evaluation} owns that convention and is where to go for it.

The bridge into the rest of Part II is kernel PCA \citep{scholkopf1998kernelpca}.
Replacing the Gram matrix in classical scaling with a kernel makes the same
eigenproblem non-linear, and Isomap --- Chapter \ref{geodesic} --- is exactly that
substitution with a kernel built from geodesic distances. Lee and Verleysen
\citep{lee2007nonlinear} set out the whole family in those terms.

\chapter{Geodesic methods}\label{geodesic}

\emph{R. Heller · Date: 2026-07-29 · Last changed: 2026-10-30}

\emph{If the data lie on a folded sheet, how do you flatten it back out?}

\section{What this chapter answers}\label{geo-question}

Isomap's premise is this book's premise. Estimate the distances measured \emph{along}
the surface, embed those instead of the distances measured \emph{through} the space
around it, and a sheet that has been bent without being stretched should come
back flat. \texttt{R/methods.R} records what every method consumes, and Isomap is the
only entry in the registry whose answer to that question is ``geodesic distance''.

That makes this the one chapter where the method's input and the answer key are
the same object. A rigid folding is an intrinsic isometry, so the geodesic
distance between two sampled points \textbf{is} the Euclidean distance between their
rows of \texttt{truth}. Not an estimate of it, and not a limit approached as the sample
grows: the exact number, in closed form, before any graph is built.

The consequence is that Isomap's error decomposes into two terms that can be
measured apart rather than argued about. A \textbf{graph term}, from approximating
geodesics by shortest paths on a \(k\)-nearest-neighbour graph. An \textbf{embedding
term}, from putting the resulting distance matrix into two dimensions with
classical scaling. Hand \texttt{cmdscale()} the exact geodesics and the second term
vanishes to floating point. Everything that is left belongs to the first.

This chapter measures that first term across the fold sweep, predicts where it
turns from a bias into a break, and reports what the break actually is on this
object --- which is not the mechanism the plan of record predicted.

Diffusion maps is here as the contrast. It sits in the same registry family and
consumes \emph{ambient} distance, and the gap between the two curves is the value of
consuming the right thing.

\section{Setup}\label{geo-setup}

\begin{Shaded}
\begin{Highlighting}[]
\NormalTok{grid   }\OtherTok{\textless{}{-}} \FunctionTok{read\_run}\NormalTok{(}\StringTok{"data/processed/benchmark{-}grid.rds"}\NormalTok{, }\StringTok{"2483a468d82d4d64"}\NormalTok{)}
\NormalTok{sweeps }\OtherTok{\textless{}{-}} \FunctionTok{read\_run}\NormalTok{(}\StringTok{"data/processed/part2{-}sweeps.rds"}\NormalTok{,  }\StringTok{"16d43ff4992c0e58"}\NormalTok{)}
\NormalTok{p }\OtherTok{\textless{}{-}} \FunctionTok{miura\_ori}\NormalTok{(}\AttributeTok{nx =} \DecValTok{6}\DataTypeTok{L}\NormalTok{, }\AttributeTok{ny =} \DecValTok{6}\DataTypeTok{L}\NormalTok{, }\AttributeTok{alpha =}\NormalTok{ pi }\SpecialCharTok{/} \DecValTok{3}\NormalTok{)}
\NormalTok{s }\OtherTok{\textless{}{-}} \FunctionTok{sample\_manifold}\NormalTok{(p, }\AttributeTok{theta =} \FloatTok{0.5}\NormalTok{, }\AttributeTok{n =}\NormalTok{ N\_DEFAULT, }\AttributeTok{seed =}\NormalTok{ BENCH\_SEEDS[}\DecValTok{1}\NormalTok{])}
\end{Highlighting}
\end{Shaded}

The main grid fixes \(k\) at \texttt{K\_DEFAULT}; \texttt{part2-sweeps.rds} is the artefact that
moves it, over 7 values at
3 fold angles and all 20 benchmark seeds.
Chapter \ref{linear} explains what \texttt{read\_run()} refuses.

\section{Background}\label{geo-background}

\subsection{Isomap, in three stages}\label{isomap-in-three-stages}

Isomap \citep{tenenbaum2000isomap} is classical multidimensional scaling
\citep{torgerson1952mds} with the distance matrix replaced. Build a
\(k\)-nearest-neighbour graph on the ambient points, weight each edge by its
ambient length, take all-pairs shortest paths, and hand the result to
\texttt{cmdscale()}. The first two stages estimate the geodesic metric; the third is the
same eigenproblem Chapter \ref{linear} already ran. Read as a kernel method
\citep{scholkopf1998kernelpca, lee2007nonlinear} it is classical scaling with a
geodesic kernel, which is exactly how Chapter \ref{linear} hands the reader over.

The estimate is justified asymptotically. Graph shortest paths converge to true
geodesics as the sample density rises, at a rate controlled by the sampling
radius and the surface's curvature \citep{bernstein2000graph}. Both stages of that
guarantee are worth naming here, because both fail on a crease pattern: the
convergence is stated for a manifold of positive reach, and a crease has reach
exactly zero (Chapter \ref{folding-geometry}), and it is asymptotic in a
density this book fixes at \texttt{N\_DEFAULT} points.

\begin{Shaded}
\begin{Highlighting}[]
\NormalTok{reg }\OtherTok{\textless{}{-}} \FunctionTok{method\_table}\NormalTok{()}
\NormalTok{fam }\OtherTok{\textless{}{-}}\NormalTok{ reg[reg}\SpecialCharTok{$}\NormalTok{family }\SpecialCharTok{==} \StringTok{"geodesic"}\NormalTok{, }\FunctionTok{c}\NormalTok{(}\StringTok{"method"}\NormalTok{, }\StringTok{"label"}\NormalTok{, }\StringTok{"consumes"}\NormalTok{, }\StringTok{"stochastic"}\NormalTok{)]}
\NormalTok{knitr}\SpecialCharTok{::}\FunctionTok{kable}\NormalTok{(fam, }\AttributeTok{row.names =} \ConstantTok{FALSE}\NormalTok{,}
             \AttributeTok{col.names =} \FunctionTok{c}\NormalTok{(}\StringTok{"registry name"}\NormalTok{, }\StringTok{"label"}\NormalTok{, }\StringTok{"consumes"}\NormalTok{, }\StringTok{"stochastic"}\NormalTok{),}
             \AttributeTok{caption =} \StringTok{"The geodesic family of the embedding registry."}\NormalTok{)}
\end{Highlighting}
\end{Shaded}

\begin{table}

\caption{\label{tab:geo-registry}The geodesic family of the embedding registry.}
\centering
\begin{tabular}[t]{l|l|l|l}
\hline
registry name & label & consumes & stochastic\\
\hline
isomap & Isomap & geodesic & FALSE\\
\hline
diffusion & Diffusion map & ambient & TRUE\\
\hline
\end{tabular}
\end{table}

Table \ref{tab:geo-registry} is the chapter in one row each. Isomap consumes
geodesic distance and is deterministic. Diffusion maps consumes \textbf{ambient}
distance and is stochastic, and both of those are easy to get wrong.

\subsection{Diffusion maps, and why it is filed here}\label{diffusion-maps-and-why-it-is-filed-here}

Diffusion maps \citep{coifman2006diffusion} puts a Gaussian kernel on the ambient
distances, normalises it into a Markov chain, and embeds the points so that
Euclidean distance in the embedding approximates \emph{diffusion distance} --- the
difference between the transition distributions of a random walk started at each
point. Diffusion distance is a geodesic-flavoured quantity, which is why the
family label is what it is, but nothing in the construction estimates a geodesic.
The kernel sees ambient distance and only ambient distance, and the registry's
\texttt{consumes} field says so. The same holds for Laplacian eigenmaps
\citep{belkin2003laplacian}, which Chapter \ref{neighbor} takes up.

Two implementation facts travel with it. It is \textbf{stochastic} ---
\texttt{diffusionMap::diffuse()} uses a randomised eigensolver --- so \texttt{embed()} refuses to
run it without a seed. And its bandwidth is set from the \(k\)-th
nearest-neighbour distance rather than the median pairwise distance, because the
outlier noise model inflates the second and not the first; \texttt{R/methods.R} records
what that cost when it was done the obvious way.

\begin{boxreport}
\textbf{Two changes from the plan of record.} The outline for this
chapter named \texttt{igraph::distances()} for the shortest paths and
locally linear embedding as a third method. Neither survived.
\texttt{igraph} is not installable in this project --- a missing system
library --- so \texttt{R/metrics.R} computes all-pairs shortest paths by
min-plus relaxation with a work list, which is Bellman--Ford rather than
Dijkstra {[}@dijkstra1959note{]} and was timed against two alternatives
at the book's sample size before it was kept. LLE consumes neighbourhood
structure, not geodesics, and is in the \hyperref[neighbor]{neighbor
embeddings chapter} with the other methods that do.
\end{boxreport}

\section{The geodesic you do not have to estimate}\label{geo-core}

\subsection{The claim, and its one condition}\label{the-claim-and-its-one-condition}

Folding preserves distance along the surface exactly (Chapter
\ref{folding-geometry}). So the length of a path on the folded sheet equals the
length of its preimage in the flat chart, and the shortest path in the chart
between two points is the straight segment between them. Therefore
\texttt{dist(s\$truth)} is the geodesic distance matrix --- not a discretisation of it.

The condition is that the straight segment stays on the paper. A Miura's
unfolded outline is not convex, so this is a real condition and not a
formality; \texttt{sample\_manifold()}'s default \texttt{boundary\ =\ FALSE} trims the sample to a
convex sub-rectangle, and the trimming is what makes the equality hold rather
than a bound.

\begin{Shaded}
\begin{Highlighting}[]
\NormalTok{dU   }\OtherTok{\textless{}{-}} \FunctionTok{reference\_dist}\NormalTok{(s, }\StringTok{"chart"}\NormalTok{)}
\NormalTok{dA   }\OtherTok{\textless{}{-}} \FunctionTok{reference\_dist}\NormalTok{(s, }\StringTok{"ambient"}\NormalTok{)}
\NormalTok{same }\OtherTok{\textless{}{-}} \FunctionTok{outer}\NormalTok{(s}\SpecialCharTok{$}\NormalTok{facet, s}\SpecialCharTok{$}\NormalTok{facet, }\StringTok{"=="}\NormalTok{)}
\NormalTok{up   }\OtherTok{\textless{}{-}} \FunctionTok{upper.tri}\NormalTok{(dU)}

\NormalTok{within  }\OtherTok{\textless{}{-}} \FunctionTok{max}\NormalTok{(}\FunctionTok{abs}\NormalTok{(dA }\SpecialCharTok{{-}}\NormalTok{ dU)[same }\SpecialCharTok{\&}\NormalTok{ up])          }\CommentTok{\# isometry, inside a facet}
\NormalTok{across  }\OtherTok{\textless{}{-}} \FunctionTok{max}\NormalTok{((dA }\SpecialCharTok{{-}}\NormalTok{ dU)[}\SpecialCharTok{!}\NormalTok{same }\SpecialCharTok{\&}\NormalTok{ up])            }\CommentTok{\# contraction, across a crease}
\NormalTok{exit    }\OtherTok{\textless{}{-}} \FunctionTok{chart\_exit\_fraction}\NormalTok{(s, p)             }\CommentTok{\# segments that leave the sheet}
\end{Highlighting}
\end{Shaded}

Three checks, and each one is a different half of the claim. Inside a facet
ambient and chart distance agree to 1.1e-15, which is rounding on a
closed form. Across a crease the ambient distance is strictly smaller --- the
largest signed difference over all 308,887
such pairs is -6.2e-07, a negative number. And
\texttt{chart\_exit\_fraction()} reports that 0 of sampled segments leave the
sheet, so the chart distance is the geodesic for every pair in this sample and
not merely a lower bound on it.

\subsection{The decomposition}\label{the-decomposition}

\begin{Shaded}
\begin{Highlighting}[]
\NormalTok{Y\_exact }\OtherTok{\textless{}{-}} \FunctionTok{cmdscale}\NormalTok{(}\FunctionTok{as.dist}\NormalTok{(dU), }\AttributeTok{k =}\NormalTok{ EMBED\_DIM)           }\CommentTok{\# exact geodesics in}
\NormalTok{dG      }\OtherTok{\textless{}{-}} \FunctionTok{reference\_dist}\NormalTok{(s, }\StringTok{"graph"}\NormalTok{, }\AttributeTok{k =}\NormalTok{ K\_DEFAULT)      }\CommentTok{\# what Isomap estimates}
\NormalTok{Y\_iso   }\OtherTok{\textless{}{-}} \FunctionTok{embed}\NormalTok{(}\StringTok{"isomap"}\NormalTok{, s, }\AttributeTok{k =}\NormalTok{ K\_DEFAULT)}

\NormalTok{err\_embed }\OtherTok{\textless{}{-}} \FunctionTok{reconstruction\_error}\NormalTok{(Y\_exact, s}\SpecialCharTok{$}\NormalTok{truth)}
\NormalTok{err\_total }\OtherTok{\textless{}{-}} \FunctionTok{reconstruction\_error}\NormalTok{(Y\_iso,   s}\SpecialCharTok{$}\NormalTok{truth)}
\NormalTok{agree     }\OtherTok{\textless{}{-}} \FunctionTok{max}\NormalTok{(}\FunctionTok{abs}\NormalTok{(}\FunctionTok{unname}\NormalTok{(Y\_iso) }\SpecialCharTok{{-}} \FunctionTok{unname}\NormalTok{(}\FunctionTok{cmdscale}\NormalTok{(}\FunctionTok{as.dist}\NormalTok{(dG), }\AttributeTok{k =}\NormalTok{ EMBED\_DIM))))}
\NormalTok{graph\_dev }\OtherTok{\textless{}{-}} \FunctionTok{sqrt}\NormalTok{(}\FunctionTok{mean}\NormalTok{(((dG }\SpecialCharTok{{-}}\NormalTok{ dU)[up])}\SpecialCharTok{\^{}}\DecValTok{2}\NormalTok{)) }\SpecialCharTok{/} \FunctionTok{sqrt}\NormalTok{(}\FunctionTok{mean}\NormalTok{((dU[up])}\SpecialCharTok{\^{}}\DecValTok{2}\NormalTok{))}
\end{Highlighting}
\end{Shaded}

Feed classical scaling the exact geodesic distances and it returns the chart, at
a Procrustes error of 6.9e-16. That
is the embedding term, and it is zero: the chart is a two-dimensional Euclidean
configuration, so double-centring its squared distances and taking two
eigenvectors inverts the construction exactly.

Isomap on the same sample scores 0.0266. It differs from
\texttt{cmdscale()} on the graph distances by 8.8e-15,
because it \emph{is} \texttt{cmdscale()} on the graph distances. So the whole of that
0.0266 is the graph term, and the graph term is small in the
distances themselves --- the shortest-path matrix departs from the exact geodesics
by 2.95 per cent in relative root-mean-square.

This is the sense in which the chapter's title is not a metaphor. Isomap's error
on a crease pattern has one cause, it is a property of a graph rather than of an
optimiser, and it can be looked at directly.

\begin{figure}

{\centering \includegraphics[width=0.9\linewidth,alt={Two scatter panels side by side, each with exact geodesic distance on the horizontal axis and the ratio of graph distance to geodesic distance on the vertical, with a reference line at one. In the left panel, at a small fold angle, the cloud spreads widely above the line at the shortest distances, reaching roughly twice the geodesic, then narrows quickly into a tight band sitting just above one for the rest of the range. In the right panel, at a large fold angle, the same narrowing happens but the band has dropped below the line: the short-distance scatter now runs from about a half to about twice, and beyond a moderate distance the band settles a little under one.}]{05-geodesic-methods_files/figure-latex/geo-graph-fig-1} 

}

\caption{Shortest-path distance on the ambient neighbourhood graph against the exact geodesic distance, as a ratio, for a lightly folded sheet and a deeply folded one. The horizontal line at one is exactness. Points above it are the discretisation bias of walking along graph edges; points below it are distances the fold has contracted. A random ten thousand pairs are drawn in each panel.}\label{fig:geo-graph-fig}
\end{figure}

Figure \ref{fig:geo-graph-fig} is the graph term seen twice. At a small fold
angle it is a \emph{positive} bias: a path made of graph edges zigzags where a
geodesic runs straight, so shortest paths come out too long, and the scatter
collapses onto a tight band just above one as the pair gets further apart and the
zigzag averages out. At a large fold angle that band has dropped below the line. Edges that cross a crease are shortened by
the \(\sin(\rho/2)\) factor Chapter \ref{folding-geometry} derives, shortest paths
route through as many creases as they can, and the estimate falls below the
truth. The two biases are different mechanisms with opposite signs, and
\texttt{geo-diagnostics} is where they cancel.

\section{Results}\label{geo-results}

\begin{Shaded}
\begin{Highlighting}[]
\NormalTok{clean }\OtherTok{\textless{}{-}} \FunctionTok{subset}\NormalTok{(grid, noise }\SpecialCharTok{==} \StringTok{"none"} \SpecialCharTok{\&}\NormalTok{ ran)}
\NormalTok{cur   }\OtherTok{\textless{}{-}} \FunctionTok{aggregate}\NormalTok{(}\FunctionTok{cbind}\NormalTok{(rmse, qnx) }\SpecialCharTok{\textasciitilde{}}\NormalTok{ method }\SpecialCharTok{+}\NormalTok{ theta, clean, mean)}
\NormalTok{band  }\OtherTok{\textless{}{-}} \FunctionTok{aggregate}\NormalTok{(rmse }\SpecialCharTok{\textasciitilde{}}\NormalTok{ method }\SpecialCharTok{+}\NormalTok{ theta, clean,}
                   \ControlFlowTok{function}\NormalTok{(x) }\FunctionTok{c}\NormalTok{(}\AttributeTok{lo =} \FunctionTok{min}\NormalTok{(x), }\AttributeTok{hi =} \FunctionTok{max}\NormalTok{(x)))}
\NormalTok{band  }\OtherTok{\textless{}{-}} \FunctionTok{data.frame}\NormalTok{(band[, }\FunctionTok{c}\NormalTok{(}\StringTok{"method"}\NormalTok{, }\StringTok{"theta"}\NormalTok{)], band}\SpecialCharTok{$}\NormalTok{rmse)}
\NormalTok{cur   }\OtherTok{\textless{}{-}} \FunctionTok{merge}\NormalTok{(cur, band, }\AttributeTok{by =} \FunctionTok{c}\NormalTok{(}\StringTok{"method"}\NormalTok{, }\StringTok{"theta"}\NormalTok{))}
\NormalTok{cur   }\OtherTok{\textless{}{-}} \FunctionTok{subset}\NormalTok{(cur, method }\SpecialCharTok{\%in\%} \FunctionTok{c}\NormalTok{(}\StringTok{"isomap"}\NormalTok{, }\StringTok{"diffusion"}\NormalTok{, }\StringTok{"pca"}\NormalTok{))}

\NormalTok{at }\OtherTok{\textless{}{-}} \ControlFlowTok{function}\NormalTok{(m, t, }\AttributeTok{col =} \StringTok{"rmse"}\NormalTok{) cur[[col]][cur}\SpecialCharTok{$}\NormalTok{method }\SpecialCharTok{==}\NormalTok{ m }\SpecialCharTok{\&}\NormalTok{ cur}\SpecialCharTok{$}\NormalTok{theta }\SpecialCharTok{==}\NormalTok{ t]}
\NormalTok{wide }\OtherTok{\textless{}{-}} \FunctionTok{reshape}\NormalTok{(cur[, }\FunctionTok{c}\NormalTok{(}\StringTok{"theta"}\NormalTok{, }\StringTok{"method"}\NormalTok{, }\StringTok{"rmse"}\NormalTok{, }\StringTok{"qnx"}\NormalTok{)],}
                \AttributeTok{idvar =} \StringTok{"theta"}\NormalTok{, }\AttributeTok{timevar =} \StringTok{"method"}\NormalTok{, }\AttributeTok{direction =} \StringTok{"wide"}\NormalTok{)}
\NormalTok{wide }\OtherTok{\textless{}{-}}\NormalTok{ wide[}\FunctionTok{order}\NormalTok{(wide}\SpecialCharTok{$}\NormalTok{theta), ]}
\NormalTok{cross\_rmse }\OtherTok{\textless{}{-}}\NormalTok{ wide}\SpecialCharTok{$}\NormalTok{theta[}\FunctionTok{which}\NormalTok{(wide}\SpecialCharTok{$}\NormalTok{rmse.isomap }\SpecialCharTok{\textless{}}\NormalTok{ wide}\SpecialCharTok{$}\NormalTok{rmse.pca)[}\DecValTok{1}\NormalTok{]]}
\NormalTok{cross\_qnx  }\OtherTok{\textless{}{-}}\NormalTok{ wide}\SpecialCharTok{$}\NormalTok{theta[}\FunctionTok{which}\NormalTok{(wide}\SpecialCharTok{$}\NormalTok{qnx.isomap  }\SpecialCharTok{\textgreater{}}\NormalTok{ wide}\SpecialCharTok{$}\NormalTok{qnx.pca)[}\DecValTok{1}\NormalTok{]]}
\NormalTok{noisy }\OtherTok{\textless{}{-}} \FunctionTok{aggregate}\NormalTok{(rmse }\SpecialCharTok{\textasciitilde{}}\NormalTok{ noise }\SpecialCharTok{+}\NormalTok{ theta, }\FunctionTok{subset}\NormalTok{(grid, method }\SpecialCharTok{==} \StringTok{"isomap"} \SpecialCharTok{\&}\NormalTok{ ran),}
\NormalTok{                   mean)}
\end{Highlighting}
\end{Shaded}

\begin{figure}

{\centering \includegraphics[width=0.9\linewidth,alt={Two panels sharing a horizontal axis, the folding parameter from zero to nearly one. Panel A, reconstruction error: the PCA curve starts at zero and rises smoothly to about 0.21; the Isomap curve starts at about 0.024, stays nearly flat through the first half and then rises to about 0.19, crossing below the PCA curve at around 0.4; the diffusion curve is nearly flat at about 0.2 across the whole range and rises at the very end. Panel B, Q-NX: PCA starts at one and falls to about 0.68; Isomap starts at about 0.94, falls slowly and crosses above PCA near 0.6; diffusion runs flat near 0.785 and falls at the end.}]{05-geodesic-methods_files/figure-latex/geo-curves-fig-1} 

}

\caption{Isomap and diffusion maps over the fold sweep, with PCA from Chapter 4 as the reference. Panel A is reconstruction error against the chart, panel B is $Q_{NX}$ at $K_{QNX}$ against the chart. Lines are means over the twenty benchmark seeds and bands in panel A are their full range. Lower is better in A, higher in B.}\label{fig:geo-curves-fig}
\end{figure}

Three readings of Figure \ref{fig:geo-curves-fig}, in order of how much they
cost the reader to accept.

\textbf{Isomap wins, eventually, and pays an entry fee for it.} On the flat sheet its
error is 0.024 where PCA's is zero: the graph term is
there at \(\theta = 0\), because a neighbourhood graph is a discretisation whether
or not anything has been folded. Isomap does not overtake PCA on reconstruction
error until \(\theta = 0.4\), and on \(Q_{NX}\) not until
\(\theta = 0.6\). From 0.4 to the end of the sweep it is the
best method in the book on this grid, and at the last fold angle it scores
0.188 against PCA's
0.212.

That the two metrics put the crossing in different places is not a discrepancy to
be resolved. Reconstruction error reads distances and \(Q_{NX}\) reads rank order,
and Chapter \ref{evaluation} is where the disagreement is the subject rather
than a nuisance.

\textbf{The curve has a shape, and the shape is the point.} Isomap's error is
essentially flat through the first half of the sweep --- it grows by a factor of
1.19 between the flat sheet and the
middle --- and then rises by a factor of
6.58 over the second
half. A benchmark reporting one fold angle would have called this method robust
or fragile depending entirely on which angle it chose.

\textbf{Diffusion maps does not improve on PCA anywhere.} Its error is close to flat
across the whole sweep, from 0.201 to
0.329, and it is the \emph{worst} of the
three at
20 of the
20 fold angles in the sweep. A method that consumes ambient
distance and embeds a random walk on it has no mechanism for undoing a fold; it
is in this chapter to show that the family label is not what does the work.

Noise moves Isomap's curve without changing it. At \(\theta = 0\) the error is
already 0.124
under ambient Gaussian noise and
0.416 under
outliers, on a sheet that is not folded at all. Displaced points are new
neighbours, and a graph method's whole input is who is a neighbour.

\section{When does the graph bridge?}\label{geo-bridge}

A \textbf{fold jump} is a \(k\)-nearest-neighbour edge between two points the sheet
holds far apart --- the classical short circuit that destroys Isomap on tightly
wound surfaces \citep{balasubramanian2002isomap}. It is the failure this chapter was
specified to predict, and the prediction can be made without running anything.

Two quantities cross. \texttt{facet\_gap(p,\ theta)} is the exact minimum ambient distance
between two facets that share no vertex: a property of the surface alone, no
sample and no density (Chapter \ref{folding-geometry} uses it to check that the
sheet never passes through itself). Against it stands the \(k\)-nearest-neighbour
radius \(r_k\) in the chart. That radius is a constant of the sweep, not a function
of \(\theta\): \texttt{sample\_manifold()} draws the chart points first and folds them
after, so for a fixed seed \texttt{s\$truth} is the \emph{same matrix} at every fold angle.
The graph acquires its first fold jump when the surface's own clearance falls
below the radius at which the sample connects points.

\begin{Shaded}
\begin{Highlighting}[]
\NormalTok{KS   }\OtherTok{\textless{}{-}} \FunctionTok{c}\NormalTok{(}\DecValTok{10}\DataTypeTok{L}\NormalTok{, }\DecValTok{20}\DataTypeTok{L}\NormalTok{, }\DecValTok{30}\DataTypeTok{L}\NormalTok{, }\DecValTok{45}\DataTypeTok{L}\NormalTok{)}
\NormalTok{s0   }\OtherTok{\textless{}{-}} \FunctionTok{sample\_manifold}\NormalTok{(p, }\AttributeTok{theta =} \DecValTok{0}\NormalTok{, }\AttributeTok{n =}\NormalTok{ N\_DEFAULT, }\AttributeTok{seed =}\NormalTok{ BENCH\_SEEDS[}\DecValTok{1}\NormalTok{])}
\NormalTok{dU0  }\OtherTok{\textless{}{-}} \FunctionTok{as.matrix}\NormalTok{(}\FunctionTok{dist}\NormalTok{(s0}\SpecialCharTok{$}\NormalTok{truth))}
\NormalTok{Lfar }\OtherTok{\textless{}{-}} \FloatTok{0.25} \SpecialCharTok{*} \FunctionTok{max}\NormalTok{(dU0)                       }\CommentTok{\# branch\_gap()\textquotesingle{}s "far apart"}
\NormalTok{d0   }\OtherTok{\textless{}{-}}\NormalTok{ dU0; }\FunctionTok{diag}\NormalTok{(d0) }\OtherTok{\textless{}{-}} \ConstantTok{Inf}
\NormalTok{r\_k  }\OtherTok{\textless{}{-}} \FunctionTok{vapply}\NormalTok{(KS, }\ControlFlowTok{function}\NormalTok{(k)}
  \FunctionTok{median}\NormalTok{(}\FunctionTok{apply}\NormalTok{(d0, }\DecValTok{1}\DataTypeTok{L}\NormalTok{, }\ControlFlowTok{function}\NormalTok{(r) }\FunctionTok{sort.int}\NormalTok{(r, }\AttributeTok{partial =}\NormalTok{ k)[k])), }\FunctionTok{numeric}\NormalTok{(}\DecValTok{1}\NormalTok{))}

\NormalTok{th\_fine }\OtherTok{\textless{}{-}} \FunctionTok{seq}\NormalTok{(}\FloatTok{0.70}\NormalTok{, }\FloatTok{0.99}\NormalTok{, }\AttributeTok{by =} \FloatTok{0.01}\NormalTok{)         }\CommentTok{\# deliberately past THETA\_GRID}
\NormalTok{gap     }\OtherTok{\textless{}{-}} \FunctionTok{vapply}\NormalTok{(th\_fine, }\ControlFlowTok{function}\NormalTok{(t) }\FunctionTok{facet\_gap}\NormalTok{(p, t)}\SpecialCharTok{$}\NormalTok{gap, }\FunctionTok{numeric}\NormalTok{(}\DecValTok{1}\NormalTok{))}
\NormalTok{gapfun  }\OtherTok{\textless{}{-}}\NormalTok{ stats}\SpecialCharTok{::}\FunctionTok{approxfun}\NormalTok{(th\_fine, gap)}
\NormalTok{predicted }\OtherTok{\textless{}{-}} \FunctionTok{vapply}\NormalTok{(r\_k, }\ControlFlowTok{function}\NormalTok{(r) }\FunctionTok{tryCatch}\NormalTok{(}
\NormalTok{  stats}\SpecialCharTok{::}\FunctionTok{uniroot}\NormalTok{(}\ControlFlowTok{function}\NormalTok{(t) }\FunctionTok{gapfun}\NormalTok{(t) }\SpecialCharTok{{-}}\NormalTok{ r, }\FunctionTok{range}\NormalTok{(th\_fine))}\SpecialCharTok{$}\NormalTok{root,}
  \AttributeTok{error =} \ControlFlowTok{function}\NormalTok{(e) }\ConstantTok{NA\_real\_}\NormalTok{), }\FunctionTok{numeric}\NormalTok{(}\DecValTok{1}\NormalTok{))}

\NormalTok{jumps }\OtherTok{\textless{}{-}} \ControlFlowTok{function}\NormalTok{(theta, seed) \{}
\NormalTok{  st }\OtherTok{\textless{}{-}} \FunctionTok{sample\_manifold}\NormalTok{(p, }\AttributeTok{theta =}\NormalTok{ theta, }\AttributeTok{n =}\NormalTok{ N\_DEFAULT, }\AttributeTok{seed =}\NormalTok{ seed)}
\NormalTok{  nn }\OtherTok{\textless{}{-}}\NormalTok{ FNN}\SpecialCharTok{::}\FunctionTok{get.knn}\NormalTok{(st}\SpecialCharTok{$}\NormalTok{X, }\AttributeTok{k =} \FunctionTok{max}\NormalTok{(KS))}\SpecialCharTok{$}\NormalTok{nn.index}
\NormalTok{  du }\OtherTok{\textless{}{-}} \FunctionTok{as.matrix}\NormalTok{(}\FunctionTok{dist}\NormalTok{(st}\SpecialCharTok{$}\NormalTok{truth))}
  \FunctionTok{vapply}\NormalTok{(KS, }\ControlFlowTok{function}\NormalTok{(k) }\FunctionTok{sum}\NormalTok{(du[}\FunctionTok{cbind}\NormalTok{(}\FunctionTok{rep}\NormalTok{(}\FunctionTok{seq\_len}\NormalTok{(}\FunctionTok{nrow}\NormalTok{(nn)), }\AttributeTok{times =}\NormalTok{ k),}
                                      \FunctionTok{as.vector}\NormalTok{(nn[, }\FunctionTok{seq\_len}\NormalTok{(k)]))] }\SpecialCharTok{\textgreater{}}\NormalTok{ Lfar),}
         \FunctionTok{numeric}\NormalTok{(}\DecValTok{1}\NormalTok{))}
\NormalTok{\}}
\NormalTok{counts }\OtherTok{\textless{}{-}} \FunctionTok{t}\NormalTok{(}\FunctionTok{vapply}\NormalTok{(th\_fine, }\ControlFlowTok{function}\NormalTok{(t)}
  \FunctionTok{rowMeans}\NormalTok{(}\FunctionTok{vapply}\NormalTok{(BENCH\_SEEDS, }\ControlFlowTok{function}\NormalTok{(sd) }\FunctionTok{jumps}\NormalTok{(t, sd), }\FunctionTok{numeric}\NormalTok{(}\FunctionTok{length}\NormalTok{(KS)))),}
  \FunctionTok{numeric}\NormalTok{(}\FunctionTok{length}\NormalTok{(KS))))}
\NormalTok{control }\OtherTok{\textless{}{-}} \FunctionTok{rowMeans}\NormalTok{(}\FunctionTok{vapply}\NormalTok{(BENCH\_SEEDS, }\ControlFlowTok{function}\NormalTok{(sd) }\FunctionTok{jumps}\NormalTok{(}\DecValTok{0}\NormalTok{, sd),}
                           \FunctionTok{numeric}\NormalTok{(}\FunctionTok{length}\NormalTok{(KS))))}
\NormalTok{measured }\OtherTok{\textless{}{-}} \FunctionTok{vapply}\NormalTok{(}\FunctionTok{seq\_along}\NormalTok{(KS), }\ControlFlowTok{function}\NormalTok{(a) \{}
\NormalTok{  i }\OtherTok{\textless{}{-}} \FunctionTok{which}\NormalTok{(counts[, a] }\SpecialCharTok{\textgreater{}=} \DecValTok{1}\NormalTok{)[}\DecValTok{1}\NormalTok{]}
  \ControlFlowTok{if}\NormalTok{ (}\FunctionTok{is.na}\NormalTok{(i)) }\ConstantTok{NA\_real\_}
  \ControlFlowTok{else} \ControlFlowTok{if}\NormalTok{ (i }\SpecialCharTok{==} \DecValTok{1}\DataTypeTok{L}\NormalTok{) th\_fine[}\DecValTok{1}\NormalTok{]}
  \ControlFlowTok{else}\NormalTok{ stats}\SpecialCharTok{::}\FunctionTok{approx}\NormalTok{(counts[(i }\SpecialCharTok{{-}} \DecValTok{1}\NormalTok{)}\SpecialCharTok{:}\NormalTok{i, a], th\_fine[(i }\SpecialCharTok{{-}} \DecValTok{1}\NormalTok{)}\SpecialCharTok{:}\NormalTok{i], }\AttributeTok{xout =} \DecValTok{1}\NormalTok{)}\SpecialCharTok{$}\NormalTok{y}
\NormalTok{\}, }\FunctionTok{numeric}\NormalTok{(}\DecValTok{1}\NormalTok{))}
\NormalTok{onset }\OtherTok{\textless{}{-}} \FunctionTok{data.frame}\NormalTok{(}\AttributeTok{k =}\NormalTok{ KS, }\AttributeTok{r\_k =}\NormalTok{ r\_k, }\AttributeTok{predicted =}\NormalTok{ predicted,}
                    \AttributeTok{measured =}\NormalTok{ measured, }\AttributeTok{at\_flat =}\NormalTok{ control)}
\end{Highlighting}
\end{Shaded}

An edge counts as a fold jump when its endpoints are more than \texttt{Lfar} apart in
the chart, which is \texttt{branch\_gap()}'s own definition of ``far apart along the
surface'' reused so that the chapter and the difficulty statistic mean the same
thing. The onset is where the \emph{expected} number of such edges reaches one,
because the first jump is a tail event and one edge in one seed is not an onset.

The fold sweep here runs past \texttt{max(THETA\_GRID)}, deliberately and uniquely in
this book. Whether the short circuit arrives before or after the benchmark's own
ceiling is exactly the question, so the measurement may not be confined to the
answer.

The criterion calibrates itself. At \(\theta = 0\) the sheet is flat, ambient and
chart distance coincide, and any edge counted there is sparse sampling rather
than folding. That control is 0 at every \(k\) tested, so the
counts below have no false positives in them.

\begin{figure}

{\centering \includegraphics[width=0.9\linewidth,alt={Two panels sharing a horizontal axis that runs from fold angle 0.7 to 0.99, with a vertical dashed line at 0.95 marking the benchmark grid's ceiling. Panel A: a falling curve, the facet clearance, from about 0.74 down to about 0.19, crossed by four horizontal dashed lines at heights 0.29, 0.41, 0.51 and 0.63 for neighbourhood sizes 10, 20, 30 and 45. The three upper lines cross the curve before the vertical marker; the lowest, for the smallest neighbourhood, crosses it after. Panel B: four rising curves of fold-jump edge counts on a logarithmic vertical axis, one per neighbourhood size, each leaving zero close to where panel A predicts and then climbing steeply; the curve for the smallest neighbourhood is flat at zero until past the vertical marker.}]{05-geodesic-methods_files/figure-latex/geo-onset-fig-1} 

}

\caption{The analytic onset against the measured one, over a fold sweep deliberately extended past the benchmark's own ceiling — the vertical dashed line at the largest fold angle the grid uses. Panel A is the exact clearance between non-touching facets as the sheet closes, with the chart $k$-nearest-neighbour radius for four neighbourhood sizes drawn as horizontal lines; each crossing is a predicted onset, marked with an open circle. Panel B is the measured number of fold-jump edges per sample, averaged over the twenty benchmark seeds, with the same predictions marked on the horizontal axis.}\label{fig:geo-onset-fig}
\end{figure}

\begin{table}

\caption{\label{tab:geo-onset-tab}Predicted and measured onset of the first fold jump.}
\centering
\begin{tabular}[t]{r|r|r|r|r}
\hline
k & chart radius & predicted onset & measured onset & edges at flat\\
\hline
10 & 0.292 & 0.974 & 0.971 & 0\\
\hline
20 & 0.413 & 0.945 & 0.944 & 0\\
\hline
30 & 0.507 & 0.909 & 0.914 & 0\\
\hline
45 & 0.629 & 0.832 & 0.855 & 0\\
\hline
\end{tabular}
\end{table}

Table \ref{tab:geo-onset-tab} is the chapter's headline promise, kept. The
analytic onset and the measured one agree to
0.023 in \(\theta\) at worst ---
under half a step of the benchmark's own grid --- over four neighbourhood sizes
spanning a factor of 4.5. The residual has the sign it should:
\texttt{facet\_gap()} is a minimum over the whole of two facets, and a finite sample does
not put points exactly at the closest-approach spot, so the measurement is late.

The row that matters is the first one, and it is a finding about the benchmark
rather than about Isomap.

\begin{boximportant}
\textbf{The committed grid never enters the short-circuit regime.} At
\(k =\) \texttt{K\_DEFAULT} --- the neighbourhood size every cell of
\texttt{benchmark-grid.rds} runs at --- the predicted onset lies past
\texttt{max(THETA\_GRID)}, and the measured count inside the sweep never
leaves zero. The classical failure that geodesic methods are known for
is not what this benchmark, as committed, measures.
\end{boximportant}

The count inside the committed sweep is
0 at every fold angle
in it, over all 20 seeds. The two boundaries are independent, which is
what makes that worth saying rather than assuming. \texttt{THETA\_GRID} stops where it does because at \(\theta = 1\) the sheet
has no extent in one direction and every method would be embedding a collapsed
object --- a reason with nothing to do with neighbourhood graphs (\texttt{R/constants.R}
states it). The short circuit at \(k =\) \texttt{K\_DEFAULT} arrives at
0.974 predicted and
0.971 measured, which is
0.021 past that ceiling. Nothing was
arranged; the grid simply stops just short.

So everything in Figure \ref{fig:geo-curves-fig} --- the whole of Isomap's
degradation from 0.024 to
0.188 across the sweep --- happened without
a single fold jump. Raising \(k\) is what crosses the boundary inside the grid, and
that is the subject of the next section.

\section{Diagnostics}\label{geo-diagnostics}

If the short circuit is absent, what is doing the damage? The \(k\) sweep answers
it, and the answer is that a neighbourhood size is not a setting to be tuned once.

\begin{Shaded}
\begin{Highlighting}[]
\NormalTok{sw   }\OtherTok{\textless{}{-}} \FunctionTok{aggregate}\NormalTok{(}\FunctionTok{cbind}\NormalTok{(rmse, qnx) }\SpecialCharTok{\textasciitilde{}}\NormalTok{ method }\SpecialCharTok{+}\NormalTok{ k\_tck }\SpecialCharTok{+}\NormalTok{ theta,}
                  \FunctionTok{subset}\NormalTok{(sweeps, ran), mean)}
\NormalTok{best }\OtherTok{\textless{}{-}} \FunctionTok{do.call}\NormalTok{(rbind, }\FunctionTok{lapply}\NormalTok{(}\FunctionTok{split}\NormalTok{(}\FunctionTok{subset}\NormalTok{(sw, method }\SpecialCharTok{==} \StringTok{"isomap"}\NormalTok{),}
                                    \FunctionTok{subset}\NormalTok{(sw, method }\SpecialCharTok{==} \StringTok{"isomap"}\NormalTok{)}\SpecialCharTok{$}\NormalTok{theta),}
                              \ControlFlowTok{function}\NormalTok{(d) d[}\FunctionTok{which.min}\NormalTok{(d}\SpecialCharTok{$}\NormalTok{rmse), ]))}
\NormalTok{dif  }\OtherTok{\textless{}{-}} \FunctionTok{subset}\NormalTok{(sw, method }\SpecialCharTok{==} \StringTok{"diffusion"}\NormalTok{)}
\NormalTok{iso  }\OtherTok{\textless{}{-}} \FunctionTok{subset}\NormalTok{(sw, method }\SpecialCharTok{==} \StringTok{"isomap"}\NormalTok{)}
\NormalTok{ths  }\OtherTok{\textless{}{-}} \FunctionTok{sort}\NormalTok{(}\FunctionTok{unique}\NormalTok{(sweeps}\SpecialCharTok{$}\NormalTok{theta))}
\NormalTok{kmax }\OtherTok{\textless{}{-}} \FunctionTok{max}\NormalTok{(sweeps}\SpecialCharTok{$}\NormalTok{k\_tck)}
\NormalTok{kbest }\OtherTok{\textless{}{-}} \ControlFlowTok{function}\NormalTok{(t) best}\SpecialCharTok{$}\NormalTok{k\_tck[best}\SpecialCharTok{$}\NormalTok{theta }\SpecialCharTok{==}\NormalTok{ t]}
\NormalTok{err   }\OtherTok{\textless{}{-}} \ControlFlowTok{function}\NormalTok{(t, k) iso}\SpecialCharTok{$}\NormalTok{rmse[iso}\SpecialCharTok{$}\NormalTok{theta }\SpecialCharTok{==}\NormalTok{ t }\SpecialCharTok{\&}\NormalTok{ iso}\SpecialCharTok{$}\NormalTok{k\_tck }\SpecialCharTok{==}\NormalTok{ k]}
\NormalTok{carry }\OtherTok{\textless{}{-}} \ControlFlowTok{function}\NormalTok{(from, to) }\FunctionTok{err}\NormalTok{(to, }\FunctionTok{kbest}\NormalTok{(from)) }\SpecialCharTok{/} \FunctionTok{err}\NormalTok{(to, }\FunctionTok{kbest}\NormalTok{(to))}
\end{Highlighting}
\end{Shaded}

The neighbourhood size that minimises Isomap's reconstruction error falls as the
sheet closes: 45 at \(\theta = 0.2\), 20 at
0.5, and 12 at 0.8. Neither end's choice survives
the other. Carrying the shallow optimum to the deepest fold multiplies the error
by 1.51; carrying the deep optimum back to the
shallowest multiplies it by 3.81. There is no one
\(k\) for this pattern, and the relative cost of the wrong one is larger at the
shallow end --- where the absolute error is smaller by an order of magnitude, so
the two ends are not simply harder and easier versions of one problem.

\begin{Shaded}
\begin{Highlighting}[]
\NormalTok{KB   }\OtherTok{\textless{}{-}} \FunctionTok{c}\NormalTok{(}\DecValTok{5}\DataTypeTok{L}\NormalTok{, }\DecValTok{10}\DataTypeTok{L}\NormalTok{, }\DecValTok{20}\DataTypeTok{L}\NormalTok{, }\DecValTok{45}\DataTypeTok{L}\NormalTok{)}
\NormalTok{bias }\OtherTok{\textless{}{-}} \FunctionTok{do.call}\NormalTok{(rbind, }\FunctionTok{lapply}\NormalTok{(ths, }\ControlFlowTok{function}\NormalTok{(t) \{}
\NormalTok{  st }\OtherTok{\textless{}{-}} \FunctionTok{sample\_manifold}\NormalTok{(p, }\AttributeTok{theta =}\NormalTok{ t, }\AttributeTok{n =}\NormalTok{ N\_DEFAULT, }\AttributeTok{seed =}\NormalTok{ BENCH\_SEEDS[}\DecValTok{1}\NormalTok{])}
\NormalTok{  du }\OtherTok{\textless{}{-}} \FunctionTok{reference\_dist}\NormalTok{(st, }\StringTok{"chart"}\NormalTok{); u }\OtherTok{\textless{}{-}} \FunctionTok{upper.tri}\NormalTok{(du)}
  \FunctionTok{do.call}\NormalTok{(rbind, }\FunctionTok{lapply}\NormalTok{(KB, }\ControlFlowTok{function}\NormalTok{(k)}
    \FunctionTok{data.frame}\NormalTok{(}\AttributeTok{theta =}\NormalTok{ t, }\AttributeTok{k =}\NormalTok{ k,}
               \AttributeTok{bias =} \FunctionTok{median}\NormalTok{(((}\FunctionTok{reference\_dist}\NormalTok{(st, }\StringTok{"graph"}\NormalTok{, }\AttributeTok{k =}\NormalTok{ k) }\SpecialCharTok{{-}}\NormalTok{ du) }\SpecialCharTok{/}\NormalTok{ du)[u]))))}
\NormalTok{\}))}
\NormalTok{zero }\OtherTok{\textless{}{-}} \FunctionTok{vapply}\NormalTok{(}\FunctionTok{split}\NormalTok{(bias, bias}\SpecialCharTok{$}\NormalTok{theta), }\ControlFlowTok{function}\NormalTok{(d)}
  \FunctionTok{tryCatch}\NormalTok{(}\FunctionTok{exp}\NormalTok{(stats}\SpecialCharTok{::}\FunctionTok{approx}\NormalTok{(d}\SpecialCharTok{$}\NormalTok{bias, }\FunctionTok{log}\NormalTok{(d}\SpecialCharTok{$}\NormalTok{k), }\AttributeTok{xout =} \DecValTok{0}\NormalTok{)}\SpecialCharTok{$}\NormalTok{y), }\AttributeTok{error =} \ControlFlowTok{function}\NormalTok{(e) }\ConstantTok{NA\_real\_}\NormalTok{),}
  \FunctionTok{numeric}\NormalTok{(}\DecValTok{1}\NormalTok{))}
\end{Highlighting}
\end{Shaded}

\begin{figure}

{\centering \includegraphics[width=0.9\linewidth,alt={Two panels, both with neighbourhood size on a logarithmic horizontal axis from five to seventy. Panel A: three U-shaped curves of reconstruction error. The lowest, at the smallest fold angle, falls steeply and bottoms out near k = 45. The middle bottoms out near k = 20 and rises again. The highest, at the largest fold angle, bottoms out near k = 12 and rises steeply after it. Panel B: three curves falling from positive values at k = 5 through a horizontal zero line to negative values at k = 45; the curve for the largest fold angle crosses zero first, near k = 7, and ends lowest, near minus 0.2.}]{05-geodesic-methods_files/figure-latex/geo-k-fig-1} 

}

\caption{Why the optimum moves. Panel A is Isomap's reconstruction error against neighbourhood size at three fold angles, from the committed sweep; the open circles mark each curve's minimum. Panel B is the median relative departure of the shortest-path distance from the exact geodesic, measured at the same three fold angles: positive is the discretisation bias of walking along edges, negative is the fold contracting them, and each curve crosses zero at a smaller neighbourhood size than the last.}\label{fig:geo-k-fig}
\end{figure}

Figure \ref{fig:geo-k-fig} puts the mechanism next to its consequence. Panel B
is the two biases of Figure \ref{fig:geo-graph-fig} as functions of \(k\). Raising
\(k\) shortens the zigzag and kills the positive bias; it also lengthens every edge,
and a longer edge crosses more creases and is contracted by more of them. The
crossing point --- where the median shortest path is neither long nor short --- falls
from \(k \approx 35\) at \(\theta = 0.2\) to
\(14\) at 0.5 and
\(6.8\) at 0.8, in the same order and over nearly the same
range as the error minima in panel A. The two do not coincide, and should not:
panel A minimises a Procrustes fit that pays for variance as well as bias.

Diffusion maps has no such trade-off. Its error falls monotonically in \(k\) at all
three fold angles, from
0.208 at
\(k = 5\) to
0.132 at
\(k = 70\) in the middle of the sweep, because \(k\) there sets a kernel
bandwidth and not a graph. It never approaches Isomap's optimum, and its best
setting is the largest one offered --- a curve still falling at the edge of the
sweep is a sweep that stopped too early, and this one is reported as such.

Two housekeeping facts from the artefacts, both of which would have been
warnings scrolling past in a long run. Isomap's \texttt{k\_effective} column equals the
requested \(k\) in every one of the 1200 main-grid rows
and every sweep row, so the neighbourhood graph never disconnected and never
needed repairing. Diffusion maps needed a widened kernel in
399 of
1200 cells and declined outright in
1; those are recorded in the
\texttt{tuning} and \texttt{status} columns rather than folded into the mean.

\section{Where it fails}\label{geo-limits}

\textbf{The predicted failure is not the observed one.} This chapter was specified to
pair an analytic short-circuit onset with a measured one, and it does --- to
0.023 in \(\theta\). What the
specification also assumed, and what is false, is that the short circuit is what
breaks Isomap on this object. At the grid's neighbourhood size it never occurs
inside the sweep, and Isomap's error still rises by a factor of
7.82 across it. The
mechanism that actually operates is edge contraction: present at every
\(\theta > 0\), graded rather than catastrophic, and invisible to any diagnostic
built to look for a bridge. A reader who arrives expecting Isomap to collapse
here will not see it collapse --- it is the best method on this grid above
\(\theta = 0.4\) --- and the reason is that the benchmark's difficulty
axis stops short of the regime that would do it.

\textbf{When the bridge does form, it is not graceful.} Shortest path takes a \emph{minimum}
over routes, so a single bad edge is not averaged against the good ones --- it is
preferred by every pair whose route it shortens.

\begin{Shaded}
\begin{Highlighting}[]
\NormalTok{z   }\OtherTok{\textless{}{-}}\NormalTok{ pair[[}\DecValTok{2}\NormalTok{]]                                  }\CommentTok{\# the deeply folded sample}
\NormalTok{dAf }\OtherTok{\textless{}{-}} \FunctionTok{as.matrix}\NormalTok{(}\FunctionTok{dist}\NormalTok{(z}\SpecialCharTok{$}\NormalTok{s}\SpecialCharTok{$}\NormalTok{X))}
\NormalTok{adj }\OtherTok{\textless{}{-}} \FunctionTok{.knn\_graph}\NormalTok{(dAf, K\_DEFAULT)                 }\CommentTok{\# reference\_dist()\textquotesingle{}s internals,}
\NormalTok{G0  }\OtherTok{\textless{}{-}} \FunctionTok{.apsp\_minplus}\NormalTok{(adj, dAf)                    }\CommentTok{\#   reached directly so that the}
\NormalTok{far }\OtherTok{\textless{}{-}} \FunctionTok{as.matrix}\NormalTok{(}\FunctionTok{dist}\NormalTok{(z}\SpecialCharTok{$}\NormalTok{s}\SpecialCharTok{$}\NormalTok{truth)) }\SpecialCharTok{\textgreater{}}\NormalTok{ Lfar          }\CommentTok{\#   graph can be edited between}

\NormalTok{worst }\OtherTok{\textless{}{-}}\NormalTok{ dAf; worst[}\SpecialCharTok{!}\NormalTok{far] }\OtherTok{\textless{}{-}} \ConstantTok{Inf}
\NormalTok{ij  }\OtherTok{\textless{}{-}} \FunctionTok{arrayInd}\NormalTok{(}\FunctionTok{which.min}\NormalTok{(worst), }\FunctionTok{dim}\NormalTok{(worst))}
\NormalTok{edges }\OtherTok{\textless{}{-}} \FunctionTok{sum}\NormalTok{(adj) }\SpecialCharTok{/} \DecValTok{2}
\NormalTok{adj[ij[}\DecValTok{1}\NormalTok{], ij[}\DecValTok{2}\NormalTok{]] }\OtherTok{\textless{}{-}}\NormalTok{ adj[ij[}\DecValTok{2}\NormalTok{], ij[}\DecValTok{1}\NormalTok{]] }\OtherTok{\textless{}{-}} \ConstantTok{TRUE}
\NormalTok{G1  }\OtherTok{\textless{}{-}} \FunctionTok{.apsp\_minplus}\NormalTok{(adj, dAf)}

\NormalTok{u2      }\OtherTok{\textless{}{-}} \FunctionTok{upper.tri}\NormalTok{(G0)}
\NormalTok{touched }\OtherTok{\textless{}{-}} \FunctionTok{mean}\NormalTok{((G1 }\SpecialCharTok{\textless{}}\NormalTok{ G0 }\SpecialCharTok{{-}}\NormalTok{ TOL}\SpecialCharTok{$}\NormalTok{default)[u2])}
\NormalTok{worst\_p }\OtherTok{\textless{}{-}} \FunctionTok{max}\NormalTok{(((G0 }\SpecialCharTok{{-}}\NormalTok{ G1) }\SpecialCharTok{/}\NormalTok{ G0)[u2])}
\NormalTok{before  }\OtherTok{\textless{}{-}} \FunctionTok{reconstruction\_error}\NormalTok{(}\FunctionTok{cmdscale}\NormalTok{(}\FunctionTok{as.dist}\NormalTok{(G0), }\AttributeTok{k =}\NormalTok{ EMBED\_DIM), z}\SpecialCharTok{$}\NormalTok{s}\SpecialCharTok{$}\NormalTok{truth)}
\NormalTok{after   }\OtherTok{\textless{}{-}} \FunctionTok{reconstruction\_error}\NormalTok{(}\FunctionTok{cmdscale}\NormalTok{(}\FunctionTok{as.dist}\NormalTok{(G1), }\AttributeTok{k =}\NormalTok{ EMBED\_DIM), z}\SpecialCharTok{$}\NormalTok{s}\SpecialCharTok{$}\NormalTok{truth)}
\end{Highlighting}
\end{Shaded}

One edge is added to a graph that already has 4,692
of them --- the single closest ambient approach between two points the chart holds
far apart, which is exactly the pair \texttt{branch\_gap()} reports. It shortens the
shortest path for
5.15 per cent of all pairs, by up to
38.9 per cent, and raises the reconstruction error from
0.0621 to 0.0767. Nothing about the sample changed;
one entry of an adjacency matrix did.

\begin{figure}

{\centering \includegraphics[width=0.9\linewidth,alt={Two square scatter panels. In panel A the embedded points form an even rectangular patch with five grey level-set paths running across it as separate parallel curves. In panel B the patch is a similar rectangle but slightly bowed at its ends and the grey paths, still separate and ordered, are unevenly spaced and drawn closer together near the middle.}]{05-geodesic-methods_files/figure-latex/geo-embed-fig-1} 

}

\caption{What Isomap returns, coloured by the exact chart coordinate along the corrugation, with grey level sets of that coordinate drawn over the points. Panel A is the lightly folded sheet and panel B the deeply folded one. Both recover the chart; the damage in B is a crowding of the level sets, not a tear.}\label{fig:geo-embed-fig}
\end{figure}

Figure \ref{fig:geo-embed-fig} is what those errors look like. Neither panel is
a failure a reader would spot by eye, which is the argument for having an answer
key: the deeply folded embedding is 2.6
times worse than the lightly folded one and looks like a slightly wobblier
version of it.

\textbf{The estimate has no guarantee here.} The convergence of graph distances to
geodesics is stated for a manifold of positive reach and a density that grows
\citep{bernstein2000graph, niyogi2008homology}. A crease has reach zero and this book
fixes \(n\) at \texttt{N\_DEFAULT}. What replaces the guarantee is the measurement above,
and that is a weaker thing than a theorem.

\textbf{One family, one size, one density.} Every onset in Table
\ref{tab:geo-onset-tab} is a statement about a six-by-six Miura at
\texttt{N\_DEFAULT} points. The radius \(r_k\) falls as \(\sqrt{k/n}\), so the same
prediction at four times the sample size moves every onset later, and the sweep
that would establish that has not been run.

\section{Reproduce this}\label{geo-reproduce}

Both artefacts are committed and neither is generated at render time.
\texttt{scripts/run-benchmark-grid.R} produces \texttt{benchmark-grid.rds} and
\texttt{scripts/run-part2-sweeps.R} produces \texttt{part2-sweeps.rds}; each declares its
artefact path in its header, and \texttt{scripts/check-artefact-producers.R} checks that
every artefact the book specifies has one. The digests in \texttt{geo-setup} are what
stop these sentences from being read against a different run.

The claims the chapter rests on are assertions in the test suite:

\begin{itemize}
\tightlist
\item
  \texttt{tests/testthat/test-folding.R} --- facet isometry across the sweep, and
  \texttt{facet\_gap()} against a brute-force point sample.
\item
  \texttt{tests/testthat/test-metrics.R} --- \texttt{reference\_dist()}'s three geometries, the
  neighbourhood graph's connectivity repair, and the shortest paths against a
  reference implementation.
\item
  \texttt{tests/testthat/test-methods.R} --- the registry contract, including that a
  stochastic method called without a seed is an error.
\end{itemize}

Run them with \texttt{Rscript\ tests/testthat.R} from the repository root.

\section{Further reading}\label{geo-reading}

Tenenbaum, de Silva and Langford \citep{tenenbaum2000isomap} is Isomap's source and
still the clearest statement of the three-stage construction. The convergence of
graph distances to geodesics is Bernstein and co-authors
\citep{bernstein2000graph}; read it together with Balasubramanian and Schwartz's
comment \citep{balasubramanian2002isomap}, which is where the topological instability
this chapter measures was first named, and with Niyogi, Smale and Weinberger
\citep{niyogi2008homology} for what a sampling guarantee needs in order to exist.

Coifman and Lafon \citep{coifman2006diffusion} introduce diffusion maps and the
diffusion distance the embedding approximates; Belkin and Niyogi
\citep{belkin2003laplacian} is the closely related Laplacian construction Chapter
\ref{neighbor} takes up. Lee and Verleysen \citep{lee2007nonlinear} set out the whole
family as one kernel eigenproblem, which is the framing that makes classical
scaling \citep{torgerson1952mds}, kernel PCA \citep{scholkopf1998kernelpca} and Isomap
three settings of one method rather than three methods. Dijkstra
\citep{dijkstra1959note} is where shortest paths on a weighted graph begin, though
\texttt{R/metrics.R} does not use his algorithm and says why.

\chapter{Neighbour embeddings}\label{neighbor}

\emph{R. Heller · Date: 2026-07-29 · Last changed: 2026-09-07}

\emph{Why do t-SNE and UMAP disagree about the same data?}

\section{What this chapter answers}\label{nbr-question}

The four methods in this chapter --- LLE, Laplacian eigenmaps, t-SNE and UMAP ---
do not try to recover the unfolding. That sentence has to come before any
number, because the rest of the chapter scores them against the unfolding and a
reader who meets the scores first will read them as a verdict on the methods
rather than on the pairing.

What each of them minimises is a discrepancy between \textbf{neighbourhoods}: which
points are near which, in the ambient cloud and in the picture. None of them
carries a term that asks for a distance to be reproduced. The registry records
this as a field rather than as a remark --- \texttt{consumes\ =\ "neighbourhood"} --- so a
later chapter cannot forget it. \texttt{reconstruction\_error()} is a Procrustes
distance to the chart, and asking these four to be small on it is asking them
for something they were never built to deliver.

Then this chapter judges them by it anyway, for two reasons. The reader will:
every one of these methods is used in practice as though its output were a map,
and a map is a claim about distance. And the benchmark can, which is the rarer
thing --- the answer key exists in closed form, so the question ``how much of the
geometry survived?'' has an exact answer rather than an argument. What comes out
is not that the methods are bad. It is that the setting on which they are most
often tuned, and the picture that tuning produces, move in the \emph{opposite}
direction to accuracy.

\section{Setup}\label{nbr-setup}

\begin{Shaded}
\begin{Highlighting}[]
\NormalTok{grid   }\OtherTok{\textless{}{-}} \FunctionTok{read\_run}\NormalTok{(}\StringTok{"data/processed/benchmark{-}grid.rds"}\NormalTok{, }\StringTok{"2483a468d82d4d64"}\NormalTok{)}
\NormalTok{sweeps }\OtherTok{\textless{}{-}} \FunctionTok{read\_run}\NormalTok{(}\StringTok{"data/processed/part2{-}sweeps.rds"}\NormalTok{,   }\StringTok{"16d43ff4992c0e58"}\NormalTok{)}
\NormalTok{budget }\OtherTok{\textless{}{-}} \FunctionTok{read\_run}\NormalTok{(}\StringTok{"data/processed/seed{-}budget.rds"}\NormalTok{,    }\StringTok{"e5f73dce72d53066"}\NormalTok{)}

\NormalTok{NBR  }\OtherTok{\textless{}{-}} \FunctionTok{c}\NormalTok{(}\StringTok{"lle"}\NormalTok{, }\StringTok{"laplacian"}\NormalTok{, }\StringTok{"tsne"}\NormalTok{, }\StringTok{"umap"}\NormalTok{)}
\NormalTok{pat  }\OtherTok{\textless{}{-}} \FunctionTok{miura\_ori}\NormalTok{(}\AttributeTok{nx =} \DecValTok{6}\DataTypeTok{L}\NormalTok{, }\AttributeTok{ny =} \DecValTok{6}\DataTypeTok{L}\NormalTok{)}
\NormalTok{cell }\OtherTok{\textless{}{-}} \FunctionTok{sample\_manifold}\NormalTok{(pat, }\AttributeTok{theta =} \FloatTok{0.5}\NormalTok{, }\AttributeTok{n =}\NormalTok{ N\_DEFAULT, }\AttributeTok{seed =}\NormalTok{ BENCH\_SEEDS[}\DecValTok{1}\NormalTok{])}
\end{Highlighting}
\end{Shaded}

Three committed artefacts. \texttt{benchmark-grid.rds} sweeps \(\theta\) at fixed
neighbourhood size; \texttt{part2-sweeps.rds} sweeps the neighbourhood size at three
fixed \(\theta\); \texttt{seed-budget.rds} is the measurement standing rule 1's seed
floor is checked against. \texttt{cell} is one sample, drawn once, used for every
figure that shows an embedding rather than a curve.

\section{Background}\label{nbr-background}

\begin{Shaded}
\begin{Highlighting}[]
\NormalTok{knitr}\SpecialCharTok{::}\FunctionTok{kable}\NormalTok{(}
  \FunctionTok{method\_table}\NormalTok{()[}\FunctionTok{method\_table}\NormalTok{()}\SpecialCharTok{$}\NormalTok{chapter }\SpecialCharTok{==} \DecValTok{6}\DataTypeTok{L}\NormalTok{,}
                 \FunctionTok{c}\NormalTok{(}\StringTok{"label"}\NormalTok{, }\StringTok{"consumes"}\NormalTok{, }\StringTok{"stochastic"}\NormalTok{)],}
  \AttributeTok{row.names =} \ConstantTok{FALSE}\NormalTok{,}
  \AttributeTok{caption =} \StringTok{"The four registry entries this chapter describes."}\NormalTok{)}
\end{Highlighting}
\end{Shaded}

\begin{table}

\caption{\label{tab:nbr-registry}The four registry entries this chapter describes.}
\centering
\begin{tabular}[t]{l|l|l}
\hline
label & consumes & stochastic\\
\hline
LLE & neighbourhood & FALSE\\
\hline
Laplacian eigenmaps & neighbourhood & FALSE\\
\hline
t-SNE & neighbourhood & TRUE\\
\hline
UMAP & neighbourhood & TRUE\\
\hline
\end{tabular}
\end{table}

\textbf{LLE} \citep{roweis2000lle} writes each point as a weighted combination of its \(k\)
nearest neighbours and then finds the low-dimensional configuration that keeps
those weights. \textbf{Laplacian eigenmaps} \citep{belkin2003laplacian} takes the bottom
eigenvectors of the graph Laplacian of a heat-kernel-weighted neighbourhood
graph. Both are spectral, deterministic, and solve a single eigenproblem.

\textbf{t-SNE} \citep{vandermaaten2008tsne} converts ambient distances into a
neighbour-probability distribution, converts embedding distances into a second
one with a heavy-tailed kernel, and minimises the Kullback--Leibler divergence
between them by gradient descent from a random initialisation. \textbf{UMAP}
\citep{mcinnes2018umap} builds a fuzzy simplicial set on the neighbourhood graph and
minimises a cross-entropy, also from a random start. Both are stochastic, and
the registry refuses to run them unseeded:

\begin{Shaded}
\begin{Highlighting}[]
\NormalTok{unseeded }\OtherTok{\textless{}{-}} \FunctionTok{tryCatch}\NormalTok{(}\FunctionTok{embed}\NormalTok{(}\StringTok{"tsne"}\NormalTok{, cell), }\AttributeTok{error =}\NormalTok{ conditionMessage)}
\FunctionTok{cat}\NormalTok{(}\FunctionTok{strwrap}\NormalTok{(unseeded, }\AttributeTok{width =} \DecValTok{72}\NormalTok{), }\AttributeTok{sep =} \StringTok{"}\SpecialCharTok{\textbackslash{}n}\StringTok{"}\NormalTok{)}
\end{Highlighting}
\end{Shaded}

\begin{verbatim}
#> tsne is stochastic and was called without a seed. Standing rule 1:
#> every stochastic result in this book is reported across BENCH_SEEDS,
#> and an unseeded fit cannot be one of them.
\end{verbatim}

That error is deliberate. Standing rule 1 says every stochastic result in this
book is reported across \texttt{BENCH\_SEEDS}, and an unseeded fit cannot be one of
them --- it is an anecdote with no way of being labelled as one afterwards. The
cheapest place to enforce a reporting rule is at the call that would break it.

\begin{boximportant}
Two neighbourhood parameters, two names, one quantity. t-SNE's is
\emph{perplexity}, a smooth measure of the effective number of
neighbours; UMAP's is \texttt{n\_neighbors}. \texttt{R/methods.R} maps
the registry's \texttt{k} onto both, so a sweep over \texttt{k} is a
sweep over perplexity for one and over \texttt{n\_neighbors} for the
other, and the curves are comparable. A previous version of
\texttt{embed\_tsne()} declared \texttt{k} and passed a fixed perplexity
--- every \texttt{k} returned the same embedding, and a flat curve would
have been read as insensitivity to neighbourhood size.
\end{boximportant}

\section{What the methods produce, and what it costs}\label{nbr-core}

\begin{Shaded}
\begin{Highlighting}[]
\NormalTok{fit }\OtherTok{\textless{}{-}} \FunctionTok{lapply}\NormalTok{(NBR, }\ControlFlowTok{function}\NormalTok{(m)}
  \FunctionTok{embed}\NormalTok{(m, cell, }\AttributeTok{k =}\NormalTok{ K\_DEFAULT,}
        \AttributeTok{seed =} \ControlFlowTok{if}\NormalTok{ (}\FunctionTok{isTRUE}\NormalTok{(METHOD\_REGISTRY[[m]]}\SpecialCharTok{$}\NormalTok{stochastic)) BENCH\_SEEDS[}\DecValTok{1}\NormalTok{]))}
\FunctionTok{names}\NormalTok{(fit) }\OtherTok{\textless{}{-}}\NormalTok{ NBR}

\NormalTok{Dchart   }\OtherTok{\textless{}{-}} \FunctionTok{reference\_dist}\NormalTok{(cell, }\StringTok{"chart"}\NormalTok{)}
\NormalTok{cell\_err }\OtherTok{\textless{}{-}} \FunctionTok{vapply}\NormalTok{(fit, reconstruction\_error, }\FunctionTok{numeric}\NormalTok{(}\DecValTok{1}\NormalTok{), }\AttributeTok{truth =}\NormalTok{ cell}\SpecialCharTok{$}\NormalTok{truth)}
\NormalTok{cell\_qnx }\OtherTok{\textless{}{-}} \FunctionTok{vapply}\NormalTok{(fit, qnx, }\FunctionTok{numeric}\NormalTok{(}\DecValTok{1}\NormalTok{), }\AttributeTok{reference =}\NormalTok{ Dchart, }\AttributeTok{K =}\NormalTok{ K\_QNX)}
\end{Highlighting}
\end{Shaded}

\begin{figure}

{\centering \includegraphics[width=0.9\linewidth,alt={Four scatter panels of the same 800 points, each with its own axes. The LLE panel is an evenly filled quadrilateral with a smooth top-to-bottom colour ramp and four level-set paths that run across it without crossing. The Laplacian eigenmaps panel is a quadrilateral whose points crowd onto its edges, leaving the interior sparse. The t-SNE panel is broken into many small separated islands with the level-set paths jumping between them. The UMAP panel is broken into curved filaments and loops with the ramp ordered left to right along them.}]{06-neighbor-embeddings_files/figure-latex/nbr-fig-embeddings-1} 

}

\caption{One folded sample, embedded four ways. Colour is the first chart coordinate and the black paths are level sets of the second, so a faithful embedding shows a smooth colour ramp crossed by non-intersecting paths. The Procrustes error against the chart is printed under each panel. The axes carry no numbers because an embedding is defined only up to a similarity and its units mean nothing.}\label{fig:nbr-fig-embeddings}
\end{figure}

Figure \ref{fig:nbr-fig-embeddings} is one cell of the grid --- a Miura at
\(\theta =\) 0.5, 800 points, no noise --- seen four ways. LLE
returns the chart to within 0.092; UMAP to
0.18; Laplacian eigenmaps to
0.22; t-SNE, at this neighbourhood size, to
0.63, which is most of the way to the score a
random configuration earns.

The \(Q_{NX}\) column tells a different story about the same four pictures.
Scored at \(K =\) 20 \textbf{against the chart}, they run
0.911, 0.757,
0.764 and 0.755. t-SNE,
last by Procrustes distance by a factor of
3.5, is \emph{ahead} of UMAP on
rank preservation. That is not a contradiction and not a defect in either
metric: one reads distances and the other reads rank order, and an embedding
that shatters the sheet into islands can keep every local rank while destroying
every distance between islands.

\subsection{The deception}\label{the-deception}

The parameter these methods are tuned on is the neighbourhood size, and tuning
is done by looking. That is the part exact truth is able to expose.

\begin{Shaded}
\begin{Highlighting}[]
\NormalTok{KS       }\OtherTok{\textless{}{-}} \FunctionTok{c}\NormalTok{(}\DecValTok{5}\DataTypeTok{L}\NormalTok{, }\DecValTok{12}\DataTypeTok{L}\NormalTok{, }\DecValTok{30}\DataTypeTok{L}\NormalTok{, }\DecValTok{70}\DataTypeTok{L}\NormalTok{)}
\NormalTok{tperp    }\OtherTok{\textless{}{-}} \FunctionTok{lapply}\NormalTok{(KS, }\ControlFlowTok{function}\NormalTok{(kk) }\FunctionTok{embed}\NormalTok{(}\StringTok{"tsne"}\NormalTok{, cell, }\AttributeTok{k =}\NormalTok{ kk,}
                                          \AttributeTok{seed =}\NormalTok{ BENCH\_SEEDS[}\DecValTok{1}\NormalTok{]))}
\NormalTok{perp\_err }\OtherTok{\textless{}{-}} \FunctionTok{vapply}\NormalTok{(tperp, reconstruction\_error, }\FunctionTok{numeric}\NormalTok{(}\DecValTok{1}\NormalTok{), }\AttributeTok{truth =}\NormalTok{ cell}\SpecialCharTok{$}\NormalTok{truth)}
\NormalTok{perp\_qnx }\OtherTok{\textless{}{-}} \FunctionTok{vapply}\NormalTok{(tperp, qnx, }\FunctionTok{numeric}\NormalTok{(}\DecValTok{1}\NormalTok{), }\AttributeTok{reference =}\NormalTok{ Dchart, }\AttributeTok{K =}\NormalTok{ K\_QNX)}
\end{Highlighting}
\end{Shaded}

\begin{figure}

{\centering \includegraphics[width=0.9\linewidth,alt={Four scatter panels of the same 800 points at perplexity 5, 12, 30 and 70. The first is a spray of small separated islands with black level-set paths crossing all over it. The second has fewer, larger islands. The third is nearly a filled quadrilateral with two indentations. The fourth is a filled quadrilateral with a smooth diagonal colour ramp and level-set paths that do not cross.}]{06-neighbor-embeddings_files/figure-latex/nbr-fig-perplexity-1} 

}

\caption{t-SNE on one sample at four perplexities. Left to right the picture loses every visible feature and gains an order of magnitude in accuracy. Colour is the first chart coordinate; the black paths are level sets of the second, and the axes are unnumbered for the same reason as in the previous figure.}\label{fig:nbr-fig-perplexity}
\end{figure}

The left-hand panel of Figure \ref{fig:nbr-fig-perplexity} is the picture
practitioners recognise and reward: crisp, well-separated islands, obvious
structure, something to point at. It scores 0.88. The
right-hand panel is a featureless filled quadrilateral with no structure to
point at whatsoever. It scores 0.089, a factor of
9.9 better, and it is featureless because
\textbf{the answer is a flat rectangular patch}. The truth here has no clusters in it. Every
island in the left-hand panel is an artefact of the setting.

\(Q_{NX}\) against the chart moves the same way --- 0.732 to
0.875 --- so this is not one metric's idiosyncrasy. What
disagrees with both is the eye. The tuning advice that circulates for these
methods is largely advice about how the picture looks
\citep{wattenberg2016tsne, kobak2019art}, and on this benchmark that advice points
away from the truth, monotonically, across the whole swept range.

\section{Results}\label{nbr-results}

\begin{Shaded}
\begin{Highlighting}[]
\NormalTok{cl }\OtherTok{\textless{}{-}}\NormalTok{ grid[grid}\SpecialCharTok{$}\NormalTok{noise }\SpecialCharTok{==} \StringTok{"none"} \SpecialCharTok{\&}\NormalTok{ grid}\SpecialCharTok{$}\NormalTok{status }\SpecialCharTok{==} \StringTok{"ok"} \SpecialCharTok{\&}
\NormalTok{           grid}\SpecialCharTok{$}\NormalTok{method }\SpecialCharTok{\%in\%} \FunctionTok{c}\NormalTok{(NBR, }\StringTok{"isomap"}\NormalTok{), ]}
\NormalTok{sp }\OtherTok{\textless{}{-}} \FunctionTok{do.call}\NormalTok{(data.frame,}
              \FunctionTok{aggregate}\NormalTok{(rmse }\SpecialCharTok{\textasciitilde{}}\NormalTok{ method }\SpecialCharTok{+}\NormalTok{ theta, cl,}
                        \ControlFlowTok{function}\NormalTok{(x) }\FunctionTok{c}\NormalTok{(}\AttributeTok{lo =} \FunctionTok{min}\NormalTok{(x), }\AttributeTok{med =} \FunctionTok{median}\NormalTok{(x),}
                                      \AttributeTok{hi =} \FunctionTok{max}\NormalTok{(x))))}
\FunctionTok{names}\NormalTok{(sp) }\OtherTok{\textless{}{-}} \FunctionTok{sub}\NormalTok{(}\StringTok{"\^{}rmse}\SpecialCharTok{\textbackslash{}\textbackslash{}}\StringTok{."}\NormalTok{, }\StringTok{""}\NormalTok{, }\FunctionTok{names}\NormalTok{(sp))}

\NormalTok{th\_max }\OtherTok{\textless{}{-}} \FunctionTok{max}\NormalTok{(THETA\_GRID)}
\NormalTok{degrade }\OtherTok{\textless{}{-}} \FunctionTok{vapply}\NormalTok{(}\FunctionTok{c}\NormalTok{(NBR, }\StringTok{"isomap"}\NormalTok{), }\ControlFlowTok{function}\NormalTok{(m) \{}
\NormalTok{  a }\OtherTok{\textless{}{-}}\NormalTok{ sp[sp}\SpecialCharTok{$}\NormalTok{method }\SpecialCharTok{==}\NormalTok{ m, ]}
\NormalTok{  a}\SpecialCharTok{$}\NormalTok{med[a}\SpecialCharTok{$}\NormalTok{theta }\SpecialCharTok{==}\NormalTok{ th\_max] }\SpecialCharTok{/}\NormalTok{ a}\SpecialCharTok{$}\NormalTok{med[a}\SpecialCharTok{$}\NormalTok{theta }\SpecialCharTok{==} \DecValTok{0}\NormalTok{]}
\NormalTok{\}, }\FunctionTok{numeric}\NormalTok{(}\DecValTok{1}\NormalTok{))}

\NormalTok{widths }\OtherTok{\textless{}{-}} \FunctionTok{transform}\NormalTok{(sp, }\AttributeTok{w =}\NormalTok{ hi }\SpecialCharTok{{-}}\NormalTok{ lo)}
\NormalTok{worst  }\OtherTok{\textless{}{-}}\NormalTok{ widths[widths}\SpecialCharTok{$}\NormalTok{method }\SpecialCharTok{==} \StringTok{"tsne"}\NormalTok{, ]}
\NormalTok{worst  }\OtherTok{\textless{}{-}}\NormalTok{ worst[}\FunctionTok{which.max}\NormalTok{(worst}\SpecialCharTok{$}\NormalTok{w), ]}

\NormalTok{out\_hi }\OtherTok{\textless{}{-}} \FunctionTok{aggregate}\NormalTok{(rmse }\SpecialCharTok{\textasciitilde{}}\NormalTok{ method,}
\NormalTok{                    grid[grid}\SpecialCharTok{$}\NormalTok{noise }\SpecialCharTok{==} \StringTok{"outlier"} \SpecialCharTok{\&}\NormalTok{ grid}\SpecialCharTok{$}\NormalTok{status }\SpecialCharTok{==} \StringTok{"ok"} \SpecialCharTok{\&}
\NormalTok{                         grid}\SpecialCharTok{$}\NormalTok{theta }\SpecialCharTok{==}\NormalTok{ th\_max, ], median)}
\NormalTok{out\_hi }\OtherTok{\textless{}{-}} \FunctionTok{setNames}\NormalTok{(out\_hi}\SpecialCharTok{$}\NormalTok{rmse, out\_hi}\SpecialCharTok{$}\NormalTok{method)}

\NormalTok{out\_hi\_mid }\OtherTok{\textless{}{-}} \FunctionTok{aggregate}\NormalTok{(rmse }\SpecialCharTok{\textasciitilde{}}\NormalTok{ method,}
\NormalTok{                        grid[grid}\SpecialCharTok{$}\NormalTok{noise }\SpecialCharTok{==} \StringTok{"outlier"} \SpecialCharTok{\&}\NormalTok{ grid}\SpecialCharTok{$}\NormalTok{status }\SpecialCharTok{==} \StringTok{"ok"} \SpecialCharTok{\&}
\NormalTok{                             grid}\SpecialCharTok{$}\NormalTok{theta }\SpecialCharTok{\textgreater{}=} \FloatTok{0.5}\NormalTok{, ], median)}
\NormalTok{out\_hi\_mid }\OtherTok{\textless{}{-}} \FunctionTok{setNames}\NormalTok{(out\_hi\_mid}\SpecialCharTok{$}\NormalTok{rmse, out\_hi\_mid}\SpecialCharTok{$}\NormalTok{method)}

\NormalTok{tsn  }\OtherTok{\textless{}{-}}\NormalTok{ sp[sp}\SpecialCharTok{$}\NormalTok{method }\SpecialCharTok{==} \StringTok{"tsne"}\NormalTok{, ]; tsn }\OtherTok{\textless{}{-}}\NormalTok{ tsn[}\FunctionTok{order}\NormalTok{(tsn}\SpecialCharTok{$}\NormalTok{theta), ]}
\NormalTok{ump  }\OtherTok{\textless{}{-}}\NormalTok{ sp[sp}\SpecialCharTok{$}\NormalTok{method }\SpecialCharTok{==} \StringTok{"umap"}\NormalTok{, ]; ump }\OtherTok{\textless{}{-}}\NormalTok{ ump[}\FunctionTok{order}\NormalTok{(ump}\SpecialCharTok{$}\NormalTok{theta), ]}
\NormalTok{gap  }\OtherTok{\textless{}{-}} \FunctionTok{abs}\NormalTok{(tsn}\SpecialCharTok{$}\NormalTok{med }\SpecialCharTok{{-}}\NormalTok{ ump}\SpecialCharTok{$}\NormalTok{med)}
\NormalTok{over }\OtherTok{\textless{}{-}} \FunctionTok{pmin}\NormalTok{(tsn}\SpecialCharTok{$}\NormalTok{hi, ump}\SpecialCharTok{$}\NormalTok{hi) }\SpecialCharTok{\textgreater{}} \FunctionTok{pmax}\NormalTok{(tsn}\SpecialCharTok{$}\NormalTok{lo, ump}\SpecialCharTok{$}\NormalTok{lo)}

\NormalTok{jump }\OtherTok{\textless{}{-}} \ControlFlowTok{function}\NormalTok{(noise) \{}
\NormalTok{  a }\OtherTok{\textless{}{-}} \FunctionTok{aggregate}\NormalTok{(qnx }\SpecialCharTok{\textasciitilde{}}\NormalTok{ theta,}
\NormalTok{                 grid[grid}\SpecialCharTok{$}\NormalTok{method }\SpecialCharTok{==} \StringTok{"lle"} \SpecialCharTok{\&}\NormalTok{ grid}\SpecialCharTok{$}\NormalTok{noise }\SpecialCharTok{==}\NormalTok{ noise }\SpecialCharTok{\&}
\NormalTok{                      grid}\SpecialCharTok{$}\NormalTok{status }\SpecialCharTok{==} \StringTok{"ok"}\NormalTok{, ], median)}
  \FunctionTok{max}\NormalTok{(}\FunctionTok{diff}\NormalTok{(a}\SpecialCharTok{$}\NormalTok{qnx[}\FunctionTok{order}\NormalTok{(a}\SpecialCharTok{$}\NormalTok{theta)]))}
\NormalTok{\}}
\NormalTok{lle\_jump }\OtherTok{\textless{}{-}} \FunctionTok{c}\NormalTok{(}\AttributeTok{none =} \FunctionTok{jump}\NormalTok{(}\StringTok{"none"}\NormalTok{), }\AttributeTok{outlier =} \FunctionTok{jump}\NormalTok{(}\StringTok{"outlier"}\NormalTok{))}
\end{Highlighting}
\end{Shaded}

The prediction, stated before the figure because a flat curve otherwise reads
as nothing happening: \textbf{these curves should be nearly flat.} Fold angle is a
difficulty parameter for a method that consumes distance, because folding
contracts ambient distance and leaves the chart alone (Chapter
\ref{folding-geometry}). It is a much weaker difficulty parameter for a method
that consumes rank order, because within a facet the fold changes no distance
at all and across a crease it changes ranks only where two facets pass close.
A method that never trusted global structure has less to lose as the global
structure degrades.

\begin{figure}

{\centering \includegraphics[width=0.9\linewidth,alt={Four line panels, one per method, sharing axes, in the order LLE, Laplacian eigenmaps, t-SNE, UMAP. The horizontal axis is the folding parameter from zero to one and the vertical axis is Procrustes error from zero to about one. LLE starts low and rises steadily, with a band that widens sharply beyond the middle of the range. Laplacian eigenmaps is almost horizontal with a narrow band. t-SNE is roughly horizontal with a band that spans most of the plot at every fold angle. UMAP is almost horizontal with a moderate band. In every panel the dashed grey Isomap line starts far below the band and rises steeply, reaching the lower edge of the band only at the right-hand end.}]{06-neighbor-embeddings_files/figure-latex/nbr-fig-theta-1} 

}

\caption{Procrustes error against the chart across the fold sweep, no noise. The band is the full spread across the twenty benchmark seeds and is the result; the solid line is its median, drawn to locate the band rather than to summarise it. The dashed grey line in each panel is Isomap's median, for scale.}\label{fig:nbr-fig-theta}
\end{figure}

Figure \ref{fig:nbr-fig-theta} is the prediction met. From flat sheet to
\(\theta =\) 0.95, Isomap's median error grows by a factor of
7.9; Laplacian eigenmaps' by
1.1, UMAP's by
1.3, t-SNE's by
1.7. Only LLE behaves like a distance method, at
5.3, which is what one would expect of a method
whose weights are least-squares reconstructions in the ambient space.

Flatness is not a compliment here. Isomap is better than all four at \emph{every}
\(\theta\) in this regime, including the last one, and it gets there by
consuming a geometry the crease pattern supplies exactly. A flat curve at a
mediocre level says the error is dominated by something the fold does not
touch --- the objective mismatch itself --- not that the method is robust.

Where the flatness does buy something is off the clean diagonal. Under the
outlier noise model at \(\theta =\) 0.95, the ordering inverts: t-SNE
lands at 0.402 and UMAP at
0.407, against 0.608 for
Isomap, 0.626 for PCA and
0.888 for LLE. A handful of displaced points
lengthens a shortest path through the neighbourhood graph and moves an ambient
distance; it changes few ranks. The methods that were never reading distances
are the ones that do not notice.

\subsection{Neighbourhood size, and where Isomap breaks}\label{neighbourhood-size-and-where-isomap-breaks}

\begin{Shaded}
\begin{Highlighting}[]
\NormalTok{sw }\OtherTok{\textless{}{-}}\NormalTok{ sweeps[sweeps}\SpecialCharTok{$}\NormalTok{status }\SpecialCharTok{==} \StringTok{"ok"} \SpecialCharTok{\&}\NormalTok{ sweeps}\SpecialCharTok{$}\NormalTok{method }\SpecialCharTok{\%in\%} \FunctionTok{c}\NormalTok{(NBR, }\StringTok{"isomap"}\NormalTok{), ]}
\NormalTok{ks }\OtherTok{\textless{}{-}} \FunctionTok{aggregate}\NormalTok{(rmse }\SpecialCharTok{\textasciitilde{}}\NormalTok{ method }\SpecialCharTok{+}\NormalTok{ k\_tck }\SpecialCharTok{+}\NormalTok{ theta, sw, median)}
\NormalTok{k\_max }\OtherTok{\textless{}{-}} \FunctionTok{max}\NormalTok{(ks}\SpecialCharTok{$}\NormalTok{k\_tck)}

\NormalTok{best\_k }\OtherTok{\textless{}{-}} \FunctionTok{do.call}\NormalTok{(rbind, }\FunctionTok{lapply}\NormalTok{(}\FunctionTok{split}\NormalTok{(ks, }\FunctionTok{list}\NormalTok{(ks}\SpecialCharTok{$}\NormalTok{method, ks}\SpecialCharTok{$}\NormalTok{theta)),}
  \ControlFlowTok{function}\NormalTok{(a) }\FunctionTok{data.frame}\NormalTok{(}\AttributeTok{method =}\NormalTok{ a}\SpecialCharTok{$}\NormalTok{method[}\DecValTok{1}\NormalTok{], }\AttributeTok{theta =}\NormalTok{ a}\SpecialCharTok{$}\NormalTok{theta[}\DecValTok{1}\NormalTok{],}
                         \AttributeTok{k\_star  =}\NormalTok{ a}\SpecialCharTok{$}\NormalTok{k\_tck[}\FunctionTok{which.min}\NormalTok{(a}\SpecialCharTok{$}\NormalTok{rmse)],}
                         \AttributeTok{penalty =}\NormalTok{ a}\SpecialCharTok{$}\NormalTok{rmse[a}\SpecialCharTok{$}\NormalTok{k\_tck }\SpecialCharTok{==}\NormalTok{ k\_max] }\SpecialCharTok{/} \FunctionTok{min}\NormalTok{(a}\SpecialCharTok{$}\NormalTok{rmse))))}
\NormalTok{iso\_k }\OtherTok{\textless{}{-}}\NormalTok{ best\_k[best\_k}\SpecialCharTok{$}\NormalTok{method }\SpecialCharTok{==} \StringTok{"isomap"}\NormalTok{, ]}
\NormalTok{iso\_k }\OtherTok{\textless{}{-}}\NormalTok{ iso\_k[}\FunctionTok{order}\NormalTok{(iso\_k}\SpecialCharTok{$}\NormalTok{theta), ]}
\NormalTok{nbr\_k }\OtherTok{\textless{}{-}}\NormalTok{ best\_k[best\_k}\SpecialCharTok{$}\NormalTok{method }\SpecialCharTok{\%in\%}\NormalTok{ NBR, ]}
\end{Highlighting}
\end{Shaded}

\begin{figure}

{\centering \includegraphics[width=0.9\linewidth,alt={Five line panels, one per method, with a logarithmic horizontal axis for neighbourhood size from 5 to 70. Each panel carries three lines, one per fold angle, distinguished by colour and by line style. In the Isomap panel the three curves are U-shaped and the minimum moves left as the fold angle increases. In the LLE, t-SNE, UMAP and Laplacian panels the curves fall steeply and then flatten, with no rise at the right-hand end.}]{06-neighbor-embeddings_files/figure-latex/nbr-fig-ksweep-1} 

}

\caption{Procrustes error against neighbourhood size, at three fold angles. Each panel has its own vertical scale, because the levels differ by more than an order of magnitude and the shape of each curve is the point. Isomap has an interior optimum that moves to smaller neighbourhoods as the sheet folds; every neighbour method is still improving at the largest size swept.}\label{fig:nbr-fig-ksweep}
\end{figure}

Figure \ref{fig:nbr-fig-ksweep} is the sharpest thing in the chapter. Isomap's
best neighbourhood size falls from 45 at \(\theta =\)
0.2, to 20 at 0.5, to
8 at 0.8. That is the short-circuit onset of
Chapter \ref{geodesic} seen from the other side: as the sheet closes, the
neighbourhood size at which the graph starts bridging between facets comes
down, and past it Isomap's error rises. Choosing 70 --- an unremarkable
choice --- costs it a factor of 1.68 over its own
optimum at the tightest fold.

Not one of the four neighbour methods has an interior optimum anywhere in this
sweep. Their worst penalty for taking the largest neighbourhood swept is
1.02, across all twelve method--angle
combinations; for most of them it is exactly one, because the largest
neighbourhood \emph{is} the best. They cannot be short-circuited by a neighbourhood
that is too large, because they never built a global object out of local ones.
This is the genuine win, and it is narrow: it is robustness to a tuning
mistake and to contamination, not accuracy.

\subsection{Do they actually disagree?}\label{do-they-actually-disagree}

The chapter's title asks why t-SNE and UMAP disagree about the same data. On
this benchmark the honest answer is that mostly they do not, and what looks
like disagreement is one draw from each of two wide distributions.

Their median errors differ by 0.048 across the sweep,
against a t-SNE seed band whose median width is
0.72. The two bands overlap at
20 of the 20 fold angles in the grid --- all of them.
Two practitioners running one t-SNE and one UMAP on the same sample will get
pictures that differ, and will usually be looking at within-method seed
variation rather than at a difference between the methods. Where the two do
separate is not in the median but in the spread: UMAP's band is
4.9 times narrower
than t-SNE's, so a UMAP run is a much better estimate of what UMAP does than a
t-SNE run is of what t-SNE does.

\section{Diagnostics}\label{nbr-diagnostics}

\subsection{Two conventions, and which claim needs which}\label{two-conventions-and-which-claim-needs-which}

Every claim above about recovery is made against the \textbf{chart} --- the answer key
--- using either the Procrustes error or \(Q_{NX}\) at \(K =\) 20. The grid
also carries trustworthiness, continuity and \(k\)-NN preservation, and those are
scored at \(k =\) 10 against \textbf{ambient} distance, which is a
different question: whether the embedding preserved the neighbourhoods a
method could actually see. Both are defensible and they are not
interchangeable. In \texttt{part2-sweeps.rds} the ambient-referenced trio is scored at
the \emph{swept} \(k\), so its rows are not comparable with each other; \(Q_{NX}\) stays
pinned at 20 against the chart and is. Every number in this chapter's
figures is one of the two chart-referenced ones.

\begin{Shaded}
\begin{Highlighting}[]
\NormalTok{Kgrid }\OtherTok{\textless{}{-}} \FunctionTok{unique}\NormalTok{(}\FunctionTok{round}\NormalTok{(}\FunctionTok{exp}\NormalTok{(}\FunctionTok{seq}\NormalTok{(}\FunctionTok{log}\NormalTok{(}\DecValTok{2}\NormalTok{), }\FunctionTok{log}\NormalTok{(}\DecValTok{400}\NormalTok{), }\AttributeTok{length.out =} \DecValTok{18}\DataTypeTok{L}\NormalTok{))))}
\NormalTok{ref   }\OtherTok{\textless{}{-}} \FunctionTok{c}\NormalTok{(fit, }\FunctionTok{list}\NormalTok{(}\AttributeTok{pca    =} \FunctionTok{embed}\NormalTok{(}\StringTok{"pca"}\NormalTok{, cell),}
                     \AttributeTok{isomap =} \FunctionTok{embed}\NormalTok{(}\StringTok{"isomap"}\NormalTok{, cell, }\AttributeTok{k =}\NormalTok{ K\_DEFAULT)))}
\NormalTok{rx }\OtherTok{\textless{}{-}} \FunctionTok{do.call}\NormalTok{(rbind, }\FunctionTok{lapply}\NormalTok{(}\FunctionTok{names}\NormalTok{(ref), }\ControlFlowTok{function}\NormalTok{(m)}
  \FunctionTok{data.frame}\NormalTok{(}\AttributeTok{method =}\NormalTok{ m, }\AttributeTok{K =}\NormalTok{ Kgrid, }\AttributeTok{rnx =} \FunctionTok{rnx}\NormalTok{(ref[[m]], Dchart, }\AttributeTok{K =}\NormalTok{ Kgrid))))}

\NormalTok{margin }\OtherTok{\textless{}{-}} \FunctionTok{min}\NormalTok{(}\FunctionTok{vapply}\NormalTok{(Kgrid, }\ControlFlowTok{function}\NormalTok{(kk) \{}
\NormalTok{  a }\OtherTok{\textless{}{-}}\NormalTok{ rx[rx}\SpecialCharTok{$}\NormalTok{K }\SpecialCharTok{==}\NormalTok{ kk, ]}
\NormalTok{  a}\SpecialCharTok{$}\NormalTok{rnx[a}\SpecialCharTok{$}\NormalTok{method }\SpecialCharTok{==} \StringTok{"pca"}\NormalTok{] }\SpecialCharTok{{-}} \FunctionTok{max}\NormalTok{(a}\SpecialCharTok{$}\NormalTok{rnx[a}\SpecialCharTok{$}\NormalTok{method }\SpecialCharTok{\%in\%}\NormalTok{ NBR])}
\NormalTok{\}, }\FunctionTok{numeric}\NormalTok{(}\DecValTok{1}\NormalTok{)))}
\NormalTok{ts }\OtherTok{\textless{}{-}}\NormalTok{ rx[rx}\SpecialCharTok{$}\NormalTok{method }\SpecialCharTok{==} \StringTok{"tsne"}\NormalTok{, ]}
\end{Highlighting}
\end{Shaded}

\begin{figure}

{\centering \includegraphics[width=0.9\linewidth,alt={A line chart with a logarithmic horizontal axis for neighbourhood size K from 2 to 400 and a vertical axis for corrected rank preservation from about 0.6 to 1. The dashed PCA curve is the highest across the whole range and is nearly flat. The dashed Isomap curve starts below LLE at the smallest K, crosses above it early, and rises to the top by the right-hand end. Four solid curves sit below: LLE is nearly flat and highest of them; t-SNE starts close to LLE at small K and falls steadily to the lowest position at K equals 400; Laplacian eigenmaps and UMAP start lowest and rise together. Each curve is labelled at its right-hand end.}]{06-neighbor-embeddings_files/figure-latex/nbr-fig-rnx-1} 

}

\caption{Rank preservation against the chart across every neighbourhood size, on one sample. The two methods that consume distance or geodesics are drawn dashed. No neighbour method reaches PCA at any neighbourhood size, and t-SNE is the only curve that falls as the neighbourhood grows.}\label{fig:nbr-fig-rnx}
\end{figure}

Figure \ref{fig:nbr-fig-rnx} puts the global loss on one axis. \(R_{NX}(K)\) is
\(Q_{NX}\) corrected for the \(K/n\) baseline, so it is comparable across
neighbourhood sizes \citep{lee2009quality}; small \(K\) is local structure and large
\(K\) is global. PCA --- a shadow, and the weakest thing in Part II by construction
--- sits above all four neighbour methods at \textbf{every} \(K\) from 2 to
400, by at least 0.032. There is no neighbourhood
size at which any of these four is the better description of the chart's rank
order.

Isomap is the exception to note rather than to bury: at the smallest
neighbourhood it sits below LLE, at 0.858
against 0.883, and only overtakes it a
few neighbours later. A geodesic method has no special claim on the very local
rank order, which is exactly the part LLE is built out of.

t-SNE is the only curve that falls: 0.846 at \(K =\)
2, 0.682 at \(K =\) 400. It
buys the local structure it is best at by spending the global structure, which
is the trade its critics describe and its defenders accept
\citep{chari2023specious, lause2024elephant}. Here the trade is priced exactly,
because the thing being spent is known.

\subsection{What twenty seeds is, and what it is not}\label{what-twenty-seeds-is-and-what-it-is-not}

\begin{Shaded}
\begin{Highlighting}[]
\NormalTok{bud }\OtherTok{\textless{}{-}}\NormalTok{ budget}\SpecialCharTok{$}\NormalTok{budget[budget}\SpecialCharTok{$}\NormalTok{budget}\SpecialCharTok{$}\NormalTok{method }\SpecialCharTok{\%in\%}\NormalTok{ NBR, ]}
\NormalTok{bud }\OtherTok{\textless{}{-}}\NormalTok{ bud[}\FunctionTok{order}\NormalTok{(}\SpecialCharTok{{-}}\NormalTok{bud}\SpecialCharTok{$}\NormalTok{seeds\_needed), }\FunctionTok{c}\NormalTok{(}\StringTok{"method"}\NormalTok{, }\StringTok{"stochastic"}\NormalTok{, }\StringTok{"sd"}\NormalTok{,}
                                       \StringTok{"seeds\_needed"}\NormalTok{)]}
\NormalTok{knitr}\SpecialCharTok{::}\FunctionTok{kable}\NormalTok{(bud, }\AttributeTok{row.names =} \ConstantTok{FALSE}\NormalTok{, }\AttributeTok{digits =} \DecValTok{4}\NormalTok{,}
             \AttributeTok{col.names =} \FunctionTok{c}\NormalTok{(}\StringTok{"method"}\NormalTok{, }\StringTok{"stochastic"}\NormalTok{, }\StringTok{"worst{-}cell sd"}\NormalTok{,}
                           \StringTok{"seeds needed"}\NormalTok{),}
             \AttributeTok{caption =} \StringTok{"Seeds required to resolve a reportable difference, from data/processed/seed{-}budget.rds."}\NormalTok{)}
\end{Highlighting}
\end{Shaded}

\begin{table}

\caption{\label{tab:nbr-seed-budget}Seeds required to resolve a reportable difference, from data/processed/seed-budget.rds.}
\centering
\begin{tabular}[t]{l|l|r|r}
\hline
method & stochastic & worst-cell sd & seeds needed\\
\hline
tsne & TRUE & 0.2873 & 3241\\
\hline
lle & FALSE & 0.2034 & 1624\\
\hline
laplacian & FALSE & 0.0446 & 79\\
\hline
umap & TRUE & 0.0438 & 76\\
\hline
\end{tabular}
\end{table}

\begin{Shaded}
\begin{Highlighting}[]
\NormalTok{declined  }\OtherTok{\textless{}{-}}\NormalTok{ grid[grid}\SpecialCharTok{$}\NormalTok{method }\SpecialCharTok{==} \StringTok{"laplacian"} \SpecialCharTok{\&}\NormalTok{ grid}\SpecialCharTok{$}\NormalTok{status }\SpecialCharTok{==} \StringTok{"declined"}\NormalTok{, ]}
\NormalTok{dec       }\OtherTok{\textless{}{-}} \FunctionTok{table}\NormalTok{(declined}\SpecialCharTok{$}\NormalTok{noise)}
\NormalTok{dec\_share }\OtherTok{\textless{}{-}} \FunctionTok{nrow}\NormalTok{(declined) }\SpecialCharTok{/} \FunctionTok{sum}\NormalTok{(grid}\SpecialCharTok{$}\NormalTok{method }\SpecialCharTok{==} \StringTok{"laplacian"}\NormalTok{)}
\end{Highlighting}
\end{Shaded}

Standing rule 1's floor is twenty seeds. \texttt{seed-budget.rds} measures what that
floor is worth: for each method it draws 40
samples per cell at three fold angles and two noise models, takes the
within-cell standard deviation of excess error over the floor, and inverts the
two-sample expression for the seeds needed to resolve a difference of
0.02 at 80\% power and the 5\% level.

None of the four passes. t-SNE would need 3241
seeds, LLE 1624, Laplacian eigenmaps
79 and UMAP
76. The floor is between one and two
orders of magnitude short for the first two. That is why the bands in Figure
\ref{fig:nbr-fig-theta} are the result and no mean line is drawn anywhere in
this chapter: at \(\theta =\) 0.4 t-SNE's twenty seeds span
0.098 to 0.875, a range of
0.777 on a metric bounded by one. A mean of that is a number
with no cell in it.

Note which methods are on that list. LLE and Laplacian eigenmaps are
deterministic: the same input gives the same output every time. What varies
across seeds for them is the \emph{sample}, redrawn from the same surface, and LLE's
sensitivity to it is larger than either stochastic method's sensitivity to its
own initialisation. ``Seed variance'' is two different quantities wearing one
name, and only for t-SNE and UMAP does it mean what the phrase suggests.

\section{Where it fails}\label{nbr-limits}

\subsection{\texorpdfstring{\texttt{preserve.seed}, and a correction to the correction}{preserve.seed, and a correction to the correction}}\label{preserve.seed-and-a-correction-to-the-correction}

\begin{Shaded}
\begin{Highlighting}[]
\NormalTok{small }\OtherTok{\textless{}{-}} \FunctionTok{sample\_manifold}\NormalTok{(pat, }\AttributeTok{theta =} \FloatTok{0.5}\NormalTok{, }\AttributeTok{n =} \DecValTok{200}\DataTypeTok{L}\NormalTok{, }\AttributeTok{seed =}\NormalTok{ BENCH\_SEEDS[}\DecValTok{1}\NormalTok{])}
\NormalTok{runs  }\OtherTok{\textless{}{-}} \ControlFlowTok{function}\NormalTok{(preserve) \{}
  \FunctionTok{set.seed}\NormalTok{(}\DecValTok{1}\DataTypeTok{L}\NormalTok{)}
  \FunctionTok{vapply}\NormalTok{(BENCH\_SEEDS, }\ControlFlowTok{function}\NormalTok{(s)}
\NormalTok{    umap}\SpecialCharTok{::}\FunctionTok{umap}\NormalTok{(small}\SpecialCharTok{$}\NormalTok{X, }\AttributeTok{n\_neighbors =}\NormalTok{ K\_DEFAULT, }\AttributeTok{n\_epochs =} \DecValTok{60}\DataTypeTok{L}\NormalTok{,}
               \AttributeTok{preserve.seed =}\NormalTok{ preserve)}\SpecialCharTok{$}\NormalTok{layout[}\DecValTok{1}\DataTypeTok{L}\NormalTok{, }\DecValTok{1}\DataTypeTok{L}\NormalTok{], }\FunctionTok{numeric}\NormalTok{(}\DecValTok{1}\NormalTok{))}
\NormalTok{\}}
\NormalTok{has\_arg  }\OtherTok{\textless{}{-}} \StringTok{"preserve.seed"} \SpecialCharTok{\%in\%} \FunctionTok{names}\NormalTok{(}\FunctionTok{formals}\NormalTok{(umap}\SpecialCharTok{::}\NormalTok{umap))}
\NormalTok{default  }\OtherTok{\textless{}{-}} \FunctionTok{formals}\NormalTok{(umap}\SpecialCharTok{::}\NormalTok{umap)}\SpecialCharTok{$}\NormalTok{preserve.seed}
\NormalTok{distinct }\OtherTok{\textless{}{-}} \FunctionTok{c}\NormalTok{(}\AttributeTok{kept =} \FunctionTok{length}\NormalTok{(}\FunctionTok{unique}\NormalTok{(}\FunctionTok{runs}\NormalTok{(}\ConstantTok{TRUE}\NormalTok{))),}
              \AttributeTok{released =} \FunctionTok{length}\NormalTok{(}\FunctionTok{unique}\NormalTok{(}\FunctionTok{runs}\NormalTok{(}\ConstantTok{FALSE}\NormalTok{))))}
\end{Highlighting}
\end{Shaded}

The natural idiom for standing rule 1 is to seed once and loop. With
\texttt{umap::umap()} that idiom silently collapses the whole loop into one run. The
pinned version takes a \texttt{preserve.seed} argument --- present:
TRUE, default: TRUE --- and that default is the mechanism.
Twenty calls with the argument at its default return
1 distinct answer; the same twenty with it switched off
return 20. The function restores R's random state on
exit, so consecutive calls all start from the same place, and the reported
seed-to-seed spread of every UMAP result in this book would have been exactly
zero --- a violation of standing rule 1 that produces no warning, no error and
numbers that look fine.

\begin{boxreport}
The provenance of this finding needs stating, because the repository
twice recorded it wrongly and in opposite directions. \texttt{PLAN.md}
S1-6 predicted the defect and named \texttt{preserve.seed} as its cause.
\texttt{R/methods.R} and \texttt{tests/testthat/test-methods.R} then
recorded that the prediction had named the wrong mechanism, on the
ground that the pinned version has no such argument, and asserted the
behaviour instead. The prediction was right: the argument is present, it
defaults to \texttt{TRUE}, and switching it off restores the variation
--- all three checked in the chunk above rather than asserted here. The
test is still a good test, because it pins the behaviour rather than an
explanation of it; its comment is wrong and the registry's is too.
\end{boxreport}

The registry uses \texttt{uwot::umap()} instead, which consumes R's random stream, so
a caller who forgets to seed gets visibly different answers rather than twenty
identical ones. The choice of implementation moves the answer, and it moves it
by changing what a seed means rather than by changing an algorithm.

\subsection{Laplacian eigenmaps declines a whole noise model}\label{laplacian-eigenmaps-declines-a-whole-noise-model}

\texttt{embed\_laplacian()} returned nothing on 400 of the
1200 cells it was offered, which is
33\% of them, and every one of those cells carries the
outlier noise model: outlier and nothing else.
The heat kernel underflows on a spread sample, rows of the weight matrix go to
zero, the graph disconnects at working precision, and the second eigenvector is
then read out of a degenerate space where every direction is equally valid. The
helper detects that and refuses. A refusal is a result and the grid records it
as one; what it is not is a score, and Chapter \ref{results} must not average
over the cells that are missing.

\subsection{LLE under outliers stops being a curve}\label{lle-under-outliers-stops-being-a-curve}

LLE never declines a cell, and that is worse than declining one. Under the
outlier model its median error over the upper half of the fold sweep is
0.832, second only to the diffusion map,
and its \(Q_{NX}\) against the chart stops behaving like a function of \(\theta\):
between adjacent grid points it steps \emph{upward} by as much as
0.18, against
0.003 on the clean sweep. Standing rule 2 asks for
a curve; what this cell array produces is a scatter that happens to be indexed
by \(\theta\). The reconstruction weights are least-squares fits in the ambient
space, so one displaced neighbour rewrites a row of the weight matrix, and the
whole embedding is a null space of that matrix. Nothing in the method notices.

\subsection{What this chapter does not establish}\label{what-this-chapter-does-not-establish}

One pattern family, one tessellation size, one ambient dimension, and
800 points. The neighbourhood sweep runs to 70, which is
9\% of the sample --- the curves are still
falling there and this chapter cannot say where they stop. Nothing here
measures whether these methods are good at the task they were built for, which
is a question about cluster recovery that a folded sheet with no clusters
cannot ask \citep{espadoto2021survey}.

\section{Reproduce this}\label{nbr-reproduce}

\begin{itemize}
\tightlist
\item
  \texttt{scripts/run-benchmark-grid.R} → \texttt{data/processed/benchmark-grid.rds}, the
  \(\theta\) sweep at fixed \(k\).
\item
  \texttt{scripts/run-part2-sweeps.R} → \texttt{data/processed/part2-sweeps.rds}, the \(k\)
  sweep at three fixed \(\theta\).
\item
  \texttt{scripts/measure-seed-budget.R} → \texttt{data/processed/seed-budget.rds}.
\item
  \texttt{tests/testthat/test-methods.R} asserts the unseeded refusal, the
  seed-to-seed variation of \texttt{Rtsne} and \texttt{uwot}, and the \texttt{umap::umap}
  behaviour above.
\end{itemize}

Each script's header declares the artefact it writes, and
\texttt{scripts/check-artefact-producers.R} maps every artefact the book reads back to
one.

\section{Further reading}\label{nbr-reading}

Roweis and Saul \citep{roweis2000lle} and Belkin and Niyogi \citep{belkin2003laplacian}
are the primary sources for the two spectral methods, and both are usually read
beside Isomap \citep{tenenbaum2000isomap}, which Chapter \ref{geodesic} covers and
which is the one method here that consumes a geometry a crease pattern supplies
exactly. Van der Maaten and Hinton \citep{vandermaaten2008tsne} and
McInnes, Healy and Melville \citep{mcinnes2018umap} are the primary sources for the
other two.

On what the pictures are worth, Wattenberg, Viégas and Johnson
\citep{wattenberg2016tsne} is the clearest demonstration that the perplexity setting
manufactures structure; Kobak and Berens \citep{kobak2019art} is the practical
counterpart. Kobak and Linderman \citep{kobak2021initialization} show that the
global structure both methods are accused of losing is substantially a matter
of initialisation rather than of the objective. Chari and Pachter
\citep{chari2023specious} press the case against reading these embeddings as maps
and Lause, Berens and Kobak \citep{lause2024elephant} answer it. For the rank-based
criteria used here, Lee and Verleysen \citep{lee2009quality} is the reference, and
Machado, Behrisch and Telea \citep{machado2025necessary} make the point this
chapter's two-metric disagreement illustrates: a projection quality metric is
necessary and not sufficient.

\chapter{Learned folds}\label{autoencoders}

\emph{R. Heller · Date: 2026-07-29 · Last changed: 2026-11-09}

\emph{Can a network learn the fold and the unfold at once?}

\textbf{Chapter status: planned.} The items below describe work still to be implemented; this chapter does not provide a trained autoencoder or benchmark results.

\begin{itemize}
\tightlist
\item
  A \texttt{torch} autoencoder with a two-unit bottleneck. (\texttt{keras} would pull TensorFlow and break clean-clone reproducibility.)
\item
  The only method here that learns an explicit inverse: encoder = low-dimensional representation, decoder = reconstruction in the input space.
\item
  Reconstruction loss is the fold--unfold identity made computable.
\item
  Compare the learned decoder against the true unfolding, not just against the input.
\end{itemize}

\part{The benchmark}\label{part-the-benchmark}

\chapter{Building crease-pattern manifolds}\label{benchmarks}

\emph{R. Heller · Date: 2026-07-29 · Last changed: 2026-09-07}

\emph{How do you build a manifold whose true geometry you already know?}

\section{What this chapter answers}\label{bench-question}

A benchmark for dimension reduction has to supply four things. Three of them are
routine. Data: points in an ambient space that a method can be handed. A
difficulty axis: some parameter that makes the same object harder without making
it a different object. An answer key: the coordinates the method should have
recovered. The Swiss roll and its relatives supply all three, and so does every
folded sheet in this book.

The fourth is the one that is almost never supplied, and it is what this chapter
builds. \textbf{How much of the error was unavoidable?} A method that scores an error
of one half on a dataset has either thrown away most of the structure or reached
the limit of what two dimensions can hold, and the score alone does not say
which. Without that number a benchmark can rank methods against each other but
cannot say whether any of them did well, and a ranking with no scale under it is
how a survey ends up reporting that every method is mediocre on every dataset
\citep{espadoto2021survey}.

The number exists and is computable in closed form, but computing it needs a
manifold whose intrinsic dimension exceeds the target. That is this chapter's
argument: the constructions in \texttt{R/constructions.R} exist to raise intrinsic
dimension while keeping the answer key exact, because a chart of dimension two
embedded into dimension two has an irreducible loss of exactly zero and nothing
to say.

The chapter also records what the benchmark does \emph{not} contain. Two pattern
families were specified for it and both were withdrawn on evidence, and they are
reported here as results --- what it takes for a pattern to qualify is worth more
to a reader than a list of the ones that did.

\section{Setup}\label{bench-setup}

\begin{Shaded}
\begin{Highlighting}[]
\NormalTok{grid }\OtherTok{\textless{}{-}} \FunctionTok{read\_run}\NormalTok{(}\StringTok{"data/processed/product{-}grid.rds"}\NormalTok{, }\StringTok{"bb8ade76d88b95a3"}\NormalTok{)}
\NormalTok{fac\_a }\OtherTok{\textless{}{-}} \FunctionTok{miura\_ori}\NormalTok{(}\DecValTok{4}\DataTypeTok{L}\NormalTok{, }\DecValTok{4}\DataTypeTok{L}\NormalTok{, }\AttributeTok{alpha =}\NormalTok{ pi }\SpecialCharTok{/} \DecValTok{3}\NormalTok{)}
\NormalTok{fac\_b }\OtherTok{\textless{}{-}} \FunctionTok{miura\_ori}\NormalTok{(}\DecValTok{3}\DataTypeTok{L}\NormalTok{, }\DecValTok{5}\DataTypeTok{L}\NormalTok{, }\AttributeTok{b =} \FloatTok{1.4}\NormalTok{, }\AttributeTok{alpha =}\NormalTok{ pi }\SpecialCharTok{/} \DecValTok{4}\NormalTok{)}
\end{Highlighting}
\end{Shaded}

\texttt{grid} is the product suite, 12960 rows produced by
\texttt{scripts/run-product-grid.R}; \texttt{fac\_a} and \texttt{fac\_b} are the two factors it is
built from. Everything else this chapter calls is in \texttt{R/}, sourced when the book
renders.

\section{Background}\label{bench-background}

\subsection{Four objects, four verbs}\label{four-objects-four-verbs}

The helper API is small on purpose. Three S3 classes carry everything between
files, each a plain list, and \texttt{R/README.md} states the contract field by field.

A \textbf{\texttt{crease\_pattern}} is a flat, unfolded pattern: vertices in the plane,
facets as vertex-index cycles, and a crease table carrying a
mountain/valley/boundary assignment. Its \texttt{vertices} matrix is the ground truth ---
not an approximation of the correct two-dimensional embedding, the embedding
itself.

A \textbf{\texttt{folded\_pattern}} is the same pattern placed in \(\mathbb{R}^3\) at folding
parameter \(\theta\), with one dihedral angle per crease derived from the
placement rather than assumed.

A \textbf{\texttt{manifold\_sample}} is \(n\) points drawn from a folded pattern: \texttt{X} in the
ambient space, \texttt{truth} in the chart, row-aligned. \texttt{X{[}i,\ {]}} and \texttt{truth{[}i,\ {]}} are
the same point seen two ways, and every measurement in this book compares an
embedding of the first against the second. The verbs run one line each.

\begin{Shaded}
\begin{Highlighting}[]
\NormalTok{p  }\OtherTok{\textless{}{-}} \FunctionTok{miura\_ori}\NormalTok{(}\AttributeTok{nx =} \DecValTok{6}\DataTypeTok{L}\NormalTok{, }\AttributeTok{ny =} \DecValTok{6}\DataTypeTok{L}\NormalTok{, }\AttributeTok{alpha =}\NormalTok{ pi }\SpecialCharTok{/} \DecValTok{3}\NormalTok{)}
\NormalTok{s  }\OtherTok{\textless{}{-}} \FunctionTok{sample\_manifold}\NormalTok{(p, }\AttributeTok{theta =} \FloatTok{0.6}\NormalTok{, }\AttributeTok{n =} \DecValTok{400}\DataTypeTok{L}\NormalTok{, }\AttributeTok{seed =}\NormalTok{ BENCH\_SEEDS[}\DecValTok{1}\NormalTok{])}
\NormalTok{Y  }\OtherTok{\textless{}{-}} \FunctionTok{embed}\NormalTok{(}\StringTok{"isomap"}\NormalTok{, s, }\AttributeTok{d =}\NormalTok{ EMBED\_DIM, }\AttributeTok{k =}\NormalTok{ K\_DEFAULT)}
\FunctionTok{data.frame}\NormalTok{(}\AttributeTok{vertices =} \FunctionTok{nrow}\NormalTok{(p}\SpecialCharTok{$}\NormalTok{vertices), }\AttributeTok{creases =} \FunctionTok{nrow}\NormalTok{(p}\SpecialCharTok{$}\NormalTok{creases),}
           \AttributeTok{points =} \FunctionTok{nrow}\NormalTok{(s}\SpecialCharTok{$}\NormalTok{X), }\AttributeTok{ambient =} \FunctionTok{ncol}\NormalTok{(s}\SpecialCharTok{$}\NormalTok{X), }\AttributeTok{chart =} \FunctionTok{ncol}\NormalTok{(s}\SpecialCharTok{$}\NormalTok{truth),}
           \AttributeTok{error =} \FunctionTok{round}\NormalTok{(}\FunctionTok{reconstruction\_error}\NormalTok{(Y, s}\SpecialCharTok{$}\NormalTok{truth), }\DecValTok{4}\NormalTok{))}
\end{Highlighting}
\end{Shaded}

\begin{verbatim}
#>   vertices creases points ambient chart  error
#> 1       49      84    400       3     2 0.0551
\end{verbatim}

Four things about that call are load-bearing and easy to get wrong.

\(\theta\) \textbf{is a parameter of the folding map, not an angle.} It runs on
\([0, 1]\), zero being the flat sheet and one the flat-folded state, and \texttt{fold()}
refuses anything outside it; the dihedral angles are derived from it. Chapter
\ref{folding-geometry} gives the derivation.

\texttt{embed()} \textbf{takes the sample object, not a matrix.} The registry holds
9 entries and
8
of them run: the autoencoder declares itself unavailable rather than failing at
fit time, and \texttt{method\_table()} prints the registry with what each method
consumes.

\textbf{Stochastic methods refuse to run unseeded.} Standing rule 1 reports every
stochastic result across \texttt{BENCH\_SEEDS}, and an unseeded fit cannot be one of
them, so \texttt{embed()} errors rather than quietly producing an anecdote.

\textbf{Sampling trims the boundary by default.} The chart is the geodesic only while
the straight chart segment between two points stays on the paper, and a Miura
unfolds to a region with a zigzag edge. \texttt{boundary\ =\ FALSE} is not tidiness --- it
is the condition under which the answer key is the answer.

\section{Constructions that keep the answer key}\label{bench-core}

\subsection{One pattern family, and two honest gaps}\label{one-pattern-family-and-two-honest-gaps}

The book ships one family that folds. That is fewer than were specified, and the
two that were withdrawn were withdrawn for different reasons, both of which are
worth more than the families would have been.

\begin{figure}

{\centering \includegraphics[width=0.9\linewidth,alt={Two side-by-side panels of flat crease patterns drawn in the plane. The left panel is a Miura-ori: a grid of congruent parallelograms whose transverse crease lines zigzag, with a jagged left and right edge. The right panel is a Yoshimura diamond tessellation: rows of triangles with horizontal, slanted and diagonal creases meeting six to a vertex. Both use two colours and two line styles for mountain and valley creases and a solid grey line for the sheet boundary.}]{08-building-benchmarks_files/figure-latex/bench-fig-families-1} 

}

\caption{Two valid flat crease patterns. Only the left one has a certified rigid folding; the right one is drawn here and never folded anywhere in this book. Mountain and valley are distinguished by colour and by line style together.}\label{fig:bench-fig-families}
\end{figure}

Figure \ref{fig:bench-fig-families} shows the Miura beside the Yoshimura. Both
panels are legitimate crease patterns. \texttt{yoshimura()} returns a \texttt{crease\_pattern}
like any other, its facets are congruent, and its interior vertices satisfy
Kawasaki's condition --- the alternating sum of the sector angles vanishes --- to
numerical noise. What it does not have is a folding.

\begin{Shaded}
\begin{Highlighting}[]
\NormalTok{yo }\OtherTok{\textless{}{-}} \FunctionTok{yoshimura}\NormalTok{(}\DecValTok{6}\DataTypeTok{L}\NormalTok{, }\DecValTok{6}\DataTypeTok{L}\NormalTok{)}
\NormalTok{kaw }\OtherTok{\textless{}{-}} \FunctionTok{max}\NormalTok{(}\FunctionTok{vapply}\NormalTok{(}\FunctionTok{interior\_vertices}\NormalTok{(yo), }\ControlFlowTok{function}\NormalTok{(v) \{}
\NormalTok{  a }\OtherTok{\textless{}{-}} \FunctionTok{sector\_angles}\NormalTok{(yo, v)}
  \FunctionTok{abs}\NormalTok{(}\FunctionTok{sum}\NormalTok{(a[}\FunctionTok{c}\NormalTok{(}\ConstantTok{TRUE}\NormalTok{, }\ConstantTok{FALSE}\NormalTok{)]) }\SpecialCharTok{{-}} \FunctionTok{sum}\NormalTok{(a[}\FunctionTok{c}\NormalTok{(}\ConstantTok{FALSE}\NormalTok{, }\ConstantTok{TRUE}\NormalTok{)]))}
\NormalTok{\}, }\FunctionTok{numeric}\NormalTok{(}\DecValTok{1}\NormalTok{)))}
\NormalTok{first\_line }\OtherTok{\textless{}{-}} \ControlFlowTok{function}\NormalTok{(e) }\FunctionTok{sub}\NormalTok{(}\StringTok{"}\SpecialCharTok{\textbackslash{}\textbackslash{}}\StringTok{. See .*"}\NormalTok{, }\StringTok{"."}\NormalTok{, }\FunctionTok{conditionMessage}\NormalTok{(e))}
\NormalTok{refusals }\OtherTok{\textless{}{-}} \FunctionTok{c}\NormalTok{(}
  \AttributeTok{yoshimura =} \FunctionTok{tryCatch}\NormalTok{(}\FunctionTok{fold}\NormalTok{(yo, }\AttributeTok{theta =} \FloatTok{0.5}\NormalTok{), }\AttributeTok{error =}\NormalTok{ first\_line),}
  \AttributeTok{waterbomb =} \FunctionTok{tryCatch}\NormalTok{(}\FunctionTok{waterbomb}\NormalTok{(), }\AttributeTok{error =}\NormalTok{ first\_line))}
\FunctionTok{cat}\NormalTok{(}\FunctionTok{strwrap}\NormalTok{(refusals, }\AttributeTok{width =} \DecValTok{76}\NormalTok{, }\AttributeTok{prefix =} \StringTok{"  "}\NormalTok{, }\AttributeTok{initial =} \StringTok{""}\NormalTok{), }\AttributeTok{sep =} \StringTok{"}\SpecialCharTok{\textbackslash{}n}\StringTok{"}\NormalTok{)}
\end{Highlighting}
\end{Shaded}

\begin{verbatim}
#> no rigid folding of the Yoshimura pattern is certified in which all three
#>   crease families fold. Two were derived and each leaves one family flat,
#>   which makes the folded object a parallelogram tessellation rather than a
#>   diamond one -- a Miura in this pattern's clothing.
#>   waterbomb() is not implemented: whether this tessellation admits a
#>   one-parameter rigid folding is an open question in this project, not a
#>   coding task.
\end{verbatim}

Kawasaki's alternating sum is satisfied at all
25 interior vertices of that patch, worst
departure 1.8e-15, and that buys
nothing. Kawasaki is a necessary condition for flat-foldability at a vertex; it
is not sufficient, and
flat-foldability of a pattern is not the same question as the existence of a
one-parameter rigid folding of a tessellation \citep{hull2020origametry, demaine2007folding}.

\textbf{The Yoshimura.} Two rigid foldings of the planar patch were derived,
implemented and verified in this repository. Each of them is exact --- facet
isometry at the same order as the Miura's, flat at \(\theta = 0\) --- and each
leaves one of the three crease families entirely unfolded, at a dihedral angle
of \(\pi\) along its whole length. A folding that leaves a crease family flat is a
folding of a coarser pattern: fuse the facets across the creases that never
move, and what remains is a parallelogram tessellation. It is a Miura wearing
the Yoshimura's crease lines. Shipping it as a second, independent family would
have been a claim about pattern variety that the geometry does not support, so
\texttt{fold()} refuses it and says why.

\textbf{The waterbomb.} \texttt{waterbomb()} stops before building anything. A degree-six
vertex with sectors of forty-five, forty-five and ninety degrees repeated
satisfies Kawasaki and is flat-foldable \emph{as an isolated vertex}, which does not
imply that the tessellation admits a one-parameter rigid folding. Whether it
does is an open question in this project, not a coding task. A continuation
moved it only under an imposed symmetry, and three things disqualified that
result: the embedding is unproved, the fold amplitude is not monotone in
\(\theta\), and degrees of freedom remain on the free boundary, so \(\theta\) does
not determine the configuration.

Both withdrawals share one moral and it generalises well past origami. \textbf{A test
that a correct implementation passes is not the same as a test that a wrong one
fails.} An isometry check cannot tell a fold from a pleat, because a pleat is a
rigid folding and passes perfectly. A first-order flex at the flat state proves
nothing either: at \(\theta = 0\) every \(z\) column of the rigidity Jacobian
vanishes, so the count of non-trivial infinitesimal flexes is
\(V - 3 = 46\) for a pattern of this size and for \emph{any}
planar pattern whatsoever, foldable or not. Each trap here was sprung, and each
was caught only by a second, independent check.

\subsection{The isometry certificate}\label{the-isometry-certificate}

What a family has to produce, then, is a closed form, a fold that reaches every
crease family, and an isometry that is exact rather than converged. The third is
certified by measurement, not by the derivation that produced it.

\begin{Shaded}
\begin{Highlighting}[]
\NormalTok{alphas }\OtherTok{\textless{}{-}} \FunctionTok{seq}\NormalTok{(MIURA\_ALPHA\_RANGE[}\DecValTok{1}\NormalTok{], MIURA\_ALPHA\_RANGE[}\DecValTok{2}\NormalTok{], }\AttributeTok{length.out =} \DecValTok{6}\DataTypeTok{L}\NormalTok{)}
\NormalTok{thetas }\OtherTok{\textless{}{-}} \FunctionTok{seq}\NormalTok{(}\DecValTok{0}\NormalTok{, }\DecValTok{1}\NormalTok{, }\AttributeTok{by =} \FloatTok{0.05}\NormalTok{)}
\NormalTok{cert }\OtherTok{\textless{}{-}} \FunctionTok{do.call}\NormalTok{(rbind, }\FunctionTok{lapply}\NormalTok{(alphas, }\ControlFlowTok{function}\NormalTok{(al) \{}
\NormalTok{  q }\OtherTok{\textless{}{-}} \FunctionTok{miura\_ori}\NormalTok{(}\DecValTok{4}\DataTypeTok{L}\NormalTok{, }\DecValTok{4}\DataTypeTok{L}\NormalTok{, }\AttributeTok{alpha =}\NormalTok{ al)}
  \FunctionTok{data.frame}\NormalTok{(}\AttributeTok{alpha =} \FunctionTok{round}\NormalTok{(al }\SpecialCharTok{*} \DecValTok{180} \SpecialCharTok{/}\NormalTok{ pi), }\AttributeTok{theta =}\NormalTok{ thetas,}
             \AttributeTok{err =} \FunctionTok{vapply}\NormalTok{(thetas, }\ControlFlowTok{function}\NormalTok{(t) }\FunctionTok{facet\_isometry\_error}\NormalTok{(q, t),}
                          \FunctionTok{numeric}\NormalTok{(}\DecValTok{1}\NormalTok{)))}
\NormalTok{\}))}
\NormalTok{worst }\OtherTok{\textless{}{-}} \FunctionTok{max}\NormalTok{(cert}\SpecialCharTok{$}\NormalTok{err)}
\end{Highlighting}
\end{Shaded}

\texttt{facet\_isometry\_error()} compares \textbf{all pairwise distances} between a facet's
vertices, folded against flat, and not merely its edges: a placement that
preserved every edge of a parallelogram while shearing it would pass an edge
test and would not be a rigid folding. Over the whole \(\theta\) sweep and the
full \(\alpha\) range the worst departure is
1.1e-15, against a contract of
1e-10. That is floating-point noise on a
closed form, not a tolerance anyone chose, and it is the certificate the product
construction below inherits.

\begin{figure}

{\centering \includegraphics[width=0.9\linewidth,alt={A line chart with a logarithmic vertical axis. The horizontal axis is the folding parameter from zero to one. Six nearly flat curves, one per opening angle and distinguished by colour, sit together in a narrow band around ten to the power minus fifteen. A horizontal dashed line far above them at ten to the power minus ten marks the contract tolerance.}]{08-building-benchmarks_files/figure-latex/bench-fig-certificate-1} 

}

\caption{Worst within-facet distance error across the folding sweep, at six opening angles spanning the range the contract asserts. The dashed line is the contract; every curve sits five orders of magnitude below it and shows no trend in either the folding parameter or the opening angle.}\label{fig:bench-fig-certificate}
\end{figure}

Figure \ref{fig:bench-fig-certificate} is that measurement across the sweep.
The flatness matters as much as the level: an error growing with \(\theta\) would
mean the closed form was an approximation that happened to be good near the flat
state, which is what a derivation by facet rotation tends to produce.

\subsection{Products: raising intrinsic dimension without losing truth}\label{products-raising-intrinsic-dimension-without-losing-truth}

A folded sheet has intrinsic dimension two and every embedding in this book has
target dimension two. That is a benchmark with no headroom, and the product
construction supplies some.

\texttt{product\_manifold(a,\ b)} takes two samples of the same size and pairs them
row-wise: the ambient coordinates concatenate to \(\mathbb{R}^{3+3}\), the charts
to \(\mathbb{R}^{2+2}\). Nothing is fitted and nothing approximated, for a reason
worth stating because everything downstream rests on it.

\textbf{A product of isometries is an isometry.} In a Riemannian product the squared
distances add: for points \((u, v)\) and \((u', v')\),
\(d^2 = d_1(u, u')^2 + d_2(v, v')^2\), and the same identity holds in the ambient
product because concatenating coordinate blocks adds squared differences. Each
factor map preserves its own term by construction, so the product map preserves
the sum. The chart of the product is therefore the product of the charts, exactly, and
its intrinsic dimension is the sum of the factors' --- \(2k\) in ambient \(3k\) for
\(k\) sheets.

\begin{Shaded}
\begin{Highlighting}[]
\NormalTok{sa }\OtherTok{\textless{}{-}} \FunctionTok{sample\_manifold}\NormalTok{(fac\_a, }\AttributeTok{theta =} \FloatTok{0.55}\NormalTok{, }\AttributeTok{n =} \DecValTok{400}\DataTypeTok{L}\NormalTok{, }\AttributeTok{seed =}\NormalTok{ BENCH\_SEEDS[}\DecValTok{1}\NormalTok{])}
\NormalTok{sb }\OtherTok{\textless{}{-}} \FunctionTok{sample\_manifold}\NormalTok{(fac\_b, }\AttributeTok{theta =} \FloatTok{0.60}\NormalTok{, }\AttributeTok{n =} \DecValTok{400}\DataTypeTok{L}\NormalTok{, }\AttributeTok{seed =}\NormalTok{ BENCH\_SEEDS[}\DecValTok{1}\NormalTok{] }\SpecialCharTok{+} \DecValTok{500}\DataTypeTok{L}\NormalTok{)}
\NormalTok{pr }\OtherTok{\textless{}{-}} \FunctionTok{product\_manifold}\NormalTok{(sa, sb)}

\NormalTok{pythag }\OtherTok{\textless{}{-}} \FunctionTok{max}\NormalTok{(}\FunctionTok{abs}\NormalTok{(}\FunctionTok{as.matrix}\NormalTok{(}\FunctionTok{dist}\NormalTok{(pr}\SpecialCharTok{$}\NormalTok{truth)) }\SpecialCharTok{{-}}
                  \FunctionTok{sqrt}\NormalTok{(}\FunctionTok{as.matrix}\NormalTok{(}\FunctionTok{dist}\NormalTok{(sa}\SpecialCharTok{$}\NormalTok{truth))}\SpecialCharTok{\^{}}\DecValTok{2} \SpecialCharTok{+} \FunctionTok{as.matrix}\NormalTok{(}\FunctionTok{dist}\NormalTok{(sb}\SpecialCharTok{$}\NormalTok{truth))}\SpecialCharTok{\^{}}\DecValTok{2}\NormalTok{)))}
\NormalTok{floors }\OtherTok{\textless{}{-}} \FunctionTok{vapply}\NormalTok{(}\DecValTok{1}\SpecialCharTok{:}\DecValTok{4}\NormalTok{, }\ControlFlowTok{function}\NormalTok{(d) }\FunctionTok{irreducible\_loss}\NormalTok{(pr, d), }\FunctionTok{numeric}\NormalTok{(}\DecValTok{1}\NormalTok{))}
\FunctionTok{data.frame}\NormalTok{(}\AttributeTok{ambient =} \FunctionTok{ncol}\NormalTok{(pr}\SpecialCharTok{$}\NormalTok{X), }\AttributeTok{chart =} \FunctionTok{ncol}\NormalTok{(pr}\SpecialCharTok{$}\NormalTok{truth),}
           \AttributeTok{pythagoras =} \FunctionTok{formatC}\NormalTok{(pythag, }\AttributeTok{format =} \StringTok{"e"}\NormalTok{, }\AttributeTok{digits =} \DecValTok{1}\NormalTok{),}
           \AttributeTok{one\_sheet\_floor =} \FunctionTok{irreducible\_loss}\NormalTok{(sa, }\DecValTok{2}\DataTypeTok{L}\NormalTok{))}
\end{Highlighting}
\end{Shaded}

\begin{verbatim}
#>   ambient chart pythagoras one_sheet_floor
#> 1       6     4    8.9e-16               0
\end{verbatim}

The Pythagorean identity holds to 8.9e-16 across all pairs, which
is the numerical form of the argument above. The last entry is the sentence the
whole construction exists to escape: a single folded sheet has a chart of
dimension two, and the smallest error any two-dimensional embedding of it could
achieve is 0. Exactly zero, in every cell, at every
\(\theta\), under every noise model.

\textbf{The bound.} For an \(n \times p\) chart \(U\) and the normalised Procrustes error
of \texttt{R/metrics.R} --- which optimises over translation, rotation, reflection and
isotropic scale \citep{schonemann1966procrustes, gower1975procrustes} --- every
\(n \times d\) configuration \(Z\) satisfies

\[\mathrm{err}(Z, U) \;\ge\; \mathrm{floor}(d) \;=\; \sqrt{\frac{\sum_{i > d}\lambda_i}{\sum_i \lambda_i}},\]

where the \(\lambda_i\) are the classical-MDS eigenvalues of the chart
\citep{torgerson1952mds}, equivalently the squared singular values of its centred
coordinates. The proof is Eckart--Young \citep{eckart1936approximation} and it is
short. As \(Z\) ranges over all configurations and the similarity group ranges
over its elements, the fitted matrix ranges over exactly the centred \(n \times p\)
matrices of rank at most \(d\); the best such matrix is the rank-\(d\) truncated
singular value decomposition; its residual is the tail sum of squared singular
values; normalising gives the expression above.

It is exact, it depends on the data alone, and no method appears anywhere in it.

\begin{figure}

{\centering \includegraphics[width=0.9\linewidth,alt={Two panels. Panel A is a bar chart of four bars, decreasing from about 0.43 to about 0.06, labelled as shares of chart variance. Panel B is a line falling from about 0.75 at target dimension one, through 0.53 and 0.25, to zero at target dimension four; black cross markers sit exactly on the line at all four points.}]{08-building-benchmarks_files/figure-latex/bench-fig-spectrum-1} 

}

\caption{The chart spectrum of one product cell, and the bound it produces. Panel A: the four eigenvalues of the four-dimensional chart, as shares of its variance. Panel B: the bound at each target dimension, as a line, with the rank-d projection of the chart drawn on top of it as crosses. The tail of A is the square of B.}\label{fig:bench-fig-spectrum}
\end{figure}

Figure \ref{fig:bench-fig-spectrum} shows why it is a bound rather than a
guess: \textbf{it is attained.} The crosses in panel B are the chart projected onto
its own top \(d\) principal directions, scored by the same function every method
in the book is scored by. They sit on the line, to
6.1e-16. A minimum, not an infimum --- which also
makes clear what the bound is not. The minimiser is constructed from the answer
key, so nothing that has not been given the answer can be expected to reach it.

The floor depends on the chart and on nothing else. That is the property that
makes it useful rather than a limitation: it does not move with \(\theta\), with
the noise model or with the sample's difficulty, so ``every method landed within
a few per cent of the floor'' is a statement about a regime rather than about a
lucky draw.

\subsection{Lifting to arbitrary ambient dimension}\label{lifting-to-arbitrary-ambient-dimension}

Real data does not arrive in three dimensions, and a method that only works when
the ambient dimension is small should be caught by that rather than flattered by
it. \texttt{lift()} maps a sample into \(\mathbb{R}^D\) through a random orthonormal
matrix.

\begin{Shaded}
\begin{Highlighting}[]
\NormalTok{lf }\OtherTok{\textless{}{-}} \FunctionTok{lift}\NormalTok{(pr, }\AttributeTok{D =} \DecValTok{50}\DataTypeTok{L}\NormalTok{, }\AttributeTok{seed =} \DecValTok{20260901}\DataTypeTok{L}\NormalTok{)}
\NormalTok{dist\_change }\OtherTok{\textless{}{-}} \FunctionTok{max}\NormalTok{(}\FunctionTok{abs}\NormalTok{(}\FunctionTok{as.matrix}\NormalTok{(}\FunctionTok{dist}\NormalTok{(lf}\SpecialCharTok{$}\NormalTok{X)) }\SpecialCharTok{{-}} \FunctionTok{as.matrix}\NormalTok{(}\FunctionTok{dist}\NormalTok{(pr}\SpecialCharTok{$}\NormalTok{X))))}
\FunctionTok{data.frame}\NormalTok{(}\AttributeTok{ambient =} \FunctionTok{ncol}\NormalTok{(lf}\SpecialCharTok{$}\NormalTok{X),}
           \AttributeTok{dist\_change =} \FunctionTok{formatC}\NormalTok{(dist\_change, }\AttributeTok{format =} \StringTok{"e"}\NormalTok{, }\AttributeTok{digits =} \DecValTok{1}\NormalTok{),}
           \AttributeTok{floor\_before =} \FunctionTok{round}\NormalTok{(}\FunctionTok{irreducible\_loss}\NormalTok{(pr, }\DecValTok{2}\DataTypeTok{L}\NormalTok{), }\DecValTok{4}\NormalTok{),}
           \AttributeTok{floor\_after  =} \FunctionTok{round}\NormalTok{(}\FunctionTok{irreducible\_loss}\NormalTok{(lf, }\DecValTok{2}\DataTypeTok{L}\NormalTok{), }\DecValTok{4}\NormalTok{))}
\end{Highlighting}
\end{Shaded}

\begin{verbatim}
#>   ambient dist_change floor_before floor_after
#> 1      50     1.8e-15         0.53        0.53
\end{verbatim}

Orthonormality is load-bearing and \texttt{lift()} asserts it rather than assuming it:
a merely random matrix would distort distances and invalidate the answer key
silently. Because \(Q^\top Q = I\) the map is exactly distance-preserving --- the
worst change over all
79800 pairs is
1.8e-15 --- and the
chart is untouched, so the floor is the same number before and after.

That decoupling is the point. Ambient dimension and intrinsic difficulty are two
axes, and a benchmark that moves them together cannot say which a method is
responding to. \texttt{lift()} moves the first alone.

\section{Results}\label{bench-results}

The product suite is 12960 rows: two arms, each swept over nine values
of its own difficulty parameter, 20 seeds, four target
dimensions and 9 registry entries at
\(n = 400\).

The second arm is a concession turned into a finding. \texttt{PROJECT\_CONCEPT.md}
listed crease-specificity among the book's differentiators, and
\texttt{product\_manifold()} imposes no family constraint: two isometric Swiss rolls
give intrinsic dimension four, an exact chart and a computable floor exactly as
two Miura sheets do. The differentiator as written is refuted, and reporting
both arms side by side is the stronger statement. What crease-specificity rests
on afterwards is zero reach \citep{federer1959curvature, niyogi2008homology} and an
answer key that is not smooth, not the bound.

\subsection{The floor bites, and nothing goes under it}\label{the-floor-bites-and-nothing-goes-under-it}

\begin{figure}

{\centering \includegraphics[width=0.9\linewidth,alt={Two panels sharing a vertical axis of normalised Procrustes error from zero to one. In each, a heavy grey line falls with target dimension one to four and the area beneath it is shaded. In the left panel, labelled crease product, three curves for ambient, geodesic and neighbourhood methods run just above the shaded region at dimensions one and two and separate from it at three and four. In the right panel, labelled roll product, the same three curves run flat near 0.95 across all four dimensions, far above a bound that falls quickly to near zero.}]{08-building-benchmarks_files/figure-latex/bench-fig-dimension-1} 

}

\caption{Achieved error against target dimension, by what each method consumes, with the irreducible-loss bound as the heavy line and the region below it shaded. Both arms have a four-dimensional chart, so the bound reaches zero at target dimension four. The shaded region is unreachable by any configuration whatsoever.}\label{fig:bench-fig-dimension}
\end{figure}

\begin{Shaded}
\begin{Highlighting}[]
\NormalTok{fits }\OtherTok{\textless{}{-}} \FunctionTok{nrow}\NormalTok{(ok)}
\NormalTok{under }\OtherTok{\textless{}{-}} \FunctionTok{sum}\NormalTok{(ok}\SpecialCharTok{$}\NormalTok{excess }\SpecialCharTok{\textless{}} \SpecialCharTok{{-}}\FloatTok{1e{-}9}\NormalTok{)}
\NormalTok{crease2 }\OtherTok{\textless{}{-}}\NormalTok{ ok[ok}\SpecialCharTok{$}\NormalTok{d }\SpecialCharTok{==} \DecValTok{2}\DataTypeTok{L} \SpecialCharTok{\&}\NormalTok{ ok}\SpecialCharTok{$}\NormalTok{family }\SpecialCharTok{==} \StringTok{"crease"}\NormalTok{, ]}
\NormalTok{roll2   }\OtherTok{\textless{}{-}}\NormalTok{ ok[ok}\SpecialCharTok{$}\NormalTok{d }\SpecialCharTok{==} \DecValTok{2}\DataTypeTok{L} \SpecialCharTok{\&}\NormalTok{ ok}\SpecialCharTok{$}\NormalTok{family }\SpecialCharTok{==} \StringTok{"roll"}\NormalTok{, ]}
\NormalTok{by\_con }\OtherTok{\textless{}{-}} \FunctionTok{aggregate}\NormalTok{(excess }\SpecialCharTok{\textasciitilde{}}\NormalTok{ consumes, crease2, mean)}
\end{Highlighting}
\end{Shaded}

Across all 11160 fits in the suite, the number that went under the floor is
0. That is a regression check and not evidence, and the distinction
matters: the bound is
\emph{proved}, so a method that beat it would mean the implementation is wrong, not
that the claim is. It is the assertion the producer exits non-zero on, and it
should be read the way a failing test is read, not the way a corroborated
conjecture is.

Figure \ref{fig:bench-fig-dimension} is what the bound buys. At target
dimension two the crease arm's floor is
0.517: the best configuration that exists still
leaves a residual half the size of the chart's own spread, before any method is
chosen. Isomap's raw error there averages
0.536, which against
zero reads as failure and against the floor is an excess of
0.019 --- it gives away
3.9 per
cent of the range actually available to it. \textbf{That contrast is this book's
thesis in one sentence.} Read the first way the table says everything failed;
read the second it says the loss belongs to the data.

Excess sorts by what a method consumes. On the crease arm at target dimension
two the mean is 0.019 for
the one geodesic method,
0.041 for the ambient ones
and 0.113 for the
neighbourhood ones. Isomap consumes exactly what a crease pattern supplies
exactly \citep{tenenbaum2000isomap}, and it is the method closest to the bound. That
ordering is measured over the whole sweep and all 20
seeds. The provisional version of it came from two values of \(\theta\) in a
quick run that was never saved and put the ambient methods first --- which is
exactly the reading standing rule 2 exists to prevent.

\subsection{Excess as a curve}\label{excess-as-a-curve}

\begin{figure}

{\centering \includegraphics[width=0.9\linewidth,alt={Two panels with independent axes. The left panel, crease product, plots excess over the floor against the folding parameter: three nearly flat curves for ambient, geodesic and neighbourhood methods sit between 0.01 and 0.10 until the folding parameter passes 0.8, where the ambient curve turns sharply upward and the geodesic one bends up more gently. The right panel, roll product, plots excess against the number of turns: three curves rise steadily together from about 0.72 to about 0.92 and stay within a hair of each other throughout.}]{08-building-benchmarks_files/figure-latex/bench-fig-excess-1} 

}

\caption{Mean excess over the floor at target dimension two, against each family's own difficulty parameter — the folding parameter on the left, the number of turns on the right. The vertical scales differ between panels and the two arms are matched on nothing, so this is a statement about the construction and not about which family is harder.}\label{fig:bench-fig-excess}
\end{figure}

Figure \ref{fig:bench-fig-excess} is the whole sweep.

\begin{Shaded}
\begin{Highlighting}[]
\NormalTok{crease\_curve }\OtherTok{\textless{}{-}} \FunctionTok{aggregate}\NormalTok{(}\FunctionTok{cbind}\NormalTok{(excess, floor) }\SpecialCharTok{\textasciitilde{}}\NormalTok{ par, crease2, mean)}
\NormalTok{roll\_curve   }\OtherTok{\textless{}{-}} \FunctionTok{aggregate}\NormalTok{(}\FunctionTok{cbind}\NormalTok{(excess, floor) }\SpecialCharTok{\textasciitilde{}}\NormalTok{ par, roll2, mean)}
\end{Highlighting}
\end{Shaded}

Two things read off it. On the crease arm the floor is
0.517 at every one of the nine folding
parameters --- spread across the sweep exactly
0 --- while the excess climbs from
0.062 to
0.146, almost all of it in the last two cells.
The bound is a property of the answer key and folding does not touch it; the
difficulty is entirely in the recovery. On the roll arm the floor \emph{does} move,
from 0.168 down to
0.068, because winding more turns lengthens the
roll's arc-length coordinate and elongates its chart. Two arms, two different
relationships between difficulty and the bound, and neither is a defect.

\subsection{Why a benchmark stuck at intrinsic dimension 2 is not a serious benchmark}\label{why-a-benchmark-stuck-at-intrinsic-dimension-2-is-not-a-serious-benchmark}

Everything above is unavailable on the book's main grid, and the reason is one
line of arithmetic. That grid's chart is two-dimensional and its target is
two-dimensional, so the \texttt{floor} column in it is identically zero --- as
\texttt{irreducible\_loss(sa,\ 2L)} returned above, for the same reason, in closed form.

That is not a defect in the grid and it is worth one sentence rather than a
footnote: \textbf{on the main grid every error is loss the method is responsible
for.} It is the cleanest possible reading, and it is available precisely
because the construction has no headroom in it.

But a benchmark that \emph{only} ever has no headroom cannot ask the question this
chapter opened with. It cannot distinguish a method that discarded structure
from one that ran out of room, because there is never any room to run out of. It
cannot show what Figure \ref{fig:bench-fig-dimension} shows --- the crease arm's
methods hugging the bound at target dimensions one and two and pulling away from
it at three and four, which says where each method's loss comes from rather than
how large it is. And it cannot report a result as a ratio, because a ratio needs
both ends fixed and one end is nailed to zero.

The classical benchmarks are all in this position. A Swiss roll, an S-curve, a
severed sphere: two-dimensional charts, two-dimensional targets, floor zero
everywhere \citep{lee2007nonlinear}. Raising intrinsic dimension is the one move that
makes the bound bite, and doing it while keeping the answer key exact is what
\texttt{product\_manifold()} is for. The product suite, not the main grid, carries the
argument.

\section{Diagnostics}\label{bench-diagnostics}

\textbf{Do not read the two arms as a difficulty comparison.} The crease arm is swept
over the folding parameter and the roll arm over turns; the two are matched on
nothing, and their floors differ by a factor of about
5 for reasons about how each
chart is parameterised rather than about how hard each object is. What the pair
supports is that the floor exists, is computable, is attained and is not beaten,
in both families. A claim that one family is harder needs matched difficulty and
does not live here.

\textbf{Rows that did not run are results.} Of the 12960 rows,
1440 are the autoencoder declaring itself
unavailable and 360 are t-SNE refusing target
dimension four --- the Barnes--Hut implementation accepts one, two or three, and
the \texttt{reason} column says so. Both are recorded rather than dropped, because a
missing row and a refused row look identical in an aggregate and mean opposite
things.

\textbf{The bound has no opinion about rank metrics.} \texttt{metric\_floor()} bounds the
normalised Procrustes error and nothing else. Trustworthiness, continuity and
\(Q_{NX}\) read rank order rather than distance, and an embedding sitting on this
floor may score better or worse on any of them than one well above it. That is
Chapter \ref{evaluation}'s subject, not a defect here.

\textbf{One stale message.} \texttt{waterbomb()}'s error text still says the book ships two
families plus one negative result; it predates the Yoshimura's withdrawal and
was never revised. The count in this chapter is the correct one --- one family,
two negative results --- and the discrepancy is left visible rather than quietly
patched, because it is exactly the kind of stale count a registry keyed by
family name is supposed to catch and did not.

\section{Where it fails}\label{bench-limits}

\textbf{A product is not a curved four-manifold.} \texttt{product\_manifold()} gives
intrinsic dimension four with an exact chart, and it gives it in the easiest
possible way: the factors are independent, their metrics do not interact, and
the chart's four coordinates come in two decoupled pairs. Nothing here models a
manifold whose curvature entangles its directions. The construction is for the
bound, and the bound needs only a chart of dimension greater than the target.

\textbf{The lift is linear.} What makes it a clean axis is also its limit. A random
orthonormal lift adds no correlated nuisance structure, no heteroscedastic noise
and no nearly redundant coordinates --- none of the pathologies that make real
high-dimensional data hard \citep{chari2023specious, lause2024elephant}. A method
that survives \texttt{lift()} has been shown insensitive to ambient dimension alone.

\textbf{The chart must be Euclidean.} The bound is an Eckart--Young residual and
needs the doubly-centred squared-distance matrix to be positive semi-definite.
It is here, by construction. A distance matrix handed in from elsewhere may not
be, and \texttt{chart\_spectrum()} refuses that case rather than reporting a tail sum
that bounds nothing.

\textbf{The answer key is the chart only inside the paper.} Chart distance equals
geodesic distance only while the straight segment between two chart points stays
on the sheet, and a Miura's outline is not convex. The producer asserts a
chart-exit fraction of zero on every cell rather than trusting the sampling
default; where that assertion fails the answer key is wrong in the direction
that flatters every method.

\textbf{One family.} A benchmark built on a single pattern family cannot support
claims about pattern variety, and this book makes none. It can support claims
about a construction, and the suite's two arms are the evidence that the
construction is what the claims are about.

\section{Reproduce this}\label{bench-reproduce}

To add a family, write the constructor in \texttt{R/patterns.R} and the folding map in
\texttt{R/folding.R}, and let \texttt{tests/testthat/test-folding.R} decide whether it
qualifies: within-facet pairwise distances preserved to
1e-10 across the sweep, flatness at
\(\theta = 0\), an independent identity that was not built into the derivation,
and a mountain/valley assignment derived from the folded geometry rather than
asserted. A family that leaves a crease family flat fails the last of those,
which is how the Yoshimura was caught.

To build a product of your own, hand \texttt{product\_manifold()} two \texttt{manifold\_sample}
objects of equal size --- they need not come from the same family, or from a
crease pattern at all --- and score against \texttt{irreducible\_loss()} rather than
against zero. \texttt{against\_floor()} returns error, floor, excess and ratio together,
which is the form every result in Part III is reported in.

The suite itself is \texttt{scripts/run-product-grid.R}, which is run neither in CI nor
at render time: run it locally, commit the \texttt{.rds}, and let \texttt{read\_run()} refuse
any chapter whose prose was written against a different one.

\section{Further reading}\label{bench-reading}

Miura's report \citep{miura1985method} introduces the one family this benchmark
folds, and Schenk and Guest \citep{schenk2013geometry} give the mechanics of folded
shells. Hull \citep{hull2020origametry} is the reference for what Kawasaki's and
Maekawa's conditions do and do not guarantee, Demaine and O'Rourke
\citep{demaine2007folding} for rigid foldability as a computational question, and
Lang \citep{lang2011origami} for the design side this chapter does not need.

On the measurement side, Eckart and Young \citep{eckart1936approximation} is the
approximation theorem the bound instantiates, Schönemann
\citep{schonemann1966procrustes} and Gower \citep{gower1975procrustes} fix the alignment
group it is stated over, and Torgerson \citep{torgerson1952mds} is where the chart
spectrum comes from. Balasubramanian and Schwartz \citep{balasubramanian2002isomap}
state the failure mode the folding parameter drives, and Lee and Verleysen
\citep{lee2007nonlinear} is the textbook treatment of the methods being measured.
Espadoto and colleagues \citep{espadoto2021survey} survey the benchmarking practice
this chapter argues with, and Machado and colleagues \citep{machado2025necessary}
make the case for evaluating the evaluators, which Chapter \ref{evaluation}
takes up.

\chapter{Ground truth and the evaluators}\label{evaluation}

\emph{R. Heller · Date: 2026-07-29 · Last changed: 2026-09-07}

\emph{Your embedding looks good. Is it correct?}

\section{What ``correct'' means here}\label{eval-question}

Machado, Behrisch and Telea built projections that score well on trustworthiness
and continuity while displaying structure the data does not have
\citep{machado2025necessary}. That is the false-\emph{positive} direction, and it needs no
ground truth: to show that a bad embedding scores high, it is enough to know
that the embedding is bad. This chapter runs the other direction. It hands each
evaluator the exact answer as a candidate embedding and asks whether the
evaluator recognises it. That cannot be done without an answer key, which is why
it is not usually done, and it is the one thing a crease pattern buys that no
real dataset can supply.

The complement matters because the two failures have different consequences. A
metric that rewards a wrong embedding wastes a reader's attention. A metric that
punishes the right one changes which method gets shipped.

``Correct'' is fixed here before anything is measured. An embedding is correct
when it \emph{is} the chart, up to a similarity of the plane. Chapter
\ref{folding-geometry} established that the flat pattern is not an approximation
to the folded sheet's intrinsic structure but that structure itself, so there is
an answer and it is exact. Everything in this chapter is a comparison against
that answer, and the object under test is not an embedding method. It is the
scoring function.

The disagreement about two-dimensional embeddings of single-cell data has
largely migrated from the embeddings to the evaluators. That framing is not this
book's: Lause, Berens and Kobak state it as their published position, concluding
that the case against those embeddings \citep{chari2023specious} rested on metrics
that ``focused on preservation of distances, where 2D PCA was unsurprisingly the
best'' \citep{lause2024elephant}. This book takes no position on what an embedding is
for. It measures how far each evaluator sits from the evaluator that has the
answer key.

Three things are established below, and one that the plan of record expected is
not. The evaluators disagree with the answer key, systematically and in a
direction that depends on which reference geometry they are given. Continuity
in particular cannot tell a perfect embedding from a plainly wrong one at any
fold angle. And the four metrics that a paper reports side by side are not four
pieces of evidence. What does \emph{not} survive is the shape of the law: the
inversion threshold is not a property of the fold angle alone, and saying so is
the point of having run the experiment rather than described it.

\section{Setup}\label{eval-setup}

\begin{Shaded}
\begin{Highlighting}[]
\NormalTok{audit }\OtherTok{\textless{}{-}} \FunctionTok{read\_run}\NormalTok{(}\StringTok{"data/processed/evaluator{-}audit.rds"}\NormalTok{, }\StringTok{"d7982cf7eb3afe7d"}\NormalTok{)}
\NormalTok{pilot }\OtherTok{\textless{}{-}} \FunctionTok{read\_run}\NormalTok{(}\StringTok{"data/processed/evaluator{-}audit{-}pilot.rds"}\NormalTok{, }\StringTok{"736906d41b87b8b3"}\NormalTok{)}
\NormalTok{grid  }\OtherTok{\textless{}{-}} \FunctionTok{read\_run}\NormalTok{(}\StringTok{"data/processed/benchmark{-}grid.rds"}\NormalTok{,  }\StringTok{"2483a468d82d4d64"}\NormalTok{)}

\NormalTok{prov }\OtherTok{\textless{}{-}} \FunctionTok{attr}\NormalTok{(audit, }\StringTok{"provenance"}\NormalTok{)}

\CommentTok{\# knitr rounds an inline numeric to options(digits), which turns 4e{-}16 into 0}
\CommentTok{\# and 0.0104 into 0.01. These return character, so a small number survives the}
\CommentTok{\# trip into prose intact.}
\NormalTok{fmt }\OtherTok{\textless{}{-}} \ControlFlowTok{function}\NormalTok{(x, }\AttributeTok{d =} \DecValTok{3}\DataTypeTok{L}\NormalTok{) }\FunctionTok{formatC}\NormalTok{(}\FunctionTok{signif}\NormalTok{(x, d), }\AttributeTok{format =} \StringTok{"fg"}\NormalTok{, }\AttributeTok{digits =}\NormalTok{ d)}
\NormalTok{sci }\OtherTok{\textless{}{-}} \ControlFlowTok{function}\NormalTok{(x, }\AttributeTok{d =} \DecValTok{2}\DataTypeTok{L}\NormalTok{) }\FunctionTok{formatC}\NormalTok{(x, }\AttributeTok{format =} \StringTok{"e"}\NormalTok{, }\AttributeTok{digits =}\NormalTok{ d }\SpecialCharTok{{-}} \DecValTok{1}\DataTypeTok{L}\NormalTok{)}

\NormalTok{p    }\OtherTok{\textless{}{-}} \FunctionTok{miura\_ori}\NormalTok{(}\AttributeTok{nx =} \DecValTok{6}\DataTypeTok{L}\NormalTok{, }\AttributeTok{ny =} \DecValTok{6}\DataTypeTok{L}\NormalTok{)}
\NormalTok{m    }\OtherTok{\textless{}{-}} \FunctionTok{sample\_manifold}\NormalTok{(p, }\AttributeTok{theta =} \FloatTok{0.8}\NormalTok{, }\AttributeTok{n =}\NormalTok{ N\_DEFAULT, }\AttributeTok{seed =}\NormalTok{ BENCH\_SEEDS[}\DecValTok{1}\NormalTok{])}
\NormalTok{U    }\OtherTok{\textless{}{-}}\NormalTok{ m}\SpecialCharTok{$}\NormalTok{truth}
\NormalTok{pca2 }\OtherTok{\textless{}{-}}\NormalTok{ stats}\SpecialCharTok{::}\FunctionTok{prcomp}\NormalTok{(m}\SpecialCharTok{$}\NormalTok{X, }\AttributeTok{rank. =}\NormalTok{ EMBED\_DIM)}\SpecialCharTok{$}\NormalTok{x[, }\FunctionTok{seq\_len}\NormalTok{(EMBED\_DIM)]}
\end{Highlighting}
\end{Shaded}

\texttt{audit} is the chapter's evidence: 19200 scored rows over
20 values of \(\theta\), 4 values of \(k\),
2 sample sizes, 20 seeds,
3 reference geometries and
2 candidate embeddings, produced by
\texttt{scripts/run-evaluator-audit.R} in 2.1 hours.
\texttt{pilot} is the narrower design that decided this chapter existed at all, and
\texttt{grid} is the main benchmark, read here for the metrics it already reports.

\section{Background}\label{eval-background}

\subsection{The headline metric: normalised Procrustes RMSE}\label{the-headline-metric-normalised-procrustes-rmse}

There is exactly one alignment convention in this book, and it lives in
\texttt{procrustes\_align()}. It is the \textbf{full} Procrustes fit
\citep{schonemann1966procrustes, gower1975procrustes}: the similarity group of
translation, rotation, \textbf{reflection} and one isotropic scale, and nothing else.

Each admission is a decision about what ``the same chart'' means.

\emph{Translation and rotation} are free because a chart has no origin and no
preferred axis. \emph{Isotropic scale} is free because two of the methods under test
have no scale to speak of --- a t-SNE embedding's extent is a function of its
perplexity and its iteration count \citep{vandermaaten2008tsne} --- and scoring them on
one would measure a setting rather than a geometry.

\emph{Reflection is the one that must be argued.} A mirrored unfolding is a correct
unfolding. Turn the flat sheet over and the crease pattern is the same pattern,
its mountains and valleys exchanged; no method in Part II is given anything from
which it could recover a handedness, and none claims to. The implementation
detail carries the whole argument: the rotation is taken as \(UV'\) from the
singular value decomposition of the cross-product matrix, with \textbf{no determinant
correction}. A version that forced \(\det R = 1\) would score an exactly mirrored
recovery as a failure.

\begin{Shaded}
\begin{Highlighting}[]
\NormalTok{rot    }\OtherTok{\textless{}{-}} \ControlFlowTok{function}\NormalTok{(a) }\FunctionTok{matrix}\NormalTok{(}\FunctionTok{c}\NormalTok{(}\FunctionTok{cos}\NormalTok{(a), }\FunctionTok{sin}\NormalTok{(a), }\SpecialCharTok{{-}}\FunctionTok{sin}\NormalTok{(a), }\FunctionTok{cos}\NormalTok{(a)), }\DecValTok{2}\DataTypeTok{L}\NormalTok{, }\DecValTok{2}\DataTypeTok{L}\NormalTok{)}
\NormalTok{moved  }\OtherTok{\textless{}{-}} \FloatTok{2.7} \SpecialCharTok{*}\NormalTok{ U }\SpecialCharTok{\%*\%} \FunctionTok{rot}\NormalTok{(}\FloatTok{1.1}\NormalTok{) }\SpecialCharTok{+} \FunctionTok{matrix}\NormalTok{(}\FunctionTok{c}\NormalTok{(}\DecValTok{4}\NormalTok{, }\SpecialCharTok{{-}}\DecValTok{9}\NormalTok{), }\FunctionTok{nrow}\NormalTok{(U), }\DecValTok{2}\DataTypeTok{L}\NormalTok{, }\AttributeTok{byrow =} \ConstantTok{TRUE}\NormalTok{)}
\NormalTok{mirror }\OtherTok{\textless{}{-}}\NormalTok{ moved }\SpecialCharTok{\%*\%} \FunctionTok{matrix}\NormalTok{(}\FunctionTok{c}\NormalTok{(}\SpecialCharTok{{-}}\DecValTok{1}\NormalTok{, }\DecValTok{0}\NormalTok{, }\DecValTok{0}\NormalTok{, }\DecValTok{1}\NormalTok{), }\DecValTok{2}\DataTypeTok{L}\NormalTok{, }\DecValTok{2}\DataTypeTok{L}\NormalTok{)}

\NormalTok{err\_moved  }\OtherTok{\textless{}{-}} \FunctionTok{reconstruction\_error}\NormalTok{(moved,  U)}
\NormalTok{err\_mirror }\OtherTok{\textless{}{-}} \FunctionTok{reconstruction\_error}\NormalTok{(mirror, U)}
\NormalTok{det\_moved  }\OtherTok{\textless{}{-}} \FunctionTok{det}\NormalTok{(}\FunctionTok{attr}\NormalTok{(}\FunctionTok{procrustes\_align}\NormalTok{(moved,  U), }\StringTok{"rotation"}\NormalTok{))}
\NormalTok{det\_mirror }\OtherTok{\textless{}{-}} \FunctionTok{det}\NormalTok{(}\FunctionTok{attr}\NormalTok{(}\FunctionTok{procrustes\_align}\NormalTok{(mirror, U), }\StringTok{"rotation"}\NormalTok{))}

\NormalTok{zb       }\OtherTok{\textless{}{-}} \FunctionTok{complex}\NormalTok{(}\AttributeTok{real =}\NormalTok{ U[, }\DecValTok{1}\DataTypeTok{L}\NormalTok{], }\AttributeTok{imaginary =}\NormalTok{ U[, }\DecValTok{2}\DataTypeTok{L}\NormalTok{])}\SpecialCharTok{\^{}}\DecValTok{2}
\NormalTok{bend     }\OtherTok{\textless{}{-}} \FunctionTok{cbind}\NormalTok{(}\FunctionTok{Re}\NormalTok{(zb), }\FunctionTok{Im}\NormalTok{(zb))}
\NormalTok{err\_exact }\OtherTok{\textless{}{-}} \FunctionTok{reconstruction\_error}\NormalTok{(U, U)}
\NormalTok{err\_bend  }\OtherTok{\textless{}{-}} \FunctionTok{reconstruction\_error}\NormalTok{(bend, U)}
\NormalTok{qnx\_bend }\OtherTok{\textless{}{-}} \FunctionTok{qnx}\NormalTok{(bend, }\FunctionTok{reference\_dist}\NormalTok{(m, }\StringTok{"chart"}\NormalTok{), }\AttributeTok{K =}\NormalTok{ K\_QNX)}
\end{Highlighting}
\end{Shaded}

A rotated, rescaled, translated copy of the chart scores
4.0e-16. Its mirror image scores 4.0e-16 ---
the same number to the last digit that means anything --- and the fitted rotation
has determinant -1 against 1 for the unmirrored copy.
Both are zero to floating point, which is what ``quotients out'' has to mean.

What is \emph{not} quotiented out is everything else, and the boundary is sharp.
Anisotropic scale and shear stay in the residual, because those are precisely the
distortions the isometry claim is about. So does a conformal bend: the map
\(z \mapsto z^2\) is a rotation and an isotropic scale at every point, so it moves
no neighbour out of any neighbourhood, while its scale factor varies across the
chart. It scores 0.255 here, and \(Q_{NX}\) at \(K\) =
20 against the chart calls the same configuration
0.963.

The score is normalised by the chart's own root-sum-square radius, which fixes
both ends. Zero is exact recovery. One is an embedding that explains none of the
chart's variance, and a random configuration reaches it to \(O(1/n)\). A bounded
metric matters because Chapter \ref{results} reports every result as an excess
over a floor, and an excess needs both ends nailed down.

\begin{figure}

{\centering \includegraphics[width=0.9\linewidth,alt={Three scatter panels of the same 800 points. Panel A is the exact chart, a filled rectangle whose colour ramps left to right with four roughly straight level-set paths running top to bottom. Panel B is the same rectangle, rotated, enlarged and flipped, so its colour ramp now runs diagonally from the dark top-right corner to the bright bottom-left one; its subtitle gives a score of four times ten to the minus sixteen. Panel C is the chart bent by a conformal map into a wedge that fans out towards the top, with the level sets splayed apart rather than parallel; it scores about one quarter.}]{09-ground-truth-evaluation_files/figure-latex/eval-fig-alignment-1} 

}

\caption{What the alignment quotients out, and what it does not. All three panels are scored against the same chart. Point colour is the first chart coordinate and the grey paths are its level sets, drawn so the structure survives greyscale printing.}\label{fig:eval-fig-alignment}
\end{figure}

Figure \ref{fig:eval-fig-alignment} is the convention in one row. Panel B is
the same chart and scores as such. Panel C is not, and the rank metrics will
disagree.

\subsection{The four metrics under audit}\label{the-four-metrics-under-audit}

Four scores are computed for every cell of the book's grid, and each is defined
here against a \textbf{reference geometry} \(D\) and a neighbourhood size \(k\). Write
\(r_{ij}\) for the rank of point \(j\) in point \(i\)'s ordering under \(D\), and
\(\hat r_{ij}\) for the same rank in the embedding.

\textbf{Trustworthiness} and \textbf{continuity} \citep{venna2001neighborhood, kaski2003trustworthiness} charge the embedding for the neighbours it invents
and the neighbours it loses:

\[M_1(k) = 1 - A(k)\sum_{i}\sum_{j \in U_k(i)}\big(r_{ij} - k\big), \qquad
  M_2(k) = 1 - A(k)\sum_{i}\sum_{j \in V_k(i)}\big(\hat r_{ij} - k\big),\]

with \(U_k(i)\) the points in the display's \(k\)-neighbourhood but not the
reference's, \(V_k(i)\) the reverse, and
\(A(k) = 2 / \big(Nk(2N - 3k - 1)\big)\) the scaling transcribed from Kaski et al.
Two things about \(A(k)\) travel with it and most re-implementations drop them.
It is the reciprocal of the largest penalty sum any embedding can produce, which
is why \(M_1 = 0\) is a floor rather than merely a small number --- and the authors
are explicit that it is a \emph{scaling}, not a proven bound on what a real
projection can attain, so a value slightly below zero is not necessarily a bug
and nothing here clamps one. Read as an information-retrieval pair they are
precision and recall of the neighbourhood \citep{venna2010information}.

\textbf{\(Q_{NX}(K)\)} is read off the co-ranking matrix of Lee and Verleysen
\citep{lee2009quality}, \(Q[k, l]\) counting the ordered pairs at rank \(k\) in the
reference and \(l\) in the embedding. Its top-left block is the fraction of
\(K\)-neighbourhoods preserved,
\(Q_{NX}(K) = \frac{1}{Kn}\sum_{k \le K}\sum_{l \le K} Q[k, l]\), which is
\(K/(n-1)\) for a random embedding rather than zero; \texttt{rnx()} applies that
correction and \texttt{auc\_qnx()} integrates it over \(K\) on a log scale.

\textbf{kNN preservation} is the mean agreement of the two \(k\)-nearest-neighbour
sets. One identity is worth stating rather than leaving to be rediscovered:
its overlap form is \emph{exactly} \(Q_{NX}(k)\). \texttt{knn\_preservation()} therefore
defaults to the Jaccard form, so that the book's four metrics are four and not
three.

Kruskal stress \citep{kruskal1964fit} is the fifth member of the usual set and is
absent here on purpose. Chapter \ref{folding-geometry} proves as a one-line
proposition that ambient-referenced stress must rank the exact unfolding below a
two-component PCA of the folded cloud, for every \(\theta > 0\); a simulation of a
theorem is not evidence. What is not entailed, and what this chapter measures,
is what the \emph{rank} metrics do, because they are invariant to a monotone
contraction of distances until neighbour identity changes.

\subsection{Three reference geometries}\label{three-reference-geometries}

Every one of the four takes a distance matrix and does not care where it came
from. That is the lever.

\begin{itemize}
\tightlist
\item
  \textbf{ambient} --- Euclidean distance in \(\mathbb{R}^3\) between the points as
  observed. This is what PCA, MDS, t-SNE and UMAP actually consume, and what
  essentially every published evaluation uses \citep{espadoto2021survey}.
\item
  \textbf{chart} --- Euclidean distance in the exact two-dimensional truth. On this
  construction that \emph{is} the geodesic distance on the folded surface, exactly
  and not as an estimate.
\item
  \textbf{graph} --- the estimate a method has to work with: a symmetric \(k\)-nearest-
  neighbour graph on the ambient points with all-pairs shortest paths, which is
  Isomap's internal geometry \citep{tenenbaum2000isomap}.
\end{itemize}

Nothing in \texttt{R/metrics.R} prefers one, and the argument is named \texttt{reference}
rather than \texttt{data} to keep that visible at the call site.

The book's own grid uses two of them at once, and the fact is worth stating
plainly because it is invisible in the output.

\begin{Shaded}
\begin{Highlighting}[]
\NormalTok{D\_amb }\OtherTok{\textless{}{-}} \FunctionTok{reference\_dist}\NormalTok{(m, }\StringTok{"ambient"}\NormalTok{)}
\NormalTok{D\_cha }\OtherTok{\textless{}{-}} \FunctionTok{reference\_dist}\NormalTok{(m, }\StringTok{"chart"}\NormalTok{)}

\NormalTok{conventions }\OtherTok{\textless{}{-}} \FunctionTok{data.frame}\NormalTok{(}
  \AttributeTok{candidate  =} \FunctionTok{c}\NormalTok{(}\StringTok{"the exact chart"}\NormalTok{, }\StringTok{"two{-}component PCA"}\NormalTok{),}
  \AttributeTok{rmse       =} \FunctionTok{c}\NormalTok{(}\FunctionTok{reconstruction\_error}\NormalTok{(U, U), }\FunctionTok{reconstruction\_error}\NormalTok{(pca2, U)),}
  \AttributeTok{qnx\_chart  =} \FunctionTok{c}\NormalTok{(}\FunctionTok{qnx}\NormalTok{(U, D\_cha, K\_QNX), }\FunctionTok{qnx}\NormalTok{(pca2, D\_cha, K\_QNX)),}
  \AttributeTok{knn\_ambient =} \FunctionTok{c}\NormalTok{(}\FunctionTok{knn\_preservation}\NormalTok{(D\_amb, U,    K\_DEFAULT),}
                  \FunctionTok{knn\_preservation}\NormalTok{(D\_amb, pca2, K\_DEFAULT)))}
\end{Highlighting}
\end{Shaded}

\texttt{grid} scores \(Q_{NX}\) at \(K\) = 20 against the
chart, and trustworthiness, continuity and kNN preservation
at \(k\) = 10 against ambient distance. Two
neighbourhood sizes and two reference geometries, in adjacent columns of one
table. Both choices are defensible and \texttt{R/constants.R} records why: the first
three ask whether the embedding preserved the neighbourhoods a method could
\emph{see}, so ambient is the honest reference; \(Q_{NX}\) is asked against the answer
key, because that is the question a benchmark with an answer key exists to be
able to ask.

The cost of leaving that unsaid is in the table. On the exact chart at
\(\theta\) = 0.8, the \(Q_{NX}\) column reads
1 and the kNN column reads
0.789. One says the answer key is perfect.
The other says it lost a fifth of its neighbourhoods. Both are correct, and
they are not answering the same question.

\section{The known-answer test}\label{eval-core}

The design is one sentence. Hand each evaluator the exact chart as though it
were a submission, alongside a two-component PCA of the folded cloud, and record
which one it prefers.

PCA is the sharp choice of wrong answer rather than a straw man. It is the wrong
chart --- it flattens the corrugation instead of unfolding it --- and it is
simultaneously the least-squares-optimal \emph{linear} reproduction of the ambient
configuration it was given \citep{pearson1901lines}. An evaluator that scores
reproduction of ambient distance is being handed the best available competitor
at exactly the thing it measures.

\begin{figure}

{\centering \includegraphics[width=0.9\linewidth,alt={Three scatter panels of 800 points each. Panel A is the exact chart: an evenly filled rectangle with a smooth left-to-right colour ramp and four level-set paths running nearly straight down it. Panel B is the two-component PCA at the chapter's working fold angle: the same rectangle, its points gathered into faint vertical bands and its level sets kinked. Panel C is the same PCA at the last fold angle of the sweep: the points are now gathered into pronounced stripes with visible gaps between them, and the level sets are strongly zigzagged. The three subtitles give scores of about zero, one tenth and one fifth.}]{09-ground-truth-evaluation_files/figure-latex/eval-fig-candidates-1} 

}

\caption{The candidates, Procrustes-aligned to the same answer key and scored against it. Point colour is the first chart coordinate and the grey paths are its level sets, which the exact chart draws as straight nested lines and a wrong chart does not. The wrongness is a function of the fold angle, which is why the audit sweeps it rather than fixing it.}\label{fig:eval-fig-candidates}
\end{figure}

Figure \ref{fig:eval-fig-candidates} is the comparison the audit runs at every
cell. The exact chart scores 5.7e-16. PCA scores
0.108 at the fold angle used throughout this chapter
and 0.21 at the last angle of the sweep, and the level sets it
draws are kinked where the creases ran. An \textbf{inversion} is an evaluator ranking
panel A below panel B or C.

Panel B is worth pausing on, because it is the reason this is a hard test rather
than an easy one. At a moderate fold a Miura is still close to planar, so its
principal-component projection is not a grotesque failure --- it is a mildly wrong
answer that looks reasonable. An evaluator that cannot separate panel A from
panel B is failing at a discrimination that matters, not at a trick.

\begin{figure}

{\centering \includegraphics[width=0.9\linewidth,alt={Four line panels, one per metric, sharing a vertical axis from zero to one labelled share of seeds inverted, against the folding parameter theta from zero to one. In the trustworthiness, kNN preservation and Q_NX panels almost every curve lies flat along zero, apart from a small bump near theta of 0.05 and, for the largest neighbourhood size, a low ragged band around one fifth that jumps to near one at the far right. The continuity panel is different: four sigmoid curves rise from zero to one, the largest neighbourhood size rising first at about theta 0.2 and the smallest last at about theta 0.9.}]{09-ground-truth-evaluation_files/figure-latex/eval-fig-inversion-1} 

}

\caption{Inversion rate against the folding parameter: the share of seeds in which the evaluator ranks the exact chart below a two-component PCA of the same cloud. Ambient reference, n = 800. Neighbourhood size is encoded by colour and line style together. The panels do not pool over sample size, because the two sample sizes are two different phenomena.}\label{fig:eval-fig-inversion}
\end{figure}

Figure \ref{fig:eval-fig-inversion} is the result at the larger sample size,
and it is four different pictures rather than one. Continuity inverts on a
clean sigmoid whose location slides left as \(k\) grows. Trustworthiness, kNN
preservation and \(Q_{NX}\) mostly do not invert at all: they rank the exact
chart first almost everywhere, with a bump near \(\theta\) = 0.05 where the
two candidates are barely distinguishable, and a low ragged band at the largest
\(k\) that only becomes a majority in the last cell of the sweep.

\begin{Shaded}
\begin{Highlighting}[]
\NormalTok{inv }\OtherTok{\textless{}{-}}\NormalTok{ audit}\SpecialCharTok{$}\NormalTok{inversions}
\NormalTok{flat\_inversions  }\OtherTok{\textless{}{-}} \FunctionTok{sum}\NormalTok{(inv}\SpecialCharTok{$}\NormalTok{inverted[inv}\SpecialCharTok{$}\NormalTok{theta }\SpecialCharTok{==} \DecValTok{0}\NormalTok{])}
\NormalTok{chart\_inversions }\OtherTok{\textless{}{-}} \FunctionTok{sum}\NormalTok{(inv}\SpecialCharTok{$}\NormalTok{inverted[inv}\SpecialCharTok{$}\NormalTok{reference }\SpecialCharTok{==} \StringTok{"chart"}\NormalTok{])}
\NormalTok{chart\_rows       }\OtherTok{\textless{}{-}} \FunctionTok{sum}\NormalTok{(inv}\SpecialCharTok{$}\NormalTok{reference }\SpecialCharTok{==} \StringTok{"chart"}\NormalTok{)}

\NormalTok{sc }\OtherTok{\textless{}{-}}\NormalTok{ audit}\SpecialCharTok{$}\NormalTok{scores}
\NormalTok{chart\_truth }\OtherTok{\textless{}{-}}\NormalTok{ sc[sc}\SpecialCharTok{$}\NormalTok{reference }\SpecialCharTok{==} \StringTok{"chart"} \SpecialCharTok{\&}\NormalTok{ sc}\SpecialCharTok{$}\NormalTok{candidate }\SpecialCharTok{==} \StringTok{"truth"}\NormalTok{, ]}
\NormalTok{truth\_exact }\OtherTok{\textless{}{-}} \FunctionTok{all}\NormalTok{(chart\_truth}\SpecialCharTok{$}\NormalTok{trust }\SpecialCharTok{==} \DecValTok{1} \SpecialCharTok{\&}\NormalTok{ chart\_truth}\SpecialCharTok{$}\NormalTok{cont }\SpecialCharTok{==} \DecValTok{1} \SpecialCharTok{\&}
\NormalTok{                     chart\_truth}\SpecialCharTok{$}\NormalTok{knn }\SpecialCharTok{==} \DecValTok{1} \SpecialCharTok{\&}\NormalTok{ chart\_truth}\SpecialCharTok{$}\NormalTok{qnx }\SpecialCharTok{==} \DecValTok{1}\NormalTok{)}
\end{Highlighting}
\end{Shaded}

Two controls make this a mechanism rather than a rate, and both come out clean.

At \(\theta\) = 0 the inversion count over the whole artefact is
0. A flat sheet contracts nothing, ambient distance \emph{is}
chart distance, and there is nothing to invert --- which is the prediction
Chapter \ref{folding-geometry}'s contraction result makes, evaluated at the one
point where the contraction vanishes.

Against the chart reference the count is 0 in
12800 rows. Not rare: never. And the reason is arithmetic rather than
luck --- handed the chart as its reference, every metric scores the chart as a
candidate at exactly one, which is checked here rather than asserted
(\texttt{truth\_exact} is TRUE). No candidate can beat a perfect score, so
no inversion is available.

The inversion is therefore not a property of the metric. It is a property of the
pair (metric, reference geometry), and the reference is doing most of the work.

An inversion is also not a \emph{bug}. It is the evaluator answering the question it
was actually asked. An ambient-referenced metric measures reproduction of
ambient distance; folding contracts ambient distance strictly; so the exact
unfolding reproduces ambient distance badly by construction, and something that
was fitted to reproduce it does better. That much is Chapter
\ref{folding-geometry}'s stress proposition and needed no experiment.

What did need one is that a \emph{rank} metric inverts at all. Rank metrics are
invariant to any monotone transformation of distances, and the argument that
they should therefore be safe here is the argument this chapter was budgeted to
test. It fails on a detail: folding is not a monotone map on ambient distance.
The contraction factor is \(\sin(\rho/2)\) \emph{per crease crossed}, so a pair
straddling a fold contracts while a pair of the same chart distance inside one
facet does not, and their order swaps. Neighbour identity therefore changes,
and the amount it changes is measurable.

\begin{Shaded}
\begin{Highlighting}[]
\NormalTok{rk }\OtherTok{\textless{}{-}} \ControlFlowTok{function}\NormalTok{(D) \{ }\FunctionTok{diag}\NormalTok{(D) }\OtherTok{\textless{}{-}} \SpecialCharTok{{-}}\ConstantTok{Inf}
                    \FunctionTok{t}\NormalTok{(}\FunctionTok{apply}\NormalTok{(D, }\DecValTok{1}\DataTypeTok{L}\NormalTok{, rank, }\AttributeTok{ties.method =} \StringTok{"first"}\NormalTok{)) }\SpecialCharTok{{-}} \DecValTok{1}\DataTypeTok{L}\NormalTok{ \}}
\NormalTok{agree }\OtherTok{\textless{}{-}} \ControlFlowTok{function}\NormalTok{(r1, r2, K) }\FunctionTok{mean}\NormalTok{(}\FunctionTok{rowSums}\NormalTok{(r1 }\SpecialCharTok{\textgreater{}=} \DecValTok{1}\DataTypeTok{L} \SpecialCharTok{\&}\NormalTok{ r1 }\SpecialCharTok{\textless{}=}\NormalTok{ K }\SpecialCharTok{\&}\NormalTok{ r2 }\SpecialCharTok{\textgreater{}=} \DecValTok{1}\DataTypeTok{L} \SpecialCharTok{\&}\NormalTok{ r2 }\SpecialCharTok{\textless{}=}\NormalTok{ K)) }\SpecialCharTok{/}\NormalTok{ K}

\NormalTok{m\_flat  }\OtherTok{\textless{}{-}} \FunctionTok{sample\_manifold}\NormalTok{(p, }\AttributeTok{theta =} \DecValTok{0}\NormalTok{, }\AttributeTok{n =}\NormalTok{ N\_DEFAULT, }\AttributeTok{seed =}\NormalTok{ BENCH\_SEEDS[}\DecValTok{1}\NormalTok{])}
\NormalTok{ra\_flat }\OtherTok{\textless{}{-}} \FunctionTok{rk}\NormalTok{(}\FunctionTok{reference\_dist}\NormalTok{(m\_flat, }\StringTok{"ambient"}\NormalTok{))}
\NormalTok{rc\_flat }\OtherTok{\textless{}{-}} \FunctionTok{rk}\NormalTok{(}\FunctionTok{reference\_dist}\NormalTok{(m\_flat, }\StringTok{"chart"}\NormalTok{))}
\NormalTok{ra }\OtherTok{\textless{}{-}} \FunctionTok{rk}\NormalTok{(D\_amb); rc }\OtherTok{\textless{}{-}} \FunctionTok{rk}\NormalTok{(D\_cha)}

\NormalTok{moved\_flat }\OtherTok{\textless{}{-}} \DecValTok{1} \SpecialCharTok{{-}} \FunctionTok{agree}\NormalTok{(ra\_flat, rc\_flat, K\_DEFAULT)}
\NormalTok{moved\_fold }\OtherTok{\textless{}{-}} \DecValTok{1} \SpecialCharTok{{-}} \FunctionTok{agree}\NormalTok{(ra, rc, K\_DEFAULT)}
\end{Highlighting}
\end{Shaded}

At \(\theta\) = 0 a point's ambient neighbours and its chart neighbours are
the same set: the share that differ is 0. At \(\theta\) =
0.8 it is 0.13. That number is the whole mechanism.
An eighth of every neighbourhood has been rewritten by the fold, the exact chart
is scored on reproducing the rewritten version, and it does not reproduce it ---
because it is right.

\subsection{Why the obvious fix fails}\label{why-the-obvious-fix-fails}

The obvious repair is to keep the metric and replace the reference: score against
estimated geodesics instead of ambient distance, which is what a practitioner
who has read Isomap would reach for. It does nothing, and it cannot.

\begin{Shaded}
\begin{Highlighting}[]
\NormalTok{D\_gra }\OtherTok{\textless{}{-}} \FunctionTok{reference\_dist}\NormalTok{(m, }\StringTok{"graph"}\NormalTok{, }\AttributeTok{k =} \FunctionTok{max}\NormalTok{(prov}\SpecialCharTok{$}\NormalTok{k))}
\NormalTok{rg    }\OtherTok{\textless{}{-}} \FunctionTok{rk}\NormalTok{(D\_gra)}

\NormalTok{Ks          }\OtherTok{\textless{}{-}}\NormalTok{ prov}\SpecialCharTok{$}\NormalTok{k}
\NormalTok{agree\_graph }\OtherTok{\textless{}{-}} \FunctionTok{vapply}\NormalTok{(Ks, }\ControlFlowTok{function}\NormalTok{(K) }\FunctionTok{agree}\NormalTok{(ra, rg, K), }\FunctionTok{numeric}\NormalTok{(}\DecValTok{1}\NormalTok{))}
\NormalTok{agree\_chart }\OtherTok{\textless{}{-}} \FunctionTok{vapply}\NormalTok{(Ks, }\ControlFlowTok{function}\NormalTok{(K) }\FunctionTok{agree}\NormalTok{(ra, rc, K), }\FunctionTok{numeric}\NormalTok{(}\DecValTok{1}\NormalTok{))}
\NormalTok{graph\_above }\OtherTok{\textless{}{-}} \FunctionTok{all}\NormalTok{(D\_gra }\SpecialCharTok{\textgreater{}=}\NormalTok{ D\_amb }\SpecialCharTok{{-}}\NormalTok{ TOL}\SpecialCharTok{$}\NormalTok{default)}

\NormalTok{key      }\OtherTok{\textless{}{-}} \FunctionTok{c}\NormalTok{(}\StringTok{"k"}\NormalTok{, }\StringTok{"theta"}\NormalTok{, }\StringTok{"n"}\NormalTok{, }\StringTok{"seed"}\NormalTok{, }\StringTok{"candidate"}\NormalTok{)}
\NormalTok{paired   }\OtherTok{\textless{}{-}} \FunctionTok{merge}\NormalTok{(sc[sc}\SpecialCharTok{$}\NormalTok{reference }\SpecialCharTok{==} \StringTok{"ambient"}\NormalTok{, ], sc[sc}\SpecialCharTok{$}\NormalTok{reference }\SpecialCharTok{==} \StringTok{"graph"}\NormalTok{, ],}
                  \AttributeTok{by =}\NormalTok{ key, }\AttributeTok{suffixes =} \FunctionTok{c}\NormalTok{(}\StringTok{".a"}\NormalTok{, }\StringTok{".g"}\NormalTok{))}
\NormalTok{gap\_graph }\OtherTok{\textless{}{-}} \FunctionTok{vapply}\NormalTok{(}\FunctionTok{names}\NormalTok{(METRIC\_LABEL), }\ControlFlowTok{function}\NormalTok{(v)}
  \FunctionTok{max}\NormalTok{(}\FunctionTok{abs}\NormalTok{(paired[[}\FunctionTok{paste0}\NormalTok{(v, }\StringTok{".a"}\NormalTok{)]] }\SpecialCharTok{{-}}\NormalTok{ paired[[}\FunctionTok{paste0}\NormalTok{(v, }\StringTok{".g"}\NormalTok{)]])), }\FunctionTok{numeric}\NormalTok{(}\DecValTok{1}\NormalTok{))}
\end{Highlighting}
\end{Shaded}

The graph reference agrees with the ambient reference on
1, 1, 1, 1 of the neighbour sets at
\(k\) of 5, 10, 20, 40 --- that is, on all of them. The chart
reference, over the same points, agrees on only
0.908, 0.873, 0.865, 0.879.

The identity is exact and has a one-line proof. Any path through a Euclidean
graph is at least as long as the straight line between its endpoints, so graph
distance is never below ambient distance (\texttt{graph\_above} is TRUE);
and the graph contains an edge to each of a point's \(k\) nearest ambient
neighbours, whose shortest path is therefore that edge. So for every
\(K \le k\) the two geometries have the \emph{same} \(K\) nearest neighbours in the
\emph{same} order. A rank metric evaluated at \(K \le k\) cannot see the difference.

The audit says the same thing across all 19200 scored rows: the largest
gap between an ambient-referenced and a graph-referenced score is
0 for continuity, kNN
preservation and \(Q_{NX}\), and 0.0022 for
trustworthiness. Trustworthiness is the only one that moves, and that too
follows: it is the one metric that asks for the reference rank of a point the
\emph{embedding} placed nearby, which can lie far outside the graph's own \(k\).

An estimated geodesic is not a different question from an ambient distance at
the scale a rank metric looks at. Only an exact chart is, and only a benchmark
with an answer key has one. This is a stronger statement than ``graph distances
are fragile'' --- the short-circuit failures of Chapter \ref{geodesic}
\citep{balasubramanian2002isomap} are about the long-range structure, and the point
here is that even a \emph{perfectly} connected graph changes nothing locally.

\section{Results}\label{eval-results}

\subsection{How optimistic, and in what units}\label{how-optimistic-and-in-what-units}

``Optimistic by X'' is not a statement the evidence supports, and the reason is
visible before any inversion is counted.

\begin{figure}

{\centering \includegraphics[width=0.9\linewidth,alt={Five line panels sharing a vertical axis from zero to one labelled score, plotted against the folding parameter theta. In the trustworthiness and continuity panels both curves are flat along the top of the axis at one across the entire range and are indistinguishable. In the kNN preservation and Q_NX panels both curves fall away as theta grows, the exact chart staying above the PCA curve and reaching about 0.58 and 0.71 at the right-hand end. In the fifth panel, normalised Procrustes RMSE, the exact chart is flat at zero while the PCA curve rises to about 0.21.}]{09-ground-truth-evaluation_files/figure-latex/eval-fig-scores-1} 

}

\caption{The two candidates scored against ambient distance across the whole fold sweep, at the book's default neighbourhood size and n = 800. All five quantities lie on the unit interval and share one axis, so the comparison is not an artefact of scaling. Candidate is encoded by colour and line style together.}\label{fig:eval-fig-scores}
\end{figure}

Figure \ref{fig:eval-fig-scores} puts all five quantities on one axis, which
they share because the normalised Procrustes score also lives on the unit
interval. Trustworthiness and continuity are two flat lines at the top of the
panel. They do not separate a configuration whose error is zero from one whose
error reaches a fifth of the chart's radius, at any fold angle in the sweep.

\begin{Shaded}
\begin{Highlighting}[]
\NormalTok{amb }\OtherTok{\textless{}{-}}\NormalTok{ inv[inv}\SpecialCharTok{$}\NormalTok{reference }\SpecialCharTok{==} \StringTok{"ambient"}\NormalTok{, ]}
\NormalTok{amb}\SpecialCharTok{$}\NormalTok{margin }\OtherTok{\textless{}{-}}\NormalTok{ amb}\SpecialCharTok{$}\NormalTok{pca2 }\SpecialCharTok{{-}}\NormalTok{ amb}\SpecialCharTok{$}\NormalTok{truth      }\CommentTok{\# positive when the evaluator prefers PCA}
\NormalTok{big }\OtherTok{\textless{}{-}} \FunctionTok{max}\NormalTok{(prov}\SpecialCharTok{$}\NormalTok{n)}

\NormalTok{optimism }\OtherTok{\textless{}{-}} \FunctionTok{do.call}\NormalTok{(rbind, }\FunctionTok{lapply}\NormalTok{(}\FunctionTok{names}\NormalTok{(METRIC\_LABEL), }\ControlFlowTok{function}\NormalTok{(mt) \{}
\NormalTok{  d }\OtherTok{\textless{}{-}}\NormalTok{ amb[amb}\SpecialCharTok{$}\NormalTok{metric }\SpecialCharTok{==}\NormalTok{ mt }\SpecialCharTok{\&}\NormalTok{ amb}\SpecialCharTok{$}\NormalTok{n }\SpecialCharTok{==}\NormalTok{ big, ]}
  \FunctionTok{data.frame}\NormalTok{(}\AttributeTok{metric         =}\NormalTok{ METRIC\_LABEL[[mt]],}
             \AttributeTok{inverted       =} \FunctionTok{mean}\NormalTok{(d}\SpecialCharTok{$}\NormalTok{inverted),}
             \AttributeTok{widest\_pca     =} \FunctionTok{max}\NormalTok{(d}\SpecialCharTok{$}\NormalTok{margin),}
             \AttributeTok{widest\_truth   =} \SpecialCharTok{{-}}\FunctionTok{min}\NormalTok{(d}\SpecialCharTok{$}\NormalTok{margin),}
             \AttributeTok{pca\_reportable =} \FunctionTok{mean}\NormalTok{(d}\SpecialCharTok{$}\NormalTok{margin  }\SpecialCharTok{\textgreater{}}\NormalTok{ MIN\_REPORTABLE),}
             \AttributeTok{truth\_reportable =} \FunctionTok{mean}\NormalTok{(}\SpecialCharTok{{-}}\NormalTok{d}\SpecialCharTok{$}\NormalTok{margin }\SpecialCharTok{\textgreater{}}\NormalTok{ MIN\_REPORTABLE))}
\NormalTok{\}))}

\NormalTok{cont\_cells      }\OtherTok{\textless{}{-}}\NormalTok{ amb[amb}\SpecialCharTok{$}\NormalTok{metric }\SpecialCharTok{==} \StringTok{"cont"} \SpecialCharTok{\&}\NormalTok{ amb}\SpecialCharTok{$}\NormalTok{n }\SpecialCharTok{==}\NormalTok{ big, ]}
\NormalTok{cont\_reportable }\OtherTok{\textless{}{-}} \FunctionTok{sum}\NormalTok{(cont\_cells}\SpecialCharTok{$}\NormalTok{margin }\SpecialCharTok{\textgreater{}}\NormalTok{ MIN\_REPORTABLE)}

\NormalTok{one\_cell }\OtherTok{\textless{}{-}} \ControlFlowTok{function}\NormalTok{(ref, cand) \{}
\NormalTok{  d }\OtherTok{\textless{}{-}}\NormalTok{ sc[sc}\SpecialCharTok{$}\NormalTok{reference }\SpecialCharTok{==}\NormalTok{ ref }\SpecialCharTok{\&}\NormalTok{ sc}\SpecialCharTok{$}\NormalTok{candidate }\SpecialCharTok{==}\NormalTok{ cand }\SpecialCharTok{\&}
\NormalTok{            sc}\SpecialCharTok{$}\NormalTok{k }\SpecialCharTok{==}\NormalTok{ K\_DEFAULT }\SpecialCharTok{\&}\NormalTok{ sc}\SpecialCharTok{$}\NormalTok{n }\SpecialCharTok{==}\NormalTok{ big, ]}
  \FunctionTok{aggregate}\NormalTok{(}\FunctionTok{cbind}\NormalTok{(trust, cont, knn, qnx) }\SpecialCharTok{\textasciitilde{}}\NormalTok{ theta, }\AttributeTok{data =}\NormalTok{ d, }\AttributeTok{FUN =}\NormalTok{ mean)}
\NormalTok{\}}
\NormalTok{ceiling   }\OtherTok{\textless{}{-}} \FunctionTok{one\_cell}\NormalTok{(}\StringTok{"ambient"}\NormalTok{, }\StringTok{"truth"}\NormalTok{)   }\CommentTok{\# what a perfect answer scores}
\NormalTok{chart\_pca }\OtherTok{\textless{}{-}} \FunctionTok{one\_cell}\NormalTok{(}\StringTok{"chart"}\NormalTok{,   }\StringTok{"pca2"}\NormalTok{)    }\CommentTok{\# what a wrong answer scores}

\NormalTok{knitr}\SpecialCharTok{::}\FunctionTok{kable}\NormalTok{(}
\NormalTok{  optimism, }\AttributeTok{digits =} \DecValTok{3}\NormalTok{, }\AttributeTok{row.names =} \ConstantTok{FALSE}\NormalTok{,}
  \AttributeTok{col.names =} \FunctionTok{c}\NormalTok{(}\StringTok{"metric"}\NormalTok{, }\StringTok{"inverted"}\NormalTok{, }\StringTok{"widest PCA lead"}\NormalTok{, }\StringTok{"widest truth lead"}\NormalTok{,}
                \StringTok{"PCA wins reportably"}\NormalTok{, }\StringTok{"truth wins reportably"}\NormalTok{),}
  \AttributeTok{caption =} \StringTok{"Four metrics, four answers. Ambient reference, at the larger sample size, over every fold angle, neighbourhood size and seed in the audit. \textquotesingle{}Reportably\textquotesingle{} means by more than the book\textquotesingle{}s pre{-}registered reportable difference."}\NormalTok{)}
\end{Highlighting}
\end{Shaded}

\begin{table}

\caption{\label{tab:eval-optimism}Four metrics, four answers. Ambient reference, at the larger sample size, over every fold angle, neighbourhood size and seed in the audit. 'Reportably' means by more than the book's pre-registered reportable difference.}
\centering
\begin{tabular}[t]{l|r|r|r|r|r}
\hline
metric & inverted & widest PCA lead & widest truth lead & PCA wins reportably & truth wins reportably\\
\hline
trustworthiness & 0.013 & 0.000 & 0.009 & 0.000 & 0.000\\
\hline
continuity & 0.418 & 0.010 & 0.000 & 0.000 & 0.000\\
\hline
kNN preservation & 0.054 & 0.029 & 0.197 & 0.003 & 0.459\\
\hline
Q\_NX & 0.056 & 0.023 & 0.157 & 0.001 & 0.329\\
\hline
\end{tabular}
\end{table}

Table \ref{tab:eval-optimism} is the honest summary and it does not reduce to a
number.

Continuity inverts on
41.8 per cent
of cells, and its widest lead over the truth anywhere at this sample size is
0.0104. It clears the
book's pre-registered reportable difference of 0.02 in
0 of 1600 cells. So the claim
the numbers support is \emph{not} ``continuity prefers a wrong embedding by X''. It is
that \textbf{continuity cannot distinguish a perfect embedding from a plainly wrong
one at all}, and reports both as excellent. The weaker version sounds bolder
and is not what happened.

Trustworthiness is inert too, and more nearly symmetrically so: it inverts on
1.3
per cent of cells and never moves by more than
0.0087
either way.

kNN preservation and \(Q_{NX}\) do the job asked of them. They rank the exact
chart first, by a margin clearing the reportable difference in
45.9
and
32.9
per cent of cells, and their widest lead for the truth is
0.197.
Ranking the right answer first is the whole of what is asked, and two of four
metrics do it.

They do it while calling the right answer mediocre, which is the second half of
the finding and the one with a practical consequence. The \texttt{ceiling} object holds
the score the answer key itself earns under the standard recipe. At the last
cell of the sweep, ambient-referenced kNN preservation at \(k\) = 10
rates the exact unfolding 0.58 and \(Q_{NX}\) rates it
0.714. Any reader shown
0.58 would call that embedding a partial failure. It
is the answer.

One more reading is available here and it corrects a tempting summary. The
result so far invites ``use the right reference geometry and the metrics behave'',
and that is only half true. Swapping ambient for the exact chart fixes the
\emph{ordering}: with the answer key as reference, nothing inverts anywhere. It does
not give trustworthiness and continuity a range to answer in. Scored against the
exact chart at the hardest cell of the sweep, the two-component PCA --- the
configuration whose true error is 0.212 ---
still earns trustworthiness 0.995 and continuity
0.995. A perfect reference geometry and a wrong
embedding, and the metric reports near-perfection. The reference decides who
wins; the metric decides whether the win means anything, and for these two it
does not.

\subsection{Four metrics, one number}\label{four-metrics-one-number}

The four are reported side by side in almost every evaluation, and read as four
independent lines of evidence. On the audit they gave four different answers. On
the benchmark grid --- nine methods, three noise models, the whole \(\theta\) sweep,
9199 scored cells --- they largely do not.

\begin{figure}

{\centering \includegraphics[width=0.9\linewidth,alt={Three scatter panels, each of roughly nine thousand semi-transparent points. The left panel plots Q_NX against kNN preservation and the cloud lies close to a rising curve, annotated Spearman 0.92. The middle panel plots trustworthiness against continuity, a looser rising cloud crammed into the top-right corner above 0.9 on both axes, annotated Spearman 0.70. The right panel plots trustworthiness against normalised Procrustes RMSE: a dense horizontal band pinned near trustworthiness one that stretches across the whole range of the error axis, annotated Spearman minus 0.79.}]{09-ground-truth-evaluation_files/figure-latex/eval-fig-agreement-1} 

}

\caption{Metric agreement across the benchmark grid. Each point is one scored cell. The two rank metrics in the left panel are computed at different neighbourhood sizes against different reference geometries and still move together; the right panel is the gap this chapter is about.}\label{fig:eval-fig-agreement}
\end{figure}

Figure \ref{fig:eval-fig-agreement}, left panel: \(Q_{NX}\) and kNN preservation
rank the grid's cells at Spearman 0.921 --- across
9199 cells, at \emph{different} neighbourhood sizes,
against \emph{different} reference geometries. They are near-redundant despite being
configured to ask two different questions, because both count neighbour-set
overlap and neighbour sets are what varies. Trustworthiness and continuity agree
at 0.695. Reporting three of these is not three pieces of
evidence; it is one piece of evidence and two restatements, and the practice of
tabling them together \citep{espadoto2021survey} hides that.

The four fall into two families rather than four instruments, and the split is
the one the definitions predict. \(Q_{NX}\) and kNN preservation count how much of
a neighbour \emph{set} survived, and their overlap forms are the same number, so
correlation near one is the expected result rather than a discovery.
Trustworthiness and continuity weight each lost or invented neighbour by \emph{how
far} it moved in ranks, which is a different functional and gives them the
looser agreement above. A reader wanting two independent readings should take
one from each family and a distance-based score alongside --- which is what
\texttt{rank\_metrics()} returns in one pass and what the book's grid records per cell.

The right panel is what a chapter on evaluators is for.
74.1 per cent of the grid's scored cells have
trustworthiness above 0.99, and 46.6 per cent of
\emph{those} have a normalised Procrustes error above 0.2 --- an embedding that has
lost a fifth of the chart. The grid's median cell scores
0.227 on Procrustes and
0.996 on trustworthiness. A metric whose median is
that close to its ceiling is not measuring the range it is being asked about.

\section{Diagnostics}\label{eval-diagnostics}

\subsection{Does k change the answer?}\label{does-k-change-the-answer}

It does, and this is where the chapter's plan of record has to be corrected
against its own artefact.

The contribution was specified as a \textbf{threshold law in \(\theta\)}: rank metrics
are invariant to monotone contraction until neighbour identity changes, so they
should invert only above a fold angle of about 0.8 while stress inverts
everywhere. The first half is right and is why the rank metrics behave at all
differently from stress. The threshold is not.

\begin{figure}

{\centering \includegraphics[width=0.9\linewidth,alt={A line chart with two nearly coincident descending curves. The horizontal axis is the neighbourhood fraction k over n on a logarithmic scale from about 0.006 to 0.1; the vertical axis is the threshold fold angle from zero to one. Both curves fall steeply from about 0.9 at the left to about 0.1 to 0.2 at the right. Each point is labelled with its neighbourhood size k; the two curves interleave, so that k of 5 at n of 400 sits at the same place as k of 10 at n of 800.}]{09-ground-truth-evaluation_files/figure-latex/eval-fig-threshold-1} 

}

\caption{Where continuity's inversion sets in. The threshold is the smallest fold angle from which the evaluator inverts on at least half of seeds and never recovers. Plotted against the neighbourhood fraction k/n, the two sample sizes fall on one curve; plotted against k alone they would not. Sample size is encoded by colour, shape and line style together.}\label{fig:eval-fig-threshold}
\end{figure}

Figure \ref{fig:eval-fig-threshold} locates it. At the smallest neighbourhood
the threshold is 0.9,
which is roughly where the plan expected it. At the largest it is
0.2 --- near the flat
sheet, fourteen grid steps earlier. A single threshold in \(\theta\) does not exist,
because \(k\) moves it across most of the available range.

What the two sample sizes do is settle what \(k\) is a proxy for. Plotted against
the neighbourhood \emph{fraction} \(k/n\) the two curves nearly coincide: at
\(k/n\) = 0.0125 the threshold is
0.8 at the smaller
sample and 0.8 at
the larger --- the same value, at half the neighbourhood size and twice the
points. The mechanism follows: a \(k\)-neighbourhood covers a chart radius
that scales as \(\sqrt{k/n}\), and continuity inverts once that radius is large
enough that a typical neighbourhood spans a crease. Fold angle and neighbourhood
fraction are two ways of buying the same thing.

So the law the chapter can state is a threshold surface in \((\theta, k/n)\) for
\textbf{continuity}, and for continuity alone: at the larger sample size no other
metric has a \(\theta_{1/2}\) at any \(k\) in the design. The coincidence with
Chapter \ref{geodesic}'s analytic short-circuit onset that the plan
anticipated is \emph{not} established here --- this artefact does not measure the
short circuit, and the two would have to be run against each other before
either could be said to predict the other.

\subsection{What to report when you have no answer key}\label{what-to-report-when-you-have-no-answer-key}

Three recommendations follow, and they cost nothing.

\textbf{Report the reference geometry and the neighbourhood size with the score.}
``Trustworthiness 0.99'' is not a result. The audit's whole content is that the
same metric on the same embedding answers differently depending on which
distance matrix it was handed, and by more than the difference between a perfect
embedding and a bad one.

\textbf{Report the score against its floor.} \texttt{metric\_floor()} returns the smallest
normalised Procrustes error \emph{any} \(d\)-dimensional configuration of these points
could achieve --- the tail of the chart's own spectrum, by Eckart--Young
\citep{eckart1936approximation}.

\begin{Shaded}
\begin{Highlighting}[]
\NormalTok{floor\_2d }\OtherTok{\textless{}{-}} \FunctionTok{metric\_floor}\NormalTok{(m, }\DecValTok{2}\DataTypeTok{L}\NormalTok{)}
\NormalTok{floor\_1d }\OtherTok{\textless{}{-}} \FunctionTok{metric\_floor}\NormalTok{(m, }\DecValTok{1}\DataTypeTok{L}\NormalTok{)}
\NormalTok{prod4    }\OtherTok{\textless{}{-}} \FunctionTok{product\_manifold}\NormalTok{(m, }\FunctionTok{sample\_manifold}\NormalTok{(p, }\AttributeTok{theta =} \FloatTok{0.5}\NormalTok{, }\AttributeTok{n =}\NormalTok{ N\_DEFAULT,}
                                                \AttributeTok{seed =}\NormalTok{ BENCH\_SEEDS[}\DecValTok{2}\NormalTok{]))}
\NormalTok{floor\_prod }\OtherTok{\textless{}{-}} \FunctionTok{metric\_floor}\NormalTok{(prod4, EMBED\_DIM)}
\NormalTok{grid\_floor }\OtherTok{\textless{}{-}} \FunctionTok{unique}\NormalTok{(grid}\SpecialCharTok{$}\NormalTok{floor[}\SpecialCharTok{!}\FunctionTok{is.na}\NormalTok{(grid}\SpecialCharTok{$}\NormalTok{floor)])}
\end{Highlighting}
\end{Shaded}

On this sample it returns 0, and the grid's own \texttt{floor} column is
0 in every cell it holds. That is worth one sentence rather than a
column: a two-dimensional chart embedded in two dimensions has nothing
unreachable in it, so every error reported in Chapter \ref{results} is loss the
method is responsible for. Force the same chart into one dimension and the floor
is 0.652; take Chapter \ref{benchmarks}'s product of two
folded sheets, an intrinsically four-dimensional object, and it is
0.649. A method scoring 0.649 has thrown
away everything on the main grid and hit the bound exactly on the product.
Reporting the number without its floor distinguishes neither.

\textbf{Report the answer key's own score.} The floor bounds the Procrustes metric
and nothing else --- there is no analogous bound for \(Q_{NX}\) or trustworthiness,
and this chapter does not supply one. What it does supply is the empirical
substitute that only a known-answer benchmark can produce: the \texttt{ceiling} object
above is the score a \emph{perfect} embedding attains under each convention. Where
that number is 0.58, an embedding scoring
0.58 has not half failed. It has succeeded, and the
shortfall belongs to the metric.

\begin{Shaded}
\begin{Highlighting}[]
\NormalTok{DU }\OtherTok{\textless{}{-}} \FunctionTok{as.matrix}\NormalTok{(stats}\SpecialCharTok{::}\FunctionTok{dist}\NormalTok{(U))}
\FunctionTok{set.seed}\NormalTok{(}\DecValTok{11}\DataTypeTok{L}\NormalTok{)}
\NormalTok{J }\OtherTok{\textless{}{-}} \FunctionTok{matrix}\NormalTok{(stats}\SpecialCharTok{::}\FunctionTok{rnorm}\NormalTok{(}\DecValTok{2}\DataTypeTok{L} \SpecialCharTok{*} \FunctionTok{nrow}\NormalTok{(U)), }\AttributeTok{ncol =} \DecValTok{2}\DataTypeTok{L}\NormalTok{)}
\NormalTok{distortion }\OtherTok{\textless{}{-}} \FunctionTok{list}\NormalTok{(}
  \StringTok{\textasciigrave{}}\AttributeTok{conformal bend}\StringTok{\textasciigrave{}} \OtherTok{=} \ControlFlowTok{function}\NormalTok{(t) \{}
\NormalTok{    z }\OtherTok{\textless{}{-}} \FunctionTok{complex}\NormalTok{(}\AttributeTok{real =}\NormalTok{ U[, }\DecValTok{1}\DataTypeTok{L}\NormalTok{], }\AttributeTok{imaginary =}\NormalTok{ U[, }\DecValTok{2}\DataTypeTok{L}\NormalTok{])}\SpecialCharTok{\^{}}\NormalTok{(}\DecValTok{1} \SpecialCharTok{+}\NormalTok{ t)}
    \FunctionTok{cbind}\NormalTok{(}\FunctionTok{Re}\NormalTok{(z), }\FunctionTok{Im}\NormalTok{(z))}
\NormalTok{  \},}
  \AttributeTok{shear  =} \ControlFlowTok{function}\NormalTok{(t) }\FunctionTok{cbind}\NormalTok{(U[, }\DecValTok{1}\DataTypeTok{L}\NormalTok{] }\SpecialCharTok{+}\NormalTok{ t }\SpecialCharTok{*}\NormalTok{ U[, }\DecValTok{2}\DataTypeTok{L}\NormalTok{], U[, }\DecValTok{2}\DataTypeTok{L}\NormalTok{]),}
  \AttributeTok{jitter =} \ControlFlowTok{function}\NormalTok{(t) U }\SpecialCharTok{+}\NormalTok{ t }\SpecialCharTok{*}\NormalTok{ J)}

\NormalTok{match\_error }\OtherTok{\textless{}{-}} \ControlFlowTok{function}\NormalTok{(f, target, }\AttributeTok{hi =} \DecValTok{6}\NormalTok{) \{          }\CommentTok{\# bisect on the strength}
\NormalTok{  lo }\OtherTok{\textless{}{-}} \FloatTok{1e{-}6}
  \ControlFlowTok{for}\NormalTok{ (i }\ControlFlowTok{in} \FunctionTok{seq\_len}\NormalTok{(}\DecValTok{50}\DataTypeTok{L}\NormalTok{)) \{}
\NormalTok{    mid }\OtherTok{\textless{}{-}}\NormalTok{ (lo }\SpecialCharTok{+}\NormalTok{ hi) }\SpecialCharTok{/} \DecValTok{2}
    \ControlFlowTok{if}\NormalTok{ (}\FunctionTok{reconstruction\_error}\NormalTok{(}\FunctionTok{f}\NormalTok{(mid), U) }\SpecialCharTok{\textless{}}\NormalTok{ target) lo }\OtherTok{\textless{}{-}}\NormalTok{ mid }\ControlFlowTok{else}\NormalTok{ hi }\OtherTok{\textless{}{-}}\NormalTok{ mid}
\NormalTok{  \}}
  \FunctionTok{f}\NormalTok{((lo }\SpecialCharTok{+}\NormalTok{ hi) }\SpecialCharTok{/} \DecValTok{2}\NormalTok{)}
\NormalTok{\}}

\NormalTok{battery }\OtherTok{\textless{}{-}} \FunctionTok{do.call}\NormalTok{(rbind, }\FunctionTok{lapply}\NormalTok{(}\FunctionTok{names}\NormalTok{(distortion), }\ControlFlowTok{function}\NormalTok{(nm) \{}
\NormalTok{  Y }\OtherTok{\textless{}{-}} \FunctionTok{match\_error}\NormalTok{(distortion[[nm]], }\FloatTok{0.39}\NormalTok{)}
  \FunctionTok{data.frame}\NormalTok{(}\AttributeTok{distortion =}\NormalTok{ nm,}
             \AttributeTok{rmse  =} \FunctionTok{reconstruction\_error}\NormalTok{(Y, U),}
             \AttributeTok{trust =} \FunctionTok{trustworthiness}\NormalTok{(DU, Y, K\_DEFAULT),}
             \AttributeTok{cont  =} \FunctionTok{continuity}\NormalTok{(DU, Y, K\_DEFAULT),}
             \AttributeTok{knn   =} \FunctionTok{knn\_preservation}\NormalTok{(DU, Y, K\_DEFAULT),}
             \AttributeTok{qnx   =} \FunctionTok{qnx}\NormalTok{(Y, DU, K\_QNX))}
\NormalTok{\}))}

\NormalTok{knitr}\SpecialCharTok{::}\FunctionTok{kable}\NormalTok{(}
\NormalTok{  battery, }\AttributeTok{digits =} \DecValTok{3}\NormalTok{, }\AttributeTok{row.names =} \ConstantTok{FALSE}\NormalTok{,}
  \AttributeTok{col.names =} \FunctionTok{c}\NormalTok{(}\StringTok{"distortion"}\NormalTok{, }\StringTok{"Procrustes RMSE"}\NormalTok{, }\StringTok{"trustworthiness"}\NormalTok{,}
                \StringTok{"continuity"}\NormalTok{, }\StringTok{"kNN preservation"}\NormalTok{, }\StringTok{"Q\_NX"}\NormalTok{),}
  \AttributeTok{caption =} \StringTok{"Three distortions of one chart, each tuned by bisection to the same normalised Procrustes error. The rank metrics do not agree about them, and the disagreement is not small."}\NormalTok{)}
\end{Highlighting}
\end{Shaded}

\begin{table}

\caption{\label{tab:eval-battery}Three distortions of one chart, each tuned by bisection to the same normalised Procrustes error. The rank metrics do not agree about them, and the disagreement is not small.}
\centering
\begin{tabular}[t]{l|r|r|r|r|r}
\hline
distortion & Procrustes RMSE & trustworthiness & continuity & kNN preservation & Q\_NX\\
\hline
conformal bend & 0.39 & 1.000 & 1.000 & 0.927 & 0.943\\
\hline
shear & 0.39 & 0.998 & 0.998 & 0.589 & 0.739\\
\hline
jitter & 0.39 & 0.889 & 0.887 & 0.059 & 0.178\\
\hline
\end{tabular}
\end{table}

And a fourth, which is really a warning about the first three. Optimism is a
band, not a number. Table \ref{tab:eval-battery} applies three
distortions to the same chart and tunes each to an identical true error of
0.39. Trustworthiness rates them
1, 0.998, 0.889. kNN preservation rates the
same three configurations 0.927, 0.589, 0.0593 ---
a spread of 0.87 at fixed true error. There
is no correction factor to apply, because there is no factor: what the rank
metrics see is the \emph{kind} of distortion, and the Procrustes error does not
determine it.

\section{Where it fails}\label{eval-limits}

\textbf{One family, one wrong candidate.} Every number above is a Miura at one
tessellation size against one competitor. PCA is the right competitor for an
ambient-referenced metric and a poor one for anything else, and an evaluator
that survives this test has survived one test. The audit's inversion is a
property of the pair, and a different wrong embedding --- a short-circuited
Isomap, say --- would produce a different table.

\textbf{Two sample sizes are not a sample-size law.} The \(k/n\) collapse in Figure
\ref{fig:eval-fig-threshold} rests on two values of \(n\) that differ by a factor
of two, and on a threshold statistic estimated from 20 seeds
per cell on a grid of 20 fold angles. It is consistent with
\(\sqrt{k/n}\) and it does not exclude much else. A third sample size is the
cheapest experiment this chapter suggests and it has not been run.

\textbf{The pilot's own headline pooled over \(n\), and should not have.} The narrow
run that decided this chapter existed reported one table in \(k\) and \(\theta\),
averaged over both of its 2 sample sizes.
On the full sweep kNN preservation inverts on
30.4
per cent of cells at the smaller sample size and
5.44 per
cent at the larger; the pooled figure is
17.9 per cent, which
describes neither. That is recorded here rather than quietly fixed, because it
is the same error this chapter warns about in other people's tables.

\textbf{Nothing here is a claim about methods.} The known-answer test scores
evaluators. Which method to use, and whether this benchmark can tell them apart
at all, is Chapter \ref{results}'s question and is settled by a pre-registered
rule rather than by these curves.

\textbf{What this chapter does not settle.} Whether an evaluator that mishandles a
crease also mishandles a real dataset is not established and cannot be, because
the property being exploited --- an exact chart --- is the property real data lacks.
The defensible reading is the disqualification argument: a metric that ranks the
right answer below a wrong one where the right answer is \emph{known} has not earned
trust where it is not. That is a necessary condition, it asserts nothing about
biology, and it is all this chapter claims.

\section{Reproduce this}\label{eval-reproduce}

The artefact is \texttt{data/processed/evaluator-audit.rds}, produced by
\texttt{scripts/run-evaluator-audit.R}, which also writes the pilot under \texttt{-\/-pilot}.
It carries its own provenance and this chapter refuses to render against a
different run.

\begin{Shaded}
\begin{Highlighting}[]
\ExtensionTok{Rscript}\NormalTok{ scripts/run{-}evaluator{-}audit.R }\AttributeTok{{-}{-}pilot}   \CommentTok{\# the narrow decision run}
\ExtensionTok{Rscript}\NormalTok{ scripts/run{-}evaluator{-}audit.R           }\CommentTok{\# the full sweep}
\end{Highlighting}
\end{Shaded}

The metric definitions are in \texttt{R/metrics.R} and every identity used above is a
test in \texttt{tests/testthat/test-metrics.R}: that reflection alone is free, which is
what forbids a determinant-one fit; that \(A(k)\) agrees with a brute-forced worst
case derived without reference to the paper; that the normalised Procrustes
statistic agrees with \texttt{vegan}'s symmetric form; that
\texttt{knn\_preservation(measure\ =\ "overlap")} is exactly \(Q_{NX}(k)\); that on a flat
sheet the chart \emph{is} the ambient geometry and on a folded one it is not; and
that \texttt{metric\_floor()} is attained by the truncated chart and nothing gets below
it. Run them with \texttt{Rscript\ tests/testthat.R}.

\section{Further reading}\label{eval-reading}

Lee and Verleysen's textbook \citep{lee2007nonlinear} and their co-ranking paper
\citep{lee2009quality} are the reference for everything in the rank-metric section,
and Kaski et al. \citep{kaski2003trustworthiness} is the primary for trustworthiness
and continuity --- with Venna and Kaski \citep{venna2001neighborhood} for priority and
the information-retrieval reading \citep{venna2010information} for why the pair is a
pair. Espadoto et al. \citep{espadoto2021survey} survey what is actually reported in
practice, which is the target of the agreement result above.

On the adversarial side, Machado et al. \citep{machado2025necessary} is the direct
complement to this chapter and should be read alongside it. Schoeneman et al.
\citep{schoeneman2017error} make the same move in a different direction: they build a
benchmark whose answer key is exact because the parameterisation is unit speed,
and Chapter \ref{comparison} measures this book's construction against theirs.

\chapter{Benchmark results}\label{results}

\emph{R. Heller · Date: 2026-07-29 · Last changed: 2026-09-07}

\emph{Which method survives --- and can the benchmark tell?}

\section{What this chapter answers}\label{results-question}

Part II fitted each method to a folded sheet one configuration at a time. This
chapter fits all of them to the whole sweep at once --- every method against every
fold angle under every noise model, twenty seeds deep --- and reads the table.

It asks two questions, in a fixed order, and the order is the argument. The
first is whether the benchmark can tell the methods apart at all: whether the
spread \emph{across} methods exceeds the spread across repeated runs \emph{within} one
method. Only if it does is the second question --- which method wins --- about
methods rather than about seeds. A benchmark that answers the second without
checking the first publishes a ranking of its own noise.

The short answer to the first is yes, everywhere on this grid, by a margin that
narrows sharply where the sheet is most folded. The answer to the second is that
it depends on \(\theta\) and on the noise model, so that ``method X is best'' is not
a well-formed sentence here unless both are stated. The rest is those two
answers with their conditions attached, and a rule --- fixed in code before the
grid existed --- for turning them into a decision.

\section{Setup}\label{results-setup}

\begin{Shaded}
\begin{Highlighting}[]
\NormalTok{grid }\OtherTok{\textless{}{-}} \FunctionTok{read\_run}\NormalTok{(}\StringTok{"data/processed/benchmark{-}grid.rds"}\NormalTok{, }\StringTok{"2483a468d82d4d64"}\NormalTok{)}
\NormalTok{prov }\OtherTok{\textless{}{-}} \FunctionTok{attr}\NormalTok{(grid, }\StringTok{"provenance"}\NormalTok{)}
\NormalTok{ok   }\OtherTok{\textless{}{-}}\NormalTok{ grid[grid}\SpecialCharTok{$}\NormalTok{ran, ]}
\end{Highlighting}
\end{Shaded}

The grid is committed, not computed at render time. \texttt{read\_run()} refuses any
file whose contents differ from the run these numbers were read off, so a
regenerated artefact is a build failure rather than a silently changed sentence.
Producing it took 4.9 hours of wall clock across
4 forked workers --- 19.6
core-hours for 9,199 fits, or about
59
core-seconds per sampled cell, metrics included.

\texttt{grid} keeps every row the producer wrote, including the ones with no metrics;
\texttt{ok} is the subset that fitted. Both are used below, and which one a number came
from is stated, because the difference between them is itself a result.

\section{What the grid holds}\label{results-background}

\begin{Shaded}
\begin{Highlighting}[]
\NormalTok{step }\OtherTok{\textless{}{-}} \FunctionTok{unique}\NormalTok{(}\FunctionTok{round}\NormalTok{(}\FunctionTok{diff}\NormalTok{(}\FunctionTok{sort}\NormalTok{(}\FunctionTok{unique}\NormalTok{(grid}\SpecialCharTok{$}\NormalTok{theta))), }\DecValTok{10}\NormalTok{))}
\NormalTok{design }\OtherTok{\textless{}{-}} \FunctionTok{data.frame}\NormalTok{(}
  \AttributeTok{Factor =} \FunctionTok{c}\NormalTok{(}\StringTok{"pattern"}\NormalTok{, }\StringTok{"$}\SpecialCharTok{\textbackslash{}\textbackslash{}}\StringTok{theta$"}\NormalTok{, }\StringTok{"noise"}\NormalTok{, }\StringTok{"method"}\NormalTok{, }\StringTok{"seed"}\NormalTok{, }\StringTok{"cells"}\NormalTok{, }\StringTok{"rows"}\NormalTok{),}
  \AttributeTok{Levels =} \FunctionTok{c}\NormalTok{(}\FunctionTok{length}\NormalTok{(}\FunctionTok{unique}\NormalTok{(grid}\SpecialCharTok{$}\NormalTok{pattern)), }\FunctionTok{length}\NormalTok{(}\FunctionTok{unique}\NormalTok{(grid}\SpecialCharTok{$}\NormalTok{theta)),}
             \FunctionTok{length}\NormalTok{(}\FunctionTok{unique}\NormalTok{(grid}\SpecialCharTok{$}\NormalTok{noise)), }\FunctionTok{length}\NormalTok{(}\FunctionTok{unique}\NormalTok{(grid}\SpecialCharTok{$}\NormalTok{method)),}
             \FunctionTok{length}\NormalTok{(}\FunctionTok{unique}\NormalTok{(grid}\SpecialCharTok{$}\NormalTok{seed)),}
             \FunctionTok{nrow}\NormalTok{(}\FunctionTok{unique}\NormalTok{(grid[, }\FunctionTok{c}\NormalTok{(}\StringTok{"pattern"}\NormalTok{, }\StringTok{"theta"}\NormalTok{, }\StringTok{"noise"}\NormalTok{, }\StringTok{"seed"}\NormalTok{)])),}
             \FunctionTok{nrow}\NormalTok{(grid)),}
  \AttributeTok{Value  =} \FunctionTok{c}\NormalTok{(}\FunctionTok{paste}\NormalTok{(}\FunctionTok{unique}\NormalTok{(grid}\SpecialCharTok{$}\NormalTok{pattern), }\AttributeTok{collapse =} \StringTok{", "}\NormalTok{),}
             \FunctionTok{sprintf}\NormalTok{(}\StringTok{"\%s to \%s in steps of \%s"}\NormalTok{, }\FunctionTok{min}\NormalTok{(grid}\SpecialCharTok{$}\NormalTok{theta), }\FunctionTok{max}\NormalTok{(grid}\SpecialCharTok{$}\NormalTok{theta), step),}
             \FunctionTok{paste}\NormalTok{(}\FunctionTok{sort}\NormalTok{(}\FunctionTok{unique}\NormalTok{(grid}\SpecialCharTok{$}\NormalTok{noise)), }\AttributeTok{collapse =} \StringTok{", "}\NormalTok{),}
             \FunctionTok{paste}\NormalTok{(}\FunctionTok{sort}\NormalTok{(}\FunctionTok{unique}\NormalTok{(grid}\SpecialCharTok{$}\NormalTok{method)), }\AttributeTok{collapse =} \StringTok{", "}\NormalTok{),}
             \FunctionTok{sprintf}\NormalTok{(}\StringTok{"\%d to \%d"}\NormalTok{, }\FunctionTok{min}\NormalTok{(grid}\SpecialCharTok{$}\NormalTok{seed), }\FunctionTok{max}\NormalTok{(grid}\SpecialCharTok{$}\NormalTok{seed)),}
             \StringTok{"pattern $}\SpecialCharTok{\textbackslash{}\textbackslash{}}\StringTok{times$ $}\SpecialCharTok{\textbackslash{}\textbackslash{}}\StringTok{theta$ $}\SpecialCharTok{\textbackslash{}\textbackslash{}}\StringTok{times$ noise $}\SpecialCharTok{\textbackslash{}\textbackslash{}}\StringTok{times$ seed"}\NormalTok{,}
             \StringTok{"one per method per cell"}\NormalTok{))}
\NormalTok{knitr}\SpecialCharTok{::}\FunctionTok{kable}\NormalTok{(design, }\AttributeTok{align =} \StringTok{"lrl"}\NormalTok{, }\AttributeTok{caption =} \StringTok{"The benchmark grid, read off the artefact rather than asserted."}\NormalTok{)}
\end{Highlighting}
\end{Shaded}

\begin{table}

\caption{\label{tab:results-design}The benchmark grid, read off the artefact rather than asserted.}
\centering
\begin{tabular}[t]{l|r|l}
\hline
Factor & Levels & Value\\
\hline
pattern & 1 & miura\\
\hline
\$\textbackslash{}theta\$ & 20 & 0 to 0.95 in steps of 0.05\\
\hline
noise & 3 & ambient, none, outlier\\
\hline
method & 9 & autoencoder, diffusion, isomap, laplacian, lle, mds, pca, tsne, umap\\
\hline
seed & 20 & 1001 to 1020\\
\hline
cells & 1200 & pattern \$\textbackslash{}times\$ \$\textbackslash{}theta\$ \$\textbackslash{}times\$ noise \$\textbackslash{}times\$ seed\\
\hline
rows & 10800 & one per method per cell\\
\hline
\end{tabular}
\end{table}

There is no pattern axis. One family folds (Chapter \ref{benchmarks}), and a
factor with a single level is not a factor. The three noise conditions are a
clean sheet, isotropic ambient jitter, and a sparse outlier contamination; the
\(\theta\) sweep stops short of the flat-folded state, where the sheet has no
extent in one direction and every method is embedding a degenerate object.

\begin{Shaded}
\begin{Highlighting}[]
\NormalTok{n\_unavail  }\OtherTok{\textless{}{-}} \FunctionTok{sum}\NormalTok{(grid}\SpecialCharTok{$}\NormalTok{status }\SpecialCharTok{==} \StringTok{"unavailable"}\NormalTok{)}
\NormalTok{n\_declined }\OtherTok{\textless{}{-}} \FunctionTok{sum}\NormalTok{(grid}\SpecialCharTok{$}\NormalTok{status }\SpecialCharTok{==} \StringTok{"declined"}\NormalTok{)}
\NormalTok{run\_rate   }\OtherTok{\textless{}{-}} \FunctionTok{tapply}\NormalTok{(grid}\SpecialCharTok{$}\NormalTok{ran, grid}\SpecialCharTok{$}\NormalTok{method, mean)}
\NormalTok{fate }\OtherTok{\textless{}{-}} \FunctionTok{factor}\NormalTok{(}\FunctionTok{ifelse}\NormalTok{(}\SpecialCharTok{!}\NormalTok{grid}\SpecialCharTok{$}\NormalTok{ran,}
                      \FunctionTok{ifelse}\NormalTok{(grid}\SpecialCharTok{$}\NormalTok{status }\SpecialCharTok{==} \StringTok{"unavailable"}\NormalTok{, }\StringTok{"declared unavailable"}\NormalTok{,}
                             \StringTok{"declined the cell"}\NormalTok{),}
                      \FunctionTok{ifelse}\NormalTok{(}\FunctionTok{is.na}\NormalTok{(grid}\SpecialCharTok{$}\NormalTok{tuning), }\StringTok{"ran as asked"}\NormalTok{, }\StringTok{"ran after retuning"}\NormalTok{)),}
               \AttributeTok{levels =} \FunctionTok{c}\NormalTok{(}\StringTok{"ran as asked"}\NormalTok{, }\StringTok{"ran after retuning"}\NormalTok{,}
                          \StringTok{"declined the cell"}\NormalTok{, }\StringTok{"declared unavailable"}\NormalTok{))}
\NormalTok{max\_floor }\OtherTok{\textless{}{-}} \FunctionTok{max}\NormalTok{(grid}\SpecialCharTok{$}\NormalTok{floor)}
\end{Highlighting}
\end{Shaded}

Of 10,800 rows, 9,199
carry metrics. The remaining 1,601
are results, not gaps: 1,200 autoencoder rows
declared unavailable by the registry until Chapter \ref{autoencoders}, and
401 rows where a method was offered a cell and returned nothing.
\texttt{report\_failures()} prints the same accounting from the producer.

The \texttt{floor} column is 0 in every row. The irreducible loss of a
\(p\)-dimensional chart embedded in \(d\) dimensions is the tail of the chart's
spectrum, and here the chart is two-dimensional and \(d\) is 2, so
there is no tail. Every error reported below is therefore loss the method is
responsible for, and ``scored against what was achievable'' and ``scored in
absolute terms'' coincide on this grid. They do not coincide on the product
construction of Chapter \ref{benchmarks}, which is the artefact built to make
the floor positive.

\begin{figure}

{\centering \includegraphics[width=0.9\linewidth,alt={Three panels, one per noise condition, each a horizontal bar per method showing 400 fits. Under no noise and under ambient noise every method except the autoencoder ran as asked, and the autoencoder bar is entirely 'declared unavailable'. Under outlier noise the Laplacian eigenmaps bar is entirely 'declined the cell', the diffusion map bar is almost entirely 'ran after retuning' with a single declined fit, and the rest ran as asked.}]{10-benchmark-results_files/figure-latex/results-fig-status-1} 

}

\caption{What happened to every fit, by method and noise condition. Counts are printed on each segment that is not a clean run, so nothing here is carried by colour alone.}\label{fig:results-fig-status}
\end{figure}

Figure \ref{fig:results-fig-status} shows where the losses sit, and they are
structured rather than scattered. Laplacian eigenmaps declined every outlier
cell; the diffusion map ran under outliers only after widening its own kernel
bandwidth, which the \texttt{tuning} column records.

\section{Where the benchmark has resolving power}\label{results-core}

A ranking is worth reading only if the differences it orders are larger than the
noise it was measured through. \texttt{R/selection.R} names that quantity and fixes it
in code:

\[R \;=\; \frac{\text{spread across methods}}{\text{spread within a method, across seeds}},\]

with both spreads taken of the same quantity --- excess over the irreducible floor
--- so that \(R\) is dimensionless. The numerator is the standard deviation of the
per-method means in a regime; the denominator is the median, over cells, of the
seed-wise standard deviation inside one cell. Above one, the methods differ by
more than a rerun of the same method does. At or below one, any ordering read
off the grid is seed noise with a rank imposed on it.

\begin{Shaded}
\begin{Highlighting}[]
\NormalTok{rp }\OtherTok{\textless{}{-}} \FunctionTok{do.call}\NormalTok{(rbind, }\FunctionTok{lapply}\NormalTok{(}\FunctionTok{sort}\NormalTok{(}\FunctionTok{unique}\NormalTok{(grid}\SpecialCharTok{$}\NormalTok{noise)), }\ControlFlowTok{function}\NormalTok{(nz)}
  \FunctionTok{do.call}\NormalTok{(rbind, }\FunctionTok{lapply}\NormalTok{(THETA\_GRID, }\ControlFlowTok{function}\NormalTok{(t) \{}
\NormalTok{    s }\OtherTok{\textless{}{-}} \FunctionTok{select\_method}\NormalTok{(grid, }\AttributeTok{theta =}\NormalTok{ t, }\AttributeTok{noise =}\NormalTok{ nz)}\SpecialCharTok{$}\NormalTok{resolving}
    \FunctionTok{data.frame}\NormalTok{(}\AttributeTok{noise =}\NormalTok{ nz, }\AttributeTok{theta =}\NormalTok{ t, }\AttributeTok{R =}\NormalTok{ s}\SpecialCharTok{$}\NormalTok{R, }\AttributeTok{between =}\NormalTok{ s}\SpecialCharTok{$}\NormalTok{between, }\AttributeTok{within =}\NormalTok{ s}\SpecialCharTok{$}\NormalTok{within)}
\NormalTok{  \}))))}
\NormalTok{worst     }\OtherTok{\textless{}{-}}\NormalTok{ rp[}\FunctionTok{which.min}\NormalTok{(rp}\SpecialCharTok{$}\NormalTok{R), ]}
\NormalTok{pooled\_sd }\OtherTok{\textless{}{-}} \FunctionTok{tapply}\NormalTok{(ok}\SpecialCharTok{$}\NormalTok{rmse, }\FunctionTok{interaction}\NormalTok{(ok}\SpecialCharTok{$}\NormalTok{method, ok}\SpecialCharTok{$}\NormalTok{theta, ok}\SpecialCharTok{$}\NormalTok{noise, }\AttributeTok{drop =} \ConstantTok{TRUE}\NormalTok{),}
\NormalTok{                    stats}\SpecialCharTok{::}\NormalTok{sd)}
\NormalTok{ends }\OtherTok{\textless{}{-}} \ControlFlowTok{function}\NormalTok{(nz) }\FunctionTok{round}\NormalTok{(rp}\SpecialCharTok{$}\NormalTok{R[rp}\SpecialCharTok{$}\NormalTok{noise }\SpecialCharTok{==}\NormalTok{ nz][}\FunctionTok{c}\NormalTok{(}\DecValTok{1}\DataTypeTok{L}\NormalTok{, }\FunctionTok{length}\NormalTok{(THETA\_GRID))], }\DecValTok{1}\NormalTok{)}
\end{Highlighting}
\end{Shaded}

Three details of that definition earn their keep. The denominator is \emph{within a
cell}, not within a method across the regime, because a method's variation from
one fold angle to the next is difficulty and belongs in the numerator. It is a
\textbf{median} rather than a mean, because one wildly unstable method would
otherwise set the noise level for all eight --- the pooled within-cell spread on
this grid has a median of 0.019 and a maximum of
0.31, and the gap between those two is t-SNE. And it
is computed on excess over the floor rather than on raw error, which on this
grid changes nothing because the floor is 0, and on Chapter
\ref{benchmarks}'s product suite changes a great deal.

\begin{figure}

{\centering \includegraphics[width=0.9\linewidth,alt={A line chart with the folding parameter on the horizontal axis and resolving power on a logarithmic vertical axis. Three lines, distinguished by colour, line style and point shape. The outlier line sits highest, near fifteen, falling gently to about nine at the right. The clean and ambient-noise lines run between three and seven across most of the range and fall towards the right-hand end, the ambient line reaching about one and a third at the largest fold angle. A dotted horizontal line marks one; no line touches it.}]{10-benchmark-results_files/figure-latex/results-fig-resolving-1} 

}

\caption{Resolving power against the folding parameter, one line per noise condition, on a logarithmic scale. The dotted line is the pre-registered floor below which the selection rule declines to rank anything.}\label{fig:results-fig-resolving}
\end{figure}

\(R\) is above the floor of 1 in all
60 regimes of Figure \ref{fig:results-fig-resolving}. The chapter was
specified to report \(R < 1\) wherever it held, and on this grid it holds nowhere;
that is the finding, and it is a more comfortable one than the specification
expected. The weakest cell is ambient noise at \(\theta = 0.95\), where
\(R = 1.33\) --- the between-method spread of
0.031 against a within-method spread of
0.023. Two ticks of \(\theta\) earlier the same condition
resolves twice as well. The outlier condition resolves best of the three, by an
order of magnitude, and that is not a compliment to the benchmark: it separates
the methods there because one of them fails completely, which is a wide gap
rather than an informative one.

The trend across the sweep is the part worth carrying forward. Every condition
loses resolution as the sheet closes --- from 6.5 to
2.3 on the clean sheet, 4.3 to
1.3 under ambient noise, 15.6 to
9.3 under outliers --- so the benchmark is least able to
discriminate exactly where its geometry is most distinctive. That is not a
paradox: at large \(\theta\) every method is failing, and methods that are all
failing look alike. It does mean that a regime chosen for being difficult is a
regime chosen for weak resolving power, and Chapter \ref{application} has to
pick one knowing that.

The number to resist is the seed count. \(R\) is a ratio of two spreads and
neither is a standard error, so running forty seeds instead of twenty estimates
both more precisely and moves the ratio nowhere; \texttt{tests/testthat/test-selection.R}
demonstrates exactly that. Raising the seed budget until something separates
would be the same abuse as raising it until a \(p\)-value falls, with the
advantage of looking like diligence. The budget was fixed in advance and is
reported in the design table above.

Two things this section does \textbf{not} establish. It does not show that a crease
pattern discriminates better than a classical benchmark: Chapter
\ref{comparison} measures that directly and finds the opposite, which is why
this book's spine is the irreducible-loss bound and not the claim that folded
sheets are a harder test. And \(R\) is a property of a regime, not of the
generator --- it is computed here at fixed \(n\), fixed \(k\) and one pattern size,
and each of those moves it.

\section{Results}\label{results-results}

\begin{figure}

{\centering \includegraphics[width=0.9\linewidth,alt={Three panels labelled linear, geodesic and neighbour, each plotting reconstruction error against the folding parameter from zero to nearly one. In the linear panel a single visible curve rises smoothly from zero to about a fifth, with a very narrow band. In the geodesic panel Isomap runs almost flat near zero until roughly two-thirds of the way along and then rises, while the diffusion map sits near a fifth throughout. In the neighbour panel LLE rises from below a tenth to above four tenths, Laplacian eigenmaps and UMAP run flat near a quarter, and t-SNE sits highest with a very wide seed band.}]{10-benchmark-results_files/figure-latex/results-fig-headline-1} 

}

\caption{Normalised Procrustes reconstruction error against the folding parameter on the clean sheet, grouped by what each method consumes. The line is the mean over twenty seeds; the band is the full seed range. PCA and classical MDS coincide exactly and plot on top of each other.}\label{fig:results-fig-headline}
\end{figure}

\begin{Shaded}
\begin{Highlighting}[]
\NormalTok{mm }\OtherTok{\textless{}{-}} \FunctionTok{aggregate}\NormalTok{(rmse }\SpecialCharTok{\textasciitilde{}}\NormalTok{ method }\SpecialCharTok{+}\NormalTok{ theta }\SpecialCharTok{+}\NormalTok{ noise, ok, mean)}
\NormalTok{mm}\SpecialCharTok{$}\NormalTok{rank  }\OtherTok{\textless{}{-}} \FunctionTok{ave}\NormalTok{(mm}\SpecialCharTok{$}\NormalTok{rmse, mm}\SpecialCharTok{$}\NormalTok{theta, mm}\SpecialCharTok{$}\NormalTok{noise, }\AttributeTok{FUN =}\NormalTok{ rank)}
\NormalTok{mm}\SpecialCharTok{$}\NormalTok{label }\OtherTok{\textless{}{-}} \FunctionTok{vapply}\NormalTok{(mm}\SpecialCharTok{$}\NormalTok{method, }\ControlFlowTok{function}\NormalTok{(m) METHOD\_REGISTRY[[m]]}\SpecialCharTok{$}\NormalTok{label, }\FunctionTok{character}\NormalTok{(}\DecValTok{1}\NormalTok{))}
\NormalTok{lead }\OtherTok{\textless{}{-}} \FunctionTok{do.call}\NormalTok{(rbind, }\FunctionTok{lapply}\NormalTok{(}\FunctionTok{split}\NormalTok{(mm, }\FunctionTok{list}\NormalTok{(mm}\SpecialCharTok{$}\NormalTok{theta, mm}\SpecialCharTok{$}\NormalTok{noise)), }\ControlFlowTok{function}\NormalTok{(s)}
  \FunctionTok{data.frame}\NormalTok{(}\AttributeTok{theta =}\NormalTok{ s}\SpecialCharTok{$}\NormalTok{theta[}\DecValTok{1}\NormalTok{], }\AttributeTok{noise =}\NormalTok{ s}\SpecialCharTok{$}\NormalTok{noise[}\DecValTok{1}\NormalTok{], }\AttributeTok{best =}\NormalTok{ s}\SpecialCharTok{$}\NormalTok{method[}\FunctionTok{which.min}\NormalTok{(s}\SpecialCharTok{$}\NormalTok{rmse)])))}
\NormalTok{cross }\OtherTok{\textless{}{-}} \FunctionTok{min}\NormalTok{(lead}\SpecialCharTok{$}\NormalTok{theta[lead}\SpecialCharTok{$}\NormalTok{noise }\SpecialCharTok{==} \StringTok{"none"} \SpecialCharTok{\&}\NormalTok{ lead}\SpecialCharTok{$}\NormalTok{best }\SpecialCharTok{==} \StringTok{"isomap"}\NormalTok{])}
\NormalTok{at }\OtherTok{\textless{}{-}} \ControlFlowTok{function}\NormalTok{(m, i, }\AttributeTok{nz =} \StringTok{"none"}\NormalTok{) \{}
\NormalTok{  mm}\SpecialCharTok{$}\NormalTok{rmse[mm}\SpecialCharTok{$}\NormalTok{method }\SpecialCharTok{==}\NormalTok{ m }\SpecialCharTok{\&}\NormalTok{ mm}\SpecialCharTok{$}\NormalTok{noise }\SpecialCharTok{==}\NormalTok{ nz }\SpecialCharTok{\&}\NormalTok{ mm}\SpecialCharTok{$}\NormalTok{theta }\SpecialCharTok{==}\NormalTok{ THETA\_GRID[i]]}
\NormalTok{\}}
\NormalTok{iso\_rank }\OtherTok{\textless{}{-}}\NormalTok{ mm}\SpecialCharTok{$}\NormalTok{rank[mm}\SpecialCharTok{$}\NormalTok{method }\SpecialCharTok{==} \StringTok{"isomap"} \SpecialCharTok{\&}\NormalTok{ mm}\SpecialCharTok{$}\NormalTok{noise }\SpecialCharTok{==} \StringTok{"outlier"} \SpecialCharTok{\&}
\NormalTok{                    mm}\SpecialCharTok{$}\NormalTok{theta }\SpecialCharTok{\%in\%}\NormalTok{ THETA\_GRID[}\FunctionTok{c}\NormalTok{(}\DecValTok{1}\DataTypeTok{L}\NormalTok{, }\FunctionTok{length}\NormalTok{(THETA\_GRID))]]}
\NormalTok{qq }\OtherTok{\textless{}{-}} \FunctionTok{aggregate}\NormalTok{(qnx }\SpecialCharTok{\textasciitilde{}}\NormalTok{ method }\SpecialCharTok{+}\NormalTok{ theta, ok[ok}\SpecialCharTok{$}\NormalTok{noise }\SpecialCharTok{==} \StringTok{"none"}\NormalTok{, ], mean)}
\NormalTok{qq }\OtherTok{\textless{}{-}}\NormalTok{ qq[}\FunctionTok{order}\NormalTok{(qq}\SpecialCharTok{$}\NormalTok{theta), ]}
\NormalTok{q\_cross }\OtherTok{\textless{}{-}} \FunctionTok{min}\NormalTok{(qq}\SpecialCharTok{$}\NormalTok{theta[qq}\SpecialCharTok{$}\NormalTok{method }\SpecialCharTok{==} \StringTok{"isomap"}\NormalTok{][}
\NormalTok{  qq}\SpecialCharTok{$}\NormalTok{qnx[qq}\SpecialCharTok{$}\NormalTok{method }\SpecialCharTok{==} \StringTok{"isomap"}\NormalTok{] }\SpecialCharTok{\textgreater{}}\NormalTok{ qq}\SpecialCharTok{$}\NormalTok{qnx[qq}\SpecialCharTok{$}\NormalTok{method }\SpecialCharTok{==} \StringTok{"pca"}\NormalTok{]])}
\NormalTok{sd\_by }\OtherTok{\textless{}{-}} \ControlFlowTok{function}\NormalTok{(m) \{}
\NormalTok{  s }\OtherTok{\textless{}{-}}\NormalTok{ ok[ok}\SpecialCharTok{$}\NormalTok{method }\SpecialCharTok{==}\NormalTok{ m, ]}
  \FunctionTok{median}\NormalTok{(}\FunctionTok{tapply}\NormalTok{(s}\SpecialCharTok{$}\NormalTok{rmse, }\FunctionTok{interaction}\NormalTok{(s}\SpecialCharTok{$}\NormalTok{theta, s}\SpecialCharTok{$}\NormalTok{noise, }\AttributeTok{drop =} \ConstantTok{TRUE}\NormalTok{), stats}\SpecialCharTok{::}\NormalTok{sd))}
\NormalTok{\}}
\end{Highlighting}
\end{Shaded}

Figure \ref{fig:results-fig-headline} is the headline on the clean sheet.
At \(\theta = 0\) the sheet \emph{is} the chart, and the linear pair
recovers it to 4.9e-16 while Isomap --- which consumes
geodesics it must first estimate --- pays 0.024 for the
estimation \citep{tenenbaum2000isomap, torgerson1952mds}. The linear error then
climbs monotonically with the fold, because folding is a strict ambient
contraction and a linear map cannot undo one, while Isomap's stays nearly flat:
its assumption is satisfied exactly on a rigid folding, and the geodesics it
estimates do not change as the sheet closes.

The two curves cross. Isomap takes the lead at \(\theta = 0.4\) ---
0.025 against 0.024 one step
earlier, 0.026 against 0.031 at the
crossing --- and holds it to the end of the sweep, finishing at
0.188 against 0.212. Averaging
over \(\theta\) would name the linear pair, because
8 of the
20 fold angles fall on the flat side of the crossing and the
errors there are small. That average is a statement about the grid's design, not
about the methods.

\textbf{The crossing moves when the metric does.} This book reports two headline
numbers per cell, deliberately: Procrustes RMSE asks whether the embedding \emph{is}
the chart, and chart-referenced \(Q_{NX}\) asks only whether it has the chart's
neighbourhoods, which any diffeomorphism of the plane preserves. Read against
\(Q_{NX}\), Isomap does not pass the linear pair until \(\theta = 0.6\) ---
0.920 against
0.917 --- four steps of
the sweep later than under RMSE. The gap between the two crossings is the
distance between ``the linear projection is still the better \emph{map}'' and ``it is
still the better \emph{neighbour list}'', and naming a winner without naming the
metric hides it.

The other five methods break rather than cross. LLE \citep{roweis2000lle} starts
respectably at 0.077 and ends at
0.439, worse than anything else on the clean sheet at the
deepest fold: its local weights are fitted to neighbourhoods that folding has
brought into contact across facets. The diffusion map and Laplacian eigenmaps
are flat and mediocre throughout, averaging
0.204 and
0.236,
which is roughly what a method that recovers the chart up to an unknown
nonlinear warping scores against a metric permitting only similarity. And t-SNE
averages 0.387
and is last at every fold angle --- not a finding about t-SNE, which was never
built to reproduce distances, but the clearest case on this grid of a method
scored against a target it does not consume.

\begin{figure}

{\centering \includegraphics[width=0.9\linewidth,alt={Three panels, one per noise condition, with the folding parameter on the horizontal axis and rank from one at the top to eight at the bottom; the outlier panel has seven ranks because Laplacian eigenmaps produced no embedding there. Under no noise and under ambient noise PCA and classical MDS alternate between ranks one and two at small fold angles and Isomap takes rank one from roughly two-fifths of the way along. Under outlier noise UMAP holds rank one across the whole sweep, the diffusion map is last throughout, and Isomap climbs only from fourth place to third at large fold angles.}]{10-benchmark-results_files/figure-latex/results-fig-order-1} 

}

\caption{Rank by mean reconstruction error against the folding parameter, one panel per noise condition. Rank one is at the top. The leader changes with the fold angle under two conditions and is a different method entirely under the third.}\label{fig:results-fig-order}
\end{figure}

Figure \ref{fig:results-fig-order} puts the three noise conditions side by side,
and the third one breaks the story the first two tell. Under sparse outliers
UMAP \citep{mcinnes2018umap} leads at every fold angle, at
0.353 against
0.422 for the
linear pair, while the diffusion map \citep{coifman2006diffusion} sits at
0.999 ---
which is the score of an embedding that explains none of the chart's variance ---
and Laplacian eigenmaps \citep{belkin2003laplacian} does not produce one at all.
Isomap's advantage at large \(\theta\) survives as a climb from
4th place to 3rd, but UMAP leads at all
20 fold angles and Isomap
never reaches the top under this contamination.

So the ordering is not a property of the methods. It is a property of
(method, \(\theta\), noise), and any of the three sentences ``the linear methods
win'', ``Isomap wins'' and ``UMAP wins'' can be read off this one grid.

\section{Diagnostics}\label{results-diagnostics}

\begin{figure}

{\centering \includegraphics[width=0.9\linewidth,alt={A line chart with the folding parameter on the horizontal axis and the spread across methods on a logarithmic vertical axis. Procrustes RMSE, chart-referenced Q-NX and kNN preservation all run between about two-tenths and four-tenths across the whole range. Trustworthiness and continuity run two orders of magnitude lower, near three thousandths, and trustworthiness rises steeply only beyond about two-thirds of the way along.}]{10-benchmark-results_files/figure-latex/results-fig-metrics-1} 

}

\caption{How far apart the metrics place the eight methods, on the clean sheet: the range of the per-method means at each fold angle, on a logarithmic scale.}\label{fig:results-fig-metrics}
\end{figure}

Four things in the table are worth naming before anyone reads a ranking off it.

\textbf{PCA and classical MDS are the same map here.} Their reconstruction errors
differ by at most 2.8e-15
across the grid and their \(Q_{NX}\) scores are bit-identical, because classical
MDS of a Euclidean distance matrix \emph{is} the principal-component projection
\citep{torgerson1952mds}. Their alternation between ranks one and two in Figure
\ref{fig:results-fig-order} is floating-point tie-breaking, not a result.

\textbf{Two of the five metrics cannot separate anything.} Figure
\ref{fig:results-fig-metrics} is the spread across methods, metric by metric.
Trustworthiness and continuity, scored at \(k = 10\) against \emph{ambient}
distance \citep{kaski2003trustworthiness}, place the eight methods within a median
range of 0.003, while
Procrustes RMSE spreads them over
0.35. Chart-referenced
\(Q_{NX}\) at \(K = 20\) \citep{lee2009quality} sits between the two. This is the
optimism Chapter \ref{evaluation} quantifies; here it is enough to note that
trustworthiness only begins to move at all beyond about two-thirds of the sweep.

\textbf{t-SNE's seed spread swallows most of the field.} Its median within-cell
standard deviation across seeds is 0.26, against
0.003 for the linear pair and
0.004 for Isomap --- a single t-SNE run on this grid is
an anecdote
\citep{vandermaaten2008tsne}, and the band in Figure \ref{fig:results-fig-headline}
is why the book reports twenty.

\textbf{Isomap never needed graph repair.} \texttt{k\_effective} equals the requested
\(k = 10\) in all 1200 fits, so the
neighbourhood graph was connected everywhere in the sweep. That is not the same
as saying no short circuit occurred \citep{balasubramanian2002isomap}: a shortcut edge
between facets connects the graph rather than breaking it, and it is the rise in
Isomap's error past \(\theta = 0.4\) that would carry it.

\section{Where it fails}\label{results-limits}

\begin{Shaded}
\begin{Highlighting}[]
\NormalTok{armB }\OtherTok{\textless{}{-}} \FunctionTok{read\_run}\NormalTok{(}\StringTok{"data/processed/e1{-}armB.rds"}\NormalTok{, }\StringTok{"cbf9ce13782a556d"}\NormalTok{)}
\NormalTok{cc   }\OtherTok{\textless{}{-}} \FunctionTok{aggregate}\NormalTok{(rmse }\SpecialCharTok{\textasciitilde{}}\NormalTok{ creases }\SpecialCharTok{+}\NormalTok{ method, armB, mean)}
\NormalTok{iso\_range }\OtherTok{\textless{}{-}} \FunctionTok{diff}\NormalTok{(}\FunctionTok{range}\NormalTok{(cc}\SpecialCharTok{$}\NormalTok{rmse[cc}\SpecialCharTok{$}\NormalTok{method }\SpecialCharTok{==} \StringTok{"isomap"}\NormalTok{]))}
\end{Highlighting}
\end{Shaded}

The grid varies four things and holds a great many fixed, and the fixed ones
bound every sentence above.

\textbf{One pattern, at one size.} A single Miura at
6 by 6 cells. Crease count is not
an axis here; it was swept separately by Chapter \ref{comparison}'s second arm,
which found Isomap ahead of the linear pair at all
3 crease counts tested at matched branch
separation, with Isomap's own error moving by 0.01 across a
sixteen-fold change in creases. That is a reason for holding it fixed, not a
licence to generalise from one size.

\textbf{One sample size, one neighbourhood, one target dimension.} Every fit used
\(n = 800\), \(k = 10\) and \(d = 2\). Chapter
\ref{evaluation} sweeps \(k\); nothing here does, and a \(k\)-sensitive method
\citep{roweis2000lle} is being reported at one setting of it.

\textbf{No autoencoder.} 1,200 rows of the grid are a
declaration rather than a measurement, and every comparison above is over eight
methods, not nine.

\textbf{Two noise models, both synthetic.} Isotropic ambient jitter and sparse
outliers are the two failure modes easiest to write down, not the two most
common in data. Neither is structured, correlated with position on the sheet, or
heteroscedastic, and the outlier result --- a complete reordering of the field ---
is a warning that the noise model may matter more than the geometry does.

\textbf{A floor of 0.} Reporting against what was achievable is
informative only when something was unachievable, and on a two-dimensional chart
in two dimensions nothing is.

\textbf{Only one pattern family folds at all}, which is Chapter
\ref{benchmarks}'s subject. This grid cannot support any claim about variety
across families, and none is made.

\section{The pre-registered selection rule}\label{results-rule}

Chapter \ref{application} has to pick a method with this grid. The value of that
choice rests entirely on the rule having been fixed before the answer was known,
so it is committed as code --- \texttt{select\_method()} in \texttt{R/selection.R}, a total
function of the grid's columns, registered before the grid it reads had ever
been generated. Prose can be read generously afterwards; a function cannot.

Three thresholds, and nowhere else: a run-rate floor of 0.9, below
which a method is not a candidate however well it scores; a reportable
difference of 0.02 normalised RMSE, below which a lead is a tie;
and a resolving-power floor of 1, below which the rule declines
and the chapter reports a null result.

\begin{Shaded}
\begin{Highlighting}[]
\NormalTok{window }\OtherTok{\textless{}{-}}\NormalTok{ THETA\_GRID[THETA\_GRID }\SpecialCharTok{\textgreater{}=}\NormalTok{ THETA\_GRID[}\DecValTok{13}\NormalTok{]]}
\NormalTok{r\_none }\OtherTok{\textless{}{-}} \FunctionTok{select\_method}\NormalTok{(grid, }\AttributeTok{theta =}\NormalTok{ window, }\AttributeTok{noise =} \StringTok{"none"}\NormalTok{)}
\NormalTok{r\_amb  }\OtherTok{\textless{}{-}} \FunctionTok{select\_method}\NormalTok{(grid, }\AttributeTok{theta =}\NormalTok{ window, }\AttributeTok{noise =} \StringTok{"ambient"}\NormalTok{)}
\NormalTok{r\_all  }\OtherTok{\textless{}{-}} \FunctionTok{select\_method}\NormalTok{(grid, }\AttributeTok{theta =}\NormalTok{ THETA\_GRID, }\AttributeTok{noise =} \FunctionTok{unique}\NormalTok{(grid}\SpecialCharTok{$}\NormalTok{noise))}
\NormalTok{class\_excess }\OtherTok{\textless{}{-}} \FunctionTok{sort}\NormalTok{(}\FunctionTok{tapply}\NormalTok{(ok}\SpecialCharTok{$}\NormalTok{rmse }\SpecialCharTok{{-}}\NormalTok{ ok}\SpecialCharTok{$}\NormalTok{floor, ok}\SpecialCharTok{$}\NormalTok{consumes, mean))}
\end{Highlighting}
\end{Shaded}

The rule does not choose the regime; that is Chapter \ref{application}'s
argument to make from its data, before it fits anything. Applied here to a
deep-fold window as an illustration, it selects \textbf{isomap} on the
clean sheet, by 0.037 at
\(R = 3.44\). Under ambient noise, in the same window,
the leading margin falls to
0.003 --- inside the reportable
threshold --- so the rule declares a tie of 3 methods and
breaks it on failures, then on assuming least, and returns
\textbf{mds} instead. Pooled over the whole grid it excludes
autoencoder and laplacian on the run-rate floor and returns
\textbf{mds} again. Three regimes, three defensible answers, one rule:
which is the point, and the reason the regime must be argued for in the open.

One clause needs a caveat. Ties break toward the method that \emph{assumes least} ---
ambient, then geodesic, then neighbourhood --- which is an argument about
assumptions, not a prediction about scores. The grid ranks the classes the other
way: mean excess over the floor is 0.22
for the geodesic class,
0.30 for ambient and
0.32 for neighbourhood,
in that order and in every noise condition
separately, and Chapter \ref{benchmarks} measures the same ordering on the
product suite. The tie-break stands, because it applies only where the benchmark
\emph{cannot separate} --- but it is not this book's estimate of which class does
better.

The declining branch did not fire anywhere on this grid. It was written to fire
and is tested to fire, and a rule that could not have declined would have
established nothing by not declining.

\section{Reproduce this}\label{results-reproduce}

The grid is regenerated with
\texttt{Rscript\ scripts/run-benchmark-grid.R\ -\/-cores\ 4} from the repository root, and
\texttt{-\/-shard\ i/n} followed by \texttt{-\/-merge} spreads it over machines. That script is the
provenance record for \texttt{data/processed/benchmark-grid.rds} and the only thing
that writes it; it refuses to merge shards built from different versions of
\texttt{R/}, and \texttt{-\/-quick} writes to a different filename so a smoke test can never be
mistaken for the book's evidence. The cell itself is \texttt{grid\_cell()} in
\texttt{R/grid.R}, shared with the Part II sweeps and the classical grid so that three
producers cannot drift apart. \texttt{tests/testthat/test-selection.R} asserts that
each of the three thresholds changes the decision, and that raising the seed
count does not.

\section{Further reading}\label{results-reading}

Espadoto and colleagues \citep{espadoto2021survey} survey the projection-quality
literature this chapter's metrics come from, and Lee and Verleysen
\citep{lee2009quality} give the co-ranking framework behind \(Q_{NX}\); Kaski and
colleagues \citep{kaski2003trustworthiness} are the source of trustworthiness and
continuity, including the scaling constant. For the methods themselves, the
original papers remain the clearest statements of what each one consumes
\citep{tenenbaum2000isomap, roweis2000lle, belkin2003laplacian, coifman2006diffusion, vandermaaten2008tsne, mcinnes2018umap}, which is the
property that predicts the regime changes above better than any of the scores
does. What none of them supplies is a scale against which a score can be called
good; that is Chapter \ref{benchmarks}'s subject, and it is why this chapter's
floor column is a placeholder rather than a measurement.

\chapter{Crease patterns vs.~the Swiss roll}\label{comparison}

\emph{R. Heller · Date: 2026-07-29 · Last changed: 2026-09-07}

\emph{Does the new benchmark actually add anything?}

\section{What this chapter answers}\label{comp-question}

Yes --- and not the thing this book set out to claim.

The claim under test was that a crease pattern is a \emph{harder} benchmark than the
Swiss roll: that folding reaches difficulty the classical surfaces cannot.
Experiment E1 was run before this chapter was drafted, precisely so that the
prose could not outlive a refutation, and it refuted it. At matched difficulty a
folded Miura is \textbf{easier} than a Swiss roll, and it tells the methods apart
about five times less well.

What survives is narrower and better founded. The Swiss roll is one fixed
geometry, and the answer key usually shipped with it is not its isometric chart.
The severed sphere has no exact answer key and, by Gauss's theorem, cannot have
one. A crease pattern is a whole family indexed by \(\theta\) whose chart is exact
at every point of it. The contribution is the answer key and the axis, not the
difficulty --- which is why Chapters \ref{benchmarks} and \ref{evaluation} are
where it pays, and this chapter is where it does not.

\section{Setup}\label{comp-setup}

\begin{Shaded}
\begin{Highlighting}[]
\NormalTok{classic }\OtherTok{\textless{}{-}} \FunctionTok{read\_run}\NormalTok{(}\StringTok{"data/processed/classic{-}grid.rds"}\NormalTok{,   }\StringTok{"d71fb3ba0ad5de68"}\NormalTok{)}
\NormalTok{crease  }\OtherTok{\textless{}{-}} \FunctionTok{read\_run}\NormalTok{(}\StringTok{"data/processed/benchmark{-}grid.rds"}\NormalTok{, }\StringTok{"2483a468d82d4d64"}\NormalTok{)}
\NormalTok{e1      }\OtherTok{\textless{}{-}} \FunctionTok{read\_run}\NormalTok{(}\StringTok{"data/processed/e1{-}controlled.rds"}\NormalTok{,  }\StringTok{"01d82cfe8f45bb44"}\NormalTok{)}
\NormalTok{armB    }\OtherTok{\textless{}{-}} \FunctionTok{read\_run}\NormalTok{(}\StringTok{"data/processed/e1{-}armB.rds"}\NormalTok{,        }\StringTok{"cbf9ce13782a556d"}\NormalTok{)}
\end{Highlighting}
\end{Shaded}

\texttt{classic-grid.rds} is the Swiss roll, the S-curve and the severed sphere scored
through the same \texttt{grid\_cell()} the crease grid uses: same methods, same metrics,
same floor. Nothing on either side is special-cased. The two E1 artefacts are the
controlled arm and the crease-count arm of the experiment that decided this
chapter's claim, and \texttt{scripts/experiment-e1.R} produced both before it was
drafted.

\section{Background}\label{comp-background}

\subsection{The Swiss roll has an exact chart, and it is usually not the one shipped}\label{the-swiss-roll-has-an-exact-chart-and-it-is-usually-not-the-one-shipped}

A Swiss roll is a rectangle of paper wound onto a spiral. It arrived with Isomap
\citep{tenenbaum2000isomap}, appeared the same year in the LLE paper beside the
S-curve \citep{roweis2000lle}, and has been the first figure of the subject ever
since \citep{lee2007nonlinear}.

It is developable, so it \emph{does} have an isometric chart, and the chart is cheap:
arc length along the spiral paired with height. For the Archimedean spiral
\(r = at\) the arc length is closed form,

\[s(t) = \frac{a}{2}\left[\,t\sqrt{1+t^{2}} + \operatorname{asinh} t\,\right],\]

which is \texttt{archimedean\_arclength()} in \texttt{R/baselines.R}. So the Swiss roll is not
short of ground truth. What it is usually short of is the \emph{right} ground truth:
the key shipped with it is the generating angle, and arc length is not
proportional to angle.

\begin{Shaded}
\begin{Highlighting}[]
\NormalTok{roll\_ang }\OtherTok{\textless{}{-}} \FunctionTok{swiss\_roll}\NormalTok{(N\_DEFAULT, }\AttributeTok{turns =} \DecValTok{2}\NormalTok{, }\AttributeTok{seed =}\NormalTok{ BENCH\_SEEDS[}\DecValTok{1}\NormalTok{],}
                       \AttributeTok{isometric =} \ConstantTok{FALSE}\NormalTok{)}
\NormalTok{eul\_ang  }\OtherTok{\textless{}{-}} \FunctionTok{swiss\_roll}\NormalTok{(N\_DEFAULT, }\AttributeTok{turns =} \DecValTok{2}\NormalTok{, }\AttributeTok{seed =}\NormalTok{ BENCH\_SEEDS[}\DecValTok{1}\NormalTok{],}
                       \AttributeTok{isometric =} \ConstantTok{FALSE}\NormalTok{, }\AttributeTok{spiral =} \StringTok{"euler"}\NormalTok{)}
\NormalTok{key\_gap  }\OtherTok{\textless{}{-}} \FunctionTok{c}\NormalTok{(}\AttributeTok{archimedean =} \FunctionTok{reconstruction\_error}\NormalTok{(roll\_ang}\SpecialCharTok{$}\NormalTok{truth, roll\_ang}\SpecialCharTok{$}\NormalTok{iso\_chart),}
              \AttributeTok{euler       =} \FunctionTok{reconstruction\_error}\NormalTok{(eul\_ang}\SpecialCharTok{$}\NormalTok{truth,  eul\_ang}\SpecialCharTok{$}\NormalTok{iso\_chart))}
\NormalTok{key\_qnx  }\OtherTok{\textless{}{-}} \FunctionTok{c}\NormalTok{(}\AttributeTok{archimedean =} \FunctionTok{qnx}\NormalTok{(roll\_ang}\SpecialCharTok{$}\NormalTok{truth, roll\_ang}\SpecialCharTok{$}\NormalTok{iso\_chart, }\AttributeTok{K =}\NormalTok{ K\_QNX),}
              \AttributeTok{euler       =} \FunctionTok{qnx}\NormalTok{(eul\_ang}\SpecialCharTok{$}\NormalTok{truth,  eul\_ang}\SpecialCharTok{$}\NormalTok{iso\_chart,  }\AttributeTok{K =}\NormalTok{ K\_QNX))}

\NormalTok{sphere   }\OtherTok{\textless{}{-}} \FunctionTok{severed\_sphere}\NormalTok{(N\_DEFAULT, }\AttributeTok{cap =} \FloatTok{0.4}\NormalTok{, }\AttributeTok{seed =}\NormalTok{ BENCH\_SEEDS[}\DecValTok{1}\NormalTok{])}
\NormalTok{stretch  }\OtherTok{\textless{}{-}}\NormalTok{ (pi }\SpecialCharTok{{-}}\NormalTok{ sphere}\SpecialCharTok{$}\NormalTok{cap) }\SpecialCharTok{/} \FunctionTok{sin}\NormalTok{(sphere}\SpecialCharTok{$}\NormalTok{cap)    }\CommentTok{\# r / (R sin(r/R)) at the rim}
\NormalTok{n\_sphere }\OtherTok{\textless{}{-}} \FunctionTok{sum}\NormalTok{(classic}\SpecialCharTok{$}\NormalTok{family }\SpecialCharTok{==} \StringTok{"severed\_sphere"} \SpecialCharTok{\&}\NormalTok{ classic}\SpecialCharTok{$}\NormalTok{ran)}
\FunctionTok{cat}\NormalTok{(}\FunctionTok{strwrap}\NormalTok{(}\FunctionTok{tryCatch}\NormalTok{(}\FunctionTok{require\_exact\_truth}\NormalTok{(sphere, }\StringTok{"a reconstruction score"}\NormalTok{),}
                     \AttributeTok{error =}\NormalTok{ conditionMessage), }\DecValTok{72}\NormalTok{), }\AttributeTok{sep =} \StringTok{"}\SpecialCharTok{\textbackslash{}n}\StringTok{"}\NormalTok{)}
\end{Highlighting}
\end{Shaded}

\begin{verbatim}
a reconstruction score needs exact ground truth, and this sample does
not have it: severed_sphere carries a conventional chart (equidistant
(a convention; no isometric chart exists)). Scoring against it measures
agreement with a convention, not recovery of a structure.
\end{verbatim}

Score the customary angle chart against the arc-length chart on one two-turn
roll and the gap is 0.567 in normalised
Procrustes RMSE, with \(Q_{NX}\) at 0.462. That
is not a rounding difference. An evaluation against the angle chart measures
recovery of a reparameterisation, and the factor by which it is wrong is a
property of the roll rather than of the method under test.

The defect is known and was fixed in 2017. Schoeneman et al.~proposed the Euler
Isometric Swiss Roll \citep{schoeneman2017error}, replacing the Archimedean spiral
with the Fresnel integrals; it has been used since \citep{mahapatra2018learning}.
Their stated reason is that curvature proportional to arc length ``ensures that
isometry is preserved''. The checkable reason is simpler: the Fresnel
parameterisation is \textbf{unit speed}, so the parameter \emph{is} arc length by
definition. Label an Euler roll by its turning angle instead and the same error
appears at 0.258 --- smaller, and not zero.

Every Swiss roll in this book is scored against arc length. Figure
\ref{fig:comp-chart} is why.

\begin{figure}

{\centering \includegraphics[width=0.9\linewidth,alt={A line chart with position by turning angle on the horizontal axis and position by arc length on the vertical axis, both scaled to the unit interval. A straight grey diagonal marks equality. A solid dark curve for the Archimedean spiral bows well below the diagonal, farthest from it around the middle of the range. A dashed lighter curve for the Euler spiral bows above the diagonal by a much smaller amount. Both curves meet the diagonal only at the two ends.}]{11-versus-swiss-roll_files/figure-latex/comp-chart-1} 

}

\caption{Position along the spiral measured two ways, on the same sample. The grey diagonal is where the two agree. The Archimedean roll's customary angle chart departs from arc length across the whole sweep; the Euler roll's departs far less, because its parameterisation is unit speed.}\label{fig:comp-chart}
\end{figure}

\subsection{The severed sphere has no answer key at all}\label{the-severed-sphere-has-no-answer-key-at-all}

Say it bluntly. A sphere has constant positive Gaussian curvature, the plane has
none, and curvature is intrinsic --- Gauss's \emph{Theorema Egregium}. No chart of any
patch of a sphere preserves distances, however small the patch and however
clever the chart. \texttt{severed\_sphere()} therefore ships a \textbf{convention}: the
azimuthal-equidistant chart about the retained pole, which gets every radial
geodesic exactly right and stretches every circle around it by
\(r / (R\sin(r/R))\) --- a factor reaching 7.04 at the rim of the
largest patch in the classic grid.

The object knows this, and the refusal printed above is \texttt{require\_exact\_truth()}
reading the flag. That guard fires only where something calls it. \texttt{grid\_cell()}
does not, so the classic grid scored the severed sphere 1280 times,
exactly as it scored the roll. Those numbers exist, they are in the committed
artefact, and this chapter reports none of them as recovery of a structure.

\section{Difficulty has two axes, not one}\label{comp-core}

\subsection{The knob has to move the right thing}\label{the-knob-has-to-move-the-right-thing}

Both families have a difficulty parameter, and both parameters have to be
checked, because the obvious version of each does nothing.

For the Swiss roll it is the turn count, and the obvious parameterisation --- let
the angle run further and the radius run with it --- adds revolutions at \emph{constant}
sheet spacing, so the \(k\)-nearest-neighbour graph stays faithful however many
turns are asked for. \texttt{swiss\_roll()} instead holds the outer radius fixed and
fits the requested revolutions inside it, so the radial gap between consecutive
sheets falls as one over the turn count. That is Balasubramanian and Schwartz's
gap-to-spacing ratio \(g/s\) \citep{balasubramanian2002isomap}, and it is the one axis
both families can be plotted against.

For a crease pattern it is \(\theta\). But folding does two things at once: it
brings distant sheets together \emph{and} it lifts the sheet out of the plane. The
second is what \texttt{non\_planarity()} measures --- the share of ambient variance no
plane captures, the third principal eigenvalue over their sum.

\begin{Shaded}
\begin{Highlighting}[]
\NormalTok{axes }\OtherTok{\textless{}{-}} \FunctionTok{local}\NormalTok{(\{}
\NormalTok{  seeds }\OtherTok{\textless{}{-}}\NormalTok{ BENCH\_SEEDS[}\DecValTok{1}\SpecialCharTok{:}\DecValTok{3}\NormalTok{]}
\NormalTok{  one }\OtherTok{\textless{}{-}} \ControlFlowTok{function}\NormalTok{(m, family, par) }\FunctionTok{data.frame}\NormalTok{(}
    \AttributeTok{family =}\NormalTok{ family, }\AttributeTok{par =}\NormalTok{ par, }\AttributeTok{np =} \FunctionTok{non\_planarity}\NormalTok{(m),}
    \AttributeTok{gs =} \FunctionTok{branch\_gap}\NormalTok{(m)}\SpecialCharTok{$}\NormalTok{ratio, }\AttributeTok{sc =} \FunctionTok{short\_circuit\_index}\NormalTok{(m))}
\NormalTok{  run }\OtherTok{\textless{}{-}} \ControlFlowTok{function}\NormalTok{(vals, family, make) }\FunctionTok{do.call}\NormalTok{(rbind, }\FunctionTok{lapply}\NormalTok{(vals, }\ControlFlowTok{function}\NormalTok{(v)}
    \FunctionTok{do.call}\NormalTok{(rbind, }\FunctionTok{lapply}\NormalTok{(seeds, }\ControlFlowTok{function}\NormalTok{(s) }\FunctionTok{one}\NormalTok{(}\FunctionTok{make}\NormalTok{(v, s), family, v)))))}
  \FunctionTok{rbind}\NormalTok{(}
    \FunctionTok{run}\NormalTok{(}\FunctionTok{sort}\NormalTok{(}\FunctionTok{unique}\NormalTok{(classic}\SpecialCharTok{$}\NormalTok{par[classic}\SpecialCharTok{$}\NormalTok{family }\SpecialCharTok{==} \StringTok{"swiss\_roll"}\NormalTok{])), }\StringTok{"Swiss roll"}\NormalTok{,}
        \ControlFlowTok{function}\NormalTok{(v, s) }\FunctionTok{swiss\_roll}\NormalTok{(N\_DEFAULT, }\AttributeTok{turns =}\NormalTok{ v, }\AttributeTok{seed =}\NormalTok{ s)),}
    \FunctionTok{run}\NormalTok{(}\FunctionTok{sort}\NormalTok{(}\FunctionTok{unique}\NormalTok{(classic}\SpecialCharTok{$}\NormalTok{par[classic}\SpecialCharTok{$}\NormalTok{family }\SpecialCharTok{==} \StringTok{"severed\_sphere"}\NormalTok{])),}
        \StringTok{"severed sphere"}\NormalTok{, }\ControlFlowTok{function}\NormalTok{(v, s) }\FunctionTok{severed\_sphere}\NormalTok{(N\_DEFAULT, }\AttributeTok{cap =}\NormalTok{ v, }\AttributeTok{seed =}\NormalTok{ s)),}
    \FunctionTok{run}\NormalTok{(}\ConstantTok{NA\_real\_}\NormalTok{, }\StringTok{"S{-}curve"}\NormalTok{, }\ControlFlowTok{function}\NormalTok{(v, s) }\FunctionTok{s\_curve}\NormalTok{(N\_DEFAULT, }\AttributeTok{seed =}\NormalTok{ s)),}
    \FunctionTok{run}\NormalTok{(THETA\_GRID, }\StringTok{"Miura"}\NormalTok{, }\ControlFlowTok{function}\NormalTok{(v, s)}
      \FunctionTok{sample\_manifold}\NormalTok{(}\FunctionTok{miura\_ori}\NormalTok{(}\DecValTok{6}\DataTypeTok{L}\NormalTok{, }\DecValTok{6}\DataTypeTok{L}\NormalTok{), }\AttributeTok{theta =}\NormalTok{ v, }\AttributeTok{n =}\NormalTok{ N\_DEFAULT, }\AttributeTok{seed =}\NormalTok{ s)))}
\NormalTok{\})}
\CommentTok{\# The S{-}curve is added by hand because it has no difficulty parameter at all;}
\CommentTok{\# aggregate() drops a group whose key is NA, and inventing a key would be worse.}
\NormalTok{cell }\OtherTok{\textless{}{-}} \FunctionTok{aggregate}\NormalTok{(}\FunctionTok{cbind}\NormalTok{(np, gs, sc) }\SpecialCharTok{\textasciitilde{}}\NormalTok{ family }\SpecialCharTok{+}\NormalTok{ par, axes, mean)}
\NormalTok{sc\_rows }\OtherTok{\textless{}{-}}\NormalTok{ axes[axes}\SpecialCharTok{$}\NormalTok{family }\SpecialCharTok{==} \StringTok{"S{-}curve"}\NormalTok{, ]}
\NormalTok{cell }\OtherTok{\textless{}{-}} \FunctionTok{rbind}\NormalTok{(cell, }\FunctionTok{data.frame}\NormalTok{(}\AttributeTok{family =} \StringTok{"S{-}curve"}\NormalTok{, }\AttributeTok{par =} \ConstantTok{NA\_real\_}\NormalTok{,}
                               \AttributeTok{np =} \FunctionTok{mean}\NormalTok{(sc\_rows}\SpecialCharTok{$}\NormalTok{np), }\AttributeTok{gs =} \FunctionTok{mean}\NormalTok{(sc\_rows}\SpecialCharTok{$}\NormalTok{gs),}
                               \AttributeTok{sc =} \FunctionTok{mean}\NormalTok{(sc\_rows}\SpecialCharTok{$}\NormalTok{sc)))}
\NormalTok{np\_top }\OtherTok{\textless{}{-}} \FunctionTok{max}\NormalTok{(cell}\SpecialCharTok{$}\NormalTok{np[cell}\SpecialCharTok{$}\NormalTok{family }\SpecialCharTok{==} \StringTok{"Miura"}\NormalTok{])}
\NormalTok{np\_bot }\OtherTok{\textless{}{-}} \FunctionTok{min}\NormalTok{(cell}\SpecialCharTok{$}\NormalTok{np[cell}\SpecialCharTok{$}\NormalTok{family }\SpecialCharTok{!=} \StringTok{"Miura"}\NormalTok{])}

\NormalTok{swept }\OtherTok{\textless{}{-}}\NormalTok{ cell[cell}\SpecialCharTok{$}\NormalTok{family }\SpecialCharTok{!=} \StringTok{"S{-}curve"}\NormalTok{, ]}
\NormalTok{ggplot2}\SpecialCharTok{::}\FunctionTok{ggplot}\NormalTok{(cell, ggplot2}\SpecialCharTok{::}\FunctionTok{aes}\NormalTok{(gs, np, }\AttributeTok{colour =}\NormalTok{ family, }\AttributeTok{shape =}\NormalTok{ family)) }\SpecialCharTok{+}
  \CommentTok{\# Finite bounds, not +/{-} Inf: under a log transform an infinite limit is NaN}
  \CommentTok{\# and the band silently does not draw.}
\NormalTok{  ggplot2}\SpecialCharTok{::}\FunctionTok{annotate}\NormalTok{(}\StringTok{"rect"}\NormalTok{, }\AttributeTok{xmin =} \FunctionTok{min}\NormalTok{(cell}\SpecialCharTok{$}\NormalTok{gs) }\SpecialCharTok{/} \DecValTok{2}\NormalTok{, }\AttributeTok{xmax =} \FunctionTok{max}\NormalTok{(cell}\SpecialCharTok{$}\NormalTok{gs) }\SpecialCharTok{*} \DecValTok{2}\NormalTok{,}
                    \AttributeTok{ymin =}\NormalTok{ np\_top, }\AttributeTok{ymax =}\NormalTok{ np\_bot, }\AttributeTok{fill =} \StringTok{"grey70"}\NormalTok{, }\AttributeTok{alpha =} \FloatTok{0.4}\NormalTok{) }\SpecialCharTok{+}
\NormalTok{  ggplot2}\SpecialCharTok{::}\FunctionTok{geom\_line}\NormalTok{(}\AttributeTok{data =}\NormalTok{ swept[}\FunctionTok{order}\NormalTok{(swept}\SpecialCharTok{$}\NormalTok{family, swept}\SpecialCharTok{$}\NormalTok{par), ],}
\NormalTok{                     ggplot2}\SpecialCharTok{::}\FunctionTok{aes}\NormalTok{(}\AttributeTok{group =}\NormalTok{ family), }\AttributeTok{linewidth =} \FloatTok{0.5}\NormalTok{) }\SpecialCharTok{+}
\NormalTok{  ggplot2}\SpecialCharTok{::}\FunctionTok{geom\_point}\NormalTok{(}\AttributeTok{size =} \FloatTok{2.3}\NormalTok{) }\SpecialCharTok{+}
\NormalTok{  ggplot2}\SpecialCharTok{::}\FunctionTok{scale\_x\_continuous}\NormalTok{(}\AttributeTok{transform =} \FunctionTok{c}\NormalTok{(}\StringTok{"log10"}\NormalTok{, }\StringTok{"reverse"}\NormalTok{)) }\SpecialCharTok{+}
  \FunctionTok{scale\_colour\_book}\NormalTok{(}\AttributeTok{begin =} \FloatTok{0.05}\NormalTok{, }\AttributeTok{end =} \FloatTok{0.78}\NormalTok{) }\SpecialCharTok{+}
\NormalTok{  ggplot2}\SpecialCharTok{::}\FunctionTok{labs}\NormalTok{(}\AttributeTok{x =} \StringTok{"gap{-}to{-}spacing ratio (log scale, harder to the right)"}\NormalTok{,}
                \AttributeTok{y =} \StringTok{"ambient non{-}planarity"}\NormalTok{, }\AttributeTok{colour =} \ConstantTok{NULL}\NormalTok{, }\AttributeTok{shape =} \ConstantTok{NULL}\NormalTok{) }\SpecialCharTok{+}
  \FunctionTok{theme\_fold}\NormalTok{()}
\end{Highlighting}
\end{Shaded}

\begin{figure}

{\centering \includegraphics[width=0.9\linewidth,alt={A scatter plot with the gap-to-spacing ratio on a reversed logarithmic horizontal axis, so difficulty increases to the right, and ambient non-planarity from zero to about 0.32 on the vertical axis. A connected run of Miura points sweeps along the bottom of the panel and never rises above about 0.06. A connected run of Swiss roll points sits in a narrow horizontal band near 0.09 to 0.13, with the S-curve as a single point beside them, and severed sphere points well above at 0.22 to 0.30. A grey shaded horizontal band separates the Miura run from everything else and contains no points.}]{11-versus-swiss-roll_files/figure-latex/comp-plane-1} 

}

\caption{The two difficulty axes. Each family sweeps its own parameter; difficulty rises to the right. The shaded band is empty, so over this design there is no setting at which a folded Miura and a classical surface are comparable on both axes at once.}\label{fig:comp-plane}
\end{figure}

The ranges are \textbf{disjoint}. Across the whole \(\theta\) sweep a six-by-six Miura
covers non-planarity from
0.000 to
0.056; the Swiss rolls in the classic grid cover
0.090 to
0.124, the S-curve sits
at 0.116, and the severed
sphere starts at
0.218 and rises. The
corridor between 0.056 and 0.090 is
empty.

Figure \ref{fig:comp-plane} is the argument. A sweep of \(\theta\) alone is not a
controlled comparison, because it moves both axes together and the families
never meet on the second. At the easy end of the \(g/s\) range a Miura is a flat
plane, which PCA recovers exactly; a Swiss roll at the same separation is still
curved, and PCA cannot.

The third column of \texttt{cell} is the reason \(g/s\) carries the first axis rather
than the more obvious candidate. \texttt{short\_circuit\_index()} asks whether the
neighbours a method would actually pick are neighbours on the surface too; it is
scale-free, bounded above by one, and comparable across families, which makes it
the better-behaved statistic of the two. It is also nearly constant over most of
this design. Across the classic grid's turn range it falls only from
1.000 to
0.703, and almost all of
that drop happens at the last two settings; across the whole \(\theta\) sweep the
Miura falls from 1.000 to
0.799; the severed sphere
never leaves
0.999. A difficulty
axis that is flat until the very end of the range cannot order the settings in
between, so the comparison is made on \(g/s\) and the short-circuit index is
reported beside it as the check that \(g/s\) is not measuring the sampler.

\subsection{The fix, and what it costs}\label{the-fix-and-what-it-costs}

\(\theta\) is not the only knob the generator has. Cell count dominates
non-planarity: a small Miura at a wide opening angle and a deep fold reaches
above anything the Swiss roll produces, while a six-by-six at the same angles
stays an order of magnitude below. Sweeping \((n_x, \alpha, \theta)\) decouples the
axes and creates the overlap a \(\theta\) sweep lacks. That is what E1's second arm
did, and it is what \texttt{e1-controlled.rds} holds.

Crease count on its own does nothing. Arm B holds \(g/s\) fixed and varies the
tessellation from nine facets to a hundred and forty-four; PCA's mean error moves
between
0.039
and
0.071,
Isomap's between
0.028
and
0.038,
and neither monotonically. Error is flat in crease count at matched separation.
Whatever makes these patterns hard, it is not that they have creases.

\subsection{One geometry against a family}\label{one-geometry-against-a-family}

This is the structural difference, and it survives everything the measurements
below do to the rest of the chapter. A Swiss roll's turn count changes how
tightly one shape is wound; it is a real knob, and \texttt{swiss\_roll()} makes it one. A
crease pattern gives a \emph{family}: the same exact chart at every \(\theta\), so a
method's score is a curve rather than a point and the shape of the curve is a
property of the method. Nothing about the classical surfaces forbids such a
sweep, and one is drawn here. What they cannot supply is the exact key at the far
end of it.

Nor have they had to break the habit that goes with a fixed geometry. A Swiss
roll's answer key is two-dimensional and its failures are legible, so the
customary report is a picture: the unrolled sheet either comes out as a rectangle
or it does not, and the reader is invited to agree. That is a serviceable way to
show a method working and a poor way to compare two that both nearly work --- which
is exactly the regime the dispute over two-dimensional embeddings lives in
\citep{chari2023specious, lause2024elephant}. A parameterised family forces the
question into a number at every setting, and the number is what the rest of this
chapter argues over.

\section{Results}\label{comp-results}

\subsection{The head-to-head, matched on separation alone}\label{the-head-to-head-matched-on-separation-alone}

Match each Swiss roll to the \(\theta\) whose measured \(g/s\) is nearest, keep the
ones landing inside the crease family's range, and score both sides from the
committed grids at the same eight methods, the same twenty seeds and no noise.

\begin{Shaded}
\begin{Highlighting}[]
\NormalTok{mi }\OtherTok{\textless{}{-}}\NormalTok{ cell[cell}\SpecialCharTok{$}\NormalTok{family }\SpecialCharTok{==} \StringTok{"Miura"}\NormalTok{, ]}
\NormalTok{ro }\OtherTok{\textless{}{-}}\NormalTok{ cell[cell}\SpecialCharTok{$}\NormalTok{family }\SpecialCharTok{==} \StringTok{"Swiss roll"}\NormalTok{, ]}
\NormalTok{ro }\OtherTok{\textless{}{-}}\NormalTok{ ro[ro}\SpecialCharTok{$}\NormalTok{gs }\SpecialCharTok{\textgreater{}=} \FunctionTok{min}\NormalTok{(mi}\SpecialCharTok{$}\NormalTok{gs) }\SpecialCharTok{\&}\NormalTok{ ro}\SpecialCharTok{$}\NormalTok{gs }\SpecialCharTok{\textless{}=} \FunctionTok{max}\NormalTok{(mi}\SpecialCharTok{$}\NormalTok{gs), ]}
\NormalTok{near }\OtherTok{\textless{}{-}} \FunctionTok{vapply}\NormalTok{(ro}\SpecialCharTok{$}\NormalTok{gs, }\ControlFlowTok{function}\NormalTok{(x) }\FunctionTok{which.min}\NormalTok{(}\FunctionTok{abs}\NormalTok{(mi}\SpecialCharTok{$}\NormalTok{gs }\SpecialCharTok{{-}}\NormalTok{ x)), }\FunctionTok{integer}\NormalTok{(}\DecValTok{1}\NormalTok{))}
\NormalTok{ro}\SpecialCharTok{$}\NormalTok{theta    }\OtherTok{\textless{}{-}}\NormalTok{ mi}\SpecialCharTok{$}\NormalTok{par[near]}
\NormalTok{ro}\SpecialCharTok{$}\NormalTok{mismatch }\OtherTok{\textless{}{-}} \FunctionTok{abs}\NormalTok{(mi}\SpecialCharTok{$}\NormalTok{gs[near] }\SpecialCharTok{{-}}\NormalTok{ ro}\SpecialCharTok{$}\NormalTok{gs) }\SpecialCharTok{/}\NormalTok{ ro}\SpecialCharTok{$}\NormalTok{gs}

\NormalTok{A }\OtherTok{\textless{}{-}}\NormalTok{ classic[classic}\SpecialCharTok{$}\NormalTok{family }\SpecialCharTok{==} \StringTok{"swiss\_roll"} \SpecialCharTok{\&}\NormalTok{ classic}\SpecialCharTok{$}\NormalTok{noise }\SpecialCharTok{==} \StringTok{"none"} \SpecialCharTok{\&}
\NormalTok{             classic}\SpecialCharTok{$}\NormalTok{ran }\SpecialCharTok{\&}\NormalTok{ classic}\SpecialCharTok{$}\NormalTok{par }\SpecialCharTok{\%in\%}\NormalTok{ ro}\SpecialCharTok{$}\NormalTok{par, ]}
\NormalTok{B }\OtherTok{\textless{}{-}}\NormalTok{ crease[crease}\SpecialCharTok{$}\NormalTok{noise }\SpecialCharTok{==} \StringTok{"none"} \SpecialCharTok{\&}\NormalTok{ crease}\SpecialCharTok{$}\NormalTok{ran }\SpecialCharTok{\&}\NormalTok{ crease}\SpecialCharTok{$}\NormalTok{theta }\SpecialCharTok{\%in\%}\NormalTok{ ro}\SpecialCharTok{$}\NormalTok{theta, ]}
\FunctionTok{stopifnot}\NormalTok{(}\FunctionTok{all}\NormalTok{(A}\SpecialCharTok{$}\NormalTok{seed }\SpecialCharTok{\%in\%}\NormalTok{ BENCH\_SEEDS), }\FunctionTok{all}\NormalTok{(B}\SpecialCharTok{$}\NormalTok{seed }\SpecialCharTok{\%in\%}\NormalTok{ BENCH\_SEEDS))}
\NormalTok{h }\OtherTok{\textless{}{-}} \FunctionTok{merge}\NormalTok{(}\FunctionTok{setNames}\NormalTok{(}\FunctionTok{aggregate}\NormalTok{(}\FunctionTok{cbind}\NormalTok{(rmse, qnx) }\SpecialCharTok{\textasciitilde{}}\NormalTok{ method, A, mean),}
                    \FunctionTok{c}\NormalTok{(}\StringTok{"method"}\NormalTok{, }\StringTok{"roll\_rmse"}\NormalTok{, }\StringTok{"roll\_qnx"}\NormalTok{)),}
           \FunctionTok{setNames}\NormalTok{(}\FunctionTok{aggregate}\NormalTok{(}\FunctionTok{cbind}\NormalTok{(rmse, qnx) }\SpecialCharTok{\textasciitilde{}}\NormalTok{ method, B, mean),}
                    \FunctionTok{c}\NormalTok{(}\StringTok{"method"}\NormalTok{, }\StringTok{"cr\_rmse"}\NormalTok{, }\StringTok{"cr\_qnx"}\NormalTok{)))}
\NormalTok{h }\OtherTok{\textless{}{-}}\NormalTok{ h[}\FunctionTok{order}\NormalTok{(h}\SpecialCharTok{$}\NormalTok{roll\_rmse), ]}
\NormalTok{rho }\OtherTok{\textless{}{-}} \FunctionTok{c}\NormalTok{(}\AttributeTok{rmse =} \FunctionTok{cor}\NormalTok{(h}\SpecialCharTok{$}\NormalTok{roll\_rmse, h}\SpecialCharTok{$}\NormalTok{cr\_rmse, }\AttributeTok{method =} \StringTok{"spearman"}\NormalTok{),}
         \AttributeTok{qnx  =} \FunctionTok{cor}\NormalTok{(h}\SpecialCharTok{$}\NormalTok{roll\_qnx,  h}\SpecialCharTok{$}\NormalTok{cr\_qnx,  }\AttributeTok{method =} \StringTok{"spearman"}\NormalTok{))}
\NormalTok{gap }\OtherTok{\textless{}{-}} \ControlFlowTok{function}\NormalTok{(d, key) \{}
\NormalTok{  w }\OtherTok{\textless{}{-}} \FunctionTok{reshape}\NormalTok{(d[d}\SpecialCharTok{$}\NormalTok{method }\SpecialCharTok{\%in\%} \FunctionTok{c}\NormalTok{(}\StringTok{"pca"}\NormalTok{, }\StringTok{"isomap"}\NormalTok{), }\FunctionTok{c}\NormalTok{(key, }\StringTok{"seed"}\NormalTok{, }\StringTok{"method"}\NormalTok{, }\StringTok{"rmse"}\NormalTok{)],}
               \AttributeTok{idvar =} \FunctionTok{c}\NormalTok{(key, }\StringTok{"seed"}\NormalTok{), }\AttributeTok{timevar =} \StringTok{"method"}\NormalTok{, }\AttributeTok{direction =} \StringTok{"wide"}\NormalTok{)}
  \FunctionTok{abs}\NormalTok{(w}\SpecialCharTok{$}\NormalTok{rmse.pca }\SpecialCharTok{{-}}\NormalTok{ w}\SpecialCharTok{$}\NormalTok{rmse.isomap)}
\NormalTok{\}}
\NormalTok{naive }\OtherTok{\textless{}{-}} \FunctionTok{mean}\NormalTok{(}\FunctionTok{gap}\NormalTok{(A, }\StringTok{"par"}\NormalTok{)) }\SpecialCharTok{/} \FunctionTok{mean}\NormalTok{(}\FunctionTok{gap}\NormalTok{(B, }\StringTok{"theta"}\NormalTok{))}
\end{Highlighting}
\end{Shaded}

That leaves 3 matched pairs, agreeing to within
6.8\% in \(g/s\).

\begin{figure}

{\centering \includegraphics[width=0.9\linewidth,alt={A dumbbell chart. Eight methods run down the vertical axis, ordered by Swiss-roll error with Isomap at the top and PCA and classical MDS at the bottom. Each method carries two markers joined by a grey segment, one for the Swiss roll and one for the matched Miura. Every Miura marker lies to the left of its Swiss-roll marker, that is at lower error. For Isomap the two nearly coincide near zero; for PCA and MDS the Miura marker sits far to the left while the Swiss-roll marker is at the extreme right.}]{11-versus-swiss-roll_files/figure-latex/comp-ranking-1} 

}

\caption{Mean reconstruction error of eight methods on a Swiss roll and on a Miura matched to it in gap-to-spacing ratio. Methods are ordered by their Swiss-roll score, and the connecting segments make the reordering visible: the two linear methods move from joint last to near the front.}\label{fig:comp-ranking}
\end{figure}

The two rankings barely correlate: Spearman 0.036 on
reconstruction error and -0.060 on \(Q_{NX}\). So the
classical benchmark \emph{does} rank methods differently --- and the order it produces
is the one the field has internalised. Isomap first, the linear methods joint
last: the picture the Swiss roll was introduced to make
\citep{tenenbaum2000isomap}, and the one against which the neighbour embeddings have
been read ever since \citep{vandermaaten2008tsne, mcinnes2018umap}.

The reordering is not a shuffle at the margin. PCA and classical MDS score
0.898 on the roll, joint last of
eight, and 0.093 on the Miura matched
to it, which is second. Isomap barely moves ---
0.049 against
0.065 --- and neither does its
\(Q_{NX}\), 0.860 against
0.864. Every method scores better
on the crease side, which is the first thing to notice: at matched separation the
folded sheet is the easier object for all eight. What changes the order is that
the linear methods gain far more from the change than the geodesic one does.

\subsection{Why that answer is the wrong one}\label{why-that-answer-is-the-wrong-one}

Because the match was on one axis and there are two. PCA does well on a folded
Miura for a reason that has nothing to do with creases: even at the largest
\(\theta\) this grid uses the sheet is still nearly planar, and a nearly planar
object is what PCA is optimal for. Read off this comparison, the Swiss roll
separates PCA from Isomap 30 times better than the crease
pattern does --- a number produced almost entirely by the axis the match failed to
control.

\subsection{The controlled comparison}\label{the-controlled-comparison}

\texttt{e1-controlled.rds} sweeps \((n_x, \alpha, \theta)\) against turns, and has an
overlap region in which both axes lie inside the joint range of the two families.

\begin{Shaded}
\begin{Highlighting}[]
\NormalTok{ctl }\OtherTok{\textless{}{-}} \FunctionTok{local}\NormalTok{(\{}
\NormalTok{  s }\OtherTok{\textless{}{-}}\NormalTok{ e1[}\FunctionTok{is.finite}\NormalTok{(e1}\SpecialCharTok{$}\NormalTok{rmse), ]; s}\SpecialCharTok{$}\NormalTok{lg }\OtherTok{\textless{}{-}} \FunctionTok{log}\NormalTok{(s}\SpecialCharTok{$}\NormalTok{gs)}
\NormalTok{  s[s}\SpecialCharTok{$}\NormalTok{lg }\SpecialCharTok{\textgreater{}=} \FunctionTok{max}\NormalTok{(}\FunctionTok{tapply}\NormalTok{(s}\SpecialCharTok{$}\NormalTok{lg, s}\SpecialCharTok{$}\NormalTok{kind, min)) }\SpecialCharTok{\&}\NormalTok{ s}\SpecialCharTok{$}\NormalTok{lg }\SpecialCharTok{\textless{}=} \FunctionTok{min}\NormalTok{(}\FunctionTok{tapply}\NormalTok{(s}\SpecialCharTok{$}\NormalTok{lg, s}\SpecialCharTok{$}\NormalTok{kind, max)) }\SpecialCharTok{\&}
\NormalTok{    s}\SpecialCharTok{$}\NormalTok{np }\SpecialCharTok{\textgreater{}=} \FunctionTok{max}\NormalTok{(}\FunctionTok{tapply}\NormalTok{(s}\SpecialCharTok{$}\NormalTok{np, s}\SpecialCharTok{$}\NormalTok{kind, min)) }\SpecialCharTok{\&}\NormalTok{ s}\SpecialCharTok{$}\NormalTok{np }\SpecialCharTok{\textless{}=} \FunctionTok{min}\NormalTok{(}\FunctionTok{tapply}\NormalTok{(s}\SpecialCharTok{$}\NormalTok{np, s}\SpecialCharTok{$}\NormalTok{kind, max)), ]}
\NormalTok{\})}
\NormalTok{w }\OtherTok{\textless{}{-}} \FunctionTok{reshape}\NormalTok{(ctl[, }\FunctionTok{c}\NormalTok{(}\StringTok{"kind"}\NormalTok{, }\StringTok{"label"}\NormalTok{, }\StringTok{"seed"}\NormalTok{, }\StringTok{"method"}\NormalTok{, }\StringTok{"rmse"}\NormalTok{)],}
             \AttributeTok{idvar =} \FunctionTok{c}\NormalTok{(}\StringTok{"kind"}\NormalTok{, }\StringTok{"label"}\NormalTok{, }\StringTok{"seed"}\NormalTok{), }\AttributeTok{timevar =} \StringTok{"method"}\NormalTok{,}
             \AttributeTok{direction =} \StringTok{"wide"}\NormalTok{)}
\NormalTok{w}\SpecialCharTok{$}\NormalTok{spread }\OtherTok{\textless{}{-}} \FunctionTok{abs}\NormalTok{(w}\SpecialCharTok{$}\NormalTok{rmse.pca }\SpecialCharTok{{-}}\NormalTok{ w}\SpecialCharTok{$}\NormalTok{rmse.isomap)}
\NormalTok{sp    }\OtherTok{\textless{}{-}} \FunctionTok{tapply}\NormalTok{(w}\SpecialCharTok{$}\NormalTok{spread, w}\SpecialCharTok{$}\NormalTok{kind, mean)}
\NormalTok{ratio }\OtherTok{\textless{}{-}}\NormalTok{ sp[[}\StringTok{"classic"}\NormalTok{]] }\SpecialCharTok{/}\NormalTok{ sp[[}\StringTok{"crease"}\NormalTok{]]}

\FunctionTok{set.seed}\NormalTok{(}\DecValTok{11}\NormalTok{)                       }\CommentTok{\# cluster bootstrap, resampling whole settings}
\NormalTok{boot }\OtherTok{\textless{}{-}} \FunctionTok{replicate}\NormalTok{(}\DecValTok{2000}\NormalTok{, \{}
\NormalTok{  r }\OtherTok{\textless{}{-}} \FunctionTok{vapply}\NormalTok{(}\FunctionTok{split}\NormalTok{(w, w}\SpecialCharTok{$}\NormalTok{kind), }\ControlFlowTok{function}\NormalTok{(d) \{}
\NormalTok{    l }\OtherTok{\textless{}{-}} \FunctionTok{sample}\NormalTok{(}\FunctionTok{unique}\NormalTok{(d}\SpecialCharTok{$}\NormalTok{label), }\FunctionTok{length}\NormalTok{(}\FunctionTok{unique}\NormalTok{(d}\SpecialCharTok{$}\NormalTok{label)), }\AttributeTok{replace =} \ConstantTok{TRUE}\NormalTok{)}
    \FunctionTok{mean}\NormalTok{(}\FunctionTok{unlist}\NormalTok{(}\FunctionTok{lapply}\NormalTok{(l, }\ControlFlowTok{function}\NormalTok{(x) d}\SpecialCharTok{$}\NormalTok{spread[d}\SpecialCharTok{$}\NormalTok{label }\SpecialCharTok{==}\NormalTok{ x])))}
\NormalTok{  \}, }\FunctionTok{numeric}\NormalTok{(}\DecValTok{1}\NormalTok{))}
\NormalTok{  r[[}\StringTok{"classic"}\NormalTok{]] }\SpecialCharTok{/}\NormalTok{ r[[}\StringTok{"crease"}\NormalTok{]]}
\NormalTok{\})}
\NormalTok{ci }\OtherTok{\textless{}{-}} \FunctionTok{quantile}\NormalTok{(boot, }\FunctionTok{c}\NormalTok{(}\FloatTok{0.025}\NormalTok{, }\FloatTok{0.975}\NormalTok{))}

\NormalTok{fit  }\OtherTok{\textless{}{-}} \FunctionTok{lm}\NormalTok{(spread }\SpecialCharTok{\textasciitilde{}} \FunctionTok{I}\NormalTok{(kind }\SpecialCharTok{==} \StringTok{"classic"}\NormalTok{), }\AttributeTok{data =}\NormalTok{ w)   }\CommentTok{\# cluster{-}robust by setting}
\NormalTok{X }\OtherTok{\textless{}{-}} \FunctionTok{model.matrix}\NormalTok{(fit); u }\OtherTok{\textless{}{-}} \FunctionTok{residuals}\NormalTok{(fit); bread }\OtherTok{\textless{}{-}} \FunctionTok{solve}\NormalTok{(}\FunctionTok{crossprod}\NormalTok{(X))}
\NormalTok{meat }\OtherTok{\textless{}{-}} \FunctionTok{Reduce}\NormalTok{(}\StringTok{\textasciigrave{}}\AttributeTok{+}\StringTok{\textasciigrave{}}\NormalTok{, }\FunctionTok{lapply}\NormalTok{(}\FunctionTok{unique}\NormalTok{(w}\SpecialCharTok{$}\NormalTok{label), }\ControlFlowTok{function}\NormalTok{(l) \{}
\NormalTok{  i }\OtherTok{\textless{}{-}}\NormalTok{ w}\SpecialCharTok{$}\NormalTok{label }\SpecialCharTok{==}\NormalTok{ l}
  \FunctionTok{tcrossprod}\NormalTok{(}\FunctionTok{crossprod}\NormalTok{(X[i, , }\AttributeTok{drop =} \ConstantTok{FALSE}\NormalTok{], u[i]))}
\NormalTok{\}))}
\NormalTok{tval }\OtherTok{\textless{}{-}} \FunctionTok{coef}\NormalTok{(fit)[}\DecValTok{2}\NormalTok{] }\SpecialCharTok{/} \FunctionTok{sqrt}\NormalTok{((bread }\SpecialCharTok{\%*\%}\NormalTok{ meat }\SpecialCharTok{\%*\%}\NormalTok{ bread)[}\DecValTok{2}\NormalTok{, }\DecValTok{2}\NormalTok{])}
\NormalTok{pval }\OtherTok{\textless{}{-}} \DecValTok{2} \SpecialCharTok{*} \FunctionTok{pt}\NormalTok{(}\SpecialCharTok{{-}}\FunctionTok{abs}\NormalTok{(tval), }\AttributeTok{df =} \FunctionTok{length}\NormalTok{(}\FunctionTok{unique}\NormalTok{(w}\SpecialCharTok{$}\NormalTok{label)) }\SpecialCharTok{{-}} \DecValTok{2}\NormalTok{)}
\NormalTok{mscore }\OtherTok{\textless{}{-}} \FunctionTok{aggregate}\NormalTok{(}\FunctionTok{cbind}\NormalTok{(rmse, qnx) }\SpecialCharTok{\textasciitilde{}}\NormalTok{ kind }\SpecialCharTok{+}\NormalTok{ method, ctl, mean)}
\NormalTok{score  }\OtherTok{\textless{}{-}} \ControlFlowTok{function}\NormalTok{(k, m, col) mscore[[col]][mscore}\SpecialCharTok{$}\NormalTok{kind }\SpecialCharTok{==}\NormalTok{ k }\SpecialCharTok{\&}\NormalTok{ mscore}\SpecialCharTok{$}\NormalTok{method }\SpecialCharTok{==}\NormalTok{ m]}
\end{Highlighting}
\end{Shaded}

9 crease settings and
5 Swiss rolls survive the
double restriction. Inside it \textbf{the ranking is identical}: Isomap first on both
families, PCA and classical MDS behind it on both. Nothing flips.

What differs is the spread. The mean gap between PCA and Isomap is
0.574 on the classical side against
0.113 on the crease side --- a ratio of
5.06, with a cluster bootstrap interval over settings of
1.7 to 9.8 and a
cluster-robust \(p\) of 0.02. The absolute scores go the same
way: PCA at 0.922 on a roll
against 0.286 on a crease pattern,
Isomap at 0.348 against
0.183.

State it in the book's own voice. \textbf{Crease patterns are not a harder benchmark;
they are an easier one that discriminates less.} A benchmark earns its keep by
telling methods apart, and on that criterion the Swiss roll wins by a factor of
about five.

What that costs a reader is worth being concrete about. A benchmark whose methods
sit 0.113 apart needs more replicates to separate
two of them than one whose methods sit 0.574
apart. Chapter \ref{results} turns that into a statistic of its own and reports
where along the \(\theta\) axis this benchmark has any resolving power at all;
read next to the ratio above, the two are the same observation from opposite
ends. It does not make the crease pattern useless --- a benchmark whose
answer key is exact can be asked questions a benchmark whose answer key is a
convention cannot, and Chapter \ref{evaluation} asks them --- but anyone choosing
between the two for a straight method comparison should choose the roll, and this
chapter is where the book says so.

\begin{figure}

{\centering \includegraphics[width=0.9\linewidth,alt={A line chart with the gap-to-spacing ratio on a reversed logarithmic horizontal axis, so difficulty increases to the right, and the mean absolute difference between PCA and Isomap reconstruction error on the vertical axis from zero to about one. The Swiss roll curve climbs from about 0.4 at the left to a peak near 0.9 in the middle, then drops almost vertically to zero and stays there across the whole right-hand half. The Miura curve is nearly flat just above zero across its entire range, never exceeding about 0.05.}]{11-versus-swiss-roll_files/figure-latex/comp-spread-1} 

}

\caption{How far apart each benchmark holds PCA and Isomap, across its whole sweep and at every seed the grids carry. Difficulty rises to the right. The Miura curve stays near the floor throughout. The Swiss roll separates the two methods widely over a middle window and then stops separating them at all, because past that window Isomap short-circuits and fails as completely as PCA.}\label{fig:comp-spread}
\end{figure}

Figure \ref{fig:comp-spread} attaches a caveat to that factor of five. The Swiss
roll's discriminating power is not monotone in its own difficulty. It reaches
0.909 in a middle window
and then collapses to
0.0003 at the tightest
rolls, where Isomap short-circuits and fails as completely as PCA does --- so past
that window the classical benchmark separates nothing either. The settings the
head-to-head matched on sit inside the peak. That is the right place to measure a
benchmark, on its best window rather than its worst, but the number is a
statement about Swiss rolls there and not about Swiss rolls generally.

\section{Diagnostics}\label{comp-diagnostics}

One objection deserves an answer before a reader raises it, and it is specific
enough that a footnote would not do.

E1's decisive arm sampled with \texttt{boundary\ =\ TRUE}, because a two-by-two Miura
leaves no room for the interior margin \texttt{sample\_manifold()} normally draws inside.
A Miura's unfolded outline is not convex, so for a pair straddling one of its
teeth the straight chart segment leaves the sheet and chart distance is a \emph{lower
bound} on the geodesic --- Chapter \ref{folding-geometry} measures this. An answer
key that understates long distances penalises the geodesic method specifically,
which is exactly the comparison being made.

\begin{Shaded}
\begin{Highlighting}[]
\NormalTok{labs }\OtherTok{\textless{}{-}} \FunctionTok{unique}\NormalTok{(ctl}\SpecialCharTok{$}\NormalTok{label[ctl}\SpecialCharTok{$}\NormalTok{kind }\SpecialCharTok{==} \StringTok{"crease"}\NormalTok{])}
\NormalTok{exit }\OtherTok{\textless{}{-}} \FunctionTok{vapply}\NormalTok{(labs, }\ControlFlowTok{function}\NormalTok{(l) \{}
\NormalTok{  q   }\OtherTok{\textless{}{-}} \FunctionTok{as.numeric}\NormalTok{(}\FunctionTok{strsplit}\NormalTok{(}\FunctionTok{sub}\NormalTok{(}\StringTok{"\^{}miura"}\NormalTok{, }\StringTok{""}\NormalTok{, l), }\StringTok{"/"}\NormalTok{)[[}\DecValTok{1}\NormalTok{]])}
\NormalTok{  pat }\OtherTok{\textless{}{-}} \FunctionTok{miura\_ori}\NormalTok{(}\FunctionTok{as.integer}\NormalTok{(q[}\DecValTok{1}\NormalTok{]), }\FunctionTok{as.integer}\NormalTok{(q[}\DecValTok{1}\NormalTok{]), }\AttributeTok{alpha =}\NormalTok{ q[}\DecValTok{2}\NormalTok{])}
  \FunctionTok{chart\_exit\_fraction}\NormalTok{(}\FunctionTok{sample\_manifold}\NormalTok{(pat, }\AttributeTok{theta =}\NormalTok{ q[}\DecValTok{3}\NormalTok{], }\AttributeTok{n =} \DecValTok{600}\DataTypeTok{L}\NormalTok{,}
                                      \AttributeTok{seed =}\NormalTok{ BENCH\_SEEDS[}\DecValTok{1}\NormalTok{], }\AttributeTok{boundary =} \ConstantTok{TRUE}\NormalTok{), pat)}
\NormalTok{\}, }\FunctionTok{numeric}\NormalTok{(}\DecValTok{1}\NormalTok{))}
\NormalTok{loo }\OtherTok{\textless{}{-}} \FunctionTok{vapply}\NormalTok{(labs, }\ControlFlowTok{function}\NormalTok{(l) \{}
\NormalTok{  k }\OtherTok{\textless{}{-}}\NormalTok{ w[w}\SpecialCharTok{$}\NormalTok{kind }\SpecialCharTok{==} \StringTok{"classic"} \SpecialCharTok{|}\NormalTok{ w}\SpecialCharTok{$}\NormalTok{label }\SpecialCharTok{!=}\NormalTok{ l, ]}
\NormalTok{  m }\OtherTok{\textless{}{-}} \FunctionTok{tapply}\NormalTok{(k}\SpecialCharTok{$}\NormalTok{spread, k}\SpecialCharTok{$}\NormalTok{kind, mean)}
\NormalTok{  m[[}\StringTok{"classic"}\NormalTok{]] }\SpecialCharTok{/}\NormalTok{ m[[}\StringTok{"crease"}\NormalTok{]]}
\NormalTok{\}, }\FunctionTok{numeric}\NormalTok{(}\DecValTok{1}\NormalTok{))}
\end{Highlighting}
\end{Shaded}

Between 0.00\% and
3.47\% of sampled pairs are affected. Dropping the
worst-affected setting moves the ratio from 5.06 to
5.11, and over the
9 leave-one-out refits it stays between
4.41 and 5.80. The
approximation biases against the finding rather than for it, and no single
setting carries it.

The second question a reader is entitled to ask is how much of the ratio is seed
noise. Every cell above is an average over seeds --- 5 in
the controlled arm, 20 in the two grids --- and the
within-family scatter is large: the standard deviation of the PCA-to-Isomap gap
across settings and seeds is 0.393
on the classical side against
0.070 on the crease side.
That scatter is why the interval quoted above comes from a bootstrap over whole
settings rather than over rows, and why the \(p\)-value is cluster-robust: seeds
within a setting are not independent draws of the quantity being estimated, and
treating them as such would report an interval several times too narrow.

\section{Where it fails}\label{comp-limits}

\textbf{The overlap is a corner.} It holds
9 crease settings,
5 rolls, three methods and two
metrics. The bootstrap interval ---
1.7 to 9.8 --- is the honest
width of the claim, and it is wide. The head-to-head above is broader in methods
and narrower in control. Neither is the study a settled answer would need.

\textbf{Under the neighbourhood metric the families nearly do collapse.} Isomap scores
0.773 on creases against
0.778 on Swiss rolls. The whole
difference between the families sits in how badly the linear methods fail, and
that is a fact about ambient non-planarity rather than about creases.

\textbf{The pre-registration got nothing right.} \texttt{PLAN.md} drafted three outcomes
before the run. The families did not collapse; they did not separate in the
registered sense, which was \emph{different rankings}, and the ranking is identical;
and arm B's null settled neither. The replacement claim is in this book because
it fits the evidence, not because its antecedent was met, and recording that is
what pre-registration is for.

\textbf{The Swiss roll is still the right tool for two jobs.} It is the clearest
picture in the subject and it belongs in a first lecture, and a pipeline that
fails on it is broken. What it cannot be is the evidence, because one fixed
geometry has nowhere to put a curve.

\section{Reproduce this}\label{comp-reproduce}

\texttt{scripts/run-classic-grid.R} writes \texttt{classic-grid.rds}, \texttt{scripts/experiment-e1.R}
the three E1 artefacts, \texttt{scripts/run-benchmark-grid.R} the crease grid; all are
committed with provenance. \texttt{tests/testthat/test-baselines.R} asserts the
arc-length closed form, the Fresnel error, the sphere's stretch factor and the
refusal printed above.

\section{Further reading}\label{comp-reading}

The Swiss roll and the S-curve enter with Isomap \citep{tenenbaum2000isomap} and LLE
\citep{roweis2000lle}; Lee and Verleysen \citep{lee2007nonlinear} is the textbook treatment
and the place to see how settled their role became. Balasubramanian and Schwartz
\citep{balasubramanian2002isomap} is one page long and is the origin of the
gap-to-spacing argument this chapter's first axis rests on.

On the answer key, Schoeneman et al. \citep{schoeneman2017error} is the paper that
noticed the Archimedean roll's chart was not isometric and proposed the Euler
replacement, and Mahapatra and Chandola \citep{mahapatra2018learning} use it. This
chapter's contribution is not the observation but its size.

On what a comparison of this kind can establish, Espadoto et al.
\citep{espadoto2021survey} survey the quantitative benchmarks and Machado et al.
\citep{machado2025necessary} argue that the projection quality metrics they rest on
are necessary but not sufficient --- a caution that applies with full force to the
five-fold separation reported above. The dispute about reading two-dimensional
embeddings at all \citep{chari2023specious, lause2024elephant} is the reason this
book measures rather than argues.

\part{Application}\label{part-application}

\chapter{From benchmark to decision}\label{application}

\emph{R. Heller · Date: 2026-07-29 · Last changed: 2026-11-09}

\emph{What happens when the data are real?}

\textbf{Chapter status: planned.} The application below has not yet been carried out. No single-cell comparison or improvement over a default method is claimed.

\begin{itemize}
\tightlist
\item
  Use the Chapter \ref{results} grid to select a method and its hyperparameters.
\item
  Apply that choice to a single-cell dataset.
\item
  Compare against the field-default choice.
\item
  Report a null result honestly if the benchmark-guided choice does not win.
\item
  Closes with visualisation and honest-reporting standards: alt text, no colour-only encoding, seed disclosure, aspect-ratio discipline.
\end{itemize}

\appendix


\chapter{Notation}\label{notation}

\emph{R. Heller · Date: 2026-07-29 · Last changed: 2026-08-30}

\section*{Symbols}\label{notation-symbols}
\addcontentsline{toc}{section}{Symbols}

{\def\LTcaptype{none} % do not increment counter
\begin{longtable}[]{@{}
  >{\raggedright\arraybackslash}p{(\linewidth - 2\tabcolsep) * \real{0.5000}}
  >{\raggedright\arraybackslash}p{(\linewidth - 2\tabcolsep) * \real{0.5000}}@{}}
\toprule\noalign{}
\begin{minipage}[b]{\linewidth}\raggedright
Symbol
\end{minipage} & \begin{minipage}[b]{\linewidth}\raggedright
Meaning
\end{minipage} \\
\midrule\noalign{}
\endhead
\bottomrule\noalign{}
\endlastfoot
\(\theta\) & Folding parameter; the difficulty axis. \(\theta = 0\) is the flat sheet \\
\(\rho\) & Dihedral angle at a crease. \(\rho = \pi\) when \(\theta = 0\); derived from \(\theta\), not a second name for it \\
\(a, b, \alpha\) & Miura-ori cell parameters \\
\(M, V\) & Mountain and valley creases \\
\(k\) & Neighbourhood size for kNN-graph methods \\
\(D\) & Ambient dimension after linear lift \\
\(\mathcal{M}\) & The folded manifold in \(\mathbb{R}^3\) \\
\(\mathcal{U}\) & The exact unfolding --- ground truth in \(\mathbb{R}^2\) \\
\(d_A\) & Ambient distance, measured in \(\mathbb{R}^3\) through the fold \\
\(d_U\) & Intrinsic (path) distance, measured along the surface --- equal to Euclidean distance in \(\mathcal{U}\) \\
\end{longtable}
}

\chapter{Datasets and codebook}\label{datasets}

\emph{R. Heller · Date: 2026-07-29 · Last changed: 2026-08-31}

\section*{Generated data}\label{datasets-generated}
\addcontentsline{toc}{section}{Generated data}

All crease-pattern manifolds are generated, not collected. Every dataset in
Parts I--III is reproducible from a pattern specification, a fold angle, a
sample size, a noise specification, and a seed.

\section*{Precomputed results}\label{datasets-precomputed}
\addcontentsline{toc}{section}{Precomputed results}

\texttt{data/processed/benchmark-grid.rds} --- the Chapter \ref{results} grid.
Provenance: \texttt{scripts/run-benchmark-grid.R}.

\texttt{data/processed/product-grid.rds} --- the product suite behind Chapter
\ref{benchmarks}'s irreducible-loss bound, in two arms: a product of two Miura
sheets and a product of two Swiss rolls. Provenance:
\texttt{scripts/run-product-grid.R}.

\texttt{data/processed/e1-difficulty.rds}, \texttt{e1-controlled.rds}, \texttt{e1-armB.rds} --- the
three arms of E1, which retired the book's original headline claim. Provenance:
\texttt{scripts/experiment-e1.R}.

\texttt{data/processed/part2-sweeps.rds} --- the same cell as the main grid, swept over
neighbourhood size instead of over noise. This is where Chapters
\ref{linear}, \ref{geodesic} and \ref{neighbor} answer how much of a
method's result is the method and how much is the \(k\) it was handed. Provenance:
\texttt{scripts/run-part2-sweeps.R}.

\texttt{data/processed/classic-grid.rds} --- the Swiss roll, the S-curve and the severed
sphere, scored by the same function, the same metrics and the same floor as the
crease patterns. Chapter \ref{comparison} needs it: a comparison in which the
two sides are scored by two code paths is not a comparison. Provenance:
\texttt{scripts/run-classic-grid.R}.

\texttt{data/processed/evaluator-audit-pilot.rds} --- the narrow pilot that decided
whether Chapter \ref{evaluation} had a subject. Provenance:
\texttt{scripts/run-evaluator-audit.R\ -\/-pilot}.

\texttt{data/processed/seed-budget.rds} --- the within-cell spread of every method, and
the seed count each one needs before a difference the size of the pre-registered
reportable difference (\texttt{MIN\_REPORTABLE}, fixed in \texttt{R/selection.R}) can be
detected. This is what standing rule 1's floor is measured against rather than
assumed. Provenance: \texttt{scripts/measure-seed-budget.R}.

Every one of these carries a provenance block recording the \texttt{R/} tree hash, the
R version, the BLAS, the package versions, the wall clock and whether the tree
was dirty when it ran. \texttt{read\_run()} refuses one that does not, refuses one
written by a \texttt{-\/-quick} smoke test, and refuses one whose contents are not the
contents the chapter around the call was written against.

\section*{External data}\label{datasets-external}
\addcontentsline{toc}{section}{External data}

One dataset, used only in Chapter \ref{application}.

\textbf{Fresh 68k PBMCs (Donor A)} --- peripheral blood mononuclear cells from a single
healthy donor, from Zheng et al. \citep{zheng2017massively}. Obtained from the 10x
Genomics dataset repository, Cell Ranger 1.1.0 processing, Chromium Universal
3' chemistry.

\texttt{scripts/prepare-single-cell.R} downloads the filtered matrix, applies the
preprocessing recorded in that script, and writes the artefact Chapter
\ref{application} reads. It is not run in CI and not run at render time, on the
same provenance convention as the benchmark grid: run it once, commit the
result, keep the script as the record of how.

\subsection*{Licence, and what may be redistributed}\label{licence-and-what-may-be-redistributed}
\addcontentsline{toc}{subsection}{Licence, and what may be redistributed}

The dataset is published under \textbf{CC BY 4.0}, which permits redistribution with
attribution. So the processed matrix \emph{may} ship in this repository, and the
constraint on whether it does is size rather than licence --- the full filtered
matrix is far too large for git, and Chapter \ref{application} does not need
it. What is committed is the reduced matrix the chapter actually reads,
produced by the script above, with the reduction stated there and in the
chapter.

\subsection*{Label provenance, and why the comparison is not circular}\label{label-provenance-and-why-the-comparison-is-not-circular}
\addcontentsline{toc}{subsection}{Label provenance, and why the comparison is not circular}

This matters more than it looks, and it was an open question until it was
checked. Chapter \ref{application}'s only external evaluation channel is the
cell-type labels. If those labels had themselves been derived by clustering one
of the embeddings under test, the comparison would be circular and the chapter
would be measuring its own assumptions.

They were not. Zheng et al.~generated reference profiles by sequencing ten
bead-enriched purified subpopulations from the same donor, separately, and
assigned each cell in the 68k population the identity of the purified
population with which it had the highest Spearman correlation. The labels come
from a separate set of experiments, and no clustering or embedding of the
matrix under test enters the assignment.

Two caveats belong in the chapter rather than in a footnote, because both bound
what a per-label result can mean.

\begin{itemize}
\tightlist
\item
  The assignment is made in gene-expression space --- the same space the
  embeddings are computed from. It is independent of the \emph{embedding}, which is
  the circularity that would matter here, but it is not an independent
  measurement of cell identity, and the chapter should not describe it as
  ground truth. The book reserves that phrase for the crease patterns, where it
  is earned.
\item
  The authors flag their own assignments as unreliable in specific places: the
  gap between the highest and second-highest correlation is small for some
  cells, cytotoxic T against NK being the example they give, and the purified
  populations are nested --- CD4+ T-helper contains all CD4+ cells --- so a
  hand-written rule reassigns cells to the next-highest correlation within a
  nested family. Any per-label accuracy in Chapter \ref{application} inherits
  that, and the pre-registered selection rule of Chapter \ref{results} is what
  keeps it from being read as a ranking.
\end{itemize}

\chapter*{Glossary}\label{glossary}
\addcontentsline{toc}{chapter}{Glossary}

\emph{R. Heller · Date: 2026-07-29 · Last changed: 2026-08-21}

\begin{description}
\tightlist
\item[\textbf{Developable surface}]
A surface that can be flattened into a plane without stretching or tearing.
\item[\textbf{Dihedral angle (\(\rho\))}]
The angle between the two facets meeting at a crease, measured in the ambient
space. A crease in an unfolded sheet has \(\rho = \pi\), not 0, and \(\rho\)
decreases as the sheet folds. It is \emph{derived} from \(\theta\) and is never a
synonym for it.
\item[\textbf{Folding parameter (\(\theta\))}]
The parameter of the folding map, and the difficulty axis of the whole book.
\(\theta = 0\) is the flat sheet, and the pattern folds monotonically as
\(\theta\) grows. It is a parameter, not an angle read off a figure; the map
\(\rho(\theta)\) from it to the dihedral angle is stated once per pattern in
Chapter \ref{folding-geometry}.
\item[\textbf{Isometry, intrinsic and ambient}]
An \emph{intrinsic} isometry preserves distances measured \textbf{along} the surface.
An \emph{ambient} isometry preserves straight-line distance in the space the
surface sits in. Rigid folding is the first and emphatically not the second:
it contracts ambient distance strictly for every \(\theta > 0\). Which of the
two a method consumes is the distinction the book turns on, and the reason
Chapter \ref{evaluation} reports two headline numbers rather than one.
\item[\textbf{Mountain / valley assignment}]
The direction of each crease. Written M and V; encoded by linetype as well as
colour in every figure.
\item[\textbf{Reach}]
The largest \(r\) such that every point within \(r\) of the surface has a unique
nearest point on it. Most manifold-learning guarantees are stated in terms of
a positive reach; a creased surface has reach exactly zero, which is what puts
these objects outside the theory rather than merely at its edge.
\item[\textbf{Rigid origami}]
Folding in which facets stay flat and all deformation happens at creases.
\item[\textbf{Short-circuit}]
Two points close in ambient space but far apart geodesically. The failure mode
that breaks kNN-graph methods.
\end{description}

\chapter*{Colophon}\label{colophon}
\addcontentsline{toc}{chapter}{Colophon}

\emph{R. Heller · Date: 2026-07-29 · Last changed: 2026-07-29}

Written in R Markdown and built with \textbf{bookdown} \citep{R-bookdown}. Typeset in
Inter and JetBrains Mono. Figures use the viridis plasma palette (option \texttt{"C"}),
the series convention for this volume.

Interface icons are Font Awesome Free 6.5.1, vendored in \texttt{style/webfonts/}
rather than loaded from a content delivery network. The icons are licensed
CC BY 4.0, the fonts under the SIL OFL 1.1, and the accompanying code under the
MIT License; see \url{https://fontawesome.com/license/free}.

\begin{verbatim}
#> - Session info ---------------------------------------------------------------
#>  setting  value
#>  version  R version 4.6.2 (2026-10-29)
#>  os       Ubuntu 24.04.5 LTS
#>  system   x86_64, linux-gnu
#>  ui       X11
#>  language (EN)
#>  collate  C.UTF-8
#>  ctype    C.UTF-8
#>  tz       UTC
#>  date     2026-11-09
#>  pandoc   3.8.3 @ /opt/hostedtoolcache/pandoc/3.8.3/x64/ (via rmarkdown)
#>  quarto   NA
#> 
#> ------------------------------------------------------------------------------
\end{verbatim}

\chapter*{About the author}\label{about-the-author}
\addcontentsline{toc}{chapter}{About the author}

\emph{R. Heller · Date: 2026-07-29 · Last changed: 2026-07-29}

R. Heller is a physician and biomedical researcher working on trauma
immunology, laboratory diagnostics, and reproducible analysis workflows.
ORCID: \href{https://orcid.org/0000-0001-8006-9742}{0000-0001-8006-9742}.

\chapter*{Citing this book}\label{citing}
\addcontentsline{toc}{chapter}{Citing this book}

\emph{R. Heller · Date: 2026-08-21 · Last changed: 2026-08-21}

This book is published under CC BY 4.0, which asks for attribution and nothing
else. Quote it, teach from it, adapt it, use its figures --- the licence permits
all of that, and a citation is the whole of what is expected in return.

Cite the book itself for the argument and the benchmark design. Cite the
repository when what you are relying on is the code: the pattern constructors,
the folding maps, the metrics, or a committed artefact under \texttt{data/processed/}.
They carry the same version number, and the repository is the only one of the
two that pins a commit.

\begin{boxempty}
Heller, R. (2026). \emph{Folding in R: Origami, Manifolds, and Dimension
Reduction.} Version 0.1.0.
\url{https://r-heller.github.io/folding-in-r/}
\end{boxempty}

The BibTeX entry is in \texttt{citation.bib}, and \texttt{CITATION.cff} carries the same
metadata in a form GitHub and Zenodo read directly:

\begin{Shaded}
\begin{Highlighting}[]
\VariableTok{@book}\NormalTok{\{}\OtherTok{heller2026folding}\NormalTok{,}
  \DataTypeTok{title}\NormalTok{     = \{Folding in R: Origami, Manifolds, and Dimension Reduction\},}
  \DataTypeTok{author}\NormalTok{    = \{Heller, R.\},}
  \DataTypeTok{year}\NormalTok{      = \{2026\},}
  \DataTypeTok{publisher}\NormalTok{ = \{Self{-}published\},}
  \DataTypeTok{url}\NormalTok{       = \{https://r{-}heller.github.io/folding{-}in{-}r/\},}
  \DataTypeTok{note}\NormalTok{      = \{Version 0.1.0\}}
\NormalTok{\}}
\end{Highlighting}
\end{Shaded}

Both files are generated from the same source and should not be edited by hand;
if a version number here disagrees with one of them, the files are right.

\chapter*{References}\label{references}
\addcontentsline{toc}{chapter}{References}

\emph{R. Heller · Date: 2026-07-29 · Last changed: 2026-07-29}

\bibliography{book.bib,packages.bib}

\end{document}
